\clearpage
\appendix
\begingroup
\makeatletter
\def\input@path{{ordinary/}{cohomology/}}
\makeatother
\newcommand{\G}{G_{\mathcal F}}
\newcommand{\Ima}{\operatorname{Im}}
\setlength{\emergencystretch}{2em}

\part*{Technical appendices}
\phantomsection
\addcontentsline{toc}{part}{Technical appendices}

The first appendices construct the prime pair and prove the
strict cohomology bound over the original field. The later
ones supply the ordinary automorphic and deformation-theoretic
inputs to the odd-prime count. In the selection of the two auxiliary primes, \(N\) denotes the conductor of the fixed primitive Dirichlet character \(\psi\) and \(q,d\) are the two auxiliary primes. The later
counting arguments specify their own moduli and let the
coefficient prime vary. Section~\ref{sec:family-window}
restricts the resulting bound to the rank used in the numerator.

\section{Selecting the two auxiliary primes}\label{sec:hosts}

Fourier projection leaves terms twisted by either of two auxiliary quadratic
characters or by their product. Their untwisted factors are the primitive characters
inducing powers of the fixed primitive character \(\psi\).
We select the auxiliary primes in prescribed progressions
using correlation estimates for the prime sums and a
simultaneous exclusion of bad pairs. For the selected pair
the unwanted modes are small throughout the interval used
in Section~\ref{sec:comparison}. Appendix~\ref{sec:cohomology}
then supplies a cohomology subinterval for this same pair. Definition~\ref{hyp:host-residue-class} specifies the congruence and size conditions on the candidate primes. The pair is selected from the two logarithmic intervals in \eqref{eq:host-pair-log-windows}.

\subsection{Correlation in a progression}

Fix the progression from which a prime will be chosen. We first average the Jacobi symbol over its candidate primes for each admissible squarefree denominator, allowing one exceptional denominator. The choice of a pair for the whole interval comes later.

\begin{paramdefinition}[Parameter range for the progression average]\label{hyp:host-residue-class}
Fix \(0<\varepsilon<1/52\) as in Proposition~\ref{prop:quadratic-prime-pair}. Let \(Y_0\) be the prime selection scale, \(J\ge1\) a fixed exponent,
and \(Z=Y_0^{400J}\). Let \(m_0,a_0\) be integers with \(8\mid m_0\) and \(\gcd(a_0,m_0)=1\). The primes under consideration form the set
\[
 \mathcal P=\{l\text{ prime}:\sqrt U<l\le U,\ l\equiv a_0\pmod{m_0}\}.
\] The admissible range of \(U\) for this modulus is
\[
 \log U\ge\tfrac14(\log Y_0)^2\log m_0,\qquad
 Y_0^{\varepsilon/6}\le\log U
       \le Y_0^{\varepsilon/5}(\log Y_0)^2 .
\]
The moduli in this range are restricted to those for which each prime divisor of \(m_0\) is at most
\(Y_0^{\varepsilon/5}\) or exceeds \(Z\).
\end{paramdefinition}
For the first selection take \(m_0=8\) and
\(\log U=Y_0^{\varepsilon/5}\); all inequalities hold for large
\(Y_0\). For the second selection \eqref{eq:host-pair-log-windows}
and the modulus calculation following it verify the same range
with the first prime already fixed.

\begin{proposition}[Quadratic correlation in one progression]\label{prop:quadratic-host-correlation}
For the data of Definition~\ref{hyp:host-residue-class}, for any fixed \(B>0\)
and all sufficiently large \(Y_0\), the relative
prime number theorem and the exceptional-character alternative give
\begin{equation}\label{eq:quadratic-corr}
       \left|\sum_{l\in\mathcal P}(l/s)\right|
       \le Y_0^{-B}\abs{\mathcal P}
\end{equation}
for all odd squarefree \(1<s\le Z^2\) coprime to \(m_0\), with at most
one exceptional value of \(s\). Here \((\cdot/s)\) is the Jacobi symbol.

\end{proposition}

\begin{proof}
\proofstep{Use the relative progression estimate.}

To obtain \eqref{eq:quadratic-corr}, use the relative progression
estimate at modulus \(m_0s\). The displayed hypotheses imply
\(\log(m_0s)=o(\log U)\) uniformly for \(s\le Z^2\).
For example, Corollary~1.4 of Thorner and Zaman applies at all
intermediate scales \(\sqrt U\le t\le U\), since
\(t\ge(m_0s)^{12}\) for sufficiently large \(Y_0\).

To pass from logarithmic weights to a count of primes, keep the exceptional character in the main term. Put \(m=m_0s\), \(V=\sqrt U\), and, for a reduced class \(b\bmod m\), set
\[
 h_b(t)=1-\xi(b)t^{\beta-1},\qquad F_b(t)=\int_0^t h_b(v)\,dv .
\]
If the modulus has no exceptional character use \(h_b=1\). For \(V\le t\le U\), Corollary~1.4 of Thorner and Zaman gives
\[
 \sum_{p\le t,\ p\equiv b\ (m)}\log p
     =\frac{F_b(t)}{\varphi(m)}(1+O(E)),\qquad
 E=e^{-c\log U/\log m}+e^{-c\sqrt{\log U}}
\]
with an absolute \(c>0\), decreased to allow the lower endpoint.
For \(\xi(b)=1\), the function \(h_b\) increases and
\(0\le F_b(t)\le t h_b(t)\) in this range. The lower inequality follows from
\(t^{1-\beta}\ge1/\beta\) for \(\log t\ge2\), since an exceptional zero has \(\beta>1/2\).
For \(\xi(b)=-1\), \(F_b(t)\le3t\) and \(h_b(t)\ge1\).

Thus partial summation gives
\[
 \#\{V<p\le U:p\equiv b\ (m)\}
 =\frac{I_b}{\varphi(m)}(1+O(E)),\qquad
 I_b=\int_V^U h_b(t)\frac{dt}{\log t}.
\]
Indeed its two endpoint errors and integral error together are
\(O(EU h_b(U)/(\varphi(m)\log U))\). This is a relative error:
the last half of the interval supplies
\[
 I_b\ge \int_{U/2}^U h_b(t)\frac{dt}{\log U}
       \ge \frac{U h_b(U)}{4\log U}.
\]
For the plus character value this follows from
\(1-e^{-z/2}\ge(1-e^{-z})/2\) with \(z=(1-\beta)\log U\).
\cite[Corollary~1.4 and (1.7)]{thorner-zaman}

Separate the two integrals in the exceptional case:
\[
 I=\int_V^U\frac{dt}{\log t},\quad I_\beta=\int_V^U t^{\beta-1}\frac{dt}{\log t},
 \qquad I_b=I-\xi(b)I_\beta .
\]

Under the Chinese remainder identification write
\(b=(a_0,c)\), \(c\in(\Z/s)^\times\), and
\(\xi(b)=\xi_0(a_0)\xi_s(c)\). The ordinary weighted
main term contains \(\sum_c(c/s)=0\). The exceptional
weighted term contains
\[
   \xi_0(a_0)\sum_{c\in(\Z/s)^\times}(c/s)\xi_s(c),
\]
which vanishes unless \(\xi_s=(\cdot/s)\). In that case
\(\sum_c\xi_s(c)=0\), so the unweighted main term is \(I/\varphi(m_0)\). Since \(t\ge\sqrt U\) throughout the integral,
\[
  0\le I_\beta/I\le U^{-(1-\beta)/2}.
\]
In both cases the sum of the errors is bounded by the
relative error factor times the unweighted main term,
because the individual progression main terms are nonnegative.

\proofstep{Use a single exceptional-character range.}

For all moduli take \(D=m_0Y_0^{800J}\). The lower bound on
\(\log U\) gives the two estimates
\[
 \log m_0\le\frac{4\log U}{(\log Y_0)^2},\qquad
 800J\log Y_0=o_J\!\left(\frac{\log U}{(\log Y_0)^2}\right).
\]
Thus, for large \(Y_0\),
\[
       \frac{\log U}{\log D}\ge\tfrac15(\log Y_0)^2 .
\]
If \(1-\beta\ge c/\log D\), the exceptional contribution therefore satisfies
\[
 I_\beta/I\le\exp\!\left(-\frac{c\log U}{2\log D}\right)
     \le \exp\!\left(-\tfrac c{10}(\log Y_0)^2\right)\ll_B Y_0^{-B}.
\]
Here \(B\) is any fixed exponent.

For nearer zeros, the Landau--Page alternative gives the same
primitive real character for all moduli up to \(D\).
Its restriction can equal the primitive quadratic character of
at most one such \(s\). Indeed, the restrictions for two squarefree
values differ on a prime dividing exactly one of them, as is seen
by taking a nonsquare there and the unit residue at the other primes.
Finally apply the same interval estimate at \(m_0\).
In the worst class the lower bound for \(I_b\) on \([U/2,U]\) gives
\[
 |\mathcal P|\gg\frac{U}{\varphi(m_0)\log U}
                 \min\{1,(1-\beta)\log U\}.
\]
Without an exceptional zero the last factor is one. With one, use
\(1-\beta\gg m_0^{-1/2}\) from \cite[(1.2), with exponent \(1/2\)]{thorner-zaman}.
Since \(\log m_0=o(\log U)\), this lower bound exceeds \(U^{9/10}\) for all sufficiently large \(Y_0\). These uses of relative prime
distribution and Page uniqueness concern the whole range of moduli.
\hfill\mbox{\cite{page,thorner-zaman}}\label{cor:cohomology-host-windows-1}

\end{proof}

\subsection{A mean square over the whole scale interval}

A bound at one scale would allow the exceptional primes to vary with the scale. To keep one set of primes throughout, we estimate the variation of the truncated sum.
For coefficients \(\abs{\kappa(n)}\le1\) put
\[
 S_l(y)=\sum_{n\le Z}\Lambda(n)\min(1,y/n)\,\kappa(n)(n/l),
 \qquad Y_0\le y\le Y_0^J .
\]
\begin{proposition}[Uniformity over the scale interval]\label{prop:quadratic-host-uniform}
For the data of Definition~\ref{hyp:host-residue-class}, fix \(\tau>0\) and a
list of bounded cardinality of coefficient sequences with
\(|\kappa(n)|\le1\) for every \(n\).
Then all but an \(O_\tau(Y_0^{-\tau})\)
fraction of \(\mathcal P\) satisfy, simultaneously for every sequence in the list, 
\[
 \lvert S_l(y)\rvert\le y^{1/2+\tau}/2\qquad(Y_0\le y\le Y_0^J).
\]

\end{proposition}

\begin{proof}
\proofstep{Truncation and coefficient bounds.}

Write \(A(t)=\sum_{j\ge2,\,p^j\le t}\log p\). Chebyshev's estimate
\(\sum_{p\le v}\log p\ll v\) gives, for \(t\ge2\),
\[
 A(t)\ll \sqrt t+t^{1/3}\log t\ll\sqrt t .
\]
Indeed the term with \(j=2\) gives the first bound and there
are at most \(\log_2 t\) remaining powers to sum. Including the
weighted tail, the unsigned contribution of higher powers is at most
\[
 A(y)+y\int_{(y,Z]}t^{-1}\,dA(t)
   =yA(Z)/Z+y\int_y^Z A(t)t^{-2}\,dt\ll\sqrt y .
\]
The prime contribution from \(p\mid m_0\) is
\(O(Y_0^{\varepsilon/5})\), since the large divisors exceed
\(Z\). For the other primes write
\(a_p=(\log p)\min(1,y/p)\kappa(p)\).
Chebyshev's bound and partial summation give
\[
 \sum_p|a_p|^2\ll
  \sum_{p\le y}(\log p)^2+y^2\sum_{p>y}\frac{(\log p)^2}{p^2}
  \ll y\log y,\qquad
 \sum_p|a_p|\ll y\log Z .
\]

\proofstep{Diagonal, exceptional pair and remaining pairs.}

On expanding \(\sum_l\abs{S_l(y)}^2\), the diagonal is bounded by
\(\abs{\mathcal P}\sum\abs{a_p}^2\). For distinct primes \(p,r\)
the product of symbols is \((l/(pr))\) times a fixed sign by
quadratic reciprocity, since the class modulo \(8\) is fixed.
The exceptional value in \eqref{eq:quadratic-corr} can correspond
to at most one unordered pair \(\{p,r\}\), which costs at most
\(\abs{\mathcal P}(\abs{a_p}^2+\abs{a_r}^2)\).
The other pairs cost \(\ll Y_0^{-B}\abs{\mathcal P}y^2(\log Z)^2\).
With \(B>J\), these estimates give
\[
 \sum_{l\in\mathcal P}\abs{S_l(y)}^2\ll\abs{\mathcal P}\,y\log y .
\]
\proofstep{Bounding the proportion of primes exceptional somewhere in a dyadic interval.}

This bound is uniform in the coefficients. Since \(yS_l'(y)\)
is the same sum restricted to \(n>y\) almost everywhere,
\(\sum_l\abs{S_l'(y)}^2\ll\abs{\mathcal P}\log y/y\).
For \(V\le W\le2V\) in the interval and arbitrary
\(w_l\in[V,W]\), the finite sum \(S_l\) is absolutely
continuous. Hence \(S_l(w_l)-S_l(V)=\int_V^{w_l}S_l'(y)\,dy\).
The square of this increment is bounded by
\((W-V)\int_V^W|S_l'(y)|^2\,dy\). Summing over \(l\)
gives
\[
 \sum_l\abs{S_l(w_l)}^2
   \le 2\sum_l\abs{S_l(V)}^2+
       2V\int_V^W\sum_l\abs{S'_l(y)}^2\,dy
   \ll\abs{\mathcal P}\,V\log(2V).
\]

Fix \(\tau>0\). For primes at which the bound
\(\abs{S_l(y)}\le y^{1/2+\tau}/2\) fails anywhere in
\([V,W]\), choose \(w_l\in[V,W]\) with \(\abs{S_l(w_l)}>w_l^{1/2+\tau}/2\). The fraction of these
primes is \(O(V^{-2\tau}\log(2V))\). Summing across dyadic blocks
shows that all but \(O_\tau(Y_0^{-\tau})\) of the primes meet the
bound for every \(y\in[Y_0,Y_0^J]\).
A union bound applies to any list of bounded cardinality.

\end{proof}

\subsection{The full smoothed sum and the shared interval}

The mean square controls a truncated polynomial. The Fourier comparison uses the full
smoothed sum and several twists at once. We now impose these requirements on the same
primes.

\begin{auxlemma}[Logarithmic derivative inside a zero-free rectangle]\label{lem:host-log-derivative}
Let \(Q\ge4\), \(3/4\le\alpha<1\), and \(0<\nu<1-\alpha\).
Suppose \(\chi\) is primitive and nonprincipal, of conductor \(f\le Q\).
If \(L(s,\chi)\) has no zero in
\(\RePart s\ge\alpha,\ |\Ima s|\le Q\), then
\[
 \left|\frac{L'}{L}(\sigma+it,\chi)\right|
   \ll(1+\nu^{-1})\log Q
     \qquad(\sigma\ge\alpha+\nu,\ |t|\le Q-1)
\]
with an absolute implied constant.
\end{auxlemma}

\begin{proof}
For \(\sigma\ge2\) the absolute Euler series gives the bound.
For \(\alpha+\nu\le\sigma\le2\) write \(s=\sigma+it\),
\(s_0=2+it\) and \(\rho=\beta+i\gamma\) for a nontrivial zero,
counted with multiplicity. Subtract the logarithmic derivatives of
the Hadamard product at \(s\) and \(s_0\). The conductor factor cancels.
The gamma factors contribute \(O(1)\), by Stirling's formula for
\(|t|\ge2\) and compactness for \(|t|\le2\). Consequently
\[
 \frac{L'}{L}(s,\chi)=O(1)+\sum_\rho
        \left(\frac1{s-\rho}-\frac1{s_0-\rho}\right).
\]
The unit-interval zero count gives
\(\#\{\rho:j\le|\gamma-t|<j+1\}\ll\log(f(3+|t|+j))\)
for integers \(j\ge0\). Both this estimate and the product are
given in \cite[Lemma~10.3 and Section~10.2]{classical-region}.

For \(j=0\) use \(|s-\rho|\ge\nu\) and \(|s_0-\rho|\ge1\).
For \(j\ge1\) use the difference of the two fractions. Its
absolute value is at most \(2/j^2\). The two contributions are
bounded, respectively, by
\[
 (1+\nu^{-1})\log(f(3+|t|)),\qquad
 \sum_{j\ge1}\frac{\log(f(3+|t|))+\log(1+j)}{j^2}
    \ll\log(f(3+|t|)).
\]
Now use \(f\le Q\) and \(|t|\le Q-1\).
\end{proof}

\begin{proposition}[Removing the truncation for a bounded list of twists]\label{prop:quadratic-host-smoothed}
Assume Definition~\ref{hyp:host-residue-class}. Fix \(\tau>0\)
and an upper bound for the cardinality of the list of primitive base
characters. Their conductors may vary subject to \(f\le F\le\sqrt U\).
There is a subset \(\mathcal P'\subseteq\mathcal P\) with
\(|\mathcal P\setminus\mathcal P'|/|\mathcal P|=o(1)\) as
\(Y_0\to\infty\), uniformly over such lists. For every \(l\in\mathcal P'\)
and each twist \(\chi\) of a listed character by \((\cdot/l)\),
\[
 \lvert G_\chi(y)\rvert\le y^{1/2+\tau}\qquad(Y_0\le y\le Y_0^J).
\]

\end{proposition}

\begin{proof}
\proofstep{Zero density is applied to primitive twists.}

For a primitive base character of conductor \(f\le F\le\sqrt U\),
the twist at \(l\in\mathcal P\) is primitive of conductor \(fl\).
Put \(Q=FU\) and \(\alpha=1-1/400\). Jutila's estimate (1.6) gives
\[
 \sum_{q\le Q}\ \sum_{\chi\bmod q}^{*}N(\alpha,Q,\chi)
     \ll (Q^3)^{2(1-\alpha)+1/100}\ll Q^{1/4} .
\]
The star denotes primitive characters and \(N(\alpha,T,\chi)\)
counts zeros with real part at least \(\alpha\) and height at most
\(T\) in absolute value. The exponent in Jutila's bound uses
\(\alpha\ge21/26\) and \(Q\ge2\), with the parameter \(k=2\).
Hence there are \(O(Q^{1/4})\) primitive characters of conductor
at most \(Q\) with a zero in that rectangle. The conductor \(fl\)
determines \(l\), as \(l>F\). Consequently the excluded primes
are \(O(U^{3/8})\), a vanishing fraction of \(\mathcal P\).
\hfill\mbox{\cite{jutila}}

\proofstep{Estimate the tail by a contour difference.}

For a remaining twist \(\chi\), write
\[
  R_\chi(z)=-\ell_\chi(1)-\sum_{n\le z}\Lambda(n)\chi(n)/n .
\]
Here \(\ell_\chi=L'/L(\,\cdot\,,\chi)\). Since \(f<l\) and the
quadratic twist is nontrivial at \(l\), the primitive twist is
nonprincipal. Lemma~\ref{lem:host-log-derivative} applies with
\(\alpha=1-1/400\) and \(0<\nu\le1/800\). Use the contour formula for \(S_\chi(z)\) on
\(\RePart s=2\) and take the difference at \(z+h\) and \(z\),
where \(0<h\le z\). Shift to \(1-1/400+\nu\) at heights
\(\pm Q/2\). The residue at \(1\) is \(-h\ell_\chi(1)\).

On the new edge \(\RePart s=\sigma_1=1-1/400+\nu\),
\[
 (z+h)^s-z^s=s\int_z^{z+h}u^{s-1}\,du,\qquad
 \int_{-Q/2}^{Q/2}\frac{dt}{|\sigma_1+it-1|}
      \ll_\nu\log Q .
\]
Since \(0<h\le z\), the first integral cancels the factor
\(s\) and leaves size at most \(h z^{\sigma_1-1}\).
The horizontal edges and the two tails on \(\RePart s=2\)
give \(O_\nu((\log Q) z^2/Q)\). Thus for
\(\sigma_1=1-1/400+\nu\)
\[
 \abs{S_\chi(z+h)-S_\chi(z)+h\ell_\chi(1)}
 \ll_\nu (\log Q)^2 h z^{\sigma_1-1}
                       +(\log Q) z^2/Q .
\]
The difference of the finite sums is
\(h\sum_{n\le z}\Lambda(n)\chi(n)/n\) plus an error from
\(z<n\le z+h\). Its modulus is at most
\(h(h+1)\log(z+h)/z\). Dividing by \(h\) bounds the tail \(R_\chi(z)\).

\proofstep{Return to the smoothed sum.}

Take \(X=Y_0^J\), \(z=X^{400}\), \(h=z/X^2\) and
\(\nu=\varepsilon/(1600J)\). Multiplication by \(X\) gives
\[
 \abs{X R_\chi(z)}
 \ll_{\varepsilon,J}(\log Q)^2Y_0^{\varepsilon/4}
      +(\log Q)X^{403}/Q+ X^{-1}\log X .
\]
Here \(Y_0^{\varepsilon/6}\le\log Q\le Y_0^{\varepsilon/4}\)
for sufficiently large \(Y_0\). The bound is eventually at most
\(Y_0^\varepsilon\). Also
\(\abs{\ell_\chi(1)}\ll_J\log Y_0\) by the definition of the
tail and the factorial identity \eqref{eq:factorial-mangoldt-bound}.

For \(z=Z\) split the defining sums at \(y\). The identity and the
two bounds, including Proposition~\ref{prop:quadratic-host-uniform}, give
\[
 \begin{aligned}
 G_\chi(y)&=-S_l(y)-yR_\chi(Z),\\
 |G_\chi(y)|&\le \tfrac12 y^{1/2+\tau}+Y_0^\varepsilon
             \le y^{1/2+\tau}\qquad (Y_0\le y\le X).
 \end{aligned}
\]
The last inequality holds for large \(Y_0\) since \(\varepsilon<1/52\). The constants and the choice of the prime apply to
every twist of the bounded list at once. To recover the explicit bound needed when the argument of the weight is shifted, use the estimate with \(0<\tau'<\min(\tau,1/2)\). The centered identity and \eqref{eq:factorial-mangoldt-bound} give, uniformly on the interval,
\[
 |\ell_\chi(1)|\le y^{-1/2+\tau'}+\sum_{n\le y}\Lambda(n)/n
              \le 4\log y
\]
for sufficiently large \(Y_0\).

\end{proof}

The remaining choice is simultaneous in the two primes. Once the first is chosen, the
twists needed for the product must also be included when choosing the second.

\begin{proposition}[A single prime pair for all primitive twists]\label{prop:common-host-prime-pair}
Fix the scale exponents \(J,\varepsilon\) of
Definition~\ref{hyp:host-residue-class}, \(\tau>0\), and a fixed
bounded list of primitive characters. For sufficiently large
\(Y_0\) there exist primes \(q\equiv3\bmod8\) and
\(d\equiv7\bmod16\), with logarithms in \eqref{eq:host-pair-log-windows}, satisfying 
\[
 q\mid d+1,\qquad p\mid d+1\quad(3\le p\le Y_0^{\varepsilon/5},\ p\text{ prime}).
\]
 The bound of
Proposition~\ref{prop:quadratic-host-smoothed} holds for the
primitive twists at \(q,d,qd\) for every \(Y_0\le y\le Y_0^J\).
\end{proposition}

\begin{proof}
\proofstep{Select the pair in order.}

For the first prime choose \(q\equiv3\bmod 8\), with
\(\log q\) in the first interval of \eqref{eq:host-pair-log-windows}.
Keeping this prime, choose \(d\equiv7\bmod16\) with \(d+1\)
divisible by \(q\) and all odd primes at most \(Y_0^{\varepsilon/5}\).
The two logarithm windows are
\begin{equation}\label{eq:host-pair-log-windows}
 \tfrac12Y_0^{\varepsilon/5}<\log q\le Y_0^{\varepsilon/5},\qquad
 \tfrac12Y_0^{\varepsilon/5}(\log Y_0)^2<\log d
           \le Y_0^{\varepsilon/5}(\log Y_0)^2 .
\end{equation}
Put \(Y=Y_0^{\varepsilon/5}\). The conditions on \(d\)
give a reduced class by the Chinese remainder theorem at
\[
 m_0=16q\prod_{3\le p\le Y}p,\qquad
 \log m_0\le (2+o(1))Y.
\]
For the second selection choose \(\log U=Y(\log Y_0)^2\).
The modulus size then gives
\[
 \tfrac14(\log Y_0)^2\log m_0
 \le (\tfrac12+o(1))Y(\log Y_0)^2\le\log U .
\]
The factor \(q\) exceeds \(Z\), and every other prime
divisor of \(m_0\) is at most \(Y\), as required in
Proposition~\ref{prop:quadratic-host-correlation}.

In the second application the base list consists of the original
list and its twists at the already selected \(q\).
If \(F_0\) is the largest original conductor, then the
largest conductor in the second list is at most \(F_0q\).
In the second selection \(F_0\) is fixed and
\[
  \log(F_0q)\le\log F_0+Y\le\tfrac12\log U
\]
for all sufficiently large \(Y_0\). The bad primes for that
list can therefore be discarded by
Proposition~\ref{prop:quadratic-host-smoothed}.
The first choice controls the twists at \(q\), and the
second choice controls the twists at \(d\) and at \(qd\).
This gives one pair for the whole interval \(Y_0\le y\le Y_0^J\).

\end{proof}

\begin{construction}
\begin{tabular}{@{}lll@{}}\toprule
Prime selected & Base list for the selection & New twists controlled\\\midrule
\(q\equiv3\bmod8\) & \(\chi\) & \(\chi\chi_q\)\\
\(d\equiv7\bmod16\) & \(\chi,\ \chi\chi_q\) &
\(\chi\chi_d,\ \chi\chi_q\chi_d\)\\\bottomrule
\end{tabular}
\caption{Selecting one ordered pair for all characters \(\chi\) in the fixed list and all scales \(Y_0\le y\le Y_0^J\). Here \(\chi_p\) is the quadratic character at \(p\), and each product denotes its inducing primitive character.}\label{fig:ordered-prime-twist-list}
\end{construction}

\begin{lemma}[Rescaling the shared interval]\label{lem:host-rescaling}
Let \(q,d\) be the pair of Proposition~\ref{prop:common-host-prime-pair}
and let \([W,W^2]\) be its cohomology interval from
Lemma~\ref{lem:exceptional-twist-window}, with \(Y_0=X^{1/J}\).
For the fixed \(c\ge1\) and \(a_*=c^{1/(2+8\varepsilon)}\), if
\(W\ge\max(a_*^3,a_*^{400})\), the dyadic frame and twist bounds transfer under the change of variables 
\[
 a_*W\le x\le (a_*W)^{3/2},\quad x'=x/a_*\quad\Longrightarrow\quad W\le x'\le W^2.
\]

\end{lemma}

\begin{proof}
\proofstep{Use the same cohomology interval after rescaling.}

Fix the common field factor \(c\ge1\) depending on \(\psi\), and
put \(a_*=c^{1/(2+8\varepsilon)}\). The cohomology window
\([W,W^2]\) will be chosen for this pair, with
\(X^{1/J}\le W\le\sqrt X\) and \(Y_0=X^{1/J}\).
On \([a_*W,(a_*W)^{3/2}]\) write \(x=a_*x'\).
If \(W\ge a_*^3\), then the upper rescaled endpoint is
\((a_*W)^{3/2}\le W^2\).

Since \(a_*\ge1\),
\[
 W\le x'=x/a_*\le x\le W^2,
 \qquad Y_0\le x\le X.
\]
Multiplication by \(c=a_*^{2+8\varepsilon}\) carries the
dyadic frame interval at \(x'\) to the frame interval at
\(x\). Moreover, if \(W\ge a_*^{400}\), then
\(x^{400}=a_*^{400}(x')^{400}\le(x')^{401}\).
In particular every integer \(1\le j<\lfloor x^{400}\rfloor\)
lies in the strict dyadic range \(1\le j<\lfloor(x')^{401}\rfloor\).
The error bound transfers because
\(\exp((x')^\varepsilon/2)\le\exp(x^\varepsilon/2)\).
The same prime pair is used before and after rescaling.
\end{proof}

\section{Local Weil--Deligne parameters and their Galois realizations}
\label{sec:local-global}

The lift construction specifies a finite type on the automorphic side. Later we will
need to recognize that same local datum in the Galois representation of an ordinary
constituent. We match the local cuspidal
and principal series data to Weil--Deligne parameters and
specify the global attachment and its hypotheses. The
ordinary-level arguments below use these compatibilities
at both the split ramified places and the coefficient prime. Proposition~\ref{prop:prescribed-local-cuspidal-types} gives the local parameter on the specified cuspidal types. Definition~\ref{hyp:local-global-ranges} specifies the places and twists at which a global representation must recover that parameter.

Let \(v\) be a finite place of a CM field \(E\).
The local assignment constructed here associates a Weil--Deligne parameter to each irreducible smooth representation of \(\mathrm{GL}_n(E_v)\). In rank one it is given by Artin reciprocity.
\begin{proposition}[Local parameters on the prescribed cuspidal types]\label{prop:prescribed-local-cuspidal-types}
At a finite place \(v\) the local assignment sends a smooth irreducible representation of \(\mathrm{GL}_n(E_v)\) to an \(n\)-dimensional Weil--Deligne parameter. On characters of \(E_v^\times\), use Artin reciprocity.
At \(E_v=\Q_p\), use induction from the degree-\(n\) unramified extension on each depth-zero cuspidal, with determinant fixed by its central character. For the regular residue character \(\theta\), the inertia representation is 
\[
 \bigoplus_{j=0}^{n-1}\theta^{p^j}.
\]
 Extend by the Artin twist to its smooth
twists, independently of the presentation. On the remaining
cuspidals use irreducible induced parameters with nonscalar
wild inertia, constructed using the first principal unit group. The full assignment extends by
the Langlands segment rule.
\end{proposition}

\begin{proof}
\proofstep{The rank-one and depth-zero normalizations.}

At a rational split completion \(E_v=\Q_p\) the other specified
cuspidals are twists of depth-zero cuspidals. The latter
are compact inductions from the finite cuspidal types of
\(\mathrm{GL}_n(\mathbf F_p)\), extended to the group generated
by the maximal compact and the center. Their finite types are
indexed by Frobenius orbits of regular characters \(\theta\) of
\(\mathbf F_{p^n}^{\times}\). The parameter is induced from the
Weil group of the degree-\(n\) unramified extension. On inertia
it is the direct sum of the characters \(\theta^{p^j}\) for
\(0\le j<n\).

Write \(F\) for the chosen Frobenius lift and \(\gamma\)
for the value of the inducing character at \(F^n\).
In a cyclic basis \(F e_j=e_{j+1}\) for \(j<n-1\)
and \(F e_{n-1}=\gamma e_0\). Thus
\[
        \det(F)=(-1)^{n-1}\gamma .
\]
For central value \(b\) at the uniformizer in the reciprocity
convention, choose \(\gamma=(-1)^{n-1}b\). This is the
depth-zero normalization of the local Langlands correspondence.
\hfill\mbox{\cite{adrian-roe}}

\proofstep{Independence of twisting presentations and segments.}

Twisting by a smooth character of \(E_v^{\times}\) twists the
parameter by its Artin image. Suppose two depth-zero cuspidals differ by a twist \(\alpha\circ\det\). Their usual Weil parameters are tame. For wild inertia \(P_v\), compatibility of the usual correspondence with twists gives
\[
 I_n=\operatorname{rec}(\pi\otimes(\alpha\circ\det))(g)
   =\operatorname{rec}_1(\alpha)(g)I_n\qquad(g\in P_v).
\]
Thus the relative character is trivial on \(1+p\mathbb Z_p\) by local reciprocity. \cite{henniart-local,adrian-roe} The corresponding
finite cuspidal types differ by the determinant character
modulo \(p\).

On the elliptic torus of the finite group, multiplication
by \(u\in\mathbf F_{p^n}^\times\) has determinant equal to
its field norm. Consequently the finite twist on the torus is
\[
       (\alpha\circ\det)(u)
          =\alpha\bigl(N_{\mathbf F_{p^n}/\mathbf F_p}(u)\bigr).
\]
Functoriality of local reciprocity gives the same pullback
on the degree-\(n\) Weil group. Thus the inertial orbits
agree in the two presentations. The central determinant gives
the same last Frobenius coefficient in the cyclic basis, which
identifies the two induced parameters.
For a general irreducible based on these cuspidals, use the unique multisegment of its Langlands quotient. For segment lengths \(s_i\) and cuspidal parameters \(W_i\) the contribution is
\[
 \bigoplus_i\bigl(W_i\otimes\operatorname{Sp}(s_i)\bigr)\otimes|\cdot|^{u_i},
\]
where \(u_i\) is the central exponent of the segment in the chosen normalization. The multisegment classification ensures independence of their order. \cite[\S9]{zelevinsky} This preserves the
dimension and determines the monodromy lengths.

\proofstep{Auxiliary parameters on the remaining cuspidals.}

For the remaining cuspidals of dimension
\(n\ge2\), choose auxiliary parameters as follows. Over the unramified extension
of degree \(n\), the first principal unit quotient gives
the additive group of its residue field. Write \(k_v\) for the residue field of \(E_v\),
and \(k_n/k_v\) for the degree-\(n\) extension. Choose \(b\in k_n\)
of degree \(n\) and use the additive character
\[
    x\longmapsto
       \exp(2\pi i\,\operatorname{Tr}_{k_n/\mathbf F_p}(bx)/p)
\]
on the first principal unit quotient. Identify \(\mathbf F_p\)
with residues in the exponent.

To extend this character to the multiplicative group, write \(F_n/E_v\) for the unramified extension, \(\varpi\) for a uniformizer in \(E_v\), and \(U_n^i=1+\varpi^i\mathcal O_{F_n}\). Reduction and Teichm\"uller lifts give
\[
 F_n^\times/U_n^2
   =\varpi^{\mathbb Z}\times [k_n^\times]\times (U_n^1/U_n^2),
 \qquad 1+\varpi x\longmapsto\bar x\quad\text{on }U_n^1/U_n^2.
\]

Extend trivially across the first two factors. Local reciprocity then gives the inducing character of the Weil group. Write \(q_v=|k_v|\) and \(\lambda_c(x)=\exp(2\pi i\operatorname{Tr}(cx)/p)\).
Nondegeneracy of the trace pairing and the degree of \(b\) give
\[
 \lambda_{b^{q_v^i}}=\lambda_{b^{q_v^j}}
 \quad\Longleftrightarrow\quad b^{q_v^i}=b^{q_v^j}
 \quad\Longleftrightarrow\quad i\equiv j\pmod n .
\]
The inducing character thus gives \(n\) distinct inertia lines. Every invariant subspace for inertia
is a sum of these lines, and Frobenius permutes them transitively,
so the induced parameter is irreducible. By local reciprocity their
restriction to wild inertia has a nonscalar element, since
the additive quotient is a \(p\)-group and the tame quotient
has prime-to-\(p\) finite quotients. 

For a cuspidal \(\tau\) of rank \(n\), write \(W_0\) for this induction. In the convention \(|\varpi|_v=q_v^{-1}\) there is a unique \(t\in\mathbb R\) such that
\[
 q_v^{-nt}|\det W_0(\operatorname{Art}_v(\varpi))|
       =|\omega_\tau(\varpi)|.
\]
Use \(W_0\otimes |\cdot|^t\). Twisting \(\tau\) by \(|\det|^u\) changes \(t\) to \(t+u\), leaving its wild inertia action unchanged.

It remains to check that this auxiliary choice cannot obscure the parameters on the specified class. On a character, or on a twist of a depth-zero cuspidal, wild inertia acts by scalars on the irreducible Weil factor. For the auxiliary induction it acts by at least two distinct characters. The uniqueness of the decomposition of a Frobenius-semisimple parameter into indecomposable summands therefore separates the two cases:
\[
 W\otimes\operatorname{Sp}(s)\simeq W'\otimes\operatorname{Sp}(s')
       \ \Longrightarrow\ s=s',\quad W\simeq W'.
\]
Indeed \(s\) is the common Jordan block length and the action on the kernel of monodromy recovers \(W\) up to the prescribed norm twist. When every constituent belongs to the specified class, the construction coincides with the ordinary local correspondence by the depth-zero formula and compatibility with twists and segments. Injectivity there follows from the local Langlands bijection. This proves the separation property used below. \cite{henniart-local,adrian-roe,zelevinsky}

\end{proof}

\subsection{Agreement for every admissible map}

It remains to see why the later local calculations do not depend on a choice among
parameter maps satisfying the comparison properties. We call a local parameter map \(R\) admissible when it preserves
dimension, satisfies the rank-one Artin formula, assigns
irreducible parameters to cuspidals, and extends by the segment
rule above. It respects real norm twists and gives the absolute
value of the central character as the absolute determinant
under reciprocity.

On representations whose cuspidal support consists of
characters and, at rational completions, twists of depth-zero
cuspidals, it also respects smooth twists and contragredients and gives
\[
        \det R(\sigma)\circ\operatorname{Art}_v
            =\omega_\sigma ,
\]
where \(\omega_\sigma\) is the central character in the fixed
reciprocity convention. It gives tame parameters on this class
when there are first congruence invariants. If one representation belongs to this class,
equality of its parameter with that of any irreducible
implies equality of the representations. At a rational completion, if the restriction to
\(\mathrm{GL}_n(\mathbb Z_p)\) contains the inflation of the
regular finite cuspidal type, the parameter satisfies the
zero-monodromy and inertia-polynomial conditions of Proposition~\ref{prop:prescribed-local-cuspidal-types}.

For the inertia formula use reciprocity on the unramified extension that realizes the residue character. The determinant identity, the separation of cuspidal parameters by wild inertia and the segment construction in Proposition~\ref{prop:prescribed-local-cuspidal-types} exhibit such a map. On the specified class contragredients and smooth twists agree with those of the ordinary correspondence. \cite{henniart-local,adrian-roe}

\begin{proposition}[Agreement in the two local comparison ranges]
Let \(R\) be an admissible map. If \(R(\sigma)\) has zero monodromy
and is a sum of lines, the rank-one Artin formula and the segment
rule determine it. At a rational split completion, if \(R(\sigma)\)
is tame with zero monodromy, the cuspidal support of \(\sigma\)
consists of characters and depth-zero cuspidals. In each of these two ranges, writing \(\operatorname{rec}_{\rm norm}\) for the usual normalized correspondence, 
\[
 R(\sigma)\simeq\operatorname{rec}_{\rm norm}(\sigma).
\]

\end{proposition}

\begin{proof}
\proofstep{Singleton segments and toral parameters.}

The global comparison is required for every admissible map.
Let \(R\) be one such map. The special parameter of a segment with cuspidal
parameter \(W\) and length \(s\) has \(\dim W\) monodromy
blocks of size \(s\). Thus in the toral comparison
\[
 N_{R(\sigma)}=0\ \Longrightarrow\ s=1\text{ on each segment},
 \qquad R(\sigma)\text{ a sum of lines}\ \Longrightarrow\ \dim W=1.
\]
The dimension law and rank-one Artin formula then determine the parameter.

\proofstep{Count tame cuspidal parameters of fixed determinant.}

At a rational split completion, suppose instead that
\(R(\sigma)\) is tame with zero monodromy. We first show
that an irreducible rank-\(r\) cuspidal \(\tau\) with tame
\(R(\tau)\) has depth zero. Fix a lift of Frobenius and
write \(b\ne0\) for its determinant on \(R(\tau)\).
In a tame irreducible Weil representation the inertia
characters form one Frobenius orbit. Each occurs once:
on one inertia eigenspace choose an eigenvector for the
Frobenius power that stabilizes it, and take the span of its
Frobenius translates. Irreducibility makes this span the
entire space. Thus the orbit has length \(r\) and is
labelled by a regular character of
\(\mathbf F_{p^r}^{\times}\), up to Frobenius.

In a cyclic basis the determinant calculation above gives
\[
      \det(XI-F)=X^r-\gamma,\qquad \gamma=(-1)^{r-1}b .
\] For each regular orbit, the
diagonal inertia action together with this cyclic Frobenius
matrix defines an irreducible tame parameter. Two such
parameters are isomorphic exactly when the inertia orbits
coincide, since rotating a cyclic basis keeps the same last
coefficient. This gives the claimed bijection at fixed \(b\).

\proofstep{Use admissible injectivity on the specified class.}

For each label there is a distinct depth-zero cuspidal
type, and compact induction with the indicated central
value gives a depth-zero cuspidal with central character
of value \(b\) at the image of the chosen Frobenius.
Let \(\mathcal D_b\) be the resulting depth-zero cuspidals
with determinant value \(b\), and \(\mathcal W_b\) the tame
irreducible parameters with that value. Their two constructions
use the same finite set of regular Frobenius orbits. The depth,
dimension and central character laws give
\[
          R:\mathcal D_b\lhook\joinrel\longrightarrow\mathcal W_b,
          \qquad |\mathcal D_b|=|\mathcal W_b|.
\]
The arrow is injective by the admissible injectivity law.
It is therefore onto.

Since \(R(\tau)\in\mathcal W_b\),
the same law, now with the other representation in
\(\mathcal D_b\), gives \(\tau\in\mathcal D_b\). Applying
this to each singleton segment puts the entire support of
\(\sigma\) in the specified class. On each of these cuspidals
the finite type law gives the inertia orbit and the central
character law fixes the Frobenius determinant. Hence the
parameter of \(\sigma\) agrees with the usual normalized
correspondence. This argument applies in particular when the
tame inertia polynomial is squarefree, as required by the
global row. The spherical condition above \(\ell\) is a
condition on the smooth local representation itself.
\hfill\mbox{\cite{adrian-roe}}\label{cor:cohomology-local-global-1}

\end{proof}

\begin{construction}
\begin{tabular}{@{}ll@{}}\toprule
Restriction on \(R(\sigma)\) & Consequence for cuspidal support\\\midrule
\(N=0\) & every segment has length one\\
\(N=0\), sum of lines & every cuspidal has rank one\\
\(N=0\), tame, at a rational split place &
characters and depth-zero cuspidals\\\bottomrule
\end{tabular}
\caption{The two local comparison ranges. The segment rule detects monodromy block sizes. For a tame irreducible factor the Frobenius orbit of its inertia character identifies its depth-zero type.}\label{fig:two-local-comparison-ranges}
\end{construction}

\begin{paramdefinition}[Local comparisons required of a single representation]\label{hyp:local-global-ranges}
Let \((\pi,\chi)\) be a regular algebraic, polarized cuspidal
pair of degree \(n\) over \(E\), and fix a prime \(\ell\) and
an identification \(\iota:\overline{\Q}_\ell\simeq\C\).
For a continuous semisimple representation
\(r\) of dimension \(n\), define the following comparison with the normalized local parameters \(W_v\) of \(\pi_v\). Away from \(\ell\) require equality of dimensions, and in
the zero-monodromy toral case compare the full parameter after
Frobenius semisimplification. The second comparison concerns
squarefree tame inertia:
\[
  E_v=\Q_p,\qquad p>16n^2\ \text{if }n>1,
  \qquad
  \det(1-X\operatorname{WD}(r|_{G_{E_v}})(i))
     =\det(1-XW_v(i))\quad(i\in I_v),
\]
with zero monodromy. Above \(\ell\), a spherical local
representation requires crystallinity and the prescribed
Hodge--Tate weights. The same toral and tame cases above \(\ell\)
require potential semistability, the Hodge--Tate weights
and the indicated Weil--Deligne comparisons. The comparison holds for \(r\) if these conditions hold at every place in the indicated ranges.
\end{paramdefinition}

The agreement just proved transfers the comparisons
for the usual correspondence to every admissible map at
the places where these comparisons are requested.

\begin{proposition}[Simultaneous Galois attachment]\label{prop:simultaneous-local-global-attachment}
For any pair and coefficient prime in
Definition~\ref{hyp:local-global-ranges}, there exists one continuous semisimple representation 
\[
 r:G_E\longrightarrow\mathrm{GL}_n(\overline{\mathbb Q}_\ell)
\]
 with the prescribed Hodge--Tate and crystallinity properties. For every admissible local map, this same representation satisfies the stated Weil--Deligne comparison at every place in its range.
\end{proposition}

\begin{proof}
\proofstep{The usual correspondence and the algebraic weights.}

Barnet-Lamb, Gee, Geraghty and Taylor give the representation
with the prescribed Hodge--Tate weights and compare its
Frobenius-semisimplified parameter away from
\(\ell\) with
\[
             \operatorname{rec}(\pi_v\otimes
                              |\det|_v^{(1-n)/2}).
\]
For a highest algebraic weight \((a_{\tau,1},\ldots,a_{\tau,n})\)
the entries of the infinitesimal character of the dual
algebraic representation, and the resulting Hodge--Tate
jumps, are
\[
 z_{\tau,j}=-a_{\tau,j}-\frac{n+1-2j}{2},
 \qquad
 \frac{n-1}{2}-z_{\tau,j}=a_{\tau,j}+n-j .
\] Purity at the two conjugate embeddings
makes their multisets complementary with respect to the
purity weight. Their Theorem~2.1.1 gives de Rhamness and
crystallinity at spherical places over \(\ell\).
\hfill\mbox{\cite{blggt}}

\proofstep{Potential semistability and the monodromy partition.}

At other places over \(\ell\) the de Rham representation is
potentially semistable by the \(p\)-adic monodromy theorem.
Theorem~1.1 of the local-global compatibility theorem gives
the semisimple Weil action and bounds the monodromy
partition by that of the normalized local parameter in
each inertial class. In a nonempty inertial class, let \(\lambda_1\) be the largest
actual Jordan block and \(\mu_1\) the largest block of the bounding
parameter. The first dominance inequality gives
\[
    \lambda_1\le\mu_1=1
      \qquad\text{if the bounding monodromy vanishes}.
\]
Thus every block has size one, and monodromy is zero. Therefore in the toral and
tame situations just described the comparison gives the
required zero monodromy even at a regular weight with no
additional regularity condition. The rank-one case also
follows from the Galois character associated to the
algebraic Hecke character by global reciprocity.
\hfill\mbox{\cite{blggt-lg}}

\end{proof}

This identifies the local Galois data that enter the cohomology calculation.
Uniformity as the level, weights and rank vary will come from the length estimates
there.

\section{A common interval for bounds on strict cohomology}\label{sec:cohomology}

The second divisibility estimate requires a weighted sum of strict cohomology lengths over the original auxiliary field. The analytic prime pair has already been fixed in Appendix~\ref{sec:hosts}. We may shorten its interval to control the twists simultaneously while retaining that same pair.
Cyclotomic towers control the lengths outside a bounded
number of exceptional twists, and a height estimate controls
those twists on the smaller interval. Let \(E\) be the auxiliary abelian CM field defined by \eqref{eq:cohom-host-field}. At \(2\) the tower gives a bound over \(E(i)\). Lemma~\ref{lem:dyadic-host} controls restriction from \(E\) to \(E(i)\), yielding a bound on the strict lengths over \(E\) before the twists are summed in Section~\ref{sec:family-window}.

Let \(x\) denote the original cohomology scale \(z\) of
\eqref{eq:cohomwin}. For \(S\) the places above \(\ell\) and the
archimedean places, put \(A_j=\Q_\ell/\Z_\ell(-j)\). The strict
length is
\[
 V(E,\ell,j)=\operatorname{length}_{\Z_\ell}
   \ker\!\left(H^1(G_{E,S},A_j)\longrightarrow
                    \prod_{v\mid\ell}H^1(E_v,A_j)\right).
\]
The fields here are totally complex. Poitou--Tate duality identifies
the length with the length of the corresponding degree-two
localization kernel for \(\Z_\ell(j+1)\).
At the original scale \(x\), take the dyadic index \(e(x)\) so that
\begin{equation}\label{eq:cohom-host-dyadic-level}
 x^{2+8\varepsilon}\le 2^{e(x)}\le2x^{2+8\varepsilon}.
\end{equation}
The auxiliary field \(E\) is cut out by the automorphisms
\(\zeta\mapsto\zeta^a\) with
\begin{equation}\label{eq:cohom-host-field}
 a\equiv\pm1\bmod 2^{e(x)+2},\qquad
 \psi(a)=1,\qquad (a/q)=(a/d)=1.
\end{equation}
Its conductor divides \(\operatorname{lcm}(2^{e(x)+2},N,q,d)\).
It is abelian and CM. In fact \(a=-1\) fails the condition at
\(q\equiv3\bmod8\), so complex conjugation acts nontrivially.

For a fixed interval exponent \(J\) and arbitrarily large
endpoints \(X\), write \([W,W^2]\) for the
subinterval of \([X^{1/J},X]\) on which the bound holds for
every finite set of primes. The proof below establishes
finiteness and the bound in \eqref{eq:cohomwin}. The cutoff
exponent \(401\) leaves room for the constant rescaling
in Section~\ref{sec:family-window}.

\subsection{Localization kernels in cyclotomic towers}

Variation of the cyclotomic layer changes the specialization of one Galois module. Describing the strict kernels as specializations locates the twists for which the usual length bound deteriorates.

Take a base \(L_0\) containing \(\mu_\ell\) for odd \(\ell\) or
\(\mu_4\) for \(\ell=2\). Let \(L_\infty/L_0\) be its cyclotomic
\(\Z_\ell\)-extension and \(L_m\) the layer of degree \(\ell^m\).
Let \(M\) be the abelian pro-\(\ell\) Galois group over \(L_\infty\)
unramified away from \(\ell\) with splitting at all places above \(\ell\).
The coefficient module \(\Q_\ell/\Z_\ell(-j)\) has trivial action
over \(L_\infty\). Denote by \(u\) the cyclotomic value of a topological
generator and by \(A\) its action by conjugation on \(M\).

\begin{lemma}[Lengths away from finitely many cyclotomic twists]\label{lem:tower-length-centres}
Suppose \(M\) is generated by \(R\ge1\) elements over \(\Z_\ell\).
Put \(j_0=\lfloor\log_\ell R\rfloor+3\) and \(c=v_\ell(u-1)\).
There are \(s_1,\ldots,s_R\in\Z_\ell\) such that for every \(m\ge0\),
integer \(j\ge1\), and \(B\ge0\) with
\(\lvert j-s_i\rvert_\ell\ge\ell^{-B}\) for all \(i\),
the strict localization kernel has finite length and
\begin{equation}\label{eq:tower-bound}
 V(L_m,\ell,j)\le R\big(\max(m,j_0)+c+B\big).
\end{equation}
The centres can be chosen once for the whole tower.
\end{lemma}

\begin{proof}
\proofstep{The coinvariant module.}
Multiplication by
\(u^{-j\ell^m}-1\ne0\) is surjective on \(\Q_\ell/\Z_\ell\).
By Lemma~\ref{lem:tower-coinvariants}, the procyclic cohomology computation and
inflation--restriction identify the global degree-one group
with the invariant homomorphisms upstairs. The analogous computation
at the places over \(\ell\) identifies the local splitting condition.
For a continuous homomorphism \(\phi\) upstairs, invariance
under the generator of \(\operatorname{Gal}(L_\infty/L_m)\)
means
\[
 \phi(A^{\ell^m}z)=u^{-j\ell^m}\phi(z)\qquad (z\in M).
\]
Equivalently \(\phi\) kills \(\big((u^j A)^{\ell^m}-1\big)M\).
The strict group is the continuous dual of that quotient.
Taking the Pontryagin dual recovers
\[
   M/\bigl((u^j A)^{\ell^m}-1\bigr)M.
\]
The inverse of the coefficient character occurs when forming
the coinvariant operator.
\proofstep{Growth above a reference layer.}
By continuity the reduction of \(A\) modulo \(\ell\)
is unipotent. Put \(j_0=\lfloor\log_\ell R\rfloor+3\) and
\(g=u^jA\). If \(m\ge j_0\), then \(D=g^{\ell^m}-1\)
maps \(M\) to \(\ell^2M\). The reduction of \(M\) has dimension at most \(R\), so
\[
   g^{\ell^k}-1\equiv(g-1)^{\ell^k}=0\pmod\ell
   \qquad(\ell^k>R).
\]
Another \(\ell\)-th power gives \(g^{\ell^{k+1}}-1\equiv0\pmod{\ell^2}\)
by the binomial formula. This applies for \(m\ge j_0\). Set \(I=DM\).

Since \(D\)
is \(\Z_\ell\)-linear, its powers satisfy
\(D^r M\subset\ell^{2(r-1)}I\) for \(r\ge1\).
Thus \(D'=(1+D)^\ell-1\) takes values in \(\ell I\), and
the terms with \(r\ge2\) take values in \(\ell^2 I\).
Reduction modulo \(\ell^2 I\) gives
\[
      D'M+\ell^2 I=\ell DM+\ell^2I=\ell I.
\]
Apply Nakayama to the finitely generated \(\Z_\ell\)-module \(\ell I/D'M\)
to get \(D'M=\ell I\). The exact sequence
\[
 0\longrightarrow I/\ell I\longrightarrow M/\ell I
       \longrightarrow M/I\longrightarrow 0
\]
then shows that the length of a finite cokernel increases by
at most \(R\), since \(I\), an image of \(M\), is generated
by at most \(R\) elements. Layers below \(j_0\) have cokernels
which are quotients of the reference cokernel.

\proofstep{Distance from a centre at the reference layer.}

Growth along the tower is now bounded in terms of a fixed reference layer. The exceptional twists arise from the eigenvalues of its cyclotomic action, so their distance to the twist under consideration enters here.

Lift \(A^{\ell^{j_0}}\) to a matrix \(A_*\) on a free module on
the \(R\) generators, and put \(w=u^{\ell^{j_0}}\).
The reference cokernel is obtained by the surjection
\[
 \operatorname{coker}(w^j A_*-1)\twoheadrightarrow \operatorname{coker}(w^j A^{\ell^{j_0}}-1\mid M).
\]
 Let \(\alpha_i\) be the eigenvalues of \(A_*\),
counted with multiplicity. Their absolute values are at most one.
If \(\abs{\alpha_i}_\ell=1\), choose \(s_i\in\Z_\ell\) minimizing
\(\abs{w^s\alpha_i-1}_\ell\), which is continuous on the compact
set \(\Z_\ell\). The ultrametric inequality and minimality give
\[
   \abs{w^j\alpha_i-1}_\ell\ge
      \abs{w^j-w^{s_i}}_\ell
      =\ell^{-c-j_0}\abs{j-s_i}_\ell ,
\]
where \(c=v_\ell(u-1)\). The identity follows from the valuation of
integer powers and then by continuity on \(\Z_\ell\).

For \(\abs{\alpha_i}_\ell<1\) the factor equals one and one can
take any \(s_i\in\Z_\ell\).
Smith normal form bounds the length of the reference cokernel by
the valuation of the determinant of \(w^j A_*-1\). Under the
separation hypothesis each of its \(R\) factors has valuation at most
\(c+j_0+B\). For \(m\le j_0\) take the quotient of this cokernel.
For larger \(m\) add at most \(R\) for each further layer. This proves
\eqref{eq:tower-bound}.
\end{proof}

\subsection{Uniform rank at the prime two}

At the prime two the original field need not contain a fourth root of unity. A uniform bound on generators upstairs is useful only after accounting for the kernel of restriction from the original field.

For the prime \(2\) start with the fixed abelian CM field obtained
from the character field by adjoining \(i\).
\begin{lemma}[A uniform number of dyadic generators]\label{lem:dyadic-generators}
Let \(d_0=[K_\psi:\Q]\), and let \(g_0\) bound the \(2\)-ranks of the
class groups in the fixed cyclotomic tower of \(K_\psi(i)\).
For the prime pair in \eqref{eq:cohom-host-field}, after adjoining
\(\sqrt{-q}\) and \(\sqrt{-d}\), the abelian
pro-\(2\) group with the unramified and local splitting conditions
in Lemma~\ref{lem:tower-length-centres} has generator rank \(R\) bounded by 
\[
 R\le4(g_0+6d_0).
\]
 This bound is uniform in the prime pair.
\end{lemma}

\begin{proof}
Ferrero--Washington's theorem gives zero Iwasawa \(\mu\)-invariant
for the cyclotomic \(\Z_2\)-extension of this abelian field
over \(\Q\). The class groups of the finite layers consequently
have bounded \(2\)-rank. The same bound \(g_0\) holds for
unramified quadratic characters of the infinite layer.
For a finite list of such characters, descend to a finite layer
after taking a sufficiently small open cyclotomic subgroup.
One can do this by first adjoining the fields of the characters
over the infinite layer, descending each finite extension to a
finite layer and then enlarging the layer until local inertia
maps trivially at the finitely many ramified places.
If \(V\) is the span of a finite independent list in the infinite layer, the restriction image from that common finite layer contains \(V\). Consequently
\[
 \dim_{\mathbf F_2} V\le
 \dim_{\mathbf F_2}\operatorname{Hom}(\operatorname{Cl}(K'),\mathbf F_2)
 \le g_0 ,
\]
where \(K'\) denotes the finite layer. Taking all finite lists proves the claimed bound.
\hfill\mbox{\cite{ferrero-washington}}

If the fixed character field has degree \(d_0\), the numbers of
primes above \(q,d\) in the fixed tower are at most \(2d_0,4d_0\).
Their decomposition subgroups in the full dyadic tower have
indices
\[
    [\Z_2^\times:\overline{\langle q\rangle}]
      =2^{v_2(q^2-1)-2}=2,\qquad
    [\Z_2^\times:\overline{\langle d\rangle}]
      =2^{v_2(d^2-1)-2}=4 .
\]
Here each projection to the sign quotient is onto, and the square
generates a subgroup of \(1+4\Z_2\) with the displayed index.
Allowing ramification at any one of these primes costs at most
one quadratic character parameter by tame inertia.
The two auxiliary quadratic extensions together have degree at
most four over the fixed tower. Apply the Schreier generator
bound to the subgroup of index at most four in the maximal
pro-\(2\) Galois group with this allowed ramification. Passing to the abelian quotient of this subgroup gives
\[
  R\le 4(g_0+2d_0+4d_0)=4(g_0+6d_0)
\]
in \eqref{eq:tower-bound}.
The splitting condition only decreases this bound. The constant
is thus fixed before the auxiliary prime pair varies.
\end{proof}

For the finite set of odd primes treated by the tower argument,
one takes the rank bound over the fixed abelian base for that
prime.

\begin{lemma}[Restriction at the prime \(2\)]\label{lem:dyadic-host}
Let \(K_\psi\) be the fixed field of \(\ker\psi\) in the
cyclotomic field of its conductor. Let \(L_0\) be the field
\(K_\psi(i,\sqrt{-q},\sqrt{-d})\). For \(E\) defined by
\eqref{eq:cohom-host-field} at dyadic level
\eqref{eq:cohom-host-dyadic-level}, the extension \(E(i)\)
is a layer \(L_m\) of degree \(2^m\) in the cyclotomic
\(\Z_2\)-extension of \(L_0\), with \(m\le e(x)+2\).
For every \(j\ge1\), the localization kernel over \(E(i)\)
controls the kernel over \(E\) by
\begin{equation}\label{eq:dyadic-host}
                  V(E,2,j)\le 1+V(L_m,2,j).
\end{equation}
\end{lemma}
\begin{proof}
\proofstep{Identifying the tower layer.}
Put \(n=e(x)+2\), \(F_n=\Q(\mu_{2^n})\), and write \(F_n^+\)
for its maximal real subfield. In the cyclotomic field containing
the four conductors, the stabilizer of a compositum is the
intersection of the stabilizers of its factors. The signs at \(q,d\)
cut out \(\Q(\sqrt{-q})\) and \(\Q(\sqrt{-d})\), since both primes
are \(3\bmod 4\). The condition \(a=\pm1\bmod 2^n\) cuts out
\(F_n^+\). Thus
\[
 E=K_\psi(\sqrt{-q},\sqrt{-d})F_n^+.
\]
For \(n\ge2\), complex conjugation acts on \(i\in F_n\) by
\(-1\), so \(F_n^+(i)=F_n\).

It follows that \(E(i)=L_0 F_n\).
Adjoining all dyadic roots to \(L_0\supset\mu_4\) gives a
\(\Z_2\)-extension. The subgroup corresponding to \(L_0F_n\)
is therefore the unique subgroup of its index \(2^m\), with
\[
        2^m=[F_n:F_n\cap L_0]\le[F_n:\Q]\le 2^n.
\]
The bound \(m=O(1+\log x)\) is uniform in the prime pairs.
The extension over \(E\) has degree one or two. For the trivial
target, an exponent congruent to \(-1\bmod 2^n\) and
\(1\bmod qd\) fixes \(E\) and moves \(i\), so the quadratic
case already occurs for the zeta application.

\proofstep{Restricting strict classes.}
Set \(A=\Q_2/\Z_2(-j)\). Let \(H^1_{\rm ur(2)}(E,A)\) denote classes unramified
outside the dyadic and infinite places, and use analogous
notation for \(L_m\). Upstairs inertia maps into downstairs
inertia. With \(w\mid v\mid2\), the two localization rows
fit in the commutative diagram
\[
 \begin{array}{ccc}
 H^1_{\rm ur(2)}(E,A)&\xrightarrow{\mathrm{loc}}&
      \displaystyle\prod_{v\mid2} H^1(E_v,A)\\
 \downarrow{\scriptstyle\mathrm{res}}&&\downarrow{\scriptstyle
                    (\mathrm{res}_{w/v})_{w\mid v}}\\
 H^1_{\rm ur(2)}(L_m,A)&\xrightarrow{\mathrm{loc}}&
      \displaystyle\prod_{w\mid2} H^1((L_m)_w,A).
 \end{array}
\]
The left arrow therefore restricts to the strict kernels. View these classes in absolute Galois cohomology
by inflation. For \(H=G_{L_m}\) and \(C=\operatorname{Gal}(L_m/E)\),
inflation--restriction gives the exact sequence
\[
 0\longrightarrow H^1(C,A^H)\xrightarrow{\mathrm{inf}} H^1(E,A)
             \xrightarrow{\mathrm{res}} H^1(L_m,A)^C .
\]
The strict restriction kernel is the intersection of the image
of inflation with the strict group over \(E\). For relative
degree one this kernel is zero.

\proofstep{Bounding the restriction kernel.}

Restriction has already placed the image in the strict kernel upstairs. The downstairs length can thus exceed the upstairs bound only by the length of the kernel of restriction.

In relative degree two put \(B=A^H\) and denote the
involution by \(\sigma\). A cocycle is determined by its value
on \(\sigma\), giving the quotient
\[
    \ker(1+\sigma:B\longrightarrow B)/(\sigma-1)B.
\]
It is killed by two, since \(2b=(1-\sigma)b\) on the numerator.
Every finitely generated subgroup of \(A\) is cyclic of
two-power order. The same property holds for any subquotient,
by lifting a finite set of generators to the subgroup before
taking the quotient. A group with this property killed by two
has at most two elements. The restriction kernel on strict
classes consequently has \(\Z_2\)-length at most one. The image of restriction on strict classes
lies in the strict group upstairs. Additivity of length gives
\eqref{eq:dyadic-host} and finiteness downstairs whenever the
upstairs length is finite. \label{cor:cohomology-cohomology-1}
\end{proof}

\begin{construction}
\begin{tabular}{@{}lll@{}}\toprule
Restriction from \(E\) to \(E(i)=L_m\) & Kernel on strict classes &
Possible length loss\\\midrule
\(i\in E\) & zero & \(0\)\\
\([E(i):E]=2\) & subgroup of \(H^1(C_2,A^{G_{L_m}})\) & \(\le1\)\\\bottomrule
\end{tabular}
\caption{Dyadic descent for a single twist with coefficients \(A=\Q_2/\Z_2(-j)\). Restriction maps strict classes upstairs. Its image has length at most \(V(L_m,2,j)\). The only additional term downstairs comes from its kernel.}\label{fig:dyadic-strict-restriction}
\end{construction}

\subsection{Moving the scale past finitely many twists}
\label{subsec:cohomology-twist-threshold}

The tower estimate can be poor near a finite set of centres. The cutoff in the twist index grows with the real scale, so shortening the scale interval can isolate the exceptions that still require a separate estimate.

\begin{lemma}[Avoiding the exceptional twists on one interval]\label{lem:exceptional-twist-window}
For a finite list of towers as in Lemma~\ref{lem:tower-length-centres},
let \(R_*\) bound the total number of centres. Fix \(0<\eta<401\)
and put
\[
 J=4(2\cdot401/\eta)^{R_*}.
\]
For \(X>1\) set \(Z_*=\lfloor X^{401}\rfloor\). Call an integer
\(1\le j\le Z_*\) exceptional if for at least one centre \(s_i\)
of an \(\ell\)-adic tower it satisfies
\(\lvert j-s_i\rvert_\ell<\ell^{-\lfloor\log_\ell Z_*\rfloor}\).
There exists \(W\ge X^{1/J}\) with \(W^2\le X\) such that, for
every real \(x\in[W,W^2]\), every exceptional \(j<x^{401}\)
satisfies \(j\le x^\eta\). Only \(R_*\) and \(\eta\) enter \(J\).
\end{lemma}

\begin{proof}
Two integers exceptional for the same centre differ by a multiple of
\(\ell^{1+\lfloor\log_\ell Z_*\rfloor}>Z_*\), so they coincide.
There are thus at most \(R_*\) exceptional integers, regardless of
where the centres lie.

Put \(u=\log x\). For an exceptional \(j\), exclude the interval
\[
 [a_j,b_j)=\left[\frac{\log j}{401},\frac{\log j}{\eta}\right).
\]
The inequalities \(x^\eta<j<x^{401}\) imply
\(a_j<u<b_j\). The left endpoint is included in the excluded
interval and the right endpoint is allowed, since there
\(j=x^\eta\). Start with
\(v=(\log X)/J\). If \([v,2v]\) meets \([a_j,b_j)\), replace \(v\)
by \(b_j\). The new \(v\) is larger, so that interval can never
interfere again. Also \(a_j\le 2v\) before the replacement, hence
\(b_j\le(2\cdot401/\eta)v\). After at most \(R_*\) replacements,
\[
 (\log X)/J\le v\le (2\cdot401/\eta)^{R_*}(\log X)/J
                   =(\log X)/4 .
\]
Take \(W=e^v\). The forbidden intervals are open at their right
endpoints: equality \(j=x^\eta\) is permitted.
\end{proof}

For the application take \(0<\varepsilon<1/52\), let
\(C_K\ge1\) be the exponent in the height estimate of
Section~\ref{sec:cohom-height}, and put \(\eta=\varepsilon/(2C_K)\).
Lemma~\ref{lem:dyadic-generators} and the fixed bounds for the
finitely many odd-prime towers give \(R_*\) independently of \(q,d\).
Thus \(J\) is available when the prime pair is chosen. The position
of \(W\) may depend on that pair.

\begin{proposition}[Sum of the dyadic lengths over the original field]\label{prop:dyadic-window-total}
For every sufficiently large selected endpoint \(X\), on the
interval from Lemma~\ref{lem:exceptional-twist-window} one has,
for each real \(x\in[W,W^2]\) and its actual host field \(E\),
\[
  \sum_{1\le j<\lfloor x^{401}\rfloor} V(E,2,j)
                  \le \exp(x^{4\varepsilon/5}).
\]
All lengths in the sum are finite. The lower threshold is uniform in
the selected prime pair at fixed character and \(\varepsilon\).
\end{proposition}

\begin{proof}
\proofstep{Exceptional twists upstairs.}
The degree-two estimate at the reference layer is
\begin{equation}\label{eq:height}
 \log\abs{H^2(G_{L,S},\Z_\ell(j+1))}
      \le \exp\bigl((2j[L:\Q])^{C_K}\log(2\abs{D_L})\bigr)
\end{equation}
for totally complex bases containing \(\mu_\ell\) for odd
\(\ell\), or \(\mu_4\) for \(\ell=2\). Apply it with \(L=L_{j_0}\).
Its degree is uniformly bounded. In proving \eqref{eq:height},
Section~\ref{sec:cohom-height} uses an auxiliary field
\(L^\dagger/L\) whose largest root level \(c^\dagger\) satisfies
\(\ell^{c^\dagger}>2j+1\) and \(r_2(L^\dagger)>2j+1\).
Its relative degree is at most \(([L:\Q]+1)(4j+4)\).
This enlargement is already accounted for in \(C_K\).
The discriminant of \(L_{j_0}\)
has logarithm \(O(1+\log(qd))\) by the conductor--discriminant
formula, and hence \(O(x^{\varepsilon/4})\) on the selected window.

Since the degree is bounded and \(j\le x^\eta\), the new
exponent is
\[
   (2j[L_{j_0}:\Q])^{C_K}\log(2|D_{L_{j_0}}|)
       \ll x^{C_K\eta+\varepsilon/4}
       \ll x^{3\varepsilon/4}.
\]
The last inequalities use the fixed choice of \(\eta\) in the
window.
By Poitou--Tate the degree-two group bounds the strict length
at \(j_0\). Lemma~\ref{lem:tower-length-centres} bounds smaller
layers by a quotient and larger layers by successive growth:
\[
 V(L_m,\ell,j)\le V(L_{j_0},\ell,j)+R\max(0,m-j_0)
       \ll 1+\log x+\exp(Cx^{3\varepsilon/4}).
\]
\proofstep{Restriction and summation.}

The exceptional and nonexceptional twists have been controlled on one upstairs interval. Restriction contributes at most one additional unit of length for each twist before the summation on the original field.

For the other twists \eqref{eq:tower-bound} gives
\(O(1+\log x)\) since \(m=O(1+\log x)\) and
\(\log Z_*\le401J\log x\). Put \(t=\lfloor x^{401}\rfloor\).
Summing \eqref{eq:dyadic-host} over the \(t-1\) integers
\(1\le j<t\) gives
\[
 (\log 2)\sum_{1\le j<\lfloor x^{401}\rfloor} V(E,2,j)
 \le (\log 2)\sum_{1\le j<\lfloor x^{401}\rfloor} V(E(i),2,j)
      +(\log 2)(t-1) .
\]
The additive term is uniform in \(q,d\). The same estimates, after dividing by the fixed positive constant \(\log 2\), bound the logarithm of one plus the unweighted sum by \(O(1+\log x+x^{3\varepsilon/4})\). This is \(o(x^{4\varepsilon/5})\). For large \(x\) the unweighted sum is therefore at most \(\exp(x^{4\varepsilon/5})\), and the weighted contribution is at most \((\log 2)\exp(x^{4\varepsilon/5})\). The same absorption applies to each of the finitely many other towers.
\end{proof}

\input{height}

\subsection{The principal term at odd primes}

For odd residue characteristic the trivial character has to be separated before decomposing into the remaining components. In the indicated parity and congruence range its strict localization kernel over the rationals vanishes.

\begin{lemma}[The principal odd-prime contribution]\label{lem:principal-odd-strict-zero}
For a prime \(\ell\ge5\) and an even integer \(j>0\) with
\(\ell-1\mid2j\), put \(T=\Q_\ell/\Z_\ell(-j)\). For classes unramified away from \(\ell\), strict localization at \(\ell\) has trivial kernel: 
\[
 \ker\!\left(H^1(G_{\mathbb Q,\{\ell,\infty\}},T)\longrightarrow H^1(\mathbb Q_\ell,T)\right)=0.
\]

\end{lemma}

\begin{proof}
Let \(A=\Q_\ell/\Z_\ell(-j)\) and
let \(S_0\) be the subgroup of degree-one cohomology
unramified outside \(\ell\) with zero localization there.
Every class has \(\ell\)-power order because a continuous
cocycle has finite image in the discrete target. Thus
it suffices to show \(S_0[\ell]=0\). For any representative \(c\) of a class killed by \(\ell\),
write \(\ell c=dm\) and take \(m'\in A\) with \(\ell m'=m\).
The cocycle \(c'=c-dm'\) has the same class and
\[
          \ell c'=dm-d(\ell m')=0.
\]
\proofstep{Trivial residual character.}
If \(\ell-1\mid j\), its action on \(A[\ell]\) is
trivial. The cocycle is an abelian homomorphism
unramified outside \(\ell\), hence trivial on the
absolute Galois group of the cyclotomic \(\ell\)-power
tower by Kronecker--Weber. Inflation from the tower
quotient has zero degree-one cohomology on \(A\).
The finite part of the quotient has order prime to \(\ell\).
For a generator \(\gamma\) of the procyclic part the scalar
\(\gamma-1\) on \(A\) is nonzero and surjective. Consequently
\[
 H^1(\overline{\langle\gamma\rangle},A)
 = A/(\gamma-1)A=A/A=0 .
\]
Taking invariants under the finite prime-to-\(\ell\) quotient is exact. The original
class is zero.
\proofstep{Quadratic residual character.}
If \(\ell-1\nmid j\), the residual action is
quadratic. Since \(j\) is even, \(\ell\equiv1\pmod4\)
and it becomes trivial over \(K=\Q(\sqrt\ell)\).
On the decomposition group at \(\ell\), the adjusted
cocycle equals \(dm_\ell\), since its local class
is zero. Choose an inertia element \(s\) acting by \(-1\) on \(A[\ell]\).
Its action minus one is a unit on \(A\), hence
\[
       (s-1)\ell m_\ell=\ell c'(s)=0\quad\Longrightarrow\quad
                  \ell m_\ell=0.
\]
The restriction to \(G_K\) thus vanishes on inertia at \(\ell\),
since \(G_K\) acts trivially on \(A[\ell]\). It gives a
homomorphism from the ideal class group of \(K\)
to \(A[\ell]\). That class group has order less
than \(\ell\).

In fact Minkowski's bound gives an
ideal of norm at most \(\sqrt\ell/2\) in each class,
and the number of ideals of norm \(n\) is at most
the divisor count of \(n\), which is at most \(n\).
For \(B=\lfloor\sqrt\ell/2\rfloor\) this gives
\[
          h_K\le \sum_{n\le B}n=B(B+1)/2<\ell.
\] The homomorphism is therefore zero.
Inflation from \(\operatorname{Gal}(K/\Q)\)
has zero degree-one cohomology since this group
has order prime to \(\ell\). This proves the vanishing of the principal-character contribution for the indicated odd primes and parity.
\label{cor:cohomology-cohomology-2}
\end{proof}

\input{odd-primes}

\section{The coefficient data at the two ordinary levels}
\label{sec:count-two-level-data}
\hypertarget{source-ordinary-level-data}{}

The odd-prime count uses ordinary eigensystems at two levels. We specify their common residual data, the ordered characters above \(\ell\), and the tame and exceptional coefficient types defining each level. At each level these choices determine the
integral spaces of forms and the Hecke localization. The next two appendices construct the Galois representations of the ordinary constituents and realize the rank-four eigenpoint in these spaces. At a tame place the coefficient type is the integral compact representation of Proposition~\ref{prop:ordinary-zero-monodromy-type}. Lemma~\ref{lem:ordinary-residual-data} gives the distinct Frobenius and ordinary inertia prescriptions at the other marked places.

Use the CM field \(L\), the coefficient DVR \(O\) with fraction
field \(E\) and residue field \(k\), and the common twist of the
fitted tower. At each selected split place \(v\) fix one of its two lifts. The local factor used by the level is then specified by
\[
 L\otimes_{L^+}L^+_v=L_w\times L_{\bar w},\qquad w\mid v\text{ chosen}.
\] Write \(S_{\mathrm t}\)
for the selected places with a compact coefficient type away from
\(\ell\), and \(S_\ell\) for the selected places over \(\ell\).
There is a separate principal congruence condition at the
exceptional place of the level. The residual datum on the level
quotient consists of global labels \(\bar d_0,\ldots,\bar d_{n-1}\)
and ordered local labels \(\bar\psi_{w,0},\ldots,\bar\psi_{w,n-1}\)
at \(S_\ell\), with \(n=4\) in the even-order case. At either level the local coefficient type must realize these labels after restriction and the common twist, on that level's quotient.

\subsection{An integral type that detects zero monodromy}

The prescribed residual characters determine the inertial labels at a tame place. To impose those labels on forms one needs an integral type whose occurrence in characteristic zero also controls the monodromy operator.

\begin{proposition}[The tame integral coefficient type]
\label{prop:ordinary-zero-monodromy-type}
For each tame place \(w\) of either actual level the Teichmüller
lifts of its prescribed characters define a type \(\tau\) with
monodromy zero. Its compact type has an integral coefficient
model with the following detection property in characteristic zero:
write \((r,N)\) for the parameter of an irreducible smooth \(\pi_w\),
and \(\sigma_\tau\) for the compact type in characteristic zero. Then
\[
 \operatorname{Hom}_{K_w}(\sigma_\tau,\pi_w)\ne0
       \quad\Longleftrightarrow\quad N=0\ \text{ and }\
                                  (r|_{I_w})^{\rm ss}\simeq\tau .
\]
The criterion allows repeated characters. It holds for every
admissible local parameter map of Section~\ref{sec:local-global}
and is invariant under unramified normalization twists.
\end{proposition}

\begin{proof}
\proofstep{The inertial type and its rational realization.}
At \(w\in S_{\mathrm t}\), lift each \(\bar d_i\) by the
multiplicative Teichmüller section \(k^\times\to O^\times\) and
restrict to inertia at \(w\). Denote the resulting characters by
\(\eta_i\). Each \(\bar u\in k^\times\) has a unique lift \([\bar u]\)
in the prime-to-\(\ell\) roots of unity of \(O\), by Hensel's lemma.
Uniqueness gives \([\bar u\bar v]=[\bar u][\bar v]\). Thus the type is
\[
  \eta_i=[\bar d_i]|_{I_w},\qquad
  \tau=\bigoplus_i\eta_i,\qquad N=0 .
\]
Each \(\eta_i\) extends to the decomposition group, with values
in the finite subgroup of \(E^\times\) given by the lifts.

It has the least
nilpotent orbit for its semisimple inertia representation. Each
partition of a multiplicity space is a column of ones. There is
an irreducible smooth type \(\sigma_\tau\) on
\(K_w=\mathrm{GL}_n(\mathcal O_{L_w})\) with the following property.
For an irreducible smooth \(\pi_w\) the detection property is the equivalence
\[
 \operatorname{Hom}_{K_w}(\sigma_\tau,\pi_w)\ne0\quad\Longleftrightarrow\quad \operatorname{WD}(\pi_w)\text{ has inertial type }\tau
\]
in the algebraic normalization. This follows from the
Schneider--Zink type theorem \cite[Theorem 3.7(1)--(2)]{shotton-types}.
Its inequality of types compares nilpotent partitions with the
semisimple inertia fixed, so for this \(\tau\) the inequality is
equality. Both implications allow repeated \(\eta_i\) and
nongeneric \(\pi_w\).

The criterion is the same for every admissible local map in the
formal chain. A parameter with monodromy zero has only length-one
segments. Write \(V_\eta\) for the space on which inertia acts through
\(\eta\). If \(\Phi_w\) is a Frobenius element, then
\[
 r(i)r(\Phi_w)v
  =r(\Phi_w)r(\Phi_w^{-1}i\Phi_w)v
  =\eta(i)\,r(\Phi_w)v
       \qquad(i\in I_w,\ v\in V_\eta).
\]
The last equality uses the extension of \(\eta\) to the Weil
group. Diagonalize Frobenius on each \(V_\eta\). Its eigenlines
are Weil lines, so the cuspidal data have rank one. The segment
rule and rank-one reciprocity give the usual local Langlands
parameter. Conversely, a usual parameter of
this inertial type also comes from singleton rank-one segments,
for which every admissible map gives the same parameter. The
criterion can be expressed by requiring \(N=0\) and, for every
inertia element \(g\),
\[
   \det(X-r(g))=\prod_{i=0}^{n-1}(X-\eta_i(g)).
\]
Equality of traces at every element of the finite common inertia
quotient determines the semisimple representation there.
An unramified normalization twist preserves both tests.

We construct an \(O\)-model that detects this condition. Work over
\(\overline E\) via the coefficient identification with \(\C\),
and let \(H=K_w/\ker(\sigma_\tau)\). The kernel is open by
intersection of the stabilizers of a finite basis. The quotient is
finite by compactness. The criterion is invariant under every
\(\gamma\in\operatorname{Gal}(\overline E/E)\). In fact, a
representation satisfying it has a Langlands quotient presentation
by singleton rank-one segments. The possible change of square root under coefficient conjugation is measured by a sign:
\[
 \gamma((\#k_w)^{1/2})=\delta(\#k_w)^{1/2},\qquad \delta\in\{1,-1\}.
\]
Hence normalized induction is carried to the normalized induction of the conjugate characters up to a common unramified quadratic character. Here \(k_w\) is the residue field of \(L_w\). Indeed the \(i\)-th diagonal coordinate of the half modulus has
exponent \((n+1-2i)/2\), \(1\le i\le n\).

If
\(\gamma(q_w^{1/2})=\delta q_w^{1/2}\) with \(\delta\in\{1,-1\}\), its
value on a uniformizer changes in that slot by
\[
 \delta^{\,n+1-2i}=\delta^{\,n-1}\qquad(1\le i\le n),
\]
so the change is the same unramified character in every slot.
The order for the unique quotient depends on precedence of
linked segments. For singleton segments the precedence relations
depend on ratios equal to powers of the absolute value, and are
preserved because the absolute value takes rational values and
the common character cancels in a ratio. Thus the quotient
still has singleton rank-one segments with the conjugate
inertia characters. All \(\eta_i\) are fixed. The inverse
conjugation proves the reverse implication.
\cite[Theorem 6.1 and the segment classification]{zelevinsky}

For each member \(\xi\) of the finite Galois orbit of
\(\sigma_\tau\), let \(e_\xi\) be its central character idempotent
in \(\overline E[H]\), given by
\[
   e_\xi=\frac{\dim\xi}{|H|}
                 \sum_{h\in H}\operatorname{tr}\xi(h^{-1})\,h.
\]
The orbit sum \(e\) belongs to \(E[H]\). Let
\[
   M=eE[H]\cap O[H].
\]
This is a free \(O\)-module spanning \(eE[H]\), since each element
of an \(E\)-basis can be scaled into \(O[H]\).

It is stable under
left multiplication by \(H\). Its dual
\(M^\vee=\operatorname{Hom}_O(M,O)\) is the integral coefficient
representation at \(w\). The fixed space of
\((M^\vee\otimes_O\overline E)\otimes_{\overline E}\pi_w\)
under the diagonal \(K_w\)-action identifies with
\(\operatorname{Hom}_{K_w}(M\otimes_O\overline E,\pi_w)\).
Over \(\overline E\) the constituents of \(eE[H]\) are the
members of the orbit. The regular representation block of \(\xi\) contains
\(\dim\xi\) copies of \(\xi\). For \(\xi=\gamma\sigma_\tau\),
\[
 \operatorname{Hom}_{K_w}(\xi,\pi_w)\ne0
 \quad\Longleftrightarrow\quad
 \operatorname{Hom}_{K_w}(\sigma_\tau,\gamma^{-1}\pi_w)\ne0.
\]
The condition on the right is the criterion for the
inverse-conjugate parameter. It has type \(\tau\) exactly
when the original parameter does. This proves detection for
\(M^\vee\). The characteristic-zero orbit block and the
intersection lattice allow a Schur index over \(E\) and allow
\(\ell\) to divide \(|H|\).

\proofstep{The integral coefficient test at all tame places.}
The full level imposes the integral types at all tame places simultaneously.

Its coefficient module is the tensor product
\[ \bigotimes_{w\in S_{\mathrm t}} M_w^\vee ,\]
where \(M_w^\vee\) is the local lattice just constructed at \(w\). The kernel contains the product of
the local open kernels. For a representation that is isotypic
with specified local factor \(\pi_w\) under each coordinate
group, the coefficient fixed space is nonzero if and only if
each local criterion holds. Take fixed vectors in the local
factors successively. The other local groups commute with each
projection, so a nonzero image is still isotypic for each
remaining group. In an \(O\)-basis, the fixed-vector equation
is the matrix equation for the tensor product on the coefficient
index.
\label{cor:ordinary-level-data-1}

\end{proof}

\subsection{The weight character and the residual Hecke values}

Ordinary localization also uses the residual spherical eigenvalues and the ordered characters on local units.

Let \(\Lambda\) be the complete local weight algebra for the level,
with finite residue field. At \(w\in S_\ell\) its specified inertia
characters are restrictions of adically continuous characters
\(\psi_{w,i}\) of the decomposition group with values in
\(\Lambda^\times\). Use the local Artin map from the abelian
Weil group to \(L_w^\times\) taking geometric Frobenius to a
uniformizer modulo units. A continuous \(\Lambda^\times\)-valued
character of the decomposition group induces a character of
\(L_w^\times\). Modulo each power of the maximal ideal this
follows from local reciprocity for finite abelian characters.
The compatible characters give the result by completeness.
Each reduced character is smooth on \(\mathcal O_{L_w}^\times\).
Write \(\epsilon_\ell\) for the \(\ell\)-adic cyclotomic character
with values mapped to \(\Lambda^\times\). On diagonal units, use
\[
   \theta((y_{w,i})_{w,i})
   =\prod_{w\in S_\ell}\prod_{i=0}^{n-1}
          (\psi_{w,i}\epsilon_\ell^{\,i})
                  (\operatorname{Art}_w^{-1} y_{w,i})^{-1},
\]
At each coefficient precision a single sufficiently deep
diagonal congruence subgroup works for all the factors.

The stabilizer proof in Paragraph~\ref{par:ordinary-functions}
applies at each exceptional prime and gives the freeness condition
for this type and weight at the auxiliary levels.

\begin{construction}
\[
\begin{array}{ll}
\text{place}&\text{coefficient datum}\\ \hline
w\in S_{\rm t}&\eta_i=[\bar d_i]|_{I_w},\quad N_{\rm mon}=0\\
w\text{ split spherical}&P_w(X)=\prod_i(X-\bar d_i(f_w))\\
w\mid\ell&\text{ordered character products and the weight character}
\end{array}
\]
\caption{The residue labels enter the ordinary level in different ways at the three kinds of place. The spherical normalization uses geometric Frobenius.}\label{fig:ordinary-coefficient-places}
\end{construction}

\begin{lemma}[The two residual coefficient prescriptions]
\label{lem:ordinary-residual-data}
At each level the residual data consist of the spherical polynomials at the unmarked split spherical places and the partial products above \(\ell\):
\[
 P_w(X)=\prod_{i=0}^{n-1}(X-\bar d_i(f_w)),\qquad
 \prod_{i=0}^{j}(\bar\psi_{w,i}\bar\epsilon_\ell^{\,i})
                 (\operatorname{Art}_w^{-1}\varpi_w)\quad(0\le j<n),
\]
respectively, with \(f_w\) the geometric Frobenius at a split spherical place. Both sets of data lift to the ordinary coefficient data, with the given weight character
and the fixed reciprocity normalization. This defines the localized
modules even when one of them is zero.
\end{lemma}

\begin{proof}
\proofstep{The spherical residual polynomial.}
At a split prime outside the allowed ramification, put
\(q_w=\# k_w\) and let \(f_w\) be geometric Frobenius on the
level quotient. Set \(P_w(X)=\prod_i(X-\bar d_i(f_w))\).
For \(1\le j\le n\) choose a lift in \(O\) of
\[
   (-1)^j q_w^{-j(j-1)/2}[X^{n-j}]P_w(X).
\]
The denominator is a unit in \(k\) since the level contains
every prime above \(\ell\). The reduction of the normalized
Hecke polynomial with these lifted values is \(P_w\).
Choices of Frobenius lift and decomposition group give the same
coefficients by unramifiedness and conjugation invariance.
\proofstep{The ordered residual values above the coefficient prime.}
The unramified spherical polynomials fix the residual Hecke values away from the coefficient prime. At the places over that prime, localization also uses the order of the diagonal characters.

At \(w\in S_\ell\) write \(\varpi_w\) for the fixed uniformizer.
The value to lift for the product through
ordered index \(j\), with \(0\le j<n\), is
\[
   \prod_{i=0}^{j}(\bar\psi_{w,i}\bar\epsilon_\ell^{\,i})
                         (\operatorname{Art}_w^{-1}\varpi_w),
\]
Local reciprocity
gives independence of the lift of \(\varpi_w\) to the Weil group.
 At each level the weight character \(\theta\) and integral
type, together with these equations, determine the ordinary data. Their
existence suffices to define the localized modules, including
the case in which a localized module vanishes. \label{cor:ordinary-level-data-2}
\end{proof}

\section{Galois representations of the contributing ordinary constituents}
\label{sec:count-two-level-attachments}
\hypertarget{source-ordinary-level-attachments}{}

The two coefficient modules determine which classical forms can occur after localization. To compare their Hecke eigenvalues with deformation rings, we must attach Galois representations to the contributing forms at the same level. Starting
with an algebraic ordinary constituent, we use the polarized
Galois attachment and the local data at the split places to
verify the required conditions. The integral representation
and its ordered flags are used when specializing the seed and
in the ensuing multiplicity arguments. To a constituent \(\sigma\), Proposition~\ref{prop:polarized-separating-constituent} attaches a representation of \(G_L\). The products in \eqref{eq:ordered-translation} determine the characters on the successive quotients of its ordinary flag, in the chosen order.

Keep one of the two levels from Section~\ref{sec:count-two-level-data}.
The hypotheses on the ground field fixing this level are
\begin{equation}\label{eq:constituent-cm-field}
 \begin{gathered}
 L=L^+K_0,\qquad K_0\text{ imaginary quadratic},\qquad 2\mid[L^+:\mathbb Q],\\
 e(w/v)=1\qquad (v\text{ finite in }L^+,\ w\mid v).
 \end{gathered}
\end{equation}
The group is the unitary group
of the standard Hermitian form and is compact at the real places.
All places of \(L^+\) above \(\ell\) and all places where the
level allows ramification are split in \(L\). Auxiliary places
used to enlarge the level are split as well. Both levels in the chosen tower meet these splitting hypotheses, so the local attachment can be made at the very places on which their Hecke operators act.

\subsection{The polarized representation of a constituent}

At a separating classical evaluation the Hecke operators give Frobenius polynomials away from the level. The polarized attachment must have these polynomials and satisfy the local conditions dictated by the coefficient type.

\begin{proposition}[The polarized representation of a classical constituent]
\label{prop:polarized-separating-constituent}
Assume \eqref{eq:constituent-cm-field}. Use the standard unitary group compact at infinity and either level of Section~\ref{sec:count-two-level-data}, with all places above \(\ell\), allowed ramification places and auxiliary places split in \(L\). A contributing classical constituent \(\sigma\) at a separating evaluation has a polarized Galois representation 
\[
 \rho_\sigma:G_L\longrightarrow\mathrm{GL}_n(\overline{\mathbb Q}_\ell)
\]
 with the prescribed classical compatibility.
The construction applies at either level with its own coefficient
data. The automorphic transfer and the Galois attachment in the
proof use the cited results with those hypotheses.
\end{proposition}

\begin{proof}
\proofstep{The constituent decomposition of the classical space.}
Let \(C\) be the tensor of the algebraic coefficient space and
the finite level coefficient spaces, and \(K\) the product of
the groups at which their transformation rules are imposed.
Realize the vector-valued forms as intertwiners from \(C^\vee\)
into the scalar automorphic space. In its discrete decomposition
\(\bigoplus_\sigma M(\sigma)\otimes\sigma\), write \(M(\sigma)\)
for the multiplicity space. At the fixed finite level and
algebraic weight the coefficient contribution is finite
dimensional and decomposes as
\[
 \bigoplus_\sigma M(\sigma)\otimes
                 \operatorname{Hom}_{K}(C^\vee,\sigma).
\]
The local level, unit-character and ordinary operators act
within each constituent, so their selected summands impose
the tests of Paragraph~\ref{par:ordinary-functions} on this
same decomposition.
\cite[Lemma 1.25]{newton-thorne-unitary}

Let \(\sigma\) have nonzero intersection. Base change to \(L\)
gives an isobaric sum of discrete conjugate self-dual automorphic
representations of general linear groups. At every split place
the local component of this sum is the component of \(\sigma\)
under the chosen identification with \(\mathrm{GL}_n(L_w)\).
At unramified places it is the unramified base change, and its
infinitesimal character is determined by \(\sigma\) at infinity.
The base-change theorem and its Galois corollary attach the continuous semisimple representation
\[
 \rho_\sigma:G_L\longrightarrow\mathrm{GL}_n(\overline{\mathbb Q}_\ell).
\] It is unramified away from \(\ell\)
where the constituent is unramified. At split places its
Frobenius-semisimplified Weil--Deligne parameter is the
algebraically normalized local parameter of the constituent,
and it is de Rham above \(\ell\).
\cite[Theorem 1.2 and Corollary 1.3]{newton-thorne-unitary}

Here is the polarization and weight calculation for this
attachment. A discrete summand of rank \(m s\) is the Speh
representation of a conjugate self-dual cuspidal representation
\(\pi_0\) of rank \(m\) and length \(s\). Its normalized parameter
has the \(s\) twists of the cuspidal parameter with exponents
\((s-1)/2-a\), \(0\le a<s\). This is the discrete spectrum
classification of M{\oe}glin and Waldspurger.

To express them
using the normalization for cuspidal rank \(m\), put
\[
     t_a=\frac{m-n}{2}+\frac{s-1}{2}-a.
\]
The corresponding Galois summands are those of
\(\pi_0|\det|^{t_a}\). Their weights are algebraic and regular:
their infinitesimal characters after the common normalization
are disjoint parts of the regular algebraic infinitesimal
character of the rank-\(n\) base change. Each summand is
essentially conjugate self-dual with a norm-character
polarization, and the cuspidal attachment theorem supplies
its de Rham weights. The shifts \(t_a\) compensate for the
difference between the rank-\(m\) and rank-\(n\)
normalizations. Hence the Hodge numbers of the full
representation at an embedding are the numbers
\((n-1)/2-z\) for the infinitesimal-character entries \(z\).
With the dual of highest weight
\(\lambda_\tau=(\lambda_{\tau,0},\ldots,\lambda_{\tau,n-1})\)
at the embedding indexed by \(\tau\), these are
\[
    h_i(\tau)=\lambda_{\tau,n-1-i}+i,\qquad 0\le i<n.
\]
They are strictly increasing in the indicated order.
\cite{moeglin-waldspurger,blggt}

Fix complex conjugation \(c\). At unramified split places
conjugation takes the rank-\(m\) Galois summand with twist
\(t_a\) to the dual of the summand with twist \(t_{s-1-a}\)
times \(\epsilon_\ell^{1-n}\). At the level of normalized Frobenius parameters, conjugation
duals the rank-\(m\) parameter and contributes the exponent
\(1-m\). Number the partner of slot \(a\) by \(a'=s-1-a\).
Compared with the dual of that partner the exponent is
\[
 (1-m)+t_a+t_{a'}=(1-m)+(m-n)=1-n .
\]
Here the shifts \(t_a\) are exponents in the normalization
of the local parameters of the algebraic rank-\(n\) summands.
The split-prime Frobenius elements outside a finite set
are dense up to conjugacy in \(G_L\).

To see this in a
finite quotient, pass to a finite Galois extension of \(L^+\)
containing the fixed field of the quotient. Any element of
the subgroup fixing \(L\) is realized by a Frobenius element
over a split place of \(L^+\), by Chebotarev and choice of the
place above it. Semisimplicity and the Frobenius identities
therefore give
\(\rho_\sigma^c\simeq \rho_\sigma^\vee\otimes
\epsilon_\ell^{1-n}\).
An intertwiner gives a nonsingular pairing \(P\) with
\[
  \rho_\sigma(g)^{\mathsf T}P\,\rho_\sigma(cgc^{-1})
      =\epsilon_\ell(g)^{1-n}P .
\]
This is the pairing required for the determinant comparison.
Write \(s_\nu\) for the Speh length of the \(\nu\)-th discrete
summand. Each of its \(s_\nu\) cuspidal twists contributes a
nonzero direct summand to the semisimple attachment. Thus
\[
 \rho_\sigma\text{ absolutely irreducible}
 \ \Longrightarrow\ \sum_\nu s_\nu=1
 \ \Longrightarrow\ \text{one discrete summand of length one}.
\]
This identifies the base change as cuspidal.

The permitted level is hyperspecial off the finite ramification
set. A nonzero contributing vector implies that the local
constituent is unramified there. The attachment consequently
factors through the same unramified Galois quotient as the
determinant ring. At \(w\in S_{\mathrm t}\), the compact type
forces the full prescribed inertia representation and \(N=0\).
Since \(w\nmid\ell\), the Weil--Deligne description then gives
the actual diagonal action on inertia by the Teichmüller lifts.
At the split auxiliary places the local compatibility applies
to the selected local factor. At unramified split places the
coefficients of the geometric Frobenius polynomial are the
Satake scalars with the powers \(q_w^{j(j-1)/2}\) of
Section~\ref{sec:count-two-level-data}. This identifies the constituent by its split Hecke action. At \(w\mid\ell\) its ordinary operators will determine the flag. \label{cor:ordinary-level-attachments-1}

\end{proof}

\subsection{The ordered flag at the separating weights}

The attachment identifies the local representation at a split place above the coefficient prime. An ordinary deformation also records an ordering of its flag. At separating weights the eigenvalues of the Jacquet operators recover the characters in that order.

For the separating classical weights of
Paragraph~\ref{par:ordinary-functions}, the smooth unit
characters at \(w\mid\ell\) are pairwise distinct. Write
\(K=L_w\), \(q=\ell^f\) and \(e=e(K/\Q_\ell)\).

\begin{lemma}[Ordering the inducing characters by the ordinary Jacquet character]\label{lem:ordinary-jacquet-reading}
Let the common eigenvector at a separating classical weight
have pairwise distinct smooth unit characters and nonzero
ordinary translation eigenvalues. Write \(v\) for this vector, \(N_B\) for the upper unipotent radical and \(q_{N_B}:\pi_w\to(\pi_w)_{N_B}\) for the Jacquet quotient. Then
\[
 q_{N_B}(v)\ne0,
\]
and the chosen Jacquet character orders the inducing
Weil characters \(\omega_i\) of the local principal series.
This series is irreducible with zero monodromy. The
weight-normalized translation eigenvalues satisfy
\eqref{eq:ordered-translation}.
\end{lemma}

\begin{proof}
Let \(v\) be the prescribed eigenvector. If its image in
unipotent coinvariants vanished, write
\[
                   v=\sum_b (u_b-1)w_b
\]
with finitely many \(u_b\) in that unipotent group. Choose a
compact averaging group containing the \(u_b\), which makes the
average of every summand zero. The formula for the translation
by a sufficiently large power of the contracting diagonal uses
such an average, after conjugation. Its value on \(v\) is
instead a nonzero power of the ordinary eigenvalue times \(v\).
This contradiction proves that the Jacquet reading is nonzero. The Jacquet subquotient theorem gives a principal
series containing the local factor. Distinct restrictions to units rule out linking: for two of the inducing characters the implication is
\[
 \chi_i|_{\mathcal O_K^\times}\ne\chi_j|_{\mathcal O_K^\times}\quad\Longrightarrow\quad\chi_i/\chi_j\ne|\cdot|^{\pm1}.
\]
Thus their principal series is irreducible.

Let \(\omega_i\), in the order of this eigenvector,
be the Weil characters in the algebraic normalization, viewed
as characters of \(K^\times\). The local comparison gives zero
monodromy and these Weil lines. Distinctness on inertia implies
that the Weil action is already Frobenius-semisimple, since
Frobenius preserves its one-dimensional inertia eigenspaces.

Put \(\alpha_i=\omega_i(\varpi_w)\), taking \(\ell\)-adic values
using the coefficient identification. Write \(b_j\) for the
eigenvalue of the weight-normalized translation with
\(j\) uniformizer entries first. The two normalizations give
\begin{equation}\label{eq:ordered-translation}
 b_j
   =\frac{q^{-j(j-1)/2}\prod_{i<j}\alpha_i}
           {\prod_{\tau:K\hookrightarrow\overline{\Q}_\ell}
                   \tau(\varpi_w)^{\sum_{i<j}\lambda_{\tau,n-1-i}}}.
\end{equation}
In the unnormalized Jacquet module, the translation sum
contributes the index \(q^{j(n-j)}\) from the upper unipotent
group. Normalizing the Jacquet action contributes the square
root of the Borel modulus \(q^{-j(n-j)/2}\). Passing from the
inducing characters to the algebraically normalized Weil lines
contributes \(q^{-j(n-1)/2}\). Their product is
\[
       q^{j(n-j)}q^{-j(n-j)/2}q^{-j(n-1)/2}
                     =q^{-j(j-1)/2}.
\]
Removing the algebraic weight divides by the remaining
factor in \eqref{eq:ordered-translation}.
\end{proof}

\begin{construction}
\[
\begin{array}{ll}
\text{ordinary datum}&\text{Weil parameter or Hodge weight}\\ \hline
b_j&\text{shifted product on the first }j\text{ Weil lines}\\
b_j/b_{j-1}&\text{next line after the corresponding shift}\\
\operatorname{gr}_i&\text{Weil character }\omega_i,\quad
       h_i(\tau)=\lambda_{\tau,n-1-i}+i
\end{array}
\]
\caption{The cumulative ordinary eigenvalues determine the successive characters of the flag at a separating evaluation. The shifts are those of \eqref{eq:ordered-translation}.}\label{fig:ordinary-ordered-flag}
\end{construction}

\begin{proposition}[The flag with the ordered local characters]
\label{prop:ordinary-ordered-separating-flag}
Use the separating evaluations and the ordered Weil lines of
Lemma~\ref{lem:ordinary-jacquet-reading}. The local Galois
representation admits an ordered flag whose successive quotients
have Weil characters \(\omega_i\) and labeled weights

\[
 h_i(\tau)=\lambda_{\tau,n-1-i}+i.
\]
 The shifted products on a
geometric Artin lift of \(\varpi_w\) equal the \(b_j\) of
\eqref{eq:ordered-translation}; dividing consecutive products reads
the individual characters.
\end{proposition}

\begin{proof}
\proofstep{Equality in weak admissibility.}
The ordinary Jacquet character determines the valuations of ordered products of Frobenius eigenvalues. Equality in weak admissibility is what turns these numerical identities into subrepresentations of the local Galois representation.

On the ordinary summand every \(b_j\) is an \(\ell\)-adic
unit. Since \(v_\ell(\tau(\varpi_w))=1/e\), this says
\[
 e\,v_\ell\left(\prod_{i<j}\alpha_i\right)
       = \sum_{\tau:K\hookrightarrow\overline{\Q}_\ell}
                      \sum_{i<j}h_i(\tau).
\]
Work with the filtered module with descent data of
\(\rho_\sigma|_{G_K}\), over a finite Galois extension \(K'/K\)
where it is crystalline and the finite inertia characters
are trivialized.
The Weil lines give ordered \(\varphi\)-stable summands
with descent data. For their first \(j\) summands the
left side times \([K':K]\) is the Newton number in the
normalization summing over all embeddings of \(K'\).
The Frobenius slopes scale by the residue degree, and the
ramification index supplies the remaining factor.
Every embedding \(\tau\) of \(K\) has \([K':K]\) extensions and the weights repeat at these extensions. Hence
\[
 \sum_{\tau' : K'\hookrightarrow\overline{\mathbb Q}_\ell}\sum_{i<j}h_i(\tau'|_K)=[K':K]\sum_\tau\sum_{i<j}h_i(\tau).
\]

Write \(D_j\) for this filtered subobject over \(K'\), and put
\(S_j=\sum_{\tau}\sum_{i<j}h_i(\tau)\), with \(\tau\) ranging
over embeddings of \(K\). At each embedding of \(K'\),
the induced filtration inherits \(j\) ambient jumps, whose sum
is at least the sum of the first \(j\). Weak admissibility gives
the other inequality:
\[
 [K':K]S_j\ \le\ t_H(D_j)\ \le\ t_N(D_j)
     =[K':K]\,e\,v_\ell\!\left(\prod_{i<j}\alpha_i\right)
     =[K':K]S_j .
\]
Equality holds at each embedding because each individual excess
above the sum of the first \(j\) jumps is nonnegative. All its subobjects inherit the weak
admissibility inequality from the ambient module.
The equivalence for weakly admissible filtered modules,
with descent by full faithfulness, therefore gives
an ordered Galois flag. The \(i\)-th successive quotient
has Weil character \(\omega_i\) and weights \(h_i(\tau)\).
\cite{colmez-fontaine}

\proofstep{The ordered characters on the decomposition group.}
Under geometric reciprocity, the rank-one quotient has
character on \(K^\times\)
\[
    a\longmapsto \omega_i(a)
                     \prod_\tau\tau(a)^{-h_i(\tau)}.
\]
This is the local class-field description of a de Rham
character with the indicated smooth Weil part and labeled
weights, where the cyclotomic character has weight \(-1\).
On units the inverse coefficient character is
\(\theta_i^{-1}(a)=\omega_i(a)
\prod_\tau\tau(a)^{-\lambda_{\tau,n-1-i}}\).
Here \(\omega_i|_{\mathcal O_K^\times}\) is the smooth
factor. The quotient agrees on inertia with the
specified character after reinstating the shift by
\(\epsilon_\ell^{\,i}\) of Section~\ref{sec:count-two-level-data}.

On a geometric Artin lift of the uniformizer, the
cyclotomic value is \(N_{K/\Q_\ell}(\varpi_w)/q\).
Since \(\prod_\tau\tau(\varpi_w)=N_{K/\mathbb Q_\ell}(\varpi_w)\),
the two contributions to the quotient multiply to
\[
 \left(\frac{N_{K/\mathbb Q_\ell}(\varpi_w)}{q}\right)^i
 \prod_\tau\tau(\varpi_w)^{-i}=q^{-i}.
\] Thus the product of these shifted
characters for \(i<j\) is \(b_j\) by
\eqref{eq:ordered-translation}. Dividing consecutive
products recovers the full ordered characters.
This supplies the local flag at the separating evaluations of the ordinary Hecke action. \label{cor:ordinary-level-attachments-2}
\end{proof}

\section{An ordinary Hecke eigenpoint from quadratic automorphic induction}
\label{sec:count-two-seed-specialization}

Knowing how to attach a representation to a constituent leaves the question of occurrence. In the even-order rank-four case, quadratic automorphic induction of the mixed-weight Hilbert eigenform, followed by the prescribed weight twist, produces an ordinary eigenpoint at each of the two levels. Proposition~\ref{prop:rank-four-ordinary-seed} identifies its Galois determinant and the ordered characters of its ordinary flag. This is the starting point for the rank-four component comparison in the even-order case. Appendix~\ref{sec:additional-rank-four} constructs the links to the component being counted.

Consider the characteristic-zero seed on one of the two links.
Make the integral weight adjustment by the prescribed scalar
twist. The resulting representation \(\rho\) over \(L\) is
absolutely irreducible. The mixed Hilbert construction and quadratic automorphic induction give a regular algebraic, conjugate self-dual cuspidal representation \(\Pi\) with
\[
 \Pi\text{ on }\mathrm{GL}_4(\mathbb A_L),\qquad \Pi^c\simeq\Pi^\vee,\qquad r_\ell(\Pi)\simeq\rho
\]
in the attachment normalization in use. The character of the twist
is locally algebraic by the parity choice of the link.

Write
\(x\) for the specialized determinant point, and \(\psi_{w,i}\)
for its ordered local characters. At every place over \(\ell\)
these characters have cyclotomic exponents \(m_{w,i}\) on open
inertia, decreasing with gaps at least two. The corresponding
Hodge numbers \(h_i=-m_{w,i}\) increase with these gaps at
each embedding. These exponents determine the weight change in both determinant comparisons. The remaining question is whether the resulting point \(x\) belongs to the localized ordinary Hecke support at the chosen level.

\begin{proposition}[Admission of the even-order seed]\label{prop:rank-four-ordinary-seed}
At either of the two link levels in Section~\ref{sec:count-two-level-data}, the determinant point
\(x\) belongs to the support of the localized ordinary Hecke
module. If \(D_{{\rm Hecke},x}\) denotes the specialized Hecke determinant, then 
\[
 D_{{\rm Hecke},x}=\det\rho.
\]
 Its ordinary characters are the ordered \(\psi_{w,i}\).
\end{proposition}

\subsection{Local vectors of the seed}

An induced representation can have the desired determinant without contributing to a given level. To prove its occurrence here, construct local vectors for the actual coefficient types, including the ordered ordinary vector at the places above the coefficient prime.

\begin{proof}[Local vectors for Proposition~\ref{prop:rank-four-ordinary-seed}]
\proofstep{The local types away from the coefficient prime.}
Away from \(\ell\) the cuspidal local--global comparison applies
to \(\rho\) and \(\Pi\). Inertia is trivial outside the finite
level set and at its exceptional place. At the other type
places, the rank-two Hilbert system has finite inertia
characters \(1,\eta\) of order prime to \(\ell\). Upon induction
the inertia image on the subgroup fixing the inducing field
is still finite of order prime to \(\ell\), and the quotient
has order dividing two. The scalar corrections on this inertia
have order prime to \(\ell\): the square of the seed correction
there has order prime to \(\ell\), and the link scalar has
only its prime-to-\(\ell\) finite part away from \(\ell\).
Hence the inertia image of \(\rho\) has finite order \(m\) prime
to \(\ell\).

For \(g\) in inertia, every eigenvalue \(\alpha\)
satisfies \(\alpha^m=1\). Reduction on this group of roots is injective:
\[
 \mu_m(\overline{\mathbb Q}_\ell)\xrightarrow{\;\sim\;}
 \mu_m(\overline{\mathbb F}_\ell),\qquad
 \alpha=[\bar\alpha].
\]
Thus the reduced characteristic polynomial determines the full multiset
of inertia roots, including their multiplicities. The residual
semisimplification is the tuple of the prescribed global
labels. It follows that the full inertial type of \(\Pi\)
at each of these places is the type of
Section~\ref{sec:count-two-level-data}, including zero monodromy.
It is spherical at the exceptional place and off the finite
level set away from \(\ell\).
\cite[Theorem 2.1.1]{blggt}

\proofstep{The Frobenius characters and zero monodromy.}
At \(w\mid\ell\), apply the exact potentially semistable
functor to the ordinary filtration. The four graded Weil
characters give the Frobenius eigenvalues \(\alpha_i\) of the
full parameter. Equality of Hodge and Newton numbers in rank
one gives, with \(e=e(L_w/\Q_\ell)\) and \(q_w=\ell^f\),
\[
 e\,v_\ell(\alpha_i)=\sum_\tau h_i(\tau),\qquad
 v_\ell(\alpha_{i+1}/\alpha_i)
    =e^{-1}\sum_\tau(h_{i+1}(\tau)-h_i(\tau))\ge2f .
\]
Write \(V_\alpha\) for a Frobenius eigenspace. The
Weil--Deligne relation is
\[
 N(V_\alpha)\subseteq V_{q_w^{\pm1}\alpha},
 \qquad v_\ell(q_w^{\pm1}\alpha)-v_\ell(\alpha)=\pm f .
\]
The sign depends on the Frobenius convention. Neither valuation
difference occurs among the roots above, so all the spaces on
the right are zero and \(N=0\). The
inertia characters extend to the Weil group, so Frobenius
preserves their isotypic spaces. Its distinct roots split the
parameter into four Weil lines.

At every embedding the Hodge numbers obey
\[
 (h_{i+1}(\tau)-(i+1))-(h_i(\tau)-i)
             =h_{i+1}(\tau)-h_i(\tau)-1\ge1 .
\]
Hence the algebraic highest weight has strict consecutive gaps.
This implies Shin regularity, which gives full Frobenius-semisimplified
local--global compatibility at \(\ell\). \cite[Theorem 1.1]{blggt-lg}
The local representation is the normalized principal series
of these four lines with the common algebraic normalization
undone. The four inducing Weil characters \(\omega_i\) are unlinked. Indeed their uniformizer values satisfy
\[
 \omega_i(\varpi_w)/\omega_j(\varpi_w)\notin\{q_w,q_w^{-1}\}\qquad(i\ne j),
\]
which rules out \(\omega_i/\omega_j=|\cdot|^{\pm1}\) and makes the normalized principal series irreducible.

Apply cuspidal descent to the standard definite unitary group.
The extension \(L=L^+K_0\) is unramified at finite places over
\(L^+\). The local calculation away from \(\ell\) and splitting
at \(\ell\) give the required ramification condition
\[
 \{v\text{ finite of }L^+: \Pi_w\text{ ramifies for some }w\mid v\}
       \subseteq \{v:v\text{ splits in }L\}.
\]
Descent therefore gives a constituent \(\sigma\) with the specified
components at split places and hyperspecial invariants at inert places.
Its archimedean weight is the one read from the labeled Hodge numbers.
\cite[Theorem 4.4]{thorne-unitary-descent}

\proofstep{A unit ordinary eigenvalue at each step.}
The chosen local types and Frobenius characters place the induced form at the correct finite level. Ordinary projection retains it only if all the required, normalized operators have unit eigenvalues.

Choose the ordered inducing datum of the ordinary quotients
at each place over \(\ell\). For sufficiently deep \(d\), write an element of the Iwahori subgroup as \(nt\bar n\) in its decomposition \(J_d=NT\bar N_d\). The diagonal character modulo that depth extends by
\[
 \chi_{J_d}(nt\bar n)=\chi_T(t) .
\] It agrees on the upper triangular
intersection with the inducing character. In the induced
model the function supported on \(B J_d\) with value one at
the identity is well-defined and has this compact character.
Here \(B\) is the upper triangular group. Identity evaluation is nonzero on this line and reads the ordered Jacquet eigenvalues of the unnormalized translations. Indeed all representatives
\(u\,\operatorname{diag}(\varpi I_j,I_{4-j})\) in the translation
sum lie in \(B\), and evaluation of the inducing relation
gives the Borel modulus and the index used in
\eqref{eq:ordered-translation}.

Let \(\lambda_j\) denote these scalars and \(U_j\) the commuting
translations. The evaluation functional \(\lambda^*\) is therefore an eigenfunctional with scalars \(\lambda_j\). For a generalized eigenvector \(v\) of \(U_j\) of eigenvalue \(b\ne\lambda_j\), choose \(a\) with \((U_j-b)^av=0\). Then
\[
 0=\lambda^*((U_j-b)^av)=(\lambda_j-b)^a\lambda^*(v)\quad\Longrightarrow\quad \lambda^*(v)=0 .
\] Thus
the summand \(V_\lambda\) at the specified tuple is nonzero.
The operators \(U_j-\lambda_j\) are commuting nilpotents there.
Successively taking kernels keeps a nonzero space, giving
\[
 0\ne\bigcap_j\ker((U_j-\lambda_j)|_{V_\lambda}).
\]
A vector in this intersection realizes the required tuple.

For a geometric Artin lift \(g_w\) of the uniformizer,
rank-one reciprocity and \eqref{eq:ordered-translation} identify
the weight-normalized eigenvalue as
\[
    \lambda_{w,j}=\prod_{i<j}(\psi_{w,i}\epsilon_\ell^{\,i})(g_w),
             \qquad v_\ell(\lambda_{w,j})=0.
\]
Indeed each continuous character of the compact decomposition
group has valuation zero.
The smooth unit characters of the vector have the required
weight after removal of the algebraic factors. This
construction allows equality of some smooth unit characters.
Taking the product of these local vectors, the type vectors
away from \(\ell\) and spherical vectors off the level gives
a classical ordinary form on the specified level. Its
spherical Hecke values are the normalized coefficients of
the Frobenius polynomials of \(\rho\).

\end{proof}

\begin{construction}
\[
\begin{array}{rcl}
\Pi\text{ on }\mathrm{GL}_4(\mathbb A_L)
    &\longrightarrow&\text{classical form with the chosen local vectors}\\
\text{classical form}&\xrightarrow{\text{weight adjustment}}&
        \text{ordinary eigenform at the link}\\
\text{Hecke eigenvalues}&\longrightarrow&
        x\text{ with determinant of }\rho\text{ and characters }\psi_{w,i}
\end{array}
\]
\caption{The specialization uses the local choices at the link where the seed is being admitted. The ordered character list is retained at that level.}\label{fig:ordinary-seed-admission}
\end{construction}

\subsection{Specialization into the same Hecke algebra}

The local vectors place the induced form at the specified level. The remaining step converts its algebraic weight into a scalar ordinary function and verifies that evaluation gives the prescribed determinant point of that level's Hecke algebra.

\begin{proof}[Completion of the proof of Proposition~\ref{prop:rank-four-ordinary-seed}]
\proofstep{The scalar ordinary function.}
Increase the coefficient DVR to contain the finitely many
values and coefficients of the form. The classical form has
a realization \(F\) as a finite adelic function with values
in \(W_\lambda\) and the fixed tame coefficient space.
It satisfies
\(F(\gamma g)=W_\lambda(\gamma)F(g)\) for rational \(\gamma\)
and has the smooth coefficient conditions on the right. It can
be defined over a finite coefficient extension because at a
fixed finite level the space is constructed on a finite
double-coset set with algebraic coefficient matrices.
The function
\[
                F_\ell(g)=W_\lambda(g_\ell)^{-1}F(g)
\]
is left invariant under the rational group, with the
displayed action on the weight factor. Project the
algebraic coefficient at each \(\ell\)-adic place to
the lowest torus weight quotient \(p_\lambda\), which is
the coinvariant line for the upper unipotent group.

On a diagonal tuple \(y\) the character is the product over embeddings and ordered diagonal entries:
\[
 y\longmapsto\prod_{\tau,i}\tau(y_i)^{\lambda_{\tau,3-i}} .
\]
The function \(p_\lambda(F_\ell(g))\) now transforms by
the inverse specialized weight \(\theta_x\).
It is invariant on the right under the integral upper
unipotents. For every fixed coefficient precision it is
invariant under a sufficiently deep lower unipotent group,
by continuity of the algebraic action and the finite level
of \(F\). Its values are bounded, using compactness at
fixed level and finitely many double cosets. Thus scaling
places the function in the integral specialized space.

\proofstep{Nonvanishing and intertwining.}
The projection is nonzero on this form. If it vanished everywhere, fix \(g\) with \(F_\ell(g)\ne0\)
and let \(C\) be a sufficiently small open compact subgroup
at \(\ell\) fixing the original smooth form. Writing \(\rho_\lambda\)
for the algebraic coefficient action, right translation would give
\[
    p_\lambda\bigl(\rho_\lambda(h)^{-1}F_\ell(g)\bigr)=0
                        \qquad(h\in C).
\]
The subgroup \(C\) is Zariski dense in the product of the local
restrictions of scalars. Thus its translates of \(p_\lambda\)
span a nonzero algebraic subrepresentation of the dual of the
irreducible coefficient representation. They span the entire dual,
contradicting \(F_\ell(g)\ne0\).
Spherical operators away from \(\ell\) commute with the
projection.

In a translation sum at \(\ell\) the upper
unipotent factors act trivially on the quotient line, and
the diagonal translation acts by the inverse algebraic
factor in \eqref{eq:ordered-translation}. This intertwines
the weight-normalized classical operators with the
operators on the scalar ordinary functions. For every ordinary translation eigenvalue \(\alpha\) on this function the valuation criterion reads
\[
 \alpha\in\mathcal O^\times\quad\Longleftrightarrow\quad v_\ell(\alpha)=0 .
\]
Thus the function lies in the ordinary subspace.

\proofstep{Integral weight specialization.}
The weight base-change comparison can be made on the
integral scalar functions. Reduction of the ordinary space
over either the weight algebra or the specialized DVR is
the ordinary summand at sufficiently deep finite level for
the reduced character. At this precision functions lift
by their values on the double cosets, which have trivial
stabilizers. The ordinary projector commutes with reduction,
and the value filtration on its image is the coefficient
filtration. This is the finite-level argument of
Paragraph~\ref{par:ordinary-functions}.

If \(M_\Lambda^{\rm ord}\) is the universal ordinary
module and \(O_x\) the specialized DVR, the map is
\[
 M_\Lambda^{\rm ord}\otimes_{\Lambda,\theta_x} O_x
                          \longrightarrow M_x^{\rm ord}.
\]
Its reduction is the identity on the ordinary summand of
the same finite scalar-function space, by the lifting
and projector comparisons. The modules are finite free
of equal rank, so the map is an isomorphism by Nakayama's
lemma. This argument uses only the reduced
weight character to choose the depth for reduction, so it
applies to the locally algebraic \(\theta_x\).

\proofstep{Identification of the specialized point.}
The ordinary eigenvalues must now recover the specified determinant point of this level.

The spherical and ordinary Hecke values of the ordinary
function reduce to the prescribed residual values of
Section~\ref{sec:count-two-level-data}. Hence it belongs to
the same residual summand. Its eigenvalues define a
character of the localized Hecke algebra. The good split Frobenius classes are dense. On them the Hecke determinant and the determinant of \(\rho\) have the same characteristic coefficients, so continuity gives
\[
 D_{\rm Hecke,x}=\det\rho .
\]
The local characters are determined by their prescribed
values on inertia and the consecutive products at a
uniformizer. Thus the eigencharacter pulls back to \(x\)
on all the determinant and ordinary-character coordinates.
The kernel \(I_i\) of the localized Hecke action lies in
\(\ker x\). This is the arithmetic point assertion used
at the actual Initial representation in the even-order case on link \(i\). The argument
applies to both exceptional places with their own compact
conditions and the same source of the seed.
\end{proof}

\section{Selected ordinary multiplicities in rank three}
\label{sec:rank-three-multiplicity}
\hypertarget{source-ordinary-rank-three-multiplicity}{}

Existence of a classical eigensystem sufficed to start the rank-four comparison. In the rank-three trace calculation the coefficient of each contributing eigensystem is its multiplicity in the specified ordinary module. Proposition~\ref{prop:ordinary-three-local-dimensions} computes the three local factors: the ordinary multiplicity above \(\ell\), the multiplicity of the prescribed tame type, and the contribution of the depth-one projector \eqref{eq:depth-one-projector-splitting} at the exceptional place.

Let \(L/F\) be the CM extension in the rank-three ordinary
setting, unramified at every finite place of \(F\), and let
\(\ell\) be the coefficient prime. Use the definite unitary group of rank three for \(L/F\), compact at each real place, with the standard maximal compact subgroups at the finite places. Keep the
chosen split places for the level and their orientations
fixed. Denote the exceptional place of \(F\) by \(v_0\) and the chosen place of \(L\) above it by \(w\). Splitting and the residue-characteristic hypothesis require
\[
 L\otimes_F F_{v_0}\simeq L_w\times L_{\bar w},\qquad L_w\simeq F_{v_0},\quad \operatorname{char} k(v_0)\ne\ell .
\] Its level is
principal congruence of depth one. At the other split tame
places the coefficient type is fixed. The places above
\(\ell\) in the ordinary charts are split.

Fix an arithmetic weight \(\lambda\), a level and coefficient
type in this setting, and an eigensystem \(x\) of the
localized ordinary Hecke module. Assume that its arithmetic
eigenspace is nonzero and that its specialized determinant
is the determinant of an absolutely irreducible
representation. The localized ordinary operators have
unit eigenvalues. We compare the multiplicity of \(x\)
with that of another such point at the same algebraic weight
exponents and level. Taking both multiplicities on the image of the same integral idempotent makes the comparison depend on dimensions at fixed level and weight.

\subsection{The stable count at a fixed weight}
\label{sec:rank-three-stable-count}

The multiplicity has a global factor determined by the archimedean weight. At a cuspidal base change, the definite group has a unique contributing representation of the prescribed weight that is spherical at the inert places. Its multiplicity can therefore be determined before imposing the ordinary and tame projectors.

\begin{proposition}[The stable count at the fixed algebraic weight]
\label{prop:ordinary-stable-constituent-count}
For rank \(1\le n\le4\), let
\(\lambda^\infty\) be a fixed dominant weight at the
infinite places. Let \(\pi\) be a regular algebraic
conjugate self-dual cuspidal representation of
\(\mathrm{GL}_n(\mathbb A_L)\), with the infinitesimal
character prescribed by \(\lambda^\infty\), spherical at
the nonsplit finite places.

At a factorizable finite
level \(K\) that is standard maximal compact at nonsplit
places, write \(\sigma_i\), \(i\in I\), for the
irreducible unitary constituents with the prescribed
weight in the finite-level spectrum. The index set is
finite and counts the summands of a direct sum of
level-fixed spaces. Define \(I_\pi\) by agreement of the
spherical Satake multisets of \(\sigma_i\) and \(\pi\) at
all but finitely many split places. Assume the split
components of \(\pi\) have vectors fixed by the
prescribed split level. Then
\[
                   |I_\pi|=a_0(\lambda^\infty)>0,
\]
where \(a_0\) is chosen for \(L,n,\lambda^\infty\)
before choosing \(\pi\), the places of the level, or
the coefficient system.

\end{proposition}

\begin{proof}
\proofstep{Identify the base change from the split Hecke operators.}
Let \(F=L^+\) and let \(\sigma_i\) be one of the summands in \(I_\pi\). Since the group is definite, its real components belong to the discrete series. It is quasi-split at the finite places in the standard self-dual lattice model: at an inert place the extension is unramified and the standard Hermitian lattice is unimodular. Its stabilizer is hyperspecial. In particular the fixed vectors used at a nonsplit place belong to
\[
                   \sigma_{i,v}^{G(\mathcal O_{F_v})}.
\] The field \(F\) has degree at least two since an imaginary quadratic extension of \(\mathbb Q\) cannot be unramified at every finite prime.

Base change of \(\sigma_i\) gives a sum \(\Pi_i\) of discrete general linear constituents. For a discrete constituent of cuspidal rank \(m\) and segment length \(s\), the twists of its \(s\) summands in the rank-\(n\) Galois normalization are
\[
                         (m-n+s-1)/2-a,\qquad 0\le a<s,
\]
as in Section~\ref{sec:count-two-level-attachments}. Take their direct sum over all constituents. At all but finitely many places of \(L\) these summands recover the local unramified parameter of \(\Pi_i\). The algebraic regularity of the twists and their polarizations follow from the fixed algebraic infinitesimal character as in that section. \cite[Corollary~5.3]{labesse-cm} \cite[Theorem~2.1.1]{blggt}

Compare with the similarly normalized attachment to \(\pi\). Agreement at the chosen half of a split place gives agreement at its conjugate by conjugate self-duality, with the same rank normalization on both sides. Hence the Frobenius polynomials agree at almost all places \(w\) split over \(F\). At an inert place above \(v\) the norm is \((\mathrm N v)^2\), so the omitted set has Dirichlet density zero:
\[
 \sum_{w\ {\rm inert}}(\mathrm Nw)^{-z}\le\sum_v(\mathrm Nv)^{-2z}=O(1)\qquad(z\downarrow1).
\] Chebotarev density and Brauer--Nesbitt therefore give
\[
          r(\Pi_i)^{\rm ss}\simeq r(\pi)^{\rm ss}.
\]
Reading the characteristic polynomials at all their common unramified primes now gives equality of the almost-everywhere local parameters of \(\Pi_i\) and \(\pi\). Uniqueness of the general linear cuspidal support, with the segment exponents included, gives \(\Pi_i=\pi\). This is the general linear rigidity used in the proof of Labesse's multiplicity formula. \cite[proof of Theorem~5.4]{labesse-cm} \cite{moeglin-waldspurger}

\proofstep{Recover the constituent from the cuspidal base change.}
At a split place, base change identifies the component at the marked half with the unitary representation. At an inert place, \(K_v\)-invariance selects the spherical representation. Spherical base change is injective there: in the standard diagonal parametrization the conjugacy class is described by unordered inverse pairs, and, in odd rank, the additional middle entry. The base-change multiset recovers these pairs. At a real place of the compact group the dominant highest weight determines the representation uniquely. Consequently any two indices in \(I_\pi\) satisfy
\[
         \sigma_{i,v}\simeq\sigma_{i',v}\ (\text{all }v)
           \quad\Longrightarrow\quad \sigma_i\simeq\sigma_{i'}.
\]

It remains to see that this representation occurs, with the required hyperspecial vectors. In the refined cuspidal trace identity in the proof of Labesse's Theorem~5.4, take the hyperspecial unit at every inert place. With the twisted intertwiner normalized to fix the spherical vector, its twisted trace on the projected line is one. At a split place transfer exchanges the two factors. On finite-dimensional images the identity
\[
 \operatorname{tr}\big((A\otimes B)\operatorname{swap}\big)=\operatorname{tr}(AB)
\]
follows by summing diagonal coefficients in a basis. Thus the matched level units have twisted trace \(\dim\pi_w^{K_v}>0\). At infinity use the discrete-series test of the prescribed weight. Once the general linear eigensystem is isolated, the twisted cuspidal term is nonzero; the discrete unitary sum therefore contains a constituent with these level and weight conditions. For the fields \(L=L^+K_0\) of \eqref{eq:constituent-cm-field}, this existence, including the hyperspecial condition, is also exactly Thorne's descent theorem. \cite[Theorem~5.4 and its proof]{labesse-cm} \cite[Theorem~4.4]{thorne-unitary-descent}

\proofstep{Distinguish multiplicity from the dimension of the fixed space.}
The unique constituent is cuspidal on the definite group and spherical at the inert places. Compactness at infinity supplies the discrete-series and ellipticity conditions in Labesse's multiplicity-one theorem. Hence its automorphic multiplicity is one. In the finite-level direct sum its contribution is one summand \(\sigma_f^K\), of possibly larger dimension. Thus, for its unique contributing isomorphism class \([\sigma]\),
\[
                |I_\pi|=m_{\rm cusp}(\sigma)=1 .
\] The local dimensions used below are computed inside that summand. \cite[Theorem~5.4]{labesse-cm}
\end{proof}

\subsection{Which constituents contribute to the ordinary system}

The stable identity counts a packet. The ordinary trace only sees those constituents in the packet that realize the chosen eigensystem with the prescribed local characters.

\begin{lemma}[Constituents with the selected ordinary eigensystem]
\label{lem:ordinary-contributing-constituents}
For the nonzero ordinary eigenspace with an absolutely irreducible
specialization, the contributing constituents are precisely the
members of \(I_{\pi_x}\) with the prescribed local ordinary character.
If \(i_*\) indexes one contributing constituent, then
\[
 i\in I_{\pi_x}\quad\Longrightarrow\quad
     \sigma_{i,v}\simeq\sigma_{i_*,v}\qquad(v\text{ split}).
\]
These split components give the ordinary line, tame type and
projected exceptional factor for each such constituent.
\end{lemma}

\begin{proof}
\proofstep{Control at a common classical depth.}
The characteristic-zero specialization at \(\lambda\)
identifies the integral ordinary fiber with a space of
classical forms with a prescribed algebraic coefficient
representation. Let \(\lambda^\infty\) be its weight
under the coefficient-field isomorphism with
\(\mathbb C\). The finite-level decomposition at this
weight may be taken at a common depth that kills the
tame type on the deeper subgroup and contains the
specified ordinary invariants. If \(c'_0\) bounds the tame-type level and \(b\) is the classical specialization's level parameter, the depth chosen for this test is
\[
 c'_0+b+1 .
\] Expansion in the
finite-level direct sum gives the direct sum of
coefficient subspaces \(Y_i\) on the classical
space. Here a constituent is first restricted to the
right level and then paired with the algebraic and
tame coefficient representations. The resulting
subspaces are invariant under the Hecke operators,
the level constraints and the exceptional-place
operator.

Write \(Z_x\) for the joint generalized eigenspace
with the prescribed ordinary conditions. The nonzero eigenspace meets at least one summand. Choose \(i_*\) and a vector with
\[
 0\ne v\in Y_{i_*}\cap Z_x .
\] The characteristic
polynomials on this summand are those of the
specialized determinant, by the unramified Hecke
polynomial identity. Equality on a dense set of
Frobenius elements gives equality everywhere,
and Brauer--Nesbitt identifies the semisimple
representation attached to the summand with the
given irreducible representation. The rank-three cuspidal transfer therefore attaches
to \(\sigma_{i_*}\) a regular algebraic conjugate
self-dual cusp \(\pi_x\) of the prescribed
infinitesimal character. It is unramified at
nonsplit places, where the unitary constituent
has maximal-compact fixed vectors. At each chosen split place \(v\), denote the prescribed group at the common depth by \(K_v\). The decomposition of the constituent at this level gives
\[
 (\sigma_{i_*,v})^{K_v}\ne0 .
\]

If \(Y_i\cap Z_x\ne0\), use the cofinite split
places outside the level and the ordinary set.
At each such place the elementary Hecke generators
act as scalars on the corresponding spherical
line. For each elementary Hecke generator \(T_{w,j}\) let
\(a_{w,j}\) be its eigenvalue in the specified system and
\(b_{w,j}\) the scalar on the constituent. A nonzero
\(v\in Y_i\cap Z_x\) satisfies
\[
 0=(T_{w,j}-a_{w,j})^N v=(b_{w,j}-a_{w,j})^N v
\]
for some \(N\). Thus \(b_{w,j}=a_{w,j}\) for every \(j\).
The Hecke polynomials agree with the same normalization.
Since the Hecke polynomial recovers the multiset
of Satake parameters, this gives \(i\in I_{\pi_x}\).

\proofstep{The converse by the split component comparison.}
One inclusion identifies the packet of every contributing form. For the converse, comparison at the split places must show that packet members with the prescribed local character really contribute to the localized module.

Conversely the split component comparison shows that membership in the packet identifies the local components, namely
\[
 i\in I_{\pi_x}\quad\Longrightarrow\quad \sigma_{i,v}\simeq\sigma_{i_*,v}\qquad(v\text{ split}).
\] In particular the spherical
eigenvalues match every tested spherical operator.
Here the sets of split level primes coincide, \(S'=S\), so no additional level variables occur. Within the classical
ordinary conditions, the only remaining test is
the ordinary joint generalized eigenspace of the
local components at the places above \(\ell\).
Its rank is at most one for the prescribed
normalization and is positive for \(i_*\), by the
ordinary calculation below.

All members of the packet have the same local components at these places. With \(r_i\) denoting the ordinary rank just bounded,
\[
 1\le r_{i_*}\le1,\qquad r_i=r_{i_*}=1\quad(i\in I_{\pi_x}).
\]
The factorization of the classical ordinary
space into local components now identifies
\(Y_i\cap Z_x\), with the exceptional projector
inserted, with the tensor product of that
ordinary line, the tame type space and the
selected exceptional space. Its dimension is
therefore independent of \(i\in I_{\pi_x}\).

The same argument with the identity in place
of the exceptional projector shows that the
tame type space and the exceptional
principal-congruence fixed space both have
positive dimension. Indeed their tensor
product with the ordinary line at \(i_*\)
contains a nonzero vector. This supplies the
nonvanishing hypothesis for projecting the
exceptional space.

\end{proof}

\subsection{The three local dimensions}

For the surviving constituents, the projected dimension is a product of local dimensions. The calculation separates the ordered ordinary line, the prescribed tame type and the exceptional split place, where the depth-one projector changes the dimension.

\begin{proposition}[The ordinary, tame and depth-one dimensions]
\label{prop:ordinary-three-local-dimensions}
For each constituent in Lemma~\ref{lem:ordinary-contributing-constituents}, the ordinary joint
eigenspaces at the places over \(\ell\) are lines. The fixed tame
type has dimension \(d_{\mathrm{tame}}>0\) independent of the point
at fixed tame labels. At the exceptional split place the prescribed
depth-one projector has rank 
\[
 D_0=(q-1)^{n-1}
\]
 on the fixed
vectors, with \(q\) the residue cardinality.
\end{proposition}

\begin{proof}
\proofstep{The depth-one projector at the exceptional place.}
At \(v_0\) write \(K_1\) for principal congruence in
\(K=\mathrm{GL}_n(\mathcal O)\), where \(\mathcal O\)
is the local integer ring. Let \(k\) have order \(q\)
and let \(U\) be the upper unitriangular group. Fix a
nontrivial additive character \(\bar\psi\) of \(k\).
The characters
\[
  \chi_t(u)=\bar\psi\!\left(\sum_{i=1}^{n-1}t_i u_{i,i+1}\right),
  \qquad t\in(k^\times)^{n-1},
\]
are distinct characters of \(U(k)\). Their projectors
sum to the convolution idempotent
\[
 e=\frac{1}{|U(k)|}\sum_{u\in U(k)}
    \prod_{i=1}^{n-1}\bigl(q\,\mathbf{1}_{u_{i,i+1}=0}-1\bigr)[u].
\]
For a simple entry \(x\in k\), additive orthogonality gives
\[
 \sum_{t\in k^\times}\bar\psi(-tx)
 =\sum_{t\in k}\bar\psi(-tx)-1=q\,\mathbf{1}_{x=0}-1.
\]

Write \(e_t=|U(k)|^{-1}\sum_u\chi_t(u)^{-1}[u]\).
Character orthogonality on the finite group gives
\begin{equation}\label{eq:depth-one-projector-splitting}
 e_t e_{t'}=\mathbf 1_{t=t'}e_t,\qquad
 e=\sum_{t\in(k^\times)^{n-1}}e_t,\qquad
 eW^{K_1}=\bigoplus_t W^{K_1}[\chi_t].
\end{equation}
Indeed the sum over the independent \(t_i\) factors into the simple-entry
sums just evaluated. The denominator \(|U(k)|=q^{n(n-1)/2}\) is
a unit in the coefficient ring at \(\ell\). The group acts through compact
lifts on \(K_1\)-fixed vectors. In the ordinary module
these translations are at a place distinct from the
spherical and ordinary operators, so \(e\) commutes
with their actions.

Each contributing component \(W\) at \(v_0\) is a
local component of a cuspidal transfer. In the Kirillov model the translates of a Whittaker functional \(\Lambda\) by the mirabolic group \(P\) separate vectors. Equivalently,
\[
 \bigcap_{g\in P}\ker\bigl(v\mapsto\Lambda(gv)\bigr)=\{0\}.
\]
\cite[\S5.2 and Corollary 8.2]{zelevinsky}
Take the additive character on the local field of
conductor the maximal ideal and use lifts of the
\(t_i\) to units. Diagonal conjugation within the
mirabolic group gives the separating functional
for this choice of \(\chi_t\).

Starting from \(0\ne w\in W^{K_1}\), choose a
translate with nonzero Whittaker evaluation. By
Iwasawa decomposition and equivariance it can be
written \(D w_1\), where \(D=\operatorname{diag}(d_i)\)
and \(w_1\in W^{K_1}\). Write \(x_{ij}(a)=1+aE_{ij}\), let \(\mathfrak p\) be the
maximal ideal and let \(\widetilde t_i\) be the unit lift of \(t_i\).

For \(a\in\mathfrak p\), the simple root group fixes \(w_1\). Thus
\[
 \Lambda(Dw_1)
 =\Lambda(Dx_{i,i+1}(a)w_1)
 =\psi(\widetilde t_i d_i d_{i+1}^{-1}a)\,\Lambda(Dw_1).
\]
The evaluation is nonzero and \(\psi\) has conductor \(\mathfrak p\);
hence \((d_i/d_{i+1})\mathfrak p\subseteq\mathfrak p\).
Consequently \(v(d_1)\ge\cdots\ge v(d_n)\). For \(i>j\), this gives
\[
 D^{-1}x_{ij}(a)D=x_{ij}(a d_j/d_i)\in K_1
       \qquad(a\in\mathfrak p).
\]
So the lower and diagonal parts of \(K_1\) fix \(Dw_1\);
the upper part preserves \(\Lambda(Dw_1)\) by its conductor. Gauss decomposition of
\(K_1\) gives these three factors, since all pivots
are units. Averaging over \(K_1\) and applying the
\(\chi_t\)-projector preserve the nonzero evaluation.

To prove uniqueness, let \(w\in W^{K_1}[\chi_t]\)
have \(\Lambda(w)=0\) and suppose \(w\ne0\).
Separation by the mirabolic translates gives a translate testing this nonzero vector, even though its original value is zero:
\[
 \Lambda(w)=0,\qquad \Lambda(gw)\ne0\quad\text{for some mirabolic }g.
\] Starting above the last row,
suppose the rows below row \(i\) are already upper
unitriangular. For \(j\le i\) and \(a\in\mathcal O\),
the conjugate of \(1+aE_{j,i+1}\) by \(g\) is upper
unipotent. Its only possibly nonzero simple entry
is \(ag_{ij}\) in position \((i,i+1)\), by the
support of \(ge_j\) and of the \((i+1)\)-st row of
\(g^{-1}\). Evaluating \(\Lambda(g x_{j,i+1}(a)w)\) in these two ways and dividing by
\(\Lambda(gw)\ne0\) gives
\[
 \psi(\widetilde t_i a g_{ij})
     =\psi(\widetilde t_i a\delta_{ij})\quad(a\in\mathcal O).
\]

As \(\widetilde t_i\) is a unit, the conductor condition gives
\(g_{ij}-\delta_{ij}\in\mathfrak p\).
Right multiplication by \(1+e_i b^{\mathsf T}\),
with \(b_j=-g_{ij}/g_{ii}\) for \(j<i\),
\(b_i=g_{ii}^{-1}-1\) and \(b_j=0\) for \(j>i\),
is a correction in \(K_1\). It makes row \(i\)
upper unitriangular and preserves the lower rows
and the nonzero evaluation. Let \(k_1\in K_1\) be the product of the successive right corrections.
The final matrix is \(u=gk_1\in U\). Since \(k_1 w=w\), this gives
\[
 0\ne\Lambda(gw)=\Lambda(gk_1w)=\Lambda(uw)
                   =\psi_U(u)\Lambda(w)=0,
\]
where \(\psi_U\) is the chosen Whittaker character on \(U\).
Thus \(\Lambda\) is injective on each nonzero space
\(W^{K_1}[\chi_t]\); each has dimension one.
Summing these dimensions by \eqref{eq:depth-one-projector-splitting} gives
\[
                         D_0=(q-1)^{n-1}.
\]

\proofstep{The fixed tame coefficient type.}
At a split tame place, the fixed coefficient type
of Section~\ref{sec:count-two-level-data} gives
zero monodromy and fixed inertia labels. These are
the same Teichm\"uller lifts for both points.
The labels extend to the Weil group, so Frobenius
preserves every inertia isotypic space. Frobenius
semisimplification expresses the parameter as a
sum of Weil characters with exactly those inertia
labels. The local components of cuspidal transfers
are generic.

By the segment classification,
zero monodromy gives singleton segments, and the
one-dimensional Weil characters give rank-one
cuspidal factors \cite[Theorem 9.7]{zelevinsky}.
Thus the local component is the irreducible
normalized principal series of the resulting
characters. For \(K_w=\mathrm{GL}_n(\mathcal O_{L_w})\) and upper
triangular subgroup \(B_w\), Iwasawa decomposition gives
\[
 \left.\operatorname{Ind}_{B_w}^{\mathrm{GL}_n(L_w)}
                      (\textstyle\bigotimes_i\chi_i)\right|_{K_w}
   \simeq
 \operatorname{Ind}_{B_w\cap K_w}^{K_w}
                         (\textstyle\bigotimes_i\chi_i|_{\mathcal O_{L_w}^\times}).
\]
The modulus factor is trivial on \(B_w\cap K_w\).
The unit characters agree by the stipulated inertia
and local reciprocity. Tensor the compact isomorphisms
and take invariants with the same fixed tame coefficient
representation. This gives the common \(d_{\mathrm{tame}}>0\).

\proofstep{The ordered uniformizer values of the ordinary line.}
At a place over \(\ell\), put \(q_\ell=\ell^f\),
where \(f\) is the residue degree. For the
unnormalized translations by
\(t_j=\operatorname{diag}(\varpi I_j,I_{n-j})\)
write \(\mu_j\) for the prescribed eigenvalues and
\(\mu_0=1\). Write
\(a_i=\prod_\tau \tau(\varpi)^{\lambda_{\tau,n-1-i}}\)
using the dominant weights with indices from zero.
Under the isomorphism from the \(\ell\)-adic
coefficient field to \(\mathbb C\), \(\mu_j\)
corresponds to \(\prod_{i<j}a_i\) times a unit. Dividing
successive products gives
\[
  \alpha_i=q_\ell^i\mu_{i+1}/\mu_i,\qquad
  v_\ell(\alpha_i)=if+v_\ell(a_i).
\]
Dominance makes \(v_\ell(a_i)\) nondecreasing, so the
valuations of the \(\alpha_i\) are strictly increasing.

\proofstep{Injectivity of the projection to the Jacquet module.}
For completeness, write \(U_0=U(\mathcal O)\),
and act on \(U_0\)-invariants by the sum of the
translations \(ut_j\) over
\(U_0/t_jU_0t_j^{-1}\). The product of these operators
is the translation sum for \(t=t_1\cdots t_n\), by
using representatives for the successive
subgroups. If \(\mathcal U_j\) is the translation sum for \(t_j\),
the product \(\prod_j\mathcal U_j\) is invertible on
the joint generalized eigenspace because each eigenvalue
\(\mu_j\) is nonzero. On the other hand any vector of
zero Jacquet image can be written
\[
                      v=\sum_m (u_m-1)w_m .
\]
Choose \(r\) with every \(u_m\in t^{-r}U_0t^r\).
Averaging over this subgroup kills each summand and hence
kills \(v\). After translation
by \(t^r\), this says that the corresponding
translation sum is zero. Consequently the
Jacquet map is injective on the refined
generalized eigenspace.

\proofstep{Uniqueness of the selected torus ordering.}
The exceptional projector and tame type have left the order of the torus characters as the remaining possible multiplicity. The ordinary eigenvalues distinguish that order.

On the normalized Jacquet module the scalar multiplying the torus action of \(t_j\) is
\[
 q_\ell^{j(n-j)}\delta_B(t_j)^{1/2}=q_\ell^{j(n-j)/2}.
\]
If the Jacquet module of an irreducible
component is nonzero, it has a character
quotient. Frobenius reciprocity embeds the
component into the corresponding principal
series. Exactness of the Jacquet functor and
the geometric lemma bound the dimension of
any refined eigenspace by the number of
permutations of the inducing characters with
the prescribed torus eigenvalues
\cite[\S2.12]{bernstein-zelevinsky}.
Write each inducing character as
\(\psi_i=\omega_i|\cdot|^{(n-1)/2}\).

For a permutation \(\sigma\), inserting the modulus into the normalized operator gives the line eigenvalue
\[
 q_\ell^{-j(j-1)/2}\prod_{i<j}\omega_{\sigma(i)}(\varpi),\qquad \delta_B(t_j)=q_\ell^{-j(n-j)} .
\]
Thus the first \(j\) selected \(\omega\)-values
at the uniformizer must have product
\(q_\ell^{j(j-1)/2}\mu_j\). Successive quotients
give the \(\alpha_i\), which are pairwise
distinct. There is at most one such
permutation, even before testing its
characters on the units. This gives the
claimed upper bound of one. Positivity on
the initial contributing constituent
forces equality there and throughout its
split component comparison.

\end{proof}

\subsection{The same projected Hecke module at both points}

The local dimensions can now be multiplied with the global packet multiplicity. Comparing two specializations requires the same projector on the same integral Hecke module, since a change of projector could change that product.

\begin{construction}
\[
\begin{array}{lr}
\text{contribution}&\text{factor in }\mu_e(x)\\ \hline
\text{global multiplicity at the fixed weight}&a_0(\lambda^\infty)\\
\text{ordered eigenspace above }\ell&1\ \text{at each place}\\
\text{tame type}&d_{\mathrm{tame}}\\
\text{exceptional depth-one projector}&D_0=(q-1)^{n-1}
\end{array}
\]
\caption{Factors of the projected multiplicity at an irreducible arithmetic point. At fixed weight and local data each contributes the same factor at both points being compared.}\label{fig:ordinary-multiplicity-factors}
\end{construction}

\begin{proposition}[Multiplicity at fixed algebraic weight]\label{prop:ordinary-multiplicity-equal}
For an irreducible arithmetic point \(x\) of the localized
ordinary support, the selected multiplicity is
\[
 \mu_e(x)=a_0(\lambda^\infty)d_{\mathrm{tame}}D_0>0 .
\]
The local type, exceptional projector, level and algebraic
weight exponents determine these factors. In particular the
multiplicity is the same at any other such point \(y\) with
those data, even if its finite-order characters differ.
\end{proposition}

\begin{proof}
\proofstep{Compatibility of control with the projector.}
Let \(M\) be the finite free localized ordinary
module over the weight ring \(\mathcal W\),
at the level used for the determinant quotient.
The determinant ring maps onto the Hecke
image acting on \(M\). The kernel is the
annihilator of \(M\). Since the projector commutes with both actions, the same splitting is defined over \(\mathcal W\) and over the determinant ring:
\[
 M=eM\oplus(1-e)M .
\]

For an arithmetic point \(x\) whose kernel contains
the Hecke ideal, specialize \(M\) by the restriction
of \(x\) to \(\mathcal W\). The arithmetic
control map is Hecke-equivariant and gives an
identification with the specified classical
ordinary space after extension of scalars.
The classical joint generalized eigenspace
is \(Z_x\), and it is nonzero for \(x\)
in the Hecke support. Write \(C\) for the Hecke-equivariant control map. Its compatibility with the integral translation sum gives
\[
 C\,(|U(k)|e)=(|U(k)|e)\,C .
\] Dividing by the unit
\(|U(k)|\) gives the equality of selected
multiplicities before and after control.
The spherical operators and the ordinary
operators generate the Hecke action. Testing
their generalized eigenvalues therefore gives
the specialized local Hecke summand, and
projecting commutes with this test.

\proofstep{Multiplication of the three local dimensions.}
The same idempotent now acts at the two specializations. The three local dimensions can therefore be multiplied with the global packet count to compare the two projected multiplicities.

Now each contributing index has projected
dimension \(d_{\mathrm{tame}}D_0\), which is
positive in view of its unprojected
nonvanishing. Proposition~\ref{prop:ordinary-stable-constituent-count}
now gives
\[
             \mu_e(x)=a_0(\lambda^\infty)
                             d_{\mathrm{tame}}D_0>0.
\]
At a second irreducible arithmetic point
\(y\) in the same Hecke support with the same algebraic
weight exponents, the stable transfer may change. The
finite-order characters can change as well. Ordinary
control uses the homomorphism from the weight ring
at each point separately, on an open subgroup where
its finite-order factors are trivial.
All three factors on the right agree.
In particular \(\mu_e(x)=\mu_e(y)\) on
the idempotent range which controls the
fixed-weight derivations. The comparison
requires a common algebraic weight, level, tame type
and exceptional projector. The constant
\(a_0\) is fixed by the field, rank and
archimedean weight. The proof allows the
depth of classical specialization and
the auxiliary trace tests to vary.
\end{proof}

\section{Ordinary components from rank-two newforms and their adjoints}
\label{sec:linked-seeds}
\hypertarget{source-ordinary-linked-seeds}{}

The rank-three trace estimate and the estimate over \(\Q(i)\)
use a comparison in ranks two and three on a common
CM field. Rank two enters because a rank-three deformation can preserve
its middle line together with the plane on its two outer coordinates.
A rank-two newform and the trace-zero adjoint of a second newform provide the two starting eigenpoints. To propagate support one must also identify the Hecke ideals containing these points. These ideals
specify the starting components for the support comparison
in Appendix~\ref{sec:linked-components}. Formula~\eqref{eq:marked-newform-choices} records the weights and characters that realize the two residual label sets on the marked field. Proposition~\ref{prop:marked-rank-two-three-seeds} identifies the resulting ordinary flags.

\subsection{Two choices of a classical newform}
\label{subsec:marked-newform-parameters}

The two newforms need different weights and nebentypus characters to give the prescribed labels over the marked field.

Write \(\epsilon\) for the \(\ell\)-adic cyclotomic character, and fix
an identification \(\iota:\overline{\Q}_\ell\simeq\C\). The marked
field data are an odd prime \(\ell\), an abelian CM extension \(L/\Q\),
an imaginary quadratic field \(K\subset L\) of character \(\eta\),
and a finite character \(\chi'\) of \(G_\Q\). The two characters
factor through \((\Z/N_0)^\times\), with \((N_0,\ell)=1\).
There is a prime \(w\nmid N_0\) with \(\chi'(w)=\eta(w)=1\) and
a nontrivial character \(\theta\bmod w\) of order \(\ell\)
trivial on \(G_L\). These requirements give
\[ \ell\mid w-1,\qquad w\ne\ell,\qquad \theta(-1)=1 . \]

For complex conjugation \(c\) the marking requires
\(\chi'(c)=(-1)^j\). The scalar character to be imposed on the frame is
\[
 \Psi(g)=\epsilon(g)^j\chi'(g) .
\] There is
\(g_*\in G_L\) with \(\bar\Psi(g_*)^2\ne1\). A bar denotes
reduction in the residue field of \(\overline{\Q}_\ell\).
In the application \(\chi'\) has order prime to \(\ell\),
\(\bar\Psi\) equals the prescribed residual scalar on the frame,
and \(\chi'\ne\eta\). The level integer \(M\ne0\) is divisible
by \(2\ell\). The extension \(L/L^+\) is unramified at all finite
places and split over the primes dividing \(M\), and
\(\chi'|_{G_L}\) is unramified outside \(M\). These are the
field and character conditions used in the two seed constructions.

We need eigenforms at each of the following parameters. In either
row \(m\equiv j\bmod(\ell-1)\) can be taken arbitrarily large.
\begin{equation}\label{eq:marked-newform-choices}
\begin{array}{c|c|c|c}
\text{seed}&k&\phi&\text{operation on the attachment}\\ \hline
2&2m+1&(\chi')^2\eta&\text{twist by }(\epsilon^m\chi')^{-1}\\
3&m+1&\chi'\eta&\text{take the trace-zero adjoint}
\end{array}
\end{equation}
In the first row \(\phi\) is odd, and in the second row it is
nontrivial by \(\chi'\ne\eta\). In both rows
\(\phi(c)=(-1)^k\) and \(\phi(w)=1\). Take the primitive character
of conductor \(f\mid N_0\) associated to \(\phi\). We have \(f>1\).
The exact conductor of \(\gamma=\phi\theta\) is \(fw\).

\begin{lemma}[Newform seeds with the prescribed congruences]
\label{lem:ordinary-newform-seeds}
For the field and character data of
Section~\ref{subsec:marked-newform-parameters}, the two rows of
\eqref{eq:marked-newform-choices} give ordinary seed representations. With
\(\Phi=\epsilon^{k-1}\phi\) and the auxiliary character \(\theta\) of
conductor \(w\), the initial two-dimensional attachment
\(\rho:G_{\mathbb Q}\to\mathrm{GL}_2(\overline{\mathbb Q}_\ell)\) satisfies
\[
       \det\rho=\Phi\theta,\qquad
       \bar\rho^{\rm ss}\simeq 1\oplus\bar\Phi.
\] The
choice at rank three is shared by both exceptional levels.
The classical existence, congruence and Galois attachment used
in the construction are the cited inputs.
\end{lemma}

\begin{proof}
\proofstep{The Eisenstein congruences for the two newform choices.}
For this application one needs the Eisenstein congruence for both
parities. Here is a calculation of its constant valid in every
admissible large weight. With \(L_k(X)=(X\,d/dX)^{k-1}(X/(1-X))\)
as in the cusp correction of Section~\ref{sec:count-three-tensor},
one has the rational identity
\[
    \sum_{h=0}^{w-1}L_k(X\zeta_w^h)=w^kL_k(X^w).
\]
It follows by summing the power series for \(|X|<1\), since only
indices divisible by \(w\) survive. For a primitive \(fw\)-th
root \(z\), put \(z_f=z^w\). The constant of the normalized
Eisenstein series at infinity is
\[
  \frac{1}{2\tau(\gamma^{-1})}
     \sum_{a\in(\Z/fw)^\times}\gamma^{-1}(a)L_k(z^a).
\]
The primitive Gauss-sum product is \(\gamma(-1)fw\), and both
factors are integral at \(\ell\). The denominator is a unit.
Each lift of a unit \(u\bmod f\), including the lift divisible by
\(w\), evaluates \(L_k\) at a nonidentity prime-to-\(\ell\) root.
Its denominator \((1-z^a)^k\) is a unit.

We may therefore replace
\(\theta^{-1}\) by 1 modulo the maximal ideal. For each \(u\)
first sum over all lifts using the identity, then remove the
one lift \(wb\) with \(wb\equiv u\bmod f\). The reduced numerator is
\[
    \bigl(w^k-\phi^{-1}(w)\bigr)
       \sum_{u\in(\Z/f)^\times}\phi^{-1}(u)L_k(z_f^u)=0
       \quad\bmod\mathfrak m .
\]
This proves the smallness needed for the integral cusp correction.
The separator of unit cusp labels constructed above has weight at
most \(8\varphi(fw)\). Completing its weight, summing with the
nebentypus and dividing by twice a unit gives the integral form
with constant one. The subtraction and simultaneous Hecke lift in
the proof of the Eigenform in the rank-four construction now apply for
\(k\ge 8\varphi(fw)+10\) of the indicated parity. In particular
the lower cutoff can be fixed before \(m\) varies on its progression.
The constants, coefficient fields and levels in this calculation
are taken for the indicated characters.

\proofstep{The ordinary Galois attachment.}
Their Galois attachments determine the marked points.

In either row let \(\rho\) be the ordinary attachment of the
resulting normalized newform of weight \(k\) and level \(fw\).
Set \(\Phi=\epsilon^{k-1}\phi\). With the arithmetic Frobenius
convention for the newform we have
\[
  \det\rho=\Phi\theta,\qquad
  \bar\rho^{\,\mathrm{ss}}=1\oplus\bar\Phi,\qquad
  \rho|_{I_p}\simeq1\oplus\phi\theta\quad(p\ne\ell).
\]
The last assertion follows from equality of the conductor with the
central character conductor at each ramified prime, as in the
prime-index eigenform argument. On \(D_\ell\) the ordinary flag has
successive characters \(\Phi\theta\nu^{-1},\nu\), where \(\nu\)
is unramified, with Frobenius the unit root of the Hecke polynomial.
The congruence \(a_\ell\equiv1\) makes \(\bar\nu=1\) on the
whole decomposition group. The inertial exponents of the flag
are \(k-1,0\). We use the absolute irreducibility and local
compatibility of the cuspidal newform attachment.
\cite{deligne-galois,carayol-local,hida-ordinary-galois}

\end{proof}

\begin{lemma}[Irreducibility of the two restricted seeds]
\label{lem:marked-newform-restriction}
For the two newforms of Lemma~\ref{lem:ordinary-newform-seeds}, write \(\rho_1,\rho_2\) for their attachments in the order of
\eqref{eq:marked-newform-choices}. Their irreducibility ranges are
\[
 \begin{aligned}
 \mathbb Q\subseteq K'\subseteq L&\quad\Longrightarrow\quad
                 \rho_1|_{G_{K'}}\text{ is irreducible},\\
 H\le G_{\mathbb Q}\text{ open}&\quad\Longrightarrow\quad
                 \rho_2|_H\text{ is irreducible}.
 \end{aligned}
\]
\end{lemma}
\begin{proof}
Suppose an open normal subgroup \(H\) has an invariant line for
\(\rho\). Its orbit under \(G_\Q\) consists of \(H\)-invariant
lines and contains at least two, by irreducibility over \(\Q\).
On \(H\cap I_\ell\) the two ordinary exponents are distinct and
the image of \(\epsilon\) is infinite. Some element thus has two
distinct eigenvalues. Its two eigenlines exhaust the orbit.
Their common stabilizer \(D\) has index two and contains \(H\).
For every \(s\notin D\), the matrix interchanges the lines and
has trace zero. The residual trace identity gives
\(\bar\Phi(s)=-1\). Applying the same identity to \(sh\) for
\(h\in D\) gives \(\bar\Phi|_D=1\). A reducible restriction therefore satisfies
\[
                  \bar\Phi|_H=1,\qquad \bar\Phi^2=1
                       \ \text{on }G_\Q .
\]

In the first row \(\bar\Phi|_{G_L}=\bar\Psi^2\), so the first
consequence at \(H=G_L\) contradicts the marker \(g_*\).
In the second row \(\bar\Phi|_{G_L}=\bar\Psi\), and the second
consequence contradicts the same marker. Passing to the normal
core of any open subgroup shows that the second attachment stays
irreducible on every open subgroup. The first assertion suffices
for every intermediate field between \(\Q\) and \(L\).

\end{proof}

\subsection{The marked representations}

The two newforms now have the required irreducibility properties on restriction. To locate their points in the marked representation rings we must also identify their ordered local characters and the effect of the common scalar twist.

\begin{construction}
\[
\begin{array}{lll}
\text{seed}&\text{ordered data on restriction}&\text{support to be met}\\ \hline
r_2&\bar\Psi,\bar\Psi^{-1}\quad\text{(residue)}&J_2\\
r_3=\operatorname{Ad}^0\rho&\alpha/\beta,1,\beta/\alpha&
        J_3\ \text{and }J_0\text{ separately}
\end{array}
\]
\caption{The two seeds before the common level twist. The adjoint column uses a diagonal rank-two parameter with characters \(\alpha,\beta\).}\label{fig:ordinary-seed-ranks}
\end{construction}

\begin{proposition}[The two polarized marked seed representations]
\label{prop:marked-rank-two-three-seeds}
Take the two attachments in Lemma~\ref{lem:ordinary-newform-seeds}
with the operations in \eqref{eq:marked-newform-choices}. Their marked
rank-two twist and rank-three adjoint are polarized and have the
ordered ordinary flags. Before the common level twist, the rank-two
restriction has determinant one, residual labels
\((\bar\Psi,\bar\Psi^{-1})\) and ordinary inertia exponents
\((m,-m)\). On a diagonal rank-two parameter the adjoint construction gives
\[
 \operatorname{Ad}^0(\alpha\oplus\beta)
      \simeq(\alpha/\beta)\oplus1\oplus(\beta/\alpha),\qquad
 \det\operatorname{Ad}^0(\alpha\oplus\beta)=1 .
\] In rank two the restriction to
\(G_L\) is irreducible. In rank three the adjoint remains irreducible
on every open subgroup of \(G_L\), using
Lemma~\ref{lem:marked-newform-restriction}.
\end{proposition}

\begin{proof}
\proofstep{The rank-two polarization and its local flag.}
For rank two let \(\Xi=\epsilon^m\chi'\) and
\(\sigma=\rho\otimes\Xi^{-1}\) on \(G_\Q\).
The untwisted determinant in this row is \(\Xi^2\eta\theta\),
hence \(\det\sigma=\eta\theta\). The rational residual
characters of \(\sigma\) are \(\bar\Psi\bar\eta,\bar\Psi^{-1}\).
On restriction to \(G_L\) the characters \(\eta,\theta\) disappear.
The ordered residual labels, tame inertia labels away from \(\ell\),
and ordinary inertia exponents of \(r_2=\sigma|_{G_L}\) are
\[
  (\bar\Psi,\bar\Psi^{-1}),\qquad
  (\chi',(\chi')^{-1}),\qquad (m,-m),
\]
where characters in the display are restricted to the relevant
groups. The ordered reductions at \(D_\ell\) use \(\bar\nu=1\).

For the pairing put \(J=\left(\begin{smallmatrix}0&1\\-1&0\end{smallmatrix}\right)\)
and \(C=\sigma(c)\). Since \(C^2=1\) and \(\det C=-1\),
\(C^tJC=-J\) gives \((JC)^t=JC\). For \(s\in G_L\),
\(\det\sigma(s)=1\), so
\[
             \sigma(s)^t(JC)\sigma(csc^{-1})=JC .
\]
Thus the pairing of the seed before the level twist is symmetric
with multiplier one.


At the level use the specified restriction \(\kappa_2=\kappa|_{G_L}\)
of the algebraic Hecke character twist on \(K\). The two labelled infinity exponents and the polarization character of the twist satisfy
\[
 (h_\tau,h_{\bar\tau})=(-1,0)\ \text{up to order},\qquad\kappa\kappa^c=\epsilon^{-1}.
\]
Its inertia away from \(\ell\) has bounded prime-to-\(\ell\) order and is trivial away from the discriminant. Here is the construction. Put
\(D=|\operatorname{disc}K|\ge3\), so \(\ell\nmid D\), and fix an
embedding \(\tau:K\hookrightarrow\C\). On the subgroup of
\((\mathcal O_K/D)^\times\) generated by rational units modulo
\(D\) and global units, prescribe \(\eta\) on the former and
\(\tau\) on the latter. Their intersection comes from \(1,-1\).
Indeed the two exceptional unit groups are the fourth and sixth
roots of unity, whose nonrational elements have coefficient
\(\pm1\) on the second integral basis element.

The sign prescriptions on their common element agree because
\[
 \eta(-1)=-1=\tau(-1).
\] The character has order
prime to \(\ell\), since \(\ell\) is odd and prime to \(D\).
It factors through the quotient by the \(\ell\)-Sylow subgroup
and extends there to a character \(e\), by extension of characters
of finite abelian groups. On ideles with unit finite components take
\(x_\infty^{-1} e(x_f\bmod D)\), using \(\tau\) for the archimedean
coordinate. This character is trivial on diagonal global units.
It extends trivially over \(K^\times\), then to all ideles by
divisibility of \(\C^\times\). The original subgroup of ideles
is open, giving continuity.

At an inert prime \((p)\) with \(p\nmid\ell D\), divide by the diagonal idele \(p\). The remaining value on a uniformizer idele is
\[
 p\,\eta(p)^{-1}=p(-1)^{-1}=-p .
\]
Thus the associated \(\ell\)-adic character satisfies
\(\kappa(t^2)=-\epsilon(t)^{-1}\) at an arithmetic Frobenius
\(t\) at such a prime, with the inverse dictated by the chosen Artin reciprocity convention. Chebotarev and continuity give the identity on
the entire nontrivial coset of \(G_K\). For \(s\in G_K\) let \(t=sc\). Since \(c^2=1\),
\[
 t^2=s(csc^{-1}),\qquad
 \epsilon(t)=-\epsilon(s),\qquad
 \kappa(s)\kappa(csc^{-1})=\kappa(t^2)=\epsilon(s)^{-1}.
\]
The two twists in the conjugate pairing therefore contribute
the multiplier \(\epsilon^{-1}\).
Base change the newform with its
algebraic twist \(\Xi^{-1}\) to \(K\), twist there by the character
of \(\kappa\), and base change to \(L\) through a cyclic tower.

At a cyclic step loss of cuspidality would give a nontrivial
self-twist by a character \(\nu\) trivial on the upper field.
Local Galois compatibility at the unramified Frobenii and
Chebotarev give an invertible intertwiner satisfying
\[
       A\rho(g)A^{-1}=\nu(g)\rho(g).
\]
For \(g\) in the upper Galois group it commutes with the
absolutely irreducible restriction, so \(A\) is scalar. Substituting \(A=cI\), with \(c\ne0\), in the intertwining identity gives 
\[
 \rho(g)=\nu(g)\rho(g)\quad\Longrightarrow\quad\nu(g)=1\qquad\text{for every }g,
\]
 since \(\rho(g)\) is invertible. This contradicts the choice of twist. The finite-place operations on
parameters are restriction and multiplication by a character.
The two distinct algebraic weights at each embedding stay distinct
under their common shift. The pairing and unramified comparison
give conjugate self-duality of the regular algebraic cusp in the chosen algebraic normalization, by strong multiplicity one.
\cite{arthur-clozel,clozel-rajan}

\proofstep{The adjoint seed in rank three.}
For rank three take the second newform. Its adjoint action on the trace-zero endomorphisms is
\[
 r_3(g)X=\rho(g)X\rho(g)^{-1},\qquad g\in G_L,\quad X\in\operatorname{End}^0(V_\rho).
\] Equivalently, this representation is \(\operatorname{Sym}^2(\rho)\otimes(\det\rho)^{-1}\). The identification
sends \(vw\) to the endomorphism
\(x\mapsto v(w\wedge x)+w(v\wedge x)\) valued in
\(V\otimes\det V\). The formula is equivariant. In a basis \(e_1,e_2\), identify
\(e_1\wedge e_2\) with \(1\). The three symmetric tensors map to
\[
 e_1^2\longmapsto 2E_{12},\qquad
 e_1e_2\longmapsto E_{22}-E_{11},\qquad
 e_2^2\longmapsto-2E_{21}.
\]
These span the trace-zero matrices since \(2\) is a unit. If the two eigenvalues on an element
are \(a,b\), the three adjoint eigenvalues are \(a/b,1,b/a\),
including multiplicities for a Jordan matrix. The stable
trace-zero lattice of an integral model therefore has the
prescribed residual characteristic polynomials.

In an upper
triangular basis for the ordinary flag, the upper root line and
the trace-zero Borel give an invariant flag with those ordered
ratios of the diagonal characters. Applying this also to the
diagonal tame inertia groups gives, respectively,
\[
 (\bar\Psi,1,\bar\Psi^{-1}),\qquad
 (\chi',1,(\chi')^{-1}),\qquad (m,0,-m).
\]
Let \(A_c=\operatorname{Ad}(\rho(c))\) on \(\mathfrak{sl}_2\).
The trace form is nondegenerate and invariant and \(A_c\)
is an isometry with \(A_c^2=1\). Thus
\(P(X,Y)=\operatorname{tr}(X A_cY)\) is symmetric.
Since \(A_c r_3(cgc^{-1})=r_3(g) A_c\), invariance gives
\(P(r_3(g)X,r_3(cgc^{-1})Y)=P(X,Y)\).
For \(r_3'=\epsilon^{-1}r_3\) on the marked level, the
same calculation gives
\[
 P(r_3'(g)X,r_3'(cgc^{-1})Y)=\epsilon(g)^{-2}P(X,Y).
\]

\proofstep{Irreducibility on every open subgroup.}
To prove absolute irreducibility of \(r_3\) on every open subgroup,
suppose a nonzero proper subspace is stable. Its dimension is one
or two, and in the latter case its trace-orthogonal line is stable.
Let \(0\ne X\in\mathfrak{sl}_2\) span the stable line.
For each \(g\) of the open subgroup there is \(a_g\ne0\)
with \(\rho(g)X\rho(g)^{-1}=a_gX\). Hence for an eigenvector
\(v\) of \(X\) of eigenvalue \(\lambda\),
\[
                   X\rho(g)v=a_g^{-1}\lambda\,\rho(g)v .
\]
The set of eigenlines is therefore invariant. It has one
or two elements: a scalar with trace zero in characteristic
zero is zero, and \(X\ne0\). The stabilizer of one eigenline
is closed by continuity and has finite index at most two,
hence is open. This contradicts irreducibility of the
original attachment on every open subgroup.

\paragraph{The rank-three cusp and its local conditions.}
Base change the second newform to \(L\) by a cyclic tower as
above. At each step a character self-twist obstructing cuspidality
would intertwine two representations identical on the next
Galois group, where the intertwiner must be scalar. The same
argument excludes any nontrivial self-twist of the resulting cusp.
Taking central characters of the rank-two self-twist gives
\[
 \pi\simeq\pi\otimes\chi\quad\Longrightarrow\quad \omega_\pi=\omega_\pi\chi^2\quad\Longrightarrow\quad\chi^2=1 .
\] Its class field character has open kernel.
On that kernel the Galois intertwiner commutes with
the absolutely irreducible restriction, hence is scalar
and forces the character to be trivial on the whole group.
Apply Gelbart--Jacquet's adjoint lift to the unitarization of
this cusp. Their cuspidality criterion now applies. The adjoint has trivial central character.

On each toral
rank-two parameter its three lines are
\[
                  \alpha/\beta,\quad 1,\quad\beta/\alpha .
\]
The two infinity weights differ by \(m\), so these three weights
are regular. In the algebraic normalization the attached
representation is \(\epsilon^{-1}r_3\). Indeed the common
rank-two scalar cancels in the ratios, and the rank-three
factor \(|\det|^{(1-3)/2}\) multiplies each parameter at a
geometric Frobenius of norm \(q\) by \(q\). This is
\(\epsilon^{-1}\) there. Equality at unramified primes identifies
the semisimple attachment by Chebotarev. The self-duality of
the adjoint and its invariance under conjugation of \(L\)
give conjugate self-duality of the cusp. \cite{gelbart-jacquet}

\proofstep{Local conditions on the attached cusps.}
We check the local size condition on the actual attached cusps
in both ranks. For \(p\nmid fw\), let \(\alpha_p,\beta_p\) be the parameters after unitary normalization. Purity for a holomorphic newform of weight \(k\ge2\) gives
\[
 |\alpha_p|=|\beta_p|=1 .
\] This statement about
complex Satake parameters includes \(p=\ell\). At \(p\mid fw\)
the conductor argument gives a principal series with distinct
unit restrictions of its two characters and zero monodromy.
If the inducing characters are \(\chi_1,\chi_2\),
unitarity identifies the multisets \(\{\chi_1,\chi_2\}\)
and \(\{\bar\chi_1^{-1},\bar\chi_2^{-1}\}\).

On units the values of \(\chi_i\) are roots of unity, so
\[
 |\chi_i(u)|=1\quad\Longrightarrow\quad\overline{\chi_i(u)}^{-1}=\chi_i(u)\qquad(u\text{ a unit}).
\] The two restrictions
to units are distinct, so the permutation identifying the two
multisets fixes both indices:
\[
 \chi_i=\bar\chi_i^{-1},\qquad
               |\chi_i(\varpi_v)|=1\quad (i=1,2).
\]
This labels the two modulus-one uniformizer values without
interchanging the two inertial characters. Restricting Weil parameters
in cyclic base change preserves these properties of absolute
values and zero monodromy. On rank three the three ratios
therefore have absolute value one on uniformizers everywhere.
On rank two an algebraic Hecke character has absolute value
a real power of the adelic norm, since the idele classes of norm
one are compact. Its twist gives a common real exponent \(t\)
for the entire cusp. The inducing characters on either cusp
have absolute value \(q_v^t\) on a uniformizer at every finite
place \(v\) for one \(t\) in the normalization in use.

In
particular their magnitudes are strictly between
\(q_v^{t-1/2}\) and \(q_v^{t+1/2}\). For any two of these inducing characters the strict bounds on magnitudes at a uniformizer give
\[
 q_v^{-1}<|\chi_i(\varpi_v)/\chi_j(\varpi_v)|<q_v.
\]
Their ratio is therefore different from \(|\cdot|^{\pm1}\), and normalized induction stays irreducible even if an inertia character becomes trivial over \(L\). This proves the toral condition at every
finite place. Strong multiplicity one gives the same condition
for every cusp with the specified unramified attachment.
\cite{deligne-weil,carayol-local,bernstein-zelevinsky,arthur-clozel,gelbart-jacquet}

At \(v\mid\ell\) the ordered Galois characters have distinct
algebraic exponents differing by multiples of \(m\) at every
embedding. After the common shift, write their exponents as \(h_i\).
At every embedding the cyclotomic ratio
\(\kappa_i/(\epsilon\kappa_j)\) has exponent
\[
                       h_i-h_j-1\in m\mathbb Z-1 .
\]
Here \(m\ge2\), so the exponent is nonzero even for \(i=j\).
Every positive power has nonzero exponent on the open
inertia subgroup killing the finite parts. This proves the potential
nonresonance of the flags in the seed condition.

\proofstep{Descent and the possible local twists.}
It remains to identify the local components of the adjoint up to the allowed twists.

The unitary spectral comparison gives another realization of the adjoint transfer. To compare it with the classical adjoint lift used here, twist the Galois representation of the second newform by an algebraic character
\(\lambda\) on \(G_L\) with
\(\lambda\lambda^c=(\epsilon^{m+1}\chi'\eta\theta)^{-1}\).
Thus the rank-two multiplier becomes \(\epsilon^{-1}\), and
the two local exponents still differ by \(m\). Twisting by a character multiplies every representing matrix by a scalar. The rank-two adjoint is therefore unchanged:
\[
 (\lambda(g)\rho(g))X(\lambda(g)\rho(g))^{-1}=\rho(g)X\rho(g)^{-1}.
\] In the unitary comparison fix one
prime above each bad split place. A single sufficiently deep
level there kills all the local smooth characters.

At the additional regular split Frobenii write \(a=\alpha/\beta\) for the input ratio. The two tests read respectively
\[
 (a,1,a^{-1})\quad\text{in rank three},\qquad a^2\quad\text{as the rank-two ratio}.
\]
The output levels and infinity types are fixed as the test
set grows. Their finite spectra therefore have a constituent
satisfying the tests simultaneously, by the finite intersection
property. In the rank-two alternative, put
\(u=\operatorname{tr}(\rho)^2/\det(\rho)-2\). The required
identity is \(u'=u^2-2\) on the regular open set after
Chebotarev. On this set it implies the character identity
\[
 \operatorname{tr}(\operatorname{Sym}^4\rho/\det(\rho)^2)+1
 =\operatorname{tr}(\operatorname{Ad}^0\rho)
    +\operatorname{tr}(\operatorname{Ad}^0\rho').
\]

To see the obstruction, take a coset of an open normal subgroup
inside the set. The identity component of the projective
Zariski closure of \(\rho\) is \(\mathrm{PGL}_2\). A proper
connected subgroup is solvable, hence fixes a line, contrary
to irreducibility on open subgroups. The five-dimensional
summand stays irreducible on the subgroup. Write \(A\) for the invertible matrix of the translating element on the five-dimensional summand. Its translated trace cannot vanish identically. Indeed vanishing on the subgroup would give, by the full matrix span,
\[
 \operatorname{tr}(AX)=0\ \text{for every }X\in M_5\quad\Longrightarrow\quad A=0,
\]
by perfectness of the trace pairing.

It cannot occur in a sum of coefficients of representations
of dimension at most three. This last assertion follows by
applying the simple block projector in the semisimple quotient
of the group algebra image on the direct sum, or by taking
the linear dual of its matrix coefficient spaces on constituent
blocks. The translated traces on reducible summands may be
computed on a filtration stable under the translating element,
by first semisimplifying as representations of the whole group.
Their restrictions to the normal subgroup are then semisimple.
In the remaining rank-three case the polarization relates the traces at \(g\) and \(g^{-1}\). Its
open-set twist test identifies the attachment with the
cyclotomic twist of the adjoint. Varying the chosen place and
using strong multiplicity one gives the local identifications
at all bad places.

\end{proof}

\subsection{Admission at the three levels}

The marked representations supply the ordered characters and determinants. To start the component comparison, their evaluations must factor through the Hecke quotients for the prescribed levels, including each of the two rank-three ideals.

Let \(A_n\) be the representation ring of the marked frame, \(B_n\) its degree-zero subalgebra and \(B_{0,n}\) the universal ordinary determinant-pair ring for the level. The determinant and the ordered characters of the scalar twist define the maps
\[
 B_{0,n}\xrightarrow{\ \Phi_n\ } B_n\lhook\joinrel\longrightarrow A_n,
\]
with \(\Phi_n\) surjective. In rank three take two levels,
each with one exceptional prime, the two primes distinct. Use
one coefficient extension containing the tame scalar data.
At rank two denote the Hecke ideal by \(J_2\), and at the two
rank-three levels by \(J_3,J_0\), with the corresponding determinant
rings and maps in rank three taken separately. The same
untwisted \(A_3,B_3\) and the same \(\kappa_3=\epsilon^{-1}\)
enter both maps. The choice of \(m\) for the adjoint is made
once for these two levels. The rank-two weight can be chosen
independently.

\begin{proposition}[Admission on all three ordinary supports]
\label{prop:ordinary-linked-seed-admission}
For each of \(J_2,J_3,J_0\), with its own quotient map
\(\Phi_n\), there is a prime \(P\) of the marked representation
ring for which
\[
 \Phi_n(J)\subseteq P\cap B_n,\qquad
 B_n/(P\cap B_n)\simeq C,\qquad
 \dim C=1+n[L^+:\mathbb Q].
\]
The rank-three seed is shared while the two rank-three quotient
maps remain separate.
\end{proposition}

\begin{proof}
\proofstep{The actual points of the three ordinary supports.}
On tame inertia away from \(\ell\) the prescribed residual
character on the frame has the reduction of \(\chi'\). Let \(m\) be the prime-to-\(\ell\) order of \(\chi'\). Since reduction is injective on \(\mu_m\),
\[
 \bar\chi'(g)=1\quad\Longrightarrow\quad\chi'(g)=1\qquad(g\text{ in tame inertia}).
\] The diagonal inertia
description shows that \(r_n\) kills the required tame kernel.
On the level chart apply the same test to the joint ordered
labels of \(\kappa_n r_n\). Their inertia orders divide the
product of the two fixed prime-to-\(\ell\) orders. The tame
kernel of this chart is again killed, including inertia at
the chosen exceptional prime, which avoids the conductors.

At \(\ell\) take the twisted flags of the constructed
representations. Their reductions have exactly the prescribed
order. These factorizations are continuous, since the maps of
the profinite groups are quotient maps. By the universal property of the determinant pair, these continuous data give a point over the coefficient integers \(O\):
\[
 x:B_{0,n}\longrightarrow O .
\]
The twisted cusp, its unramifiedness at the exceptional prime,
and its regular ordinary nonresonant flag put \(J\) in \(\ker x\)
by the classical point theorem for this Hecke level. This
reasoning applies to each of \(J_2,J_3,J_0\) and gives an
irreducible untwisted specialization in the frame in each case.

\proofstep{The dimension needed at the starting components.}
Each starting point now lies on the required Hecke support. Propagation also needs a component of sufficient dimension through each point; this is the remaining condition to check.

We need a component of large dimension. Choose a minimal prime \(Z\) over \(J\) lying in \(\ker x\), and pass to the generic field of this component:
\[
 J\subseteq Z\subseteq\ker x,\qquad C=B_{0,n}/Z,\qquad E=\overline{\operatorname{Frac}(C)} .
\]
The universal determinant gives a semisimple representation
\(R\) over \(E\). At \(x\), choose \(n^2\) group elements whose
images span the full matrix algebra of the seed. The determinant
of the matrix of traces of pairwise products is nonzero, by the
nondegenerate trace pairing. Those traces lift to the same traces
in \(C\), so the matrices span \(\operatorname{Mat}_n(E)\)
for \(R\) as well.

In particular \(R\) is absolutely irreducible.
The kernel of \(U=\kappa_n^{-1}R\) is the common equality set
\[
                \operatorname{tr}U(hg)=\operatorname{tr}U(h)
                        \qquad\text{for all }h .
\]
This follows by pairing \(U(g)-1\) against the full span. It
proves the kernel closed in the level Galois group, since
the traces come from the continuous determinant and scalar.
Conjugate duality makes the kernel invariant under the
involution. The same trace test after multiplying a putative
relation by all \(U(h)\) turns the ordinary product and
exceptional unramified determinant relations into matrix
relations.

If a remaining marked tame inertia element \(g\) has trivial
untwisted residual label, some fixed prime-to-\(\ell\) power
of \(g\) kills the scalar and is in the tame kernel of the level.
Thus for some positive integer \(a\) prime to \(\ell\), the test on the scalar and on the tame kernel gives
\[
 U(g)^a=1,\qquad(a,\ell)=1 .
\]
Each eigenvalue \(\alpha\) reduces to one in an integral
extension at a prime over the residual point, since the residual
polynomial is \((X-1)^n\). Factor \(\alpha^a-1\). The second factor reduces to \(a\ne0\) and hence is a unit:
\[
 0=(\alpha-1)\underbrace{(1+\alpha+\cdots+\alpha^{a-1})}_{\text{unit}}\quad\Longrightarrow\quad \alpha=1 .
\] Separability of
\(X^a-1\) gives \(U(g)=1\). Consequently \(U\) descends to the
frame with its polarization and its ordered flag relations.
Putting it in the marked frame normal form gives a map
\(\xi:A_n\to E\) whose restriction to \(B_n\) takes values in
\(C\).

Indeed the degree-zero matrix products and flag
characters in this normal form belong to the trace and
ordinary-character algebra, as in the frame construction.
Retwisting recovers the generic pair on \(B_{0,n}\), so by
universality the composite \(B_{0,n}\to B_n\to C\) is the
quotient map. For \(P=\ker\xi\) we have proved
\[
       \Phi_n(J)\subseteq P\cap B_n,\qquad
       B_n/(P\cap B_n)\simeq C,\qquad
       \dim C=1+n[L^+:\Q].
\]
For the dimension, \(B_{0,n}/J\) is finite and torsion-free
over the power series weight ring of relative dimension
\(n[L^+:\Q]\) over a DVR. A minimal prime is associated,
so it cannot contain a nonzero weight scalar. Its domain
quotient is finite integral over the weight ring, which gives
the displayed equality. We have a point of the required
dimension for each of the three actual ideals.
\end{proof}

\section{Taylor--Wiles tests on a chain of curve points}
\label{sec:marked-patching}
\hypertarget{source-ordinary-marked-patching}{}

To move between components of ordinary Hecke support, we shall compare a chain of
residual curve points at one marked level. Choosing Taylor--Wiles primes separately at
each point would leave us without a common comparison.
We work with the normalized representation over the
characteristic-\(\ell\) power-series ring and test its
off-diagonal entries and local flags. The tests give the
patching hypotheses used in the component comparison of
Appendix~\ref{sec:linked-components}. Lemma~\ref{hyp:marked-square-zero-data} fixes the curve representation, its cyclotomic kernel and its coefficient field before any test elements or auxiliary places are chosen.

\begin{established}[Integral representatives at a normalized ordinary curve]\label{hyp:marked-square-zero-data}
Use a fitted point of Proposition~\ref{prop:fitted-points-three-tests}
at one of the three marked levels. Let \(G\) denote the CM subgroup,
\(H=\ker(\epsilon|_G)\) its full cyclotomic kernel, \(c\) the chosen
complex conjugation and \(\sigma h=chc^{-1}\). Let \(A=K[[t]]\)
and \(E=K((t))\) at the normalized curve, with \(K\) finite of
odd characteristic \(\ell\). In the polarized coordinates the
curve representation satisfies \(r(\sigma h)=\mu(h)r(h)^{-\mathsf T}\),
where \(\mu\) is the level multiplier. The span \(S=\operatorname{span}_E r(H)\) and the off-diagonal module
\(M\) are
\[
 (S,M)=\begin{cases}
 (M_n(E),0),&\text{type I},\\
 (M_a(E)\oplus M_b(E),\text{the off-diagonal blocks}),
      &\text{type II},\quad a+b=n .
 \end{cases}
\]

Use the conjugation action on \(M\) and \(-\mathsf T\) at \(c\).
The subspace of \(H^1\) whose classes have an upper triangular
representative in each ordered ordinary flag and are unramified on
the prescribed inertia away from \(\ell\) has a finite basis
represented by integral cocycles \(\kappa_1,\ldots,\kappa_g\) with
\[
 \kappa_j(c)=0,\qquad
 \kappa_j(\sigma h)=-\kappa_j(h)^{\mathsf T}.
\]
One can choose the representatives and their local changes of
frame over a single denominator.
\end{established}
\begin{proof}
The two span formulas follow from the irreducible, pairwise
inequivalent constituents on \(H\) in
Proposition~\ref{prop:fitted-points-three-tests}. In type II, if \(V_1,V_2\) denote these constituents, their matrix
coefficient modules give
\[
 E[r(H)]=\operatorname{End}_E(V_1)\oplus\operatorname{End}_E(V_2),
 \qquad M=\operatorname{Hom}_E(V_1,V_2)\oplus\operatorname{Hom}_E(V_2,V_1).
\]

The marked
auxiliary place in Proposition~\ref{prop:separated-marked-host}
supplies an element of \(G\) with nontrivial residual cyclotomic
character. For the finitely ramified real and CM groups and their
open subgroups \(U\), class field theory gives
\(\#\operatorname{Hom}_{\rm cont}(U,\mathbf F_\ell)<\infty\):
only the finite class group, the units at finitely many finite
places and finitely many uniformizers contribute modulo \(\ell\).
After passing to an open subgroup acting trivially on the residual
lattice, the action and integral cocycles factor through its
maximal pro-\(\ell\) quotient. 
For that quotient \(P=U^{(\ell)}\), finite generation follows from
\[
      \dim_{\mathbf F_\ell}P/\overline{P^\ell[P,P]}
       =\dim_{\mathbf F_\ell}\operatorname{Hom}_{\rm cont}(U,\mathbf F_\ell)<\infty .
\] Evaluating cocycles on
these generators and a finite set of coset representatives
bounds the dimension of the full cohomology, hence of the
subspace cut out by the local conditions.

For any cocycle \(z\), the identity \(c^2=1\) gives
\(c z(c)=-z(c)\). Subtract the coboundary of \(-z(c)/2\) to
obtain a representative vanishing at \(c\). The cocycle identity
at \(chc^{-1}\) then gives the transpose formula. Its continuous
image on the compact Galois group is bounded. Choose such
representatives for a basis, together with the finitely many
local changes permitted by the definition of the ordinary
subspace. A common power of \(t\) clears their denominators
without changing independence over \(E\). Boundaries between
normalized representatives have value zero at \(c\), so their
defining matrices satisfy \(Y=-Y^{\mathsf T}\).
\end{proof}

\subsection{The paired determinant test}
\label{subsec:det-test}

Begin at one curve point. The determinant must detect the cocycles together with the
underlying representation, so over \(E\) we form the finite algebra
\[
 B=S\oplus M^g,\qquad
 (x,p)(y,q)=(xy,xq+py).
\]
The second summand has square zero. The homomorphism
\[
 R(h)=\bigl(r(h),(\kappa_j(h)r(h))_j\bigr)
\]
takes values in \(B^\times\) by the cocycle rule. Entrywise transpose
on each component is an anti-involution \(\tau\) of \(B\), and
\(R(\sigma h)=\mu(h)\tau(R(h))^{-1}\). All these formulas are defined
on the evident free \(A\)-lattices.

\begin{lemma}[Spanning the marked square-zero algebra]\label{lem:marked-square-zero-span}
For the normalized cocycles and block decomposition of
Lemma~\ref{hyp:marked-square-zero-data}, one has
\[
   B=\operatorname{span}_E\{R(h):h\in H\}.
\]
\end{lemma}

\begin{proof}
\proofstep{Projection onto the two radical directions.}
The span of \(R(H)\) is a subalgebra
mapping onto \(S\). Its intersection \(W\) with \(M^g\) is an
\(S\)-bimodule, since the ideal has square zero. In the one-block
case there is no second summand. In the two-block case write
\(M=M_{12}\oplus M_{21}\). These are distinct simple split
\(S\)-bimodules. The central idempotents separate directions. The matrix units in each block then recover all of a submodule from its coefficient subspace:
\[
 W_{ij}=M_{ij}\otimes_E U_{ij}\subseteq M_{ij}\otimes_E E^g,\qquad U_{ij}\subseteq E^g .
\] If \(W\) is proper, take a nonzero row \(d=(d_j)\) annihilating
one of the \(U_{ij}\). For a lift \((x,(p_j))\) in the span of
\(R(H)\), define \(D(x)=\sum_j d_j p_j^{ij}\). Changing the lift
adds an element of \(M_{ij}\otimes U_{ij}\), which this row kills.
The square-zero multiplication rule gives
\[
                  D(xy)=xD(y)+D(x)y.
\]
Thus \(D\) is an \(E\)-linear derivation into the \(ij\)-corner.
\proofstep{The coboundary forced by a missing direction.}
For \(e_i\) the central idempotent of
the target block, put \(Y_{ij}=D(e_i)\). On the \(ij\)-corner the left action of \(e_i\) is one
and its right action is zero. Differentiating \(e_i x=xe_i\)
therefore gives
\[
                  Y_{ij}x + D(x) = xY_{ij}.
\] It follows
that the \(ij\)-part of \(\sum_j d_j\kappa_j(h)\), for \(h\in H\),
is \(r(h)Y_{ij}r(h)^{-1}-Y_{ij}\). Polarization at \(\sigma h\)
gives the other part with matrix \(-Y_{ij}^{\mathsf T}\). The
difference from this skew coboundary vanishes on \(H\). Its
values on \(G\) are then \(H\)-invariant by normality, and an
off-diagonal matrix commuting with \(r(H)\) is zero by the
two block idempotents. At \(c\) the difference is zero by
normalization and skew symmetry. This contradicts independence
of the \(\kappa_j\) and proves the spanning assertion.

\end{proof}

\paragraph{The twisted dual and the trace test.}
To read the determinant test as a condition on a linear functional, identify \(B^\vee=\operatorname{Hom}_E(B,E)\) with pairs
\(z=(\phi,(\lambda_j)_j)\) by the perfect pairing
\[
 \langle z,(x,p)\rangle=\operatorname{tr}(\phi x)
     +\sum_j\operatorname{tr}(\lambda_jp_j),
 \qquad \phi\in S,\quad\lambda_j\in M .
\]
The dual bimodule acts by \((a z b)(d)=z(bda)\).
Thus the trace pairing exchanges the two off-diagonal
directions. In this pairing \(\phi\) has coefficient \(\eta\), and \(\lambda_j\) has coefficient
\(\varepsilon_j'\), paired with \(\varepsilon_j\) by
\(\varepsilon_i\varepsilon_j'=\delta_{ij}\eta\).
The action on this dual is conjugation multiplied by the scalar twist:
\[
 h\cdot z=\chi(h)R(h)zR(h)^{-1}.
\]

In evaluating this action, conjugate the argument \((x,p)\)
of the functional by \(R(h)^{-1}\). Factoring
\(R(h)=(1,\kappa(h))(r(h),0)\) shows that the radical
part of the argument acquires the commutator
\(x\kappa_j(h)-\kappa_j(h)x\) before conjugation by
\(r(h)^{-1}\). Moving it across the trace gives the
commutator in the first coordinate below. The second coordinates are
\(\chi(h)r(h)\lambda_jr(h)^{-1}\). Its first coordinate is
\[
 \chi(h)\left(r(h)\phi r(h)^{-1}
       +\sum_j[\kappa_j(h),r(h)\lambda_jr(h)^{-1}]\right).
\]
Here \([P,Q]=PQ-QP\). Conjugation at \(c\) acts as
\(s\tau^*\), where \(\tau^*\) is transpose under the pairing
and \(s=-\chi(c)\). In particular \(s^2=1\).
For a twisted cocycle with \(z(c)=0\), the test at \(h\in H\) is
\[
 \mathcal L_h(z)=\langle z(h),R(h)\rangle
  =\operatorname{tr}(\phi(h)r(h))
     +\sum_j\operatorname{tr}(\lambda_j(h)\kappa_j(h)r(h)).
 \tag{10}\label{eq:local-patch-10}
\]
This is the \(\eta\)-coefficient of the infinitesimal trace.

\begin{lemma}[Paired trace detection]\label{lem:paired-trace-detection}
With the curve and block data of Lemma~\ref{hyp:marked-square-zero-data},
let \(z\) be a cocycle with values in the \(\chi\)-twisted dual and
\(z(c)=0\). The detection implication over \(E\) is
\[
  (\forall h\in H,\ \mathcal L_h(z)=0)\quad\Longrightarrow\quad [z]=0
\]
in first cohomology with these coefficients. For an integral cocycle
one can clear the boundary matrix by a power of \(t\).
\end{lemma}

\begin{proof}
\proofstep{From the vanishing trace to a derivation.}
Suppose all the tests vanish. For \(h\in H\), set
\(D_h=z(h)R(h)\in B^\vee\). Then
\(D_{hh'}=D_hR(h')+R(h)D_{h'}\).
If \(\sum_i a_iR(h_i)=0\), evaluate \(\sum_i a_i D_{h_i h'}\)
at the identity of \(B\). The vanishing of the tests and the
relation imply
\[
       \left\langle\sum_i a_i D_{h_i},R(h')\right\rangle=0 .
\]
Lemma~\ref{lem:marked-square-zero-span} gives
\(\sum_i a_i D_{h_i}=0\).
There is consequently a derivation \(D:B\to B^\vee\)
with \(D(R(h))=D_h\).

\proofstep{The action of the cyclotomic quotient.}
For \(v\in G\), put \(u_v=R(v)\) and \(w_v=D(u_v)u_v^{-1}\).
Write \(\delta_Y(x)=xY-Yx\). Applying the derivation rule to
\(u_v x u_v^{-1}\) gives
\[
 u_v D(u_v^{-1}x u_v)u_v^{-1}=D(x)+\delta_{w_v}(x).
\]
The cocycle rule on \(v h v^{-1}\), with \(h\in H\), gives
\[
 D(x)=\chi(v)u_v D(u_v^{-1}x u_v)u_v^{-1}
                                  -\delta_{z(v)}(x)
 \quad (x=R(v h v^{-1})).
\]
By linearity this holds on \(B\). Choosing \(v\) with
\(\chi(v)\ne1\) shows that \(D\) is inner. Subtract its twisted
coboundary from \(z\). The resulting cocycle vanishes on \(H\),
and takes values in the invariant subspace of \(B^\vee\) under
\(R(H)\). Such a functional satisfies \(xY=Yx\) for every
\(x\in B\) by the span.

On this subspace the action of \(G\)
is therefore the scalar \(\chi\). Since \(G/H\) is abelian,
the cocycle \(z_0\) on it satisfies
\[
 (\chi(v)-1)z_0(v')=(\chi(v')-1)z_0(v).
\]
The same choice of \(v\) makes \(z_0\) a coboundary. In the final
boundary matrix replace \(Y\) by \((Y+s\tau^*(Y))/2\).
The polarized cocycle rule and the formula for \(R(\sigma h)\)
imply that the new matrix gives the same boundary on \(G\),
and it is fixed by the action at \(c\). Thus it also gives the
boundary on the full real group. This proves the test after
inverting \(t\), or, integrally, after multiplying the cocycle
by a power of \(t\) to clear the matrix denominators.

\end{proof}

\begin{construction}
\[
\begin{array}{rl}
1&\text{Fix the finite chain of marked curve points.}\\
2&\text{Choose one finite set of tests in the cyclotomic kernel.}\\
3&\text{Fix the coefficient depth }m.\\
4&\text{At each point test at the base and at } N\gg_m1.\\ \hline
 &\#\{\text{tested dual classes at the point}\}\le |K|^{C},
      \quad C\text{ independent of }m .
\end{array}
\]
\caption{Order of choices for the paired determinant tests. The finite test set works across the fixed chain, while the bound on size may depend on the point.}\label{fig:marked-test-quantifiers}
\end{construction}

\subsection{Common tests and the finite-depth dual bound}

At a single point we have a trace test. Its use along a chain requires two kinds of
uniformity. The tests must be chosen in common, and their error at each point must
stay bounded as the torsion depth grows.

\begin{proposition}[Uniform bounds for the tested dual classes]
\label{prop:marked-uniform-dual-tests}
For the fixed finite chain, a common finite choice of tests in the
cyclotomic kernel detects the stated paired dual classes. For each
point there is a constant \(C\) independent of \(m\) so that the tested dual cohomology group \(H_{\rm test}\) at depth \(m\) satisfies 
\[
 |H_{\rm test}|\le|K|^C\qquad\text{at the base and at every auxiliary level }N\gg_m1.
\]
 The finite set
of tests is chosen before the auxiliary levels.
\end{proposition}

\begin{proof}
\proofstep{A fixed generating set for the finite-depth cocycles.}
For the standard adjoint module, the odd module and the paired
dual, the action on an open normal subgroup has pro-\(\ell\)
image. For these finite-free \(A\)-modules one can take the
kernel of the reduction of the action. At all finite
truncations the kernel of reduction consists of successive
additive matrix groups. At depth \(m\) write \(V_m\) for the additive coefficient module and \(\rho_m\) for its action. The cocycle equation is precisely the multiplication rule 
\[
 (z(xy),\rho_m(xy))=(z(x)+\rho_m(x)z(y),\rho_m(x)\rho_m(y)).
\]
 The image on the chosen subgroup therefore lies in an extension of two finite \(\ell\)-groups. The maximal pro-\(\ell\) quotient of the
open subgroup has finitely generated Frattini quotient by the
stated Galois finiteness, and hence is topologically finitely
generated. Adding lifts of the finite quotient gives a fixed
finite generating set on which all the cocycles are
determined, integrally and modulo every \(t^m\).

\proofstep{Finite test functionals on the common dense set.}
On a finite chain of points, regular semisimplicity on \(H\) is
an open dense condition for each point by the nice-point
hypotheses. Its openness follows by reducing the nonzero
discriminant at sufficiently high depth. The intersection is
dense, and by continuity a cocycle vanishing under \eqref{eq:local-patch-10} there
passes every test on \(H\). Evaluation at a fixed finite generating set \(g_1,\ldots,g_d\) identifies a cocycle by its values. For the coefficient space \(V\) at this point the injection is
\[
 Z^1(G,V)\hookrightarrow V^d,\qquad z\longmapsto(z(g_1),\ldots,z(g_d)).
\] Finitely many of the tests span the same space of
linear forms as the tests on the entire intersection. Take
the union of these finite sets for all points, say
\(h_1,\ldots,h_q\).

Split the characteristic polynomials of these elements by
finite extensions of each curve field, using the integer
rings for the extensions. The support kernels are unchanged
because the ring maps are injective. The rings are complete
equicharacteristic DVRs with finite residue fields. They can
be enlarged again until the residue fields all have the same
cardinality, as follows. If their degrees over \(\mathbf F_\ell\) are \(d_1,\ldots,d_s\), take unramified extensions to the field of cardinality 
\[
 \ell^{\operatorname{lcm}(d_1,\ldots,d_s)}.
\]
 Each coefficient extension is finite
free over the preceding DVR.

For a finite free extension \(A\subset A'\) of these
DVRs, choose a basis \(b_i\) of \(A'\) over \(A\).
An \(A'\)-valued cochain is a unique sum \(z=\sum_i b_i z_i\)
with continuous \(A\)-valued cochains \(z_i\). For an \(A\)-linear
equation \(\mathcal T\) in the cocycle and the matrices imposing its local conditions,
\[
       \mathcal T(z)=\sum_i b_i \mathcal T(z_i).
\]
Hence the equations cut out the scalar extensions of their
kernels. The image of the boundary map is also the scalar
extension of its image, and so is the cohomology quotient.
After inverting the uniformizer this preserves the chosen odd
cohomology basis. In particular the new tests still give precisely
the coboundaries over the new fraction fields.

\proofstep{Frobenius approximation at every fixed depth.}
At depth \(N\) use the finite image of the product
of the current marking quotient, the CM sign, and
\(\epsilon_\ell\bmod\ell^N\). The image of \(h_j\) is
\[
            (h_j\bmod V_N,\ 1,\ 1),
\]
since \(h_j\) is in the positive full cyclotomic kernel.
Thus the simultaneous congruence is realized in that
finite quotient. Chebotarev for its inverse class gives
geometric Frobenius in the desired class. Excluding the
fixed level and previously chosen primes still leaves a
choice of distinct split places. The cyclotomic coordinate specifies the norm of every chosen place: 
\[
 Q_{j,N}\equiv1\pmod{\ell^N}.
\]
 The auxiliary
level retains the chosen decomposition groups and tame
inertia generators.

Passing to the open subgroup cutting
out the residual diagonal characters and the mod-\(\ell\)
cyclotomic character, the point representations have
unipotent residual image. A finite group consisting of
unipotent matrices is an \(\ell\)-group by Cauchy's theorem.
Its congruence kernels and the cocycle semidirect images
are likewise \(\ell\)-groups at finite depth. This is the
factorization through the pro-\(\ell\) quotient used in
comparing the level kernels. Fix \(m\). The coefficient actions and base cocycles at precision \(t^m\) are recorded by finitely many maps to finite groups. If \(N_i\) is the open normal kernel for the \(i\)-th such map, they all factor through
\[
 G/\bigcap_i N_i\ \hookrightarrow\ \prod_i G/N_i .
\] Frobenius agreement thus holds eventually
for all of them at once.

\proofstep{The Smith bounds for the tested complex.}
Here is the integral size comparison. Use evaluation on the
finite generating set to write the entire cocycle and test
system as linear equations on \(A^d\), where \(d\) here denotes
the number of scalar generator coordinates. On a free
profinite group the tuples define cocycles by the
semidirect product universal property. The relations of its
quotient, the condition at \(c\) and \eqref{eq:local-patch-10} give rows over
\(A\). The submodule of the dual of \(A^d\) spanned by the
rows is finitely generated, so the same finite matrix
\(F\) gives the system at every truncation.

The matrix
\(\Delta\) of the coboundaries from the \(c\)-fixed subspace
has image in its kernel. That subspace commutes with
reduction since the eigenprojections of the involution
use division by two. Over the fraction field, passing the tests is equivalent to being a boundary. Thus 
\[
 \operatorname{im}(\Delta\otimes_A E)=\ker(F\otimes_A E).
\]
 If their common dimension is \(j\), Smith
normal form gives
\[
        |\ker F_m|\le |K|^{mj+C_1},\qquad
        |\operatorname{im}\Delta_m|\ge |K|^{mj-C_2}
\]
with constants given by valuations of the nonzero
elementary divisors. This bounds the size of the tested
cohomology by \(|K|^{C_1+C_2}\) for all \(m\).

\proofstep{Descent of the tested auxiliary classes.}
A dual class at an auxiliary level is tested if it
restricts to a coboundary on each new decomposition
group. In that case its value at tame inertia is zero,
because the coefficient action is pulled back from the
base level. Its cocycle descends to the base: the
auxiliary kernel is generated, after passing to the
pro-\(\ell\) quotient, by the normal closure of those tame
generators. On the auxiliary kernel the action on the coefficients is trivial. Thus a cocycle that vanishes on a generator vanishes on its conjugates by the cocycle rule, and then on the closed subgroup they generate by continuity:
\[
 z(ghg^{-1})=g\,z(h)\qquad(h\text{ in the kernel}).
\] Normalize
the cocycle at \(c\) by subtracting the coboundary of
\(-z(c)/2\).

At each new Frobenius the normalized value
is still a local coboundary, whose product trace against
the Frobenius matrix vanishes by cyclicity. At sufficiently
large \(N\) the agreement modulo \(t^m\) transports this
vanishing to \eqref{eq:local-patch-10} on the chosen \(h_j\). Boundaries descend
with it, so the bound for the tested cohomology applies
at the auxiliary level. The dual bound is uniform in the sense \(\exists C\,\forall m\,\forall N\gg_m 1\). The constant is
allowed to depend on the fixed curve point and finite tests.

\end{proof}

\subsection{The first-order bound}

We next apply the common tests to first-order deformations. Passing from normalized
cocycles to their quotient by boundaries is essential for the dimension count in the
completed local ring.

\begin{proposition}[First-order cohomology at a marked point]
\label{prop:marked-first-order-bound}
With the standard adjoint and chosen local conditions at a marked
point, let \(Z^1_{m,\mathcal L}\) be the normalized cocycles at depth \(m\) satisfying these conditions and \(B_m^1\) their normalized coboundaries. Then
\[
                       \bigl|Z^1_{m,\mathcal L}/B_m^1\bigr|\le |K|^{m(nd+nq)+C},
\]
where \(d=[L^+:\mathbb Q]\) and \(q\) is the number of auxiliary
places in this comparison. The constant is independent of \(m\)
and the bound holds on the shared auxiliary system for \(N\gg_m1\).
\end{proposition}

\begin{proof}
\proofstep{Local dimensions at the places above the coefficient prime.}
Let \(d=[L^+:\Q]\) again and let \(d_v\) be the local degree
at \(v\mid\ell\). Let \(d'_v\) be the rank of the pro-\(\ell\)
inertial units after removing their torsion. The local
abelianized quotient for the diagonal cocycles then has
rank \(d'_v+1\), including the unramified generator.
Here \(\sum_v d_v=\sum_v d'_v=d\), by local class field
theory on the chosen split places. Put
\[
 \nu=\binom{a}{2}+\binom{b}{2}+gab
\]
with \(a=n,b=0\) in type I. The primal first-order
coefficient module is \(B\) under conjugation by \(R\),
with \(-\tau\) at \(c\). Its rank is \(n+2\nu\).
At depth \(m\) the perfect pairing with the twisted dual is the composite 
\[
 (z,b)\longmapsto \langle z,b\rangle\longmapsto [t^{m-1}]\langle z,b\rangle\longmapsto \operatorname{Tr}_{K/\mathbf F_\ell}\!\left([t^{m-1}]\langle z,b\rangle\right).
\]

At \(v\mid\ell\), make the local odd coboundary changes
used in the choice of \(\kappa_j\) so that \(R\) is
upper triangular in the ordered local flag. The
local upper triangular coefficient subspace has
rank \(n+\nu\). The \(n\) diagonal cocycles kill the
local torsion prescribed by the level. Each of the
other positions, including the \(g\) copies for every
mixed-block position, is filtered by its distance
above the diagonal. For positions at a fixed diagonal distance, the higher-distance contributions vanish in the associated quotient. The action on the coordinate \((i,j)\) becomes
\[
 x_{ij}\longmapsto\psi_{v,i}(g)\,x_{ij}\,\psi_{v,j}(g)^{-1}\qquad(i<j).
\] The local
nonresonance test excludes the cyclotomic character
for these ratios. The finite module with this
character at depth \(m\) has order \(|K|^m\).

Tate local duality bounds its second cohomology by
\(|K|^{O(1)}\), since some ratio divided by the
cyclotomic character differs from one by a fixed
nonzero element of \(A\). The local Euler characteristic
formula and the formula for the size of coboundaries
then give \(|K|^{m(1+d_v)+O(1)}\) for its first
cocycles. Multiplying along the successive kernels gives
\[
       |Z_v|\le
       |K|^{m[n(d'_v+1)+\nu(1+d_v)]+C_v}.
\]
The image of \(Z_v\) in local cohomology quotients
by at least the coboundaries of the upper triangular
coefficient subspace. The upper triangular subspace has \(|K|^{m(n+\nu)}\) elements. Dividing by its coboundaries in the Greenberg--Wiles local term gives 
\[
 \frac{|L_v|}{|H_v^0|}\le \frac{|Z_v|}{|K|^{m(n+\nu)}}.
\]


\proofstep{The unramified and auxiliary factors.}
At each fixed non-\(\ell\) place of the unramified local
condition, the local term is one. At the new places
and at the distinguished test place the primal local
condition is full. Local Euler characteristic in
residue characteristic different from \(\ell\) makes
the local term the size of the local twisted-dual
invariants, and the dual local class must vanish.
At the auxiliary places the twist is trivial in
characteristic \(\ell\).

Write \(T\) for the regular Frobenius of the base curve. Its simple spectrum and the disjointness of the spectra in the two blocks give 
\[
 \dim_E\{X\in S:XT=TX\}=n,\qquad\{X\in M:XT=TX\}=0.
\]
 In the paired dual,
an invariant at the fraction field therefore has
\(\lambda_j=0\) for every \(j\) and at most \(n\)
free diagonal coordinates. For \(N\) sufficiently large as a function of \(m\), the finite-depth Frobenius equations agree with the fixed ones. The at most \(n\) free coordinates and the bounded Smith factors give the estimate
\[
 |K|^{mn}\,|K|^{C_j}=|K|^{mn+C_j} .
\]

\proofstep{Invariant spaces at the distinguished and real places.}
At the distinguished place the decomposition subgroup \(D_v\) must satisfy the intertwining test 
\[
 \bigl\{Y:\chi(h)r(h)Y=Yr(h)\text{ for every }h\in D_v\bigr\}=0.
\]
 This holds
over \(E\) by clearing denominators in \(Y\). Applying it
first to the off-diagonal projections and then to the
diagonal projection gives zero invariants for the paired
dual over \(E\), and a bounded finite-depth term.
For primal global invariants, commutation
with the span \(B\) forces centrality, while the
real action requires skew symmetry. The center consists of block scalars fixed by transpose. Since \(\ell\ne2\), the real skew condition gives 
\[
 Z\in Z(B),\quad Z=-Z^{\mathsf T}\quad\Longrightarrow\quad 2Z=0\quad\Longrightarrow\quad Z=0.
\]
 Hence the integral comparison contributes only bounded torsion at finite depth. At every real place
the primal invariant module has at least
\(\nu\) independent coordinates at each depth.
Indeed if \(c'\) is its conjugation, the matrix
\(R(c'c)\) is symmetric by the polarization
relation and \(\epsilon(c'c)=1\). Multiplication
by this matrix embeds the skew matrices in the
\(c'\)-fixed space.

\proofstep{The Greenberg--Wiles product.}
Greenberg--Wiles duality, with the tested dual bound,
gives the estimate in cohomology. The exponent of
\(|K|\) that multiplies \(m\) is
\[
  \sum_{v\mid\ell}
    \bigl(n(d'_v+1)+\nu(1+d_v)-(n+\nu)\bigr)
     +nq-\nu d=nd+nq .
\]
The primal cocycles normalized at \(c\) modulo their
normalized coboundaries give the same cohomology,
again by division by two. Thus the bound is
\(|K|^{m(nd+nq)+C}\) for all \(m\) and all sufficiently
large \(N\) at that \(m\). This estimate holds on the same auxiliary systems for every point in the chain.

\end{proof}

\subsection{The odd obstruction bound}
\label{subsec:odd-obstruction}

The first-order estimate counts directions. To bound the relations among them we must
also control the obstruction to lifting; it is the odd part that enters the two-block
calculation.

\begin{proposition}[Uniform control of the odd obstruction module]
\label{prop:marked-odd-obstruction-bound}
Use the fixed normalized point of Lemma~\ref{hyp:marked-square-zero-data}.
The module \(U_m\) is relative second cohomology of the odd matrix
module over \(K[[t]]/(t^m)\), with the local projections by the upper
triangular part at \(v\mid\ell\) and the unramified local condition
away from \(\ell\). A cocycle includes a trivialization of each
projected local class. There is a constant \(C\) independent of \(e,m\)
such that for every \(e\le m\)
\[
                  |U_m/t^e U_m|\le |K|^{ge+C}.
\]
In particular the quotient by the first \(g\) elementary-divisor
summands is killed by \(t^C\). Among any \(g+1\) classes there is
a combination with a unit coefficient whose multiple by \(t^C\)
vanishes.
\end{proposition}

\begin{proof}
\proofstep{The obstruction module at a fixed depth.}
Keep the fixed point and let \(A_e=A/(t^e)\).
Use the odd matrix module \(M_e\) over \(A_e\),
with conjugation by \(r\) and with \(-\mathsf T\)
at \(c\). At a place \(v\mid\ell\) the local
coefficient quotient \(Q_{v,e}\) is by the
upper triangular part in the ordered flag. Its
rank over \(A_e\) is \(ab\), one lower entry for each
pair of indices from different blocks. At the other
finite places the quotient is \(M_e\) itself and
the local condition is unramified.

In particular at a place above \(\ell\) the localization condition on the projected class \(\bar z\) is
\[
 \operatorname{res}_v[\bar z]=0\qquad(v\mid\ell).
\] Let \(S_e\)
denote the resulting first cohomology group with
these conditions. Let \(U_e\) denote second
cohomology with these local conditions.
Represent a class of \(U_e\) by a global
two-cocycle and a local one-cochain trivializing
its projection at every place. One quotients by the
changes from global one-cochains and local
one-cocycles with the allowed local restrictions.
Let \(r_e^2\) denote restriction followed by projection
in second cohomology, and set
\[
 a_e\colon H^1(G^+,M_e)\longrightarrow
             \bigoplus_v H^1(G_v,Q_{v,e})/L_{v,e}.
\]
Then the obstruction group sits in
\[
 0\longrightarrow\operatorname{coker}a_e\longrightarrow U_e
 \longrightarrow\ker r^2_e\longrightarrow0.
 \tag{11}\label{eq:local-patch-11}
\]

In \eqref{eq:local-patch-11} the local indices are precisely the primes above \(\ell\)
and the fixed tame primes dividing the level cutoff away from \(\ell\).
The distinguished place and the new Taylor--Wiles places belong to the
global ramification set. Only the new inertia evaluations enter \eqref{eq:local-patch-12}
in comparing with the base. Here \(G^+\) denotes the real-field group
\(G_{L^+,S'}\), the Galois group of the maximal extension unramified
outside the level set \(S'\) and the archimedean places. The subgroup \(L_{v,e}\) is the group unramified on the prescribed tame inertia kernel away from
\(\ell\) and zero above \(\ell\). The sequence follows
by first sending the cocycle to its global class.
If that class is zero one subtracts a global
coboundary and is left with the indicated local
classes modulo global restrictions.

At the base level,
\[
                         |S_e|\le |K|^{ge+C_0}.
 \tag{12}\label{eq:local-patch-12}
\]
Use evaluations on finitely many generators as for
the paired dual. For each local condition add the
coordinates of a local boundary matrix at the
places above \(\ell\), and require zero restrictions
on the prescribed inertias away from \(\ell\), where
the odd module is unramified. One obtains a
finite matrix of linear equations by taking the
row span, including all relations and local
equations. If \(F\) denotes this matrix, its nonzero Smith factors contribute a bounded correction: 
\[
 \log_{|K|}|\ker(F\bmod t^e)|=e\dim_E\ker(F\otimes_A E)+O_F(1).
\]
 Project to the generator
coordinates and divide by global boundaries.
The kernel of projection is given by the local
invariance equations on the added matrices,
so its count also has the corresponding
dimension estimate. This gives \eqref{eq:local-patch-12}, since the
resulting fraction-field quotient has dimension
\(g\).

At the auxiliary level, evaluate each class on
the new tame generators. Boundaries vanish there.
The Frobenius relation, the trivial inertia action
and \(Q_{j,N}\equiv1\pmod\ell\) make the
tame value fixed by conjugation under the
auxiliary Frobenius. Fix \(m\) and take \(N\) large enough for that depth. For every \(e\le m\) the tame value \(Z\) then obeys the same commutation equation as at the test element: 
\[
 Zr(h_j)=r(h_j)Z\pmod {t^e}.
\]
 The commutant of
\(r(h_j)\) has no off-diagonal part over \(E\)
because the full eigenvalue list is distinct.
By Smith normal form the cardinality of this evaluation image satisfies the bound
\[
 |\operatorname{im}({\rm ev})|\le |K|^{C_{\rm aux}}\qquad(e\le m).
\] In the kernel
the cocycle is zero on the new inertia and
descends to the base, with its old local
conditions. This proves \eqref{eq:local-patch-12} at the auxiliary
level with the additional constant \(C_{\rm aux}\).

Two more invariant groups have uniformly bounded
order. For \(H^0(G_v,Q_{v,e})\) above \(\ell\),
filter the lower entries so that the induced
characters are the lower ordered ratios. On a lower-entry graded piece, an invariant \(x_{ij}\) satisfies 
\[
 (\psi_{v,i}(h)\psi_{v,j}(h)^{-1}-1)x_{ij}=0\quad\text{for every }h.
\]
 Distinctness of the ordinary characters forces \(x_{ij}=0\) over the fraction field. Successive application to the filtration kills every invariant. The
same finite-row and Smith argument bounds
the finite-depth invariants. For the
cyclotomic twisted dual of \(M_e\), its invariants
on \(H\) vanish over \(E\). The twist is trivial
there, and conjugation invariants on the odd
matrices vanish by the block idempotents in
the span of \(r(H)\). The matrix trace pairing
identifies the odd module with its \(A_e\)-dual.

For \(X,Y\) in the odd matrix module the induced finite-field
pairing is
\[
 (X,Y)\longmapsto
 \operatorname{Tr}_{K/\mathbf F_\ell}\!\left([t^{e-1}]
                    \operatorname{tr}(XY)\right).
\]
Indeed a coefficient at degree \(j<e\) in a nonzero trace pairing
is detected by multiplication of the other vector by \(t^{e-1-j}\)
and a scalar of \(K\), using nondegeneracy of the field trace.
It gives a fixed
bound \(|K|^{C_d}\) for the global Tate dual
invariants. This remains valid at all auxiliary
levels because the coefficient action is
pulled back by a surjection. In particular
\[
          |\operatorname{coker} r^2_e|\le |K|^{C_d}.
 \tag{13}\label{eq:local-patch-13}
\]
Indeed take first the arithmetic level defining the global and local
groups. The finite set \(S'\) contains every place over \(\ell\) and
\(\ell\) is odd.

For the finite discrete \(\ell\)-primary coefficients
here global \(H^3\) vanishes, even in the presence of real places.
Local \(H^3\) at the finite places vanishes as well. Thus for any short exact sequence of these finite discrete coefficient modules, the end of the cohomology sequence is
\[
 H^2(G,A)\longrightarrow H^2(G,B)\longrightarrow H^2(G,C)\longrightarrow H^3(G,A)=0
\]
for \(0\to A\to B\to C\to0\). This is the required right exactness.
Poitou--Tate bounds the cokernel of
global localization with full local coefficients
by the global Tate dual invariants. Drop the components outside the fixed indices and identify each retained local group with the decomposition group at the chosen split prime of the CM field. For any retained local quotient \(T\twoheadrightarrow T'\) with kernel \(T_0\), local cohomological dimension two gives the exact segment 
\[
 H^2(D_v,T)\longrightarrow H^2(D_v,T')\longrightarrow H^3(D_v,T_0)=0.
\]

The level map from the arithmetic group at \(V_N\) to the group
for this count is a continuous bijection by the level-map
surjectivity and the injectivity in the arithmetic group comparison. Compactness and Hausdorffness make its inverse
continuous. Pullback of cochains thus identifies the global cohomology
and commutes with the fixed local restrictions, the signed real
involutions and the cyclotomic twist.  They precede passage to
the relative pro-\(\ell\) quotients for the patch.

\proofstep{The global and local Euler factors.}
\proofstep{Bounding the auxiliary odd classes by duality.}
The global Euler characteristic gives
\[
 |H^1(G^+,M_e)|=
    |H^0(G^+,M_e)|\,|H^2(G^+,M_e)|\,|K|^{eabd}.
\]
Each real conjugation has \(ab\) fixed coordinates
on \(M_e\). In the normalized basis its action
is minus transpose. At another real conjugation
multiplication by the corresponding invertible
symmetric block-diagonal matrix identifies the
fixed space with the skew off-diagonal matrices.
The local Euler term at \(v\mid\ell\) for
\(Q_{v,e}\) is \(|K|^{eab d_v}\) in the numerator
of \(|H^1|/(|H^0||H^2|)\). At the other
places \(|H^1_{\rm ur}|=|H^0|\), by taking the prime-to-\(\ell\) inertia invariants
and evaluating on Frobenius, as explained below. Taking orders
in \eqref{eq:local-patch-11}, using \(\sum d_v=d\) and dropping
the global \(H^0\) denominator, gives
\[
 |U_e|\le |S_e|\,
    \prod_{v\mid\ell}|H^0(G_v,Q_{v,e})|\,
                   |\operatorname{coker} r^2_e|.
 \tag{14}\label{eq:local-patch-14}
\]
In particular the two Euler contributions cancel
even if the finite-depth local quotients have
nonzero invariants.

To control elementary divisors, use both coefficient
sequences for \(e\le m\),
\[
 0\to A_{m-e}\xrightarrow{t^e} A_m\to A_e\to0,
 \qquad
 0\to A_e\xrightarrow{t^{m-e}} A_m\to A_{m-e}\to0.
\]
They stay exact on the odd modules and their local
quotients, since the local upper triangular kernel
has a free complement in its local frame. Write
\(i:U_{m-e}\to U_m\) for the first inclusion map
and \(r:U_m\to U_{m-e}\) for the second reduction
map. Multiplication and the first reduction satisfy 
\[
 i\circ r=t^e\operatorname{id}_{U_m},\qquad \ker(U_m\longrightarrow U_e)\subseteq i(U_{m-e}).
\]


To see this, subtract a lift of the coboundary
expression after reduction, global and local.
At a fixed tame place let \(D_u,I_u\) be its absolute
decomposition and inertia groups, and
\(J_u=I_u\cap\bigcap_i\ker\bar\chi_i\) for the residual diagonal
characters pulled back to \(D_u\). Its image on the base level is the
inertia kernel prescribed for the point. Thus \(J_u\) acts trivially on
the odd modules at every depth. The inertia quotient embeds into a product of multiplicative groups of the finite residual field \(k\). Therefore, for the number \(a\) of factors in the embedding,
\[
 |P_u|\mid (|k|-1)^a,\qquad (|P_u|,\ell)=1 .
\] A cocycle zero on \(J_u\) descends to
\(D_u/J_u\). Its restriction to \(P_u\) is a boundary by averaging.

Subtract that boundary and obtain a cocycle on
\(D_u/I_u\simeq\widehat{\mathbb Z}\) with values in the
\(P_u\)-invariants. Invariants lift under a surjection by the same
averaging. Lift the value on Frobenius, which gives a continuous
cocycle since the cyclic subgroup generated in the finite affine
group is finite. Add back a lift of the boundary to lift the original
one-cocycle exactly. This calculation also identifies the local
\(H^1\) unramified on \(J_u\) with Frobenius coinvariants on
the \(P_u\)-invariants. On the finite invariant coefficient module \(M\), let \(T=\mathrm{Frob}-1\).
Rank counting for this endomorphism gives
\[
 |\operatorname{coker}T|=\frac{|M|}{|\operatorname{im}T|}
                         =|\ker T|.
\]
Thus the Frobenius coinvariants and invariants have the same order. At the
\(\ell\)-adic places the local allowed classes
are boundaries and their boundary matrices
lift. The remaining cochains are then in the
coefficient kernel and hence give the inverse
image under \(i\).

For the second sequence,
\(|\operatorname{coker}r|\) is bounded by
\(|\operatorname{coker}r^2_e|\). Lift the global class
using surjectivity on global second cohomology.
The defects of its local nullhomotopies now have
coefficients in \(A_e\). For the finite, discrete coefficient modules a section of the coefficient surjection is continuous as a map of sets. A local one-cochain \(b\) therefore has a continuous lift by composition:
\[
 \tilde b=s\circ b,\qquad \operatorname{pr}\circ\tilde b=b .
\]
The defects are second cohomology classes in
the local quotients at depth \(e\). Changing the
global lift changes them by the image of
\(r^2_e\). If the defect in its cokernel vanishes,
adjust the global lift accordingly and use the
local cochains of the corrections; this produces
a lift in \(U_m\). The same comparison on
differences shows that equal defects yield the
same element of \(\operatorname{coker}r\).
This is the local restriction cokernel bound,
using global cohomological dimension two.

\proofstep{The quotient at a smaller depth and the elementary divisors.}
Consequently
\[
 |U_m/t^eU_m|\le |U_e|\,|\operatorname{coker} r^2_e|
                  \le |K|^{ge+C},\qquad e\le m.
 \tag{15}\label{eq:local-patch-15}
\]
Here \(C\) is the sum of the constants in \eqref{eq:local-patch-12},
the auxiliary correction, the local invariant
bounds and twice the constant in \eqref{eq:local-patch-13}. It is
chosen before \(m\) or \(e\). Write the elementary
divisors of \(U_m\) as \(A/(t^{h_i})\) in
nonincreasing order, with \(h_i\le m\).
If there is a \((g+1)\)-st divisor, use its own exponent as the test depth in \eqref{eq:local-patch-15}:
\[
 e=h_{g+1}\quad\Longrightarrow\quad h_{g+1}\le C .
\]
Thus the quotient by the first \(g\) summands is
killed by \(t^C\).

For any \(g+1\) obstruction
classes, their projections to the first \(g\)
summands admit a linear relation with at least
one unit coefficient. For example lift to a
matrix of \(g\) rows over \(A\), take a
fraction-field linear dependence and rescale to
integral coefficients with a unit. Multiplying
this combination by \(t^C\) kills the full
obstruction. For \(g=0\) \eqref{eq:local-patch-15} directly gives the
same annihilation bound.  Its obstruction classes
are the odd two-cocycles with their local
trivializations. Their coboundaries are precisely
the global one-cochain changes with the same
local conditions, as in \eqref{eq:local-patch-11}.

\end{proof}

\input{patching-marking}
\input{patching-matrices}
\input{patching-cotangent}
\input{closed-link}

\subsection{Recovering the augmentation support}
\label{subsec:marked-augmentation}

We have compared primes above each curve point. To obtain a statement about ordinary
Hecke support, we must now follow them back to the common level by augmentation.

\begin{proposition}[Support transport across a marked chain]\label{prop:marked-chain-support}
Fix a finite chain of normalized curve points at the same
ordinary level, with the auxiliary system of
Proposition~\ref{prop:marked-profinite-system} and the local
comparison of Proposition~\ref{prop:marked-stable-prime-inclusion}. Suppose there is a stable prime of the
ordinary Hecke support below the initial curve kernel.
If successive curves share a thick prime as in the
completed local comparison, write \(Q_{\rm end}\) for the last stable or thick target prime
and \(J_{\rm ord}\) for the image of the ordinary Hecke ideal at the target. Then
\[
                      J_{\rm ord}\subseteq Q_{\rm end}.
\] The augmentation comparison is made at the
target, where the tame variables vanish.
\end{proposition}

\begin{proof}
\proofstep{The common augmented comparison.}
Use the ordinary finite-level modules of
Section~\ref{par:ordinary-patching} with the same
choice of auxiliary places for the finite chain.
Write \(J\subset D\) for the annihilator-support
ideal in the base determinant ring, and let
\(D^{\rm sp}\) be the finite free splitting algebra
adjoining the ordered Frobenius roots. A common
coefficient extension and the variable specializing
to \(t\) are included in the base weight algebra.
They are also included in the base module \(X\)
after extension to \(D^{\rm sp}\). For every prime
of that algebra its support is the inverse image
of \(V(J)\). Indeed, let \(D\to S\) be the base change, \(\mathfrak q\in\operatorname{Spec}S\) and \(\mathfrak p=\mathfrak q\cap D\). The fibre of the finite module is
\[
 (X\otimes_D S)\otimes_S k(\mathfrak q)
  \simeq (X\otimes_D k(\mathfrak p))
                    \otimes_{k(\mathfrak p)} k(\mathfrak q).
\]
A field extension preserves nonzero vector spaces. Applying Nakayama at each of the two primes identifies the support test with the ideal test:
\[
 \mathfrak q\in\operatorname{Supp}(X\otimes_D S)\quad\Longleftrightarrow\quad J\subseteq\mathfrak p .
\]

\proofstep{The fixed marked limit for the chosen finite chain.}
The common finitely generated limit group gives
a coefficient ring \(R_0\) acting on these modules
at each precision. Its augmentation maps to the
base splitting algebra. The limiting relation ideal
\(I\) includes the group, polarization, ordered
ordinary, old inertia and auxiliary root and inertia
relations. Write \(\mathbf y\) for the \(nq\) tame
parameters and \(f\) for the specializations of the
two-map product. At every curve point \(f\) is a unit
after localization. Its factors reduce there to
differences of distinct roots since the new norms
tend to one.

At each fixed tame degree and coefficient precision
the modules at sufficiently deep levels are free
over the truncated coefficient and tame algebra.
Let \(u\) and \(v\) be the limiting comparison maps between the augmentation of the patched module and the base module \(X\), in opposite directions. Their composites specialize to multiplication by the same \(f\), so on inverting \(f\),
\[
 (f^{-1}v)u=1,\qquad u(f^{-1}v)=1 .
\]
This identifies the augmentation after localization with \(X[1/f]\).
The tame graded pieces are the corresponding
free monomial pieces over that augmentation.
This is the same module construction as for
one arc.

In particular approximation of a power
of each relation by the finite-level annihilators
makes that relation nilpotent on the patched
module. Here both the tame degree and coefficient
precision are fixed before taking the auxiliary
level depth. One may quotient by the relations
for the support comparison. More explicitly,
if \(U\) is the finite patched module over the
completed localized ring and
\(I\subseteq\sqrt{\operatorname{Ann}U}\), the
finite module \(U/IU\) has the same support by
Nakayama. Equivalently the comparison may be
made at each support prime before imposing the
relations. At augmentation primes the identity
\[
  \mathfrak q\in\operatorname{Supp}U
  \quad\Longleftrightarrow\quad
  \mathfrak q/(\mathbf y)\in
       \operatorname{Supp}(U/(\mathbf y)U)
 \tag{16}\label{eq:local-patch-16}
\]
follows by local Nakayama at a prime containing
\((\mathbf y)\). Apply \eqref{eq:local-patch-16} to the original patched
module, whose augmentation is \(X[1/f]\); this
avoids requiring its quotient by \(I\) to have
the original graded pieces.

\proofstep{Support after completed localization.}
For completed localization at the \(i\)-th
curve kernel, the completed weight ring has
dimension \(1+nd\). Indeed its coefficient and
weight algebra before localization includes
the uniformizer, one parameter specializing
to \(t\), and \(nd\) weight parameters. The
evaluation quotient has dimension one and
the localized weight ring is regular. The
selected completed summand of the augmentation
module is finite free over this completed weight
ring. On completing the finite algebra acting faithfully on \(X\), let \(e_i\) be its mutually orthogonal local idempotents. They exhibit the summands by
\[
 1=\sum_i e_i,\qquad \widehat X=\bigoplus_i e_i\widehat X .
\]
The factor module is a direct summand of a finite
free module over a local ring, hence is free.
Quasi-regularity of \(\mathbf y\) means that the
graded pieces of the completed patched module are
the monomial pieces on this free augmentation.

Write \(W\) for the completed weight ring, \(Y=(y_1,\ldots,y_{nq})\), and \(U\) for the completed patched summand. Lift a \(W\)-basis of \(U/YU\), of size \(d_0\), to \(U\). The resulting map
\[
  W[[\mathbf y]]^{d_0}\longrightarrow U
\]
induces the isomorphism on each graded piece
\[
 \bigoplus_{|\alpha|=h}W^{d_0}\mathbf y^\alpha
      \xrightarrow{\ \sim\ }Y^hU/Y^{h+1}U\qquad(h\ge0)
\]
by quasi-regularity. Induction on \(h\) gives the isomorphism modulo \(Y^{h+1}\), and completeness gives it on the full modules. The support
algebra embeds in the endomorphisms of this
finite free module. It is finite torsion-free over
the regular power series ring, of dimension
\(1+nd+nq\).

Let \(S\) be the regular power series ring and \(C\) this finite support algebra. Fix a minimal prime \(P\) of \(C\). If \(0\ne s\in P\cap S\), its image in the zero-dimensional Noetherian local ring \(C_P\) is nilpotent. Clearing a denominator gives
\[
               u s^a=0\quad\text{for some }u\in C\setminus P,\ a\ge1 .
\]
Multiplication by \(s\) on the free module is injective and hence is injective on its endomorphisms and on \(C\). This would force \(u=0\). Therefore \(S\hookrightarrow C/P\) is finite and integral. Going-up lifts prime chains of \(S\), and incomparability makes the contractions of strict chains in \(C/P\) strict. It follows that
\[
 P\cap S=0,\qquad \dim C/P=\dim S=1+nd+nq.
\]

The same statement holds for the minimal support
primes after the relations are imposed, since
the underlying support is unchanged.

\proofstep{Movement across a link and augmentation at the target.}
For clarity, the support is moved across a link
as follows. Take the stable start prime
on the augmentation and lift it to the completion
at the starting point. It lies in the patched
support by \eqref{eq:local-patch-16}. Choose a minimal support prime
below it, which is stable because stability is
preserved under generization. Sections~\ref{sec:patching-cotangent}
and~\ref{subsec:closed-link} put this prime below a lift of the
outgoing thick prime. Let \(R\) be the common tame-adic patch ring. Imposing the tame parameters is precisely the quotient
\[
 \pi:R\twoheadrightarrow R/(\mathbf y),\qquad\ker\pi=(\mathbf y).
\] If \(\mathfrak t\) is the outgoing prime contracted to \(R\), then \((\mathbf y)\subseteq\mathfrak t\).

For the next curve kernel \(\mathfrak k'\) in \(R\), the link on the augmentation gives
\[
  \mathfrak p_{\rm st}\cap R\ \subseteq  \mathfrak t=\pi^{-1}(\mathfrak t/(\mathbf y))\ \subseteq\ \mathfrak k' ,
\]
where \(\mathfrak p_{\rm st}\) is the minimal support prime chosen in the starting completion. Lift this latter
contraction by faithful flatness to the next
completed localization. Its support membership
comes from base change of the common patched
module. The stable prime in this step need not
contain \((\mathbf y)\). After iterating, apply the
local comparison to the last target and use
\eqref{eq:local-patch-16} there, where the target does contain
\((\mathbf y)\). Contracting to the original
determinant ring gives the inclusion of \(J\)
at that target. This is the support step in
the finite-chain comparison. It uses the
finite-level relation and comparison maps at
the actual Hecke levels and the local
cotangent and dimension implication at each
point.
\hfill\mbox{\cite{milne-duality}}\label{cor:ordinary-marked-patching-1}
\end{proof}

\section{Propagation of ordinary Hecke support between irreducible components}
\label{sec:linked-components}
\hypertarget{source-ordinary-linked-components}{}

The ordinary eigenpoints of Appendix~\ref{sec:linked-seeds}
supply the initial components. To reach the desired type-I points we pass support along the prescribed links. In rank
three one such passage can meet the reducible locus with irreducible outer plane; the
patching tests also cover this case. The resulting support statement is the input to
descent over the marked field and to the integral length
comparison. For a Hecke ideal \(J\), its image \(M_n\) in the degree-zero ring \(B_n\) controls Proposition~\ref{prop:ordinary-component-propagation} at each level: 
\[
 P\text{ is good for }J\quad\Longleftrightarrow\quad M_n\subseteq P\cap B_n.
\]


Keep the frame and the three levels of Section~\ref{sec:linked-seeds}.
Put \(d=[L^+:\Q]\), \(d_v=[L_v^+:\Q_\ell]\) for \(v\mid\ell\),
and \(D=nd+1\) in rank \(n\). We give the component comparison
with its arithmetic inputs stated explicitly. On a frame chart \(R\) write \(A=(a_{ij})\) for a frame matrix. Conjugation by the
torus gives the entry degrees from the two diagonal vectors:
\[
 r(t)A r(t)^{-1}=(t^{u_i-u_j}a_{ij})_{ij},\qquad
 \mathbf u=(1,-1)\ \text{or}\ (1,0,-1).
\]
Let \(I_{\rm up}\) and \(I_{\rm low}\) be the ideals killing,
respectively, the entries below and above the diagonal.

In
rank three the odd ideal \(I_{\rm odd}\) kills the four
entries joining the middle index to either outer index.
The complement of the triangular loci is the open subset
\[
 \operatorname{Spec}(R)\setminus
              \bigl(V(I_{\rm up})\cup V(I_{\rm low})\bigr).
\]
Within it the exceptional locus in rank three is \(V(I_{\rm odd})\). A point of type I
lies off both triangular loci and, in rank three, off the
exceptional locus. The marked normal form identifies this
with the absolutely irreducible locus. A point of type E
lies off both triangular loci but on the exceptional locus
in rank three. Its outer block is irreducible. For a Hecke ideal \(J\) let \(M_n=\Phi_n(J)B_n\).
Goodness is defined by the preimage of its vanishing locus:
\[
      P\in\operatorname{Good}(J)
                  \ \Longleftrightarrow\ M_n\subseteq P\cap B_n .
\]
Goodness specializes because contraction preserves inclusion: 
\[
 P\subseteq Q,\quad M_n\subseteq P\cap B_n\quad\Longrightarrow\quad M_n\subseteq Q\cap B_n.
\]


\paragraph{The prime conditions at the three levels.}
The support movement in Proposition~\ref{prop:marked-chain-support}
is used only with the completed local inclusions of
Appendix~\ref{sec:patching-cotangent}. To use these inclusions in the component graph we need two conclusions on the level
determinant ring. Its curve points take
values in \(k[[t]]\) with \(k\) finite of characteristic \(\ell\).
The point conditions include regular semisimple detection on the
cyclotomic kernel, the prescribed residual constituents, distinct
ordinary characters on the local test subgroups, and the local and
global restrictions of the marked level. In determinant terminology
a stable prime has absolutely irreducible determinant over an
algebraic closure of its fraction field. At a thick prime in this
application the semisimplification has an outer two-dimensional
constituent and a middle character. Write \(V_{\rm out}\) for the outer plane and \(\chi_{\rm mid}\) for the middle character. The ordinary Cayley--Hamilton algebra admits a nonsplit extension with semisimplification 
\[
 V^{\rm ss}\simeq V_{\rm out}\oplus\chi_{\rm mid},\qquad \dim V_{\rm out}=2.
\]

Let \(Q\) denote the kernel of a nice type-I curve point
in the Hecke locus. The first patching implication is
\begin{equation}\label{eq:ordinary-component-stable-test}
              P\subseteq Q\quad\Longrightarrow\quad J\subseteq P .
\end{equation}
For the second implication write \(Q_0,\ldots,Q_s\) for
the kernels of nice type-II curve points. Suppose
\(Z_i\subseteq Q_{i-1}\cap Q_i\) is thick for each
\(1\le i\le s\) and \(J\subseteq P_0\subseteq Q_0\) for
a stable \(P_0\). Then the comparison gives
\begin{equation}\label{eq:ordinary-component-thick-test}
 P\subseteq Q_s,\quad P\ \text{stable or thick}
                  \quad\Longrightarrow\quad J\subseteq P .
\end{equation}
This includes the chain with \(s=0\).
One common Taylor--Wiles system is needed for the whole list.
Section~\ref{sec:marked-patching} constructs its common tests,
gives the uniform first-order and odd obstruction bounds, and
recovers the Hecke support on the augmentation fiber. Sections~\ref{sec:patching-cotangent}
and~\ref{subsec:closed-link} give the completed comparison, including
specialization of thick primes. We verify the point conditions and
identify the ordinary modules on the three levels at the curve
construction below.

The frame constructed in Section~\ref{sec:marked-hosts}
supplies the abelian quotient bound, the local inertia generators
and the ordinary Selmer inequality \eqref{eq:fitted-relative-euler}.
The locally split nonzero residual cocycles in the outer and
adjacent weights, the bound \(C_0\), and the lower bounds on
\(d,d_v\) all apply to this frame.

\subsection{The two-dimensional corner}

We start in rank two because the outer plane will enter both rank-three comparisons.
Write \(\lambda_n\) for the ideal of \(B_n\) generated by the
uniformizer and the differences \(\psi_{v,i}(g)-1\) on inertia
with trivial residual scalar. The finite weight condition on
an ideal \(I\) in these charts is
\(\mathfrak m_{B_n}\subseteq\sqrt{I+\lambda_n}\).
It holds for \(M_n\). For the weight transport let \(q_0\)
be the residue cardinality. Raising inertial labels to
\(q_0-1\) kills their residual values.

Write \(h=q_0-1\). For units \(a,c\) with the same residue the geometric factor in \(a^h-c^h\) is a unit, because
\[
 \sum_{i=0}^{h-1}a^{h-1-i}c^i\equiv h\bar a^{h-1}\ne0 .
\] Thus equality of the powers in a
prime quotient recovers equality of the units. For the
scalar twist use its fixed inertial values in the weight
specialization of the level ring. The finite algebra over
that ring has Artinian special fiber. Its map onto the
specialized degree-zero quotient forces every prime there
to be maximal.

\begin{lemma}[Outer-corner support and the weight fibre]
\label{lem:rank-three-outer-corner}
For the indicated rank-two corner of the rank-three chart, let \(I_3\) be the ideal obtained from \(M_2\). For every type-E point
\(P\) satisfying the specified outer-block condition,
\[
                       I_3\subseteq P\cap B_3 .
\] It satisfies the finite weight-fibre condition for
\(\lambda_3\).
\end{lemma}

\begin{proof}
\proofstep{The completed rank-two comparison.}
First carry out the type-I propagation below in rank two,
where only the first patching input is needed. Quotient
\(A_3\) by the odd entries and let \(C\) be the image of
\(B_3\) there. Write \(p:B_3\to C\). The outer block
\(\sigma\) and the middle entry \(u\) in this quotient are
individually multiplicative. The outer block inherits the
transpose relation by restriction of the antidiagonal matrix
to the outer indices. It has ordinary characters \(\psi_0,\psi_2\).
To verify this integrally, at an ordinary place choose the marked
element \(g_1\) for which the middle residual value differs
from both outer values. Put
\[
 P_x=\sigma(x)-\psi_0(x),\quad
 T=\sigma(g_1)-\psi_1(g_1),\quad Q_y=\sigma(y)-\psi_2(y),
\]
where scalars multiply the identity. The determinant of \(T\)
is a unit by the marked unit condition. The rank-three ordinary
product on the outer block gives \(P_xTQ_y=0\) for all \(x,y\).

Taking \(y=g_1\) gives
\(P_xT^2=(\psi_2(g_1)-\psi_1(g_1))P_xT\).
Cayley--Hamilton in dimension two gives
\[
  \det(T)P_x =
    \bigl(\operatorname{tr}T-\psi_2(g_1)+\psi_1(g_1)\bigr)P_xT .
\]
Multiplication by \(Q_y\) yields \(P_xQ_y=0\). On the middle
entry of the triple product at \((g_1,y,g_1)\) the first and
last factors are units, giving \(u(y)=\psi_1(y)\) locally.
The corner therefore defines a map from the rank-two frame
whose degree-zero part is \(\beta:B_2\to C\). Put
\[
                        I_3=p^{-1}\bigl(\beta(M_2)C\bigr).
\]
At a type-E prime a nonzero upper and lower entry must both
lie in the outer block. Its contraction in rank two is type I,
hence good for \(M_2\). If \(P\) is the type-E prime, the rank-two contraction contains \(M_2\); taking the inverse image in the degree-zero ring gives 
\[
 I_3=p^{-1}(\beta(M_2)C)\subseteq P\cap B_3.
\]


\proofstep{The resulting weight-fibre bound on the outer corner.}
The ideal \(I_3\) also has finite weight fiber. A prime above
\(I_3+\lambda_3\) gives a prime of \(C\) whose rank-two
contraction contains \(M_2+\lambda_2\). All rank-two degree-zero
values then reduce to coefficients. Consider the middle character
\(u\) in the resulting local domain in characteristic \(\ell\).
It reduces to one, is invariant under the frame involution, and
is one on the indicated inertia kernels. Modulo each power of the
maximal ideal its continuous image is a finite \(\ell\)-group in
the principal units. The abelian generator condition of the frame
gives a fixed positive exponent \(h\) for such finite quotients
once the prescribed inertia and involution differences are killed.
Separatedness turns the finite-depth identities into an equality in the domain.

For its maximal ideal \(\mathfrak m\),
\[
 u(g)^h-1\in\bigcap_{b\ge1}\mathfrak m^b=0 .
\] Its \(\ell\)-part
can be dropped by injectivity of Frobenius. The remaining geometric
sum is a unit since \(u(g)\) reduces to one. Thus \(u=1\).
In any degree-zero monomial an odd factor vanishes and all middle
diagonal factors can be removed. The remaining monomials come
from rank two, and the middle ordinary character is one. These
generators and the coefficient reduction force the prime to be
maximal. This gives the required weight-fiber control on the
exceptional locus using only the rank-two Hecke ideal.

\end{proof}

\subsection{Curves and the component graph}

To connect components through curves, work locally at a closed point on the residual
lower triangular locus off the upper triangular locus. Use its complete pointed chart
\(R\), flat over the frame ring, and contract primes of \(R\)
to \(B_n\). Call a chart prime nice if its quotient is
one-dimensional of characteristic \(\ell\), it lies off both
triangular loci, and its distinct ordinary indices give
different characters on each distinguished local subgroup \(N_v\).
The following estimates govern the choice. In the chart the
upper triangular ideal is the unit ideal and the opposite
triangular quotient has dimension at most
\[
                   (n-1)d+2+3C_0+5 < nd .
\]
For the latter estimate cut by the uniformizer, \(1+C_0+d\)
outer diagonal values at abelian generators and \(d\) values for
each middle diagonal index.

On the triangular quotient the
diagonal characters determine the characteristic polynomials.
The cut forces the degree-zero maximal ideal and leaves chart
fiber dimension at most \(3C_0+5\). For a prime off both
triangular loci above \(I\) with finite weight fiber, coincidence
of two characters on \(N_v\) in characteristic \(\ell\) gives
\[
                    \dim(R/P)+d_v\le nd+5 .
\]
The weight count drops the \(d_v\) generators of one of the
equal labels, retaining one distinguished value, the other
generators and the coefficient uniformizer. The radical test
bounds the dimension by the number of cuts. The off-triangular
framing has at most one transcendental parameter over degree
zero. The large-frame condition gives \(d_v>5\).

\begin{proposition}[A nice specialization in the pointed chart]\label{prop:nice-ordinary-specialization}
Use a complete pointed chart at a closed point of the
residual lower triangular locus off the upper triangular locus.
Let \(P\) have quotient dimension at least \(D\), with
contraction to \(B_n\) containing an ideal \(I\) of finite
weight fibre. There exists a nice prime \(Q\) with
\[
                    P\subseteq Q,\qquad \operatorname{char}\kappa(Q)=\ell .
\] The resulting curve admits
an embedding into a complete DVR with finite residue field.
\end{proposition}

\begin{proof}
\proofstep{Producing a curve with the required specializations.}
Let \(P\) have quotient dimension at least \(D\), with contraction
containing such an \(I\). Imposing the uniformizer leaves a
component of dimension at least \(nd\), by the principal ideal
theorem and catenarity. The bounds just proved keep its generic
point off the opposite triangular locus and all local coincidence
loci. In that component choose a product \(f\) of one nonvanishing
element from each bad ideal, times a nonzero element of the maximal
ideal.

Let \(A\) be the local domain of the chosen component, and write \(R\twoheadrightarrow A\) for the quotient of the pointed chart. Lemma~\ref{lem:local-curve-avoidance} applied to \(0\ne f\in\mathfrak m_A\) gives a prime \(\mathfrak q\) avoiding \(f\) with one-dimensional quotient. Pull it back to \(Q\) in \(R\). Then
\[
 P+(\varpi)\subseteq Q,\qquad R/Q\simeq A/\mathfrak q,
 \qquad f\notin\mathfrak q .
\]
The first inclusion gives the specialization and its residue characteristic. The last condition keeps its generic point away from every excluded locus, since the factor of \(f\) chosen in the corresponding ideal survives there. The complete local domain of the
resulting curve embeds into a complete DVR \(k[[t]]\) with
\(k\) finite. A parameter gives a finite extension of a
one-variable power series ring over the coefficient field. Normalizing in the finite fraction field
extension gives the DVR.

\end{proof}

\input{marked-points}

\paragraph{Contraction to a level.} Now contract to a level. A type-I generization gives
a stable determinant, since the representation
is absolutely irreducible and the level group maps onto the
frame group. For a thick generization use the nonsplit
ordinary realization of its block determinant. Its group algebra
map kills the Cayley--Hamilton relations by the characteristic
polynomial. The ordered products, before and after the involution,
map to scalar multiples of the ordinary and reversed ordinary
products of the extension. The tame inertia relations follow
from the marked prime-to-\(\ell\) power test. Adic continuity
kills the closure of the ideal of relations. Thus the nonsplit
realization factors through the ordinary Cayley--Hamilton algebra,
while its determinant is reducible. Contraction transports
both prime conditions, including the links that must be thick.

\begin{construction}
\begin{tikzpicture}[>=Stealth,every node/.style={align=center,inner sep=2pt}]
\node (two) at (0,0) {Rank-two comparison on \(M_2\)};
\node (input) at (0,-.85) {Outer-corner support};
\node (three) at (-2.3,-1.9) {\(M_3\) and its quotient};
\node (zero) at (2.3,-1.9) {\(M_0\) and its quotient};
\draw[->] (two)--(input);
\draw[->] (input)--(three);
\draw[->] (input)--(zero);
\end{tikzpicture}
\caption{Propagation uses the rank-two comparison first. Its outer-corner input is applied to two different rank-three support problems.}\label{fig:ordinary-three-levels}
\end{construction}

\paragraph{Which components can link curves.}
Which components do these links reach from the seed?

\begin{proposition}[Propagation from a seed across ordinary components]
\label{prop:ordinary-component-propagation}
In the selected ordinary frame, given the seed of Appendix~\ref{sec:linked-seeds} of the required dimension
and the completed comparisons
\eqref{eq:ordinary-component-stable-test} and~\eqref{eq:ordinary-component-thick-test},
each type-I prime \(P\) is good for its corresponding level ideal, meaning 
\[
 M_n\subseteq P\cap B_n.
\]

The application to \(M_2\) supplies the outer-corner input for the
applications to \(M_3\) and \(M_0\), which keep their separate
quotient maps.
\end{proposition}

\begin{proof}
\proofstep{Relative dimension and the presentation of the chart.}
The ungraded chart has a presentation over a complete DVR in
\(N\) variables with \(N-D\) relations, padding with zeros if
necessary. At a rank-three point on the exceptional locus, use
conjugation by the sign on the middle coordinate to grade the
presentation by parity. There are \(a\) even and \(g\) odd
variables, at most \(a-D\) even and \(g\) odd relations. The relation count uses the relative obstruction group.
For a small extension with kernel \(I\), let \(c(g,h)\) be
the \(I\)-valued error in multiplying two matrix lifts.
For each ordinary place let \(b_v(g)\) be the error in the
lifted local flag, read in the adjoint quotient by
\(\mathfrak b_v\). With the cone differential
\(d(c,b)=(dc,\operatorname{res}(c)-db)\), associativity
and the local multiplication law give
\[
      dc=0,\qquad db_v=\operatorname{res}_v(c).
\]

Write \(\alpha\) for the one-cochain changing the matrix
lifts and \(\beta_v\) for the zero-cochain changing the flag at
\(v\), in the cone convention just specified. The new errors are
\[
 c'=c+d\alpha,\qquad
 b'_v=b_v+\operatorname{res}_v(\alpha)-d\beta_v .
\]
Thus a zero obstruction class permits both errors to be removed
in the same lift. The dual
of the minimal relation space injects into this obstruction
space. Indeed if a nonzero functional gave zero obstruction,
push out along it and lift a section to the power series variables.
Two lifts of the variables then differ in the square-zero
kernel killed by the maximal ideal. Every minimal relation has
zero relative linear part modulo that ideal, so its value is
unchanged. A relation on which the functional is nonzero would
have both zero and nonzero value. Since \(2\) is invertible, an involution \(s\) admits the complementary projections
\[
 e_+=(1+s)/2,\qquad e_-=(1-s)/2,\qquad e_+e_-=0.
\]
Applying the argument to the images gives both eigenspace assertions.

In the ordinary tangent and Euler characteristic formulas of
the marked frame, the finite-place contribution is
\(\sum_{v\mid\ell}d_v\dim\mathfrak b_v\), from which one subtracts the
sum of the dimensions fixed by real conjugations. Here
\(\dim\mathfrak b_v=n(n+1)/2\). Multiplication by the
polarization matrix identifies the fixed matrices at a real
place with skew matrices, of dimension \(n(n-1)/2\).

Since \(\sum_{v\mid\ell}d_v=d\) and \(L^+\) has \(d\) real
places, counting the extra frame direction gives
\[
 d\frac{n(n+1)}2-d\frac{n(n-1)}2+1=nd+1=D .
\]
In rank three the Borel dimensions are \(4\) even and \(2\) odd:
the even part consists of the three diagonals and the outer root.
The real skew dimensions are \(1\) even and \(2\) odd, from the
outer pair and the two adjacent pairs respectively. The extra
frame direction is even. Thus the two lower bounds for relative
tangent dimension minus obstruction dimension are
\[
        4d-d+1=D,\qquad 2d-2d=0.
\]

\proofstep{Connectedness and the odd incidence equations.}
In rank two, or outside the exceptional locus of rank three,
Grothendieck connectedness applied to the ungraded presentation
joins the minimal components in a graph with intersections
of dimension at least \(D\). Each component has dimension at
least \(D+1\). At the remaining rank-three centers write \(y\)
for the odd variables. If there are odd variables, pad the odd
relations to a square system \(H y=0\). Each odd relation lies
in the ideal \((y)\) by parity, and minimality permits the
entries of \(H\) to be taken in the maximal ideal. Impose
\(\det H=0\) in the chart. Every type-I prime contains the extra equation. Indeed modulo such a prime at least one entry of \(y\) survives, whereas
\[
 0=\operatorname{adj}(H)Hy=(\det H)y.
\]
The quotient is a domain, so \(\det H\) vanishes there. The odd dimension on this pointed chart is positive in the exceptional case.

To see the same intersection bound, use an independent
\(g\times g\) matrix \(X\) and vector \(z\). The incidence
locus \(Xz=0,\det X=0\) has irreducible support of codimension
\(g\). On \(z_i\ne0\) solve for the \(i\)-th column, with
all other coordinates free, and these charts overlap. At \(z=0\), choose a nonzero vector \(v\in\ker X\) for the singular residue matrix. The line is a specialization from the open locus, since
\[
 X(tv)=0,\qquad tv\ne0\ (t\ne0),\qquad (tv)|_{t=0}=0 .
\] The homogeneous prime radical gives a domain after
completion: the initial homogeneous pieces of two nonzero
elements still have nonzero product. Keep the \(a\) even and
\(g\) odd variables of the original chart separate from the
\(g^2+g\) new variables \((X,z)\). The new incidence equations
have codimension \(g\) in the new variables. Thus the ambient
irreducible complete local germ has dimension
\[
             1+a+g+(g^2+g)-g=1+a+g+g^2.
\]

In this germ impose \(X-H=0\), \(z-y=0\) and the at most
\(a-D\) even relations. Their number is at most
\(g^2+g+a-D\), so each component has dimension at least
\(D+1\). The complete local connectedness theorem with the
same equation count gives connectedness in dimension \(D\).
\cite{stacks-connectedness}

\proofstep{The nonsplit extension at an exceptional branch.}
The minimal primes of this smaller locus are off both triangular
loci by the bounds above. An exceptional prime there has \(y=0\)
and \(H\) singular. Choose a nonzero vector in its kernel
over the fraction field and clear denominators. At \(y=0\) an odd variation annihilates the even relations. For the odd ones the linearization is
\[
 d(Hy)=(dH)y+H\,dy=H\,dy=0 .
\] This gives a derivation
of the original chart, zero on even entries and ordinary characters.
If \(r_{\rm out},r_{\rm mid}\) are the two retained blocks,
let \(u(g)\) be the derivative of the column from the middle
line to the outer block. The even derivatives vanish, so the
linearized multiplication law in that column is
\[
   u(gh)=r_{\rm out}(g)u(h)+u(g)r_{\rm mid}(h).
\]
Consequently \(\left(\begin{smallmatrix}r_{\rm out}(g)&u(g)\\
0&r_{\rm mid}(g)\end{smallmatrix}\right)\) is a representation with
the original block determinant.

The same linearization of the
two ordered ordinary products gives the prescribed local tests. A splitting would give an
invariant eigenvector with middle coordinate 1. Write a splitting vector as \((u,1)\) in the outer-plane
and middle-line decomposition. Its equations at the mark
\(g_*\) and its polarized mate \(g_*^*\) have a common
off-diagonal column and the same middle value, while the
two distinct outer eigenvalues are exchanged. Write \(D_*,D_*'\) for the two outer diagonal matrices, \(c\) for
the common column and \(m\) for the common middle value. The two splitting
equations give
\[
 D_*u+c=mu=D_*'u+c
 \quad\Longrightarrow\quad (D_*-D_*')u=0 .
\]
Both diagonal entries of the difference are nonzero, forcing \(u=0\). The splitting equation at all
\(g\) then says the differentiated column vanishes. Polarization also kills the opposite differentiated
row, contradicting the nonzero odd derivation. This proves
thickness at the exceptional branches used by the graph.

\proofstep{Propagation from a seed.}
Suppose a chart has a good type-I prime with quotient dimension
at least \(D\). Write \(P\) for this good prime. Using \(I=M_n\), choose the minimal prime \(P_{\min}\) and the nice specialization \(Q\) so that
\[
 P_{\min}\subseteq P\subseteq Q,\qquad \dim(\text{chart}/P)\ge D .
\] The minimal prime is
type I by generization, and is a vertex of the component graph,
since every type-I prime contains the additional chart equation
when one is imposed. If the nice point is type I, the first
prime condition makes that vertex good. If the nice point is
exceptional, the original good prime serves as the stable
generization for the length-zero second prime condition. This
again makes the vertex good.

\proofstep{The induction along the component graph.}
Along a path of adjacent vertices keep the following information.
A type-I vertex reached is good. At a thick vertex keep a
nice exceptional point containing it, linked by a finite
sequence of nice exceptional points to one having a good type-I
generization. In the chart ring \(A\), adjacency of successive vertices \(P,P'\) supplies a common specialization of controlled dimension: 
\[
 V\supseteq P+P',\qquad\dim A/V\ge D.
\]
 Choose the nice specialization \(Q\) above \(V\); then
\[
                       P+P'\subseteq V\subseteq Q .
\]

If \(P\) is good type I, use \(M_n\) in the specialization
and apply the first or the length-zero second condition at
\(Q\), according to its type. If \(P_1\) is any type-I generization of this \(Q\), either implication gives 
\[
 P_1\subseteq Q\quad\Longrightarrow\quad M_n\subseteq P_1\cap B_n.
\]
 If \(P\) is thick, \(V\)
contains the corner ideal \(I_3\). Use \(I_3\) in the
specialization: \(Q\) still contains the odd ideal, so is
exceptional.
Write \(Q_{\rm old}\) for the stored nice point on the current thick vertex. The common vertex joins the new exceptional point to the chain since 
\[
 P\subseteq Q_{\rm old}\cap Q.
\]

The second prime condition now makes every type-I generization
of the new point good. If the next vertex is type I, this proves
the assertion for it. If it is thick, specialize \(V\) using
its \(I_3\) as well. A previous vertex of type I gives a good type-I generization of this new exceptional point. When the previous vertex is
thick the stored chain extends through that common vertex.
This proves the invariant without an assumption that every
thick component belongs to the rank-three Hecke locus.
In rank two all the vertices and nice points used are type I.

\proofstep{Passage between charts.}
It remains to pass between charts. Write \(\pi_0\) for the projection to the degree-zero spectrum, \(x_0\) for its residual point and \(Z_{\rm up}\) for the upper triangular locus. The core used to join the charts is 
\[
 C=\pi_0^{-1}(x_0)\setminus Z_{\rm up}.
\]
 Its relevant points admit smooth lower
triangular realizations over their residue fields. Minimal
frame primes are homogeneous. A homogeneous type-I generic
point specializes to this core. One chooses a lower entry nonzero
there, uses the chart obtained by retaining its direction under
the torus, and specializes its degree-zero coefficients to the
closed point.  Applying this to a homogeneous lift of a seed contraction gives a chart prime \(\mathfrak q\) with quotient dimension
\[
 \dim(\text{chart}/\mathfrak q)=\dim\bigl(B_n/(P\cap B_n)\bigr)+1 .
\] The extra parameter is the one-dimensional torus orbit.

The core is preconnected. In rank two the sole lower entry is
a twisted cocycle subject to linear local and marked equations.
Interpolation over \(k[t]\) to a fixed nonzero solution gives
a common generization, still off the upper triangular locus.
In rank three write \(y\) for the adjacent lower entry and
\(a\) for the outer lower entry. The other adjacent entry
is given by the involution.

Write \(\partial a\) for the linear part of the outer
cocycle equation and \(Q(y)\) for its quadratic term.
A pure outer solution \(b\) satisfies \(\partial b=0\),
so
\[
        \partial(t^2 a+(1-t)b)=t^2 Q(y)=Q(ty).
\]
This is the lower multiplication law for the interpolant
\((ty,t^2a+(1-t)b)\). The local flag coordinates scale
in degrees one and two, and adding a pure outer solution adds
its outer flag coordinate. Pick \(b\ne0\) as supplied by the
frame. The kernel of the map into \(k[t]\) is a common
generization of \((0,b)\) and the given point. To see density
off the exceptional locus one needs a solution with \(y\ne0\).

For the frame's nonzero adjacent cocycle \(u\), put \(P=\bar\Psi\)
and use the conventions \(u(gh)=P(h)u(g)+u(h)\) and
\(u(g^*)=u(g)/P(g)\). Set, pointwise,
\[
 y=u+\delta(P-1),\qquad
 a=\frac{u^2}{2P}-\delta u(1+P^{-1})
                +\delta^2\bigl(1-(P+P^{-1})/2\bigr).
\]
Here \(\delta(P(g_*)^2-1)=-u(g_*)(P(g_*)-1)\), with
\(g_*\) the marked element for the normal form. Substitution
gives the marked equality for \(y\), the quadratic multiplication
law for \(a\) and the required local coboundary forms, using
\(u|_{G_v}=\alpha_v(P-1)\). If \(y=0\), substitution in the formula for \(u(g_*^*)\) gives
\[
 2\delta\bigl(P(g_*)^2-1\bigr)=0 .
\]
This forces \(u=0\), a contradiction. With one such solution
\((y_0,a_0)\), the following curve exhibits a pure outer point as a specialization: 
\[
 (y(t),a(t))=(t y_0,t^2a_0+a),\qquad (y(0),a(0))=(0,a).
\]
 Since \(y_0\ne0\), the generic point lies outside the exceptional locus.

\proofstep{From the open core to all required components.}
Divide the finitely many components meeting the open type-I
locus according as their generic points are good or bad.
The high-dimensional seed in its pointed chart gives a good
component by the local propagation. Every component under
consideration meets the core by marked specialization.
The good and bad unions are closed and cover the core, because
off the exceptional locus its closed points lie on type-I
components by the ungraded chart calculation, and these
points are dense there and in the core.

By the Jacobson
specialization property of this fiber, a closed set
meeting it contains a closed chart center there, also after
restriction to an open subset of the core.
Write \(C_{\rm good}\) and \(C_{\rm bad}\) for the two closed subsets of the core \(C\). Preconnectedness gives 
\[
 C=C_{\rm good}\cup C_{\rm bad},\quad C_{\rm good},C_{\rm bad}\ne\varnothing\quad\Longrightarrow\quad C_{\rm good}\cap C_{\rm bad}\ne\varnothing.
\]
 Choose a closed chart center in the
intersection. Flatness lifts both generic points by going
down, giving two type-I minimal chart primes. Local propagation
from the good one reaches the other, contrary to the latter's
contraction. Thus every type-I prime is good. This argument
first applies to \(M_2\), then to each of \(M_3\) and \(M_0\)
on the same rank-three frame.

\end{proof}

\subsection{Two levels and the local annihilator}

Support at the two levels supplies local information at different exceptional places.
To compare the resulting cuspidal attachments, let \(\xi:A_3\to\overline{\Q}_\ell\) be a continuous integral
type-I point with arithmetic ordinary characters on the first
level. At the other level identify the local groups via the
marked place, in the order prescribed by the common frame.
The map of pairs on each level gives the characters of precisely
the same twisted frame point. Thus local algebraicity and the
strict inequalities of its weights transfer at each embedding.

Both determinant points lie on their respective Hecke supports.
Classicality gives cuspidal attachments with the prescribed
local compatibility away from the respective exceptional places.
Both pullbacks give the same representation \(\epsilon^{-1}r_\xi\). Thus for a good Frobenius \(g\) their common characteristic polynomial is
\[
 \det\bigl(1-X\epsilon(g)^{-1}r_\xi(g)\bigr).
\]
Agreement at almost all unramified places identifies the cusps by strong multiplicity one. At
each exceptional prime the other level gives the missing local
comparison.

\paragraph{Weights for fixed-weight adjoint vanishing.}
The comparison at the two levels will be used locally on a component. Fix an integral
type-I point of the full frame ring \(A=A_n\)
with kernel \(q\). Lemma~\ref{lem:ordinary-annihilator-derivative}
will show that \(a\in M_n A\) vanishes in \(A_q\).
To supply the local hypothesis for the density argument on a component,
choose one embedding \(\tau_v:L_v\hookrightarrow\overline{\Q}_\ell\)
for each marked place over \(\ell\). At each marked place and \(j=0,1,2\), prescribe the weight according to the chosen embedding by
\[
 h_j(\tau)=\begin{cases}2j&\tau=\tau_v,\\ j&\tau\ne\tau_v.\end{cases}
\] The indices increase strictly at every
embedding. The local degree is greater than five, so choose
a nonzero trace-zero \(x\in L_v\), sufficiently small that
the exponential and logarithm are inverse on its neighborhood
of zero. Then \(u=\exp(x)\) is a nontorsion unit with norm one.

Take an inertial Artin lift of \(u\). Its cyclotomic value is
one, and the \(j\)-th ordinary character there has value
\[
   \zeta_j\,N_{L_v/\Q_\ell}(u)^{-j}\tau_v(u)^{-j}
      =\zeta_j\tau_v(u)^{-j},\qquad \zeta_j\ \text{of finite order}.
\]
Every ratio for distinct indices has infinite order on this
single element. The argument works for any finite parts in the
specialization. Density at this weight therefore supplies the
local hypothesis of fixed-weight adjoint vanishing at an
integral type-I point on each component. At
rank three the two-level comparison supplies its finite-place
compatibilities. We give the commutative algebra in detail.

\begin{lemma}[Differentiation on the ordinary support]
\label{lem:ordinary-annihilator-derivative}
With the finite weight quotient and type-I density at these
two levels, an ordinary Hecke relation \(a\) restricted to the
chart is annihilated by an element outside the prime of the
type-I point. Writing \(\xi\) for evaluation at the point and \(D_\xi\) for any
point derivation, this gives
\[
                      \xi(a)=0,\qquad D_\xi(a)=0
\]
also at a point with resonant weights.
\end{lemma}

\begin{proof}
\proofstep{Heights in a flat pointed chart.}
Use a flat pointed chart \(R\) containing a prime \(Q\) over
\(q\), minimal in that fiber, with \(\dim(R/Q)=1\). This chart
is obtained by specializing an integral model of the point
and completing at its residual center. Every prime \(P\subset Q\)
contracts to a type-I prime of the frame, so its degree-zero
contraction contains \(M_n\). That quotient is finite over an
ordinary weight quotient and has dimension at most \(D\).
The chart dimension formula adds at most one for the torus
fiber, giving \(\dim(R/P)\le D+1\). Extend any chain ending at \(Q\) by the specialization to the closed point. The resulting dimension bound is
\[
 \operatorname{ht}Q+1\le D+1 .
\]

In the presentation of \(R\) with \(N\) variables over a
DVR and \(N-D\) relations, the lift of \(Q\) has height \(N\)
by catenarity and \(\dim(R/Q)=1\). Its height after imposing
\(N-D\) equations is at least \(D\). Hence equality holds.
At the lifted prime the relations have their full possible
height in a regular local ring and form a regular sequence.
The quotient local ring \(R_Q\) is therefore Cohen--Macaulay.
For the flat local map \(A_q\to R_Q\), the depth and dimension
formulas with the zero-dimensional closed fibre give
\[
 \operatorname{depth}A_q=\operatorname{depth}R_Q
      =\dim R_Q=\dim A_q .
\]
In particular \(A_q\) is Cohen--Macaulay and has no embedded
associated primes. The dimension identity is
\(\operatorname{ht}q=\operatorname{ht}Q\).

\proofstep{Generic reducedness along every component.}
For each minimal prime \(P_0\subset q\) of \(A\), going down
and taking a minimal prime in the chart lifts it to a minimal
chart prime below \(Q\). Every such chart component has
dimension at least \(D+1\), by the presentation and the
height theorem. Its frame contraction has degree-zero quotient
dimension at least \(D\), by the bound of one on the torus
fiber. If \(P\) is the frame contraction and \(B_n\) its degree-zero ring, the fiber bound and the Hecke upper bound now agree: 
\[
 D\le\dim B_n/(P\cap B_n)\le D.
\]
 Apply density on this component and lift using the
graded frame to a type-I point \(\xi'\) on \(P_0\) with
the weights just prescribed.

At \(\xi'\) the ordinary adjoint
vanishing says every fixed-weight deformation is inner.
In the marked basis the remaining infinitesimal change of
frame has dimension at most one. The extra inertial weight directions contribute at most \(nd\), so the bound for the full tangent space is
\[
 \dim T_{\rm full}\le 1+nd=D .
\] The height lower
bound at \(q'=\ker\xi'\) is \(D\) by the same chart
presentation. Over the characteristic-zero coefficient field
the residue field extension is separable and formally smooth,
so every cotangent functional extends to a coefficient-relative
point derivation. Together the tangent and height bounds give 
\[
 D\le\dim A_{q'}\le\operatorname{edim} A_{q'}\le D.
\]
 Hence \(A_{q'}\) is regular local and a domain; localization at its minimal prime gives the field \(A_{P_0}\).

\proofstep{Differentiating the annihilator relation at the original point.}
Now if \(\operatorname{Ann}_A(a)\subseteq q\), take a prime
minimal over that annihilator inside \(q\). It is associated
to the cyclic submodule \(Aa\), hence to \(A\). It must
therefore be a minimal prime \(P_0\) by the unmixedness at
\(q\). Goodness at \(P_0\) puts \(a\) in \(P_0\). In the
field \(A_{P_0}\) it vanishes. This supplies an annihilator
outside \(P_0\), a contradiction. We obtain the asserted
\(s\notin q\). Since \(s\notin q\), we have \(\xi(s)\ne0\). Evaluation first gives \(\xi(a)=0\); differentiation then gives 
\[
 0=D_\xi(sa)=\xi(s)D_\xi(a)+\xi(a)D_\xi(s)=\xi(s)D_\xi(a).
\]
 Thus every point derivation kills \(a\), even if the original point has resonant
weights.  At that point Section~\ref{sec:rank-three-multiplicity}
supplies the selected multiplicity used on the determinant
quotient in the fixed-weight comparison.

\end{proof}

\paragraph{Exclusion of an adjacent cyclotomic ratio.}
We now use the bound on fixed-inertia derivatives to rule out a cyclotomic ratio for
the marked points specified below.

\begin{paramdefinition}[The points tested for an adjacent cyclotomic ratio]\label{hyp:cyclotomic-ratio-point}
On the rank-three type-I chart of
Proposition~\ref{prop:ordinary-component-propagation}, let
\(X_e\), \(e\ge1\), denote the field points having the finite
inertia parts and local degree bounds of
Proposition~\ref{prop:separated-marked-host}, absolutely irreducible
on each open global subgroup, and potentially crystalline at the
marked \(\ell\)-adic places with ordered exponents \((e,0,-e)\).
At such a point fix the inertia characters when taking derivatives,
as in Lemma~\ref{lem:ordinary-annihilator-derivative}.
\end{paramdefinition}
For the application to the even level, the restriction and
attachment arguments proving Lemma~\ref{lem:rank-three-host-restriction}
and Proposition~\ref{prop:ordinary-host-attachment} supply these
properties before applying the ratio test.

\begin{proposition}[Exclusion of cyclotomic resonance]
\label{prop:ordinary-no-cyclotomic-resonance}
At a point of Definition~\ref{hyp:cyclotomic-ratio-point}, writing \(D_v\) for the marked local group, any two distinct ordinary flag positions satisfy 
\[
 (\psi_i/\psi_j)|_{D_v}\ne\epsilon|_{D_v}\qquad(i\ne j).
\]
 The two estimates concern the same
fixed-inertia tangent space: the global calculation bounds its
rank by one, while the proposed ratio forces rank at least two.
\end{proposition}

\begin{proof}
\proofstep{The upper bound on fixed-inertia derivations.}
First consider the upper estimate at the point \(x\) of the
level determinant ring. A derivation induced from the frame
with fixed inertia annihilates \(J\) by the local annihilator
and kills the power series weight ring \(\Lambda\) by its
generators. Let \(J'\) be the reduced support ideal of the
projected module. A trace Gram minor \(b_0\) nonzero at \(x\)
restricts the support to its irreducible locus. On this open,
dense good arithmetic points of the ordinary Hecke component
have positive selected multiplicity. Thus \(b_0a\in J\) for
\(a\in J'\), by reducedness and density in the Hecke quotient.
The local annihilator gives \(D(b_0a)=0\). Since \(a(x)=0\) and \(b_0(x)\ne0\), 
\[
 0=D(b_0a)=b_0(x)D(a)+a(x)D(b_0)=b_0(x)D(a),
\]
 so \(D(a)=0\).

\proofstep{Separating the eigenvalues in the fixed-weight fibre.}
On the finite free projected weight module choose a linear
combination \(t\) of Hecke generators primitive for its
generic reduced algebra and separating the systems in the
weight fiber at \(x\). Finitely many proper linear conditions
on the coefficients suffice for both separations in
characteristic zero. Let \(t\) be this separating combination of Hecke generators. The characteristic polynomial in the generic algebra has the factorization
\[
 P_t(X)=\prod_i g_i(X)^{m_i},\qquad \gcd(g_i,g_j)=1\quad(i\ne j),
\]
with the \(g_i\) distinct monic irreducibles. They lie in \(\Lambda[T]\) by integral closedness.
Let \(k_i\) be the multiplicity of \(x(t)\) as a root of
the specialization of \(g_i\). If \(k_i>0\), lying over in
the corresponding finite integral component gives a point
with these values of \(\Lambda,t\). Separation identifies it
with \(x\).

On that component choose a good arithmetic point
\(y\) with the same algebraic exponents and with
\(b_0g_i'(t)\prod_{j\ne i}g_j(t)\ne0\). This uses density
at those exponents with varying finite-order parts. The
displayed element is generically nonzero there by separability.
The selected joint dimension at \(y\) is at most the root
multiplicity \(m_i\), and at \(x\) it is the whole root
multiplicity since \(t\) separates at \(x\). Consequently
\[
       \sum_i m_i k_i=\mu_e(x)=\mu_e(y)\le m_i
                              \qquad(k_i>0).
\]
Some \(k_i\) is positive. Exactly one is positive and it is
one. The square-free product \(\prod g_i(t)\) belongs to \(J'\)
and has a simple root at \(x(t)\). Differentiating gives
\(D(t)=0\). The component ideal for that root is killed by
\(D\) as well, since multiplying each element by
\(\prod_{j\ne i}g_j(t)\), a unit at \(x\), puts it in \(J'\).

For \(a\) in the component algebra define
\[
          P_a(T)=\operatorname{Tr}\bigl(a\,g_i(T)/(T-t)\bigr)
                     \in\Lambda[T].
\]
Here the trace is from the generic component field, and
integrality of the coefficients follows from monic division
by \(T-t\) and integral closedness. At \(T=t\) the identity
embedding contributes \(a g_i'(t)\) and all other embeddings
contribute zero. Differentiate this identity at \(x\). The coefficients and \(t\) have zero derivative. Differentiating the surviving term gives 
\[
 g_i'(x(t))D(a)=0,\qquad g_i'(x(t))\ne0,
\]
 and hence \(D(a)=0\). Thus the derivation
vanishes on the whole determinant ring and on its image \(B_3\).

At a type-I frame point there are nonzero specialized entries
of both signs and one of odd weight. Products give two
monomials \(p_+,p_-\) of weights \(1,-1\) nonzero there.
A derivation vanishing on \(B_3\) is determined by its value
on \(p_+\). Write \(\xi\) for the frame point above \(x\). If \(D(p_+)=0\), applying \(D\) to \(p_+p_-\) gives \(D(p_-)=0\). For an entry \(u\) of weight \(j\ne0\), take \(p=p_-\) if \(j>0\) and \(p=p_+\) if \(j<0\). Its balanced product then gives 
\[
 0=D(up^{|j|})=\xi(p)^{|j|}D(u),\qquad \xi(p)\ne0.
\]
 All generating entries
and ordinary characters have zero derivative. This proves the
upper bound one for the dimension of fixed-inertia derivations.

\proofstep{The lower bound from a cyclotomic ratio.}
For the lower bound, equality on an open inertia subgroup forces equality of cyclotomic exponents since \(\epsilon\) has infinite image there:
\[
 e(w_i-w_j)=1,\qquad (w_1,w_2,w_3)=(1,0,-1).
\] Thus \(e=1\) and
the two indices are adjacent. Normalize the adjacent
two-dimensional subquotient by its second character. Its
extension cocycle \(b\) takes values in \(E(1)\), with \(E\)
a finite coefficient field. Its class belongs to the
Bloch--Kato finite subspace. Indeed the subquotient and its
characters become crystalline after finite extension and
restriction on cohomology with crystalline period coefficients
is injective by averaging in characteristic zero. By Kummer
theory this finite subspace consists of local units, annihilated
under cup product by the unramified line in \(H^1(F_v,E)\).
For a nonzero unramified additive character \(u\) take a
cochain satisfying
\[
      \beta(st)=\beta(s)+\epsilon(s)\beta(t)-b(s)u(t).
\]
Scaling \(u,\beta\) together makes them integral by compactness.
\cite{bloch-kato-finite}

Write \(A_k=\mathcal O_E/\ell^k\), \(q_k=|A_k|\), and work
in a stable ordinary integral basis. At each \(v\mid\ell\)
the local Euler characteristic and Tate duality give at least
\(q_k^{6(1+d_v)}\) upper triangular adjoint cocycles before
fixing inertia. At each of the \(d_v\) pro-\(\ell\) unit generators, three diagonal coordinates are fixed. Thus the maximum cost is
\[
 (q_k^3)^{d_v}=q_k^{3d_v}.
\]
The marked local fields have no \(\ell\)-power roots of unity.
At the distinguished place impose equal variations on the
adjacent diagonals. The resulting module has rank five and
projects equivariantly onto \(A_k(1)\) via the adjacent entry,
so the dual invariant factor has size at least \(q_k\). Only
two inertia characters must be killed. With the two diagonal variations equal, the lower bound for the number of solutions simplifies to
\[
 q_k^{5(1+d_v)+1-2d_v}=q_k^{6+3d_v}.
\]

One further integral cocycle has unramified diagonal difference
\(u\) on the resonant pair and zero difference on the other
adjacent pair. Its resonant entry is supplied by \(\beta\). Then
the remaining error is a two-cocycle in the outer root
\(\kappa_0/\kappa_2\). This differs from \(\epsilon\)
since the other adjacent ratio is nontrivial at the marked
regular element. With \(\theta=\kappa_0/\kappa_2\), local duality identifies the obstruction space by 
\[
 H^2(G_v,E(\theta))^\vee\simeq H^0(G_v,E(\epsilon\theta^{-1}))=0,
\]
 since \(\epsilon\theta^{-1}\) is nontrivial. Correct the outer entry and clear denominators
once.

Multiples of this integral cocycle at depth \(k\)
supply at least \(q_k/|\mathcal O_E/(c)|\) different cosets
of the equal-diagonal solutions, where \(c\ne0\) is a fixed
value of its diagonal difference. Adding all full-matrix
coboundaries gives another factor at least \(q_k^3\) at
every place, since distinct residual eigenvalues at the
marked local element force the generator of any coboundary
that stays upper triangular to be upper triangular. Thus
the allowed local cocycles \(Z_v\) satisfy
\[
  \prod_{v\mid\ell}|Z_v|
        \ \ge\ q_k^{9\,\#\{v\mid\ell\}+3d+1}/C
\]
with \(C\) independent of \(k\) for this point.

\proofstep{The global cardinality estimate and contradiction.}
The global Selmer cardinality inequality for these local
conditions gives
\[
 |Z_{\rm glob}|\, q_k^{9\,\#\{v\mid\ell\}}
             \prod_{v\mid\infty}|H^0_v|
      \ \ge\ q_k^9\prod_{v\mid\ell}|Z_v| .
\]
The dual global invariant factor is trivial here. The
residual frame separates each matrix unit's adjoint character
from \(\bar\epsilon\), using \(1,\bar\Psi^{\pm1}\) on the
positive component and \(\eta\bar\Psi^{\pm2}\) on the outer
roots. Nakayama shows that the images of
\(\operatorname{Ad}\rho(s)-\epsilon(s)\) span the truncated
module, so no equivariant dual functional survives. At each real place the fixed matrices identify with the skew matrices; there are three free entries.

Taking the product over the \(d\) real places gives
\[
 \prod_{v\mid\infty}|H_v^0|=(q_k^3)^d=q_k^{3d}.
\] Before normalization the global cocycle bound is \(q_k^{10}/C\). Imposing zero
at conjugation costs six scalar entries by its square relation.
For \(Y=z(g_1)\rho(g_1)\) the remaining equations set the
off-diagonal part of \(Y-JY^tJ\) to zero, with \(J\) the
antidiagonal matrix. The two outer corners already vanish
and the other four entries form two opposite pairs. Thus, writing \(Z_k\) and \(Z_k^{\rm norm}\) for the cocycles before and after imposing the two normalizations, respectively, we obtain
\[
 |Z_k^{\rm norm}|\ge q_k^{-6-2}|Z_k|
              \ge q_k^{-8}q_k^{10}/C=q_k^2/C .
\]

Each normalized cocycle deforms the representation by
\((1+\delta z(s))\rho(s)\), \(\delta^2=0\), with the
prescribed ordinary conditions, and hence induces a distinct
fixed-inertia point derivation into \(A_k\). The relative
cotangent module after imposing those inertia equations is
finite over \(\mathcal O_E\). Write \(M\) for this module and \(T\) for its torsion submodule. If \(\operatorname{rk}M\le1\), the elementary divisor decomposition over the coefficient DVR gives
\[
  q_k^2/C\le |\operatorname{Hom}_{\mathcal O_E}(M,A_k)|
      \le q_k |T| .
\]
Here \(C\) and \(T\) are fixed as \(k\) varies. Since \(q_k\to\infty\), the rank is at least two. This contradicts
the upper bound on the very same framed tangent space and
excludes the cyclotomic equality. The constants in this
depth argument can depend on the chosen point and the field.

At a type-I point of the even-level Field used in the rank-three trace estimate over the
imaginary quadratic field \(K\), the marking for this use takes
\(M=\ell N|\operatorname{disc}K|\). On complex conjugation \(c\) the two parity factors cancel:
\[
 (\epsilon^e\psi)(c)=(-1)^e\psi(-1)=(-1)^{2e}=1 .
\] Every prime ramifying in the quadratic
character divides \(M\). Factor the integral point through
the integers of a finite coefficient extension. Its ordinary
flag pulls back to the marked groups and its restriction to
the open CM subgroups is irreducible by the type-I host
hypothesis.

Given the automorphic attachment on this field, the preceding result supplies the unlinked local flag at the even level. The resulting local nonresonance and the three distinct weights give the required freedom of the ordinary characters. The original packet level in this branch
is \(N\), while \(M\) is the ramification cutoff of the comparison.
For the auxiliary field over \(\mathbb Q(i)\), use the same local exclusion after lifting to that field in the rigidity and trace bounds. Thus the local calculation is shared
even though the host attachment and the two-dimensional
constituent of the other Case of the rank-three trace estimate require separate
arguments.
\end{proof}

\section{The finite algebra at the marked point}
\label{sec:marked-hosts}

The support comparison in Section~\ref{sec:linked-components} took us to an abelian CM
extension \(L\). We now need a point there that will contribute to a length computed
over the original coefficient DVR \(\mathcal O\).
Two base changes intervene. The residual curve
points have to descend from the enlarged frame, and a point of
the graded algebra in characteristic zero has to give an integral
representation to which the automorphy comparison applies.
The marked point algebra gives the integral model used for the
descent and for the length bound in the next two appendices. The separated frame of Proposition~\ref{prop:separated-marked-host} carries the mark \eqref{eq:tag-h1}. Its degree-zero algebra over the original coefficient DVR is the object of Proposition~\ref{prop:marked-invariant-algebra}.

\paragraph{The fields and the mark.}
The field and its mark are specified by characters. With
\(e\ge1\) and \(\psi\) of even modulus \(N\) satisfying \(\psi(-1)=(-1)^e\), form the characters
\[
 \Psi=\epsilon^e\psi,\qquad P=\bar\Psi,\qquad \Psi(c)=1 .
\] Here \(c\) denotes complex conjugation.
Let \(\eta\) cut out the imaginary quadratic field \(K\). The
linking cutoff is \(\ell N|D_K|\) in the general-data case and
\(\ell N\) for \(K=\Q(i)\). Indeed an inertia group on which
\(\eta\) is nontrivial lies over a divisor of \(D_K\), and the
Gaussian character has conductor four. Let \(B\) be the chosen
place over \(\ell\). The marking uses
\[
       g_*\in D_B\cap G_L,\qquad P(g_*)^2\ne1 .
       \tag{H1}\label{eq:tag-h1}
\]
The level of the final packet will be read from the inertia of
the original representation.

\begin{proposition}[A separated marked frame]\label{prop:separated-marked-host}
Under the general-data or Gaussian host hypotheses, the
construction yields a totally real abelian \(L^+\) and its
CM extension \(L=L^+K\). The extension \(L/L^+\) splits
at the places above the primes dividing \(M\) and is unramified
at every finite place. At every place \(v\mid\ell\) of \(L^+\),
\[
                     [L_v^+:\mathbb Q_\ell]>3C_0+7.
\]
The distinguished mark satisfies
\textup{\eqref{eq:tag-h1}}. Both signed residual extension classes with
characters \(\eta P^{-a}\), \(a=1,2\), restrict nontrivially
to the marked frame with its prescribed inertia conditions.
\end{proposition}

\begin{proof}
\proofstep{Matching the quadratic characters locally.}
Here are the choices in the construction of the frame underlying
\eqref{eq:tag-h1}. Write \(K=\Q(\sqrt{-D})\), \(D>0\), and denote the linking
cutoff by \(M\). Choose a prime \(p>4D+M\) congruent to \(-1\)
modulo \(4M\). The real quadratic field
\(\Q(\sqrt{Dp})\) has the same local character as \(K\) on
decomposition groups above \(M\). Their defining square classes have ratio \(Dp/(-D)=-p\). For an odd prime \(v\mid M\), the congruence \(-p\equiv1\pmod v\) lifts the simple root \(1\) of \(X^2+p\) to \(\mathbb Z_v\). If \(2\mid M\), then \(8\mid4M\). Put \(u=-p\equiv1\pmod8\) and start with \(x_3=1\).

For \(n\ge3\),
\[
 (x_n+e2^{n-1})^2\equiv x_n^2+e2^n x_n\pmod {2^{n+1}} .
\]
If \(x_n^2\equiv u\pmod {2^n}\), exactly one of \(e=0,1\) corrects the next bit because \(x_n\) is odd. The resulting sequence converges to a square root in \(\mathbb Z_2^\times\).
This real quadratic field also ramifies at \(p\), where \(K\)
is unramified. Denote its fixing group by \(H_0\). The general
frame data include an abelian quotient of \(\ell\)-power order
with kernel \(H_b\) on which the away-from-\(\ell\) inertia
character has the same kernel as its reduction. On the Gaussian
branch \(H_b\) kills the \(\ell\)-power part of the character
after removing the factor at the place \(\ell\).

On inertia at \(\ell\), the condition \(P=\bar\epsilon^j\) determines the reduction of the finite factor. Use the Teichm\"uller lift
\[
 \chi'=[P\bar\epsilon^{-j}],\qquad\bar\epsilon^j\bar\chi'=P .
\]
The finite characters \(\chi'\) and \(\eta\) factor through
roots of unity of order dividing some \(M^a\ell\). Both are
unramified at \(\ell\), so its \(\ell\)-power factor can be
removed from that modulus by cyclotomic inertia. Choose a
second prime \(w\) congruent to one modulo both the resulting
modulus and \(M^{a+1}\ell\), outside the divisors of those
integers. The degree-\(\ell\) subfield of \(\Q(\mu_w)\) gives
the order-\(\ell\) character \(\theta\) of the seed construction.

For \(m\ge1\), let \(H_w\) fix that subfield, and \(H_m\) fix the
degree-\(\ell^m\) layer in \(\Q(\mu_{\ell^{m+1}})\). Put
\[
       U_m=H_0\cap H_w\cap H_b\cap H_m,\qquad
       G_L=U_m\cap\ker\eta .
       \tag{H2}\label{eq:tag-h2}
\]
All quotients defining \(U_m\) are abelian and complex conjugation
lies in each kernel. Thus \(U_m=G_{L^+}\), \(L=L^+K\). The
local comparison for \(H_0\) gives splitting of \(L/L^+\) over
\(M\). Outside \(M\), \(\eta\) is unramified, so \(L/L^+\) is
everywhere unramified at finite places. The layer at \(w\)
kills \(\theta\) on \(G_L\).

\proofstep{Selecting the real abelian layers.}
For the order of the choices in \eqref{eq:tag-h2}, let \(E\) be fixed by
\[
       H_0\cap H_w\cap H_b\cap\ker\eta\cap\ker P
                    \cap\ker\bar\epsilon .
\]
It is abelian over \(\Q\) and contains \(\mu_\ell\). On adjoining
the \(m\)-th layer the corresponding intersection fixes
\(E(\mu_{\ell^{m+1}})\). The bounded \(\ell\)-rank of the class
groups in this tower follows from Ferrero--Washington for the
cyclotomic tower of an abelian field and the Iwasawa module
structure theorem \cite[Chapter~13]{washington-cyclo}. Let \(E_\infty\) be their union and \(X\) its unramified abelian pro-\(\ell\) Galois module. The vanishing of \(\mu\) makes \(X\) finitely generated
over \(\mathbb Z_\ell\). Above a sufficiently deep layer \(E_m\),
one of the primes over \(\ell\) is totally ramified in
\(E_\infty/E_m\), since cyclotomic inertia has open image.

Write \(H_m^{\rm ur}\) for the unramified abelian
\(\ell\)-class field of \(E_m\). Every nontrivial finite intermediate extension in this part of the tower is ramified at the totally ramified prime over \(\ell\). Since \(H_m^{\rm ur}\) is unramified there,
\[
 H_m^{\rm ur}\cap E_\infty=E_m .
\] Unramifiedness survives base
change, so restriction gives
\[
 X\twoheadrightarrow\operatorname{Gal}
            (H_m^{\rm ur}E_\infty/E_\infty)
      \simeq\operatorname{Gal}(H_m^{\rm ur}/E_m).
\]
The \(\ell\)-rank is bounded by the number of
\(\mathbb Z_\ell\)-generators of \(X\). The finite
set of earlier layers can be included in the same bound.
By unramified class field theory this bounds
the dimension of the unramified mod-\(\ell\) character space by
a single \(C_0\) for all \(m\). Only now take
\(\ell^m>3C_0+7\). Inertia at \(\ell\) surjects onto the
degree-\(\ell^m\) layer, so every relevant local degree is at
least \(\ell^m\).

To construct the element in \eqref{eq:tag-h1}, choose \(x\in D_B\cap\ker\eta\) with \(P(x)^2\ne1\) as in the local separation hypothesis.
It fixes \(\sqrt{Dp}\) by the local comparison. Raise \(x\)
to a common \(\ell\)-power killing the three remaining finite
quotients in \(U_m\). Write \(d=\operatorname{ord}(P(x)^2)>1\). Since \(\gcd(d,\ell)=1\),
\[
 \operatorname{ord}\bigl(P(x)^{2\ell^a}\bigr)=d.
\]
Thus the resulting power \(x_*=x^{\ell^a}\) satisfies
\[
 x_*\in D_B\cap G_L,\qquad
 P(x_*)^2=(P(x)^2)^{\ell^a}\ne1 .
\]
The second equality uses the prime-to-\(\ell\) order of \(P(x)^2\),
on whose cyclic subgroup the indicated power is an automorphism.

\proofstep{Local separation by inertia and Frobenius in the general and Gaussian cases.}
Here are the two separation tests used above. For the General case of the rank-three trace estimate, inertia at the ramified prime of \(K\) away from
\(\ell N\) gives an element \(s\) with
\(P(s)=\bar\epsilon(s)=1\), \(\eta(s)=-1\). On inertia at \(\ell\) the residual finite factor has square one and \(\eta=1\). Thus here
\[
 P^2=\bar\epsilon^{2e}\ne1,\qquad \ell-1\nmid 2e .
\]

On the Gaussian branch let \(2^c\) be the
conductor of the dyadic component \(\psi_2\). Here \(c\ge4\).
Its image on units congruent to one modulo four is a nontrivial
cyclic 2-group. Some such unit maps to \(-1\). Multiplying
\(-1\) by that unit if \(\psi_2(-1)=-1\) gives a unit
congruent to \(-1\) modulo four in the kernel of \(\psi_2\).
Its image under dyadic reciprocity gives the required \(s\).
In either case \(\chi'(s)=1\) and \(\eta(s)=-1\).

For the other Gaussian test, the base data give an integer
\(o\ge1\) such that the odd-prime components of the
prime-to-\(\ell\) part of \(\psi\) have order dividing \(2o\).
For each exceptional prime and for \(p=\ell\) the dyadic conductor bound reads
\[
 8op^2<2^c .
\] These primes are odd. Choose a
Frobenius lift \(h\) acting by \(p\) on prime-to-\(p\) roots
and with \(\bar\epsilon(h)=1\). At \(p=\ell\) use cyclotomic
inertia to arrange the latter condition. At the other primes
it follows from \(p\equiv1\pmod\ell\). If \(P(h)^4=1\),
reduction is faithful on prime-to-\(\ell\) roots of unity,
so raising to the power \(4o\) gives
\(\psi_2(p)^{4o}=1\).

Let \(s_2\) and \(d\) be the two valuations that determine the effect of raising the odd prime \(p\) to successively doubled exponents:
\[
 s_2=v_2(p^2-1),\qquad d=v_2(o).
\]
For \(a\in1+4\mathbb Z_2\) one has \(v_2(a+1)=1\), so
\[
 v_2(a^2-1)=v_2(a-1)+v_2(a+1)=v_2(a-1)+1 .
\]
The geometric sum for an odd power is a unit. Starting
from \(a=p^2\) and using \(4o=2\cdot2^{d+1}(o/2^d)\)
therefore gives
\[
               v_2(p^{4o}-1)=v_2(p^2-1)+1+d .
\]
The same order formula in the finite dyadic quotients
shows that \(p^{4o}\) generates the subgroup congruent to
one modulo \(2^{s_2+1+d}\) in each such quotient. This forces
\[
 2^c\le2^{s_2+1+d}\le 2o(p^2-1),
\] a contradiction.
Since \(x=h^2\), multiplicativity of all three characters gives
\[
 \eta(x)=1,\qquad P(x)^2=P(h)^4\ne1,\qquad\bar\epsilon(x)=1 .
\] At \(p=\ell\) this is the starting
element for \eqref{eq:tag-h1}.

\proofstep{The signed residual cocycles on the marked frame.}
One use of these tests is the construction of the two
residual extension classes on the marked frame. For
\(a=1,2\) put \(\kappa=\eta P^{-a}\). This character is
odd, and differs from \(\bar\epsilon\) on \(s\) above.
It is nontrivial on \(U_m\cap D_B\) by \eqref{eq:tag-h1}. Let \(S\)
consist of \(\ell,w\) and the ramified primes of \(\kappa\).
For the Selmer group of \(k(\kappa)\), the local subspace in \(H^1(\mathbb Q_v,k(\kappa))\) is 
\[
 L_v=\begin{cases}H^1(\mathbb Q_v,k(\kappa)),&v=w,\\ 0,&v\in S\setminus\{w\},\\ H^1_{\rm unr}(\mathbb Q_v,k(\kappa)),&v\notin S.\end{cases}
\]
 The dimension difference
with the dual Selmer group is one by the Greenberg--Wiles
formula. Indeed both global invariant spaces vanish, and
the real invariant space vanishes. At \(w\) both \(\kappa\) and \(\bar\epsilon\) are trivial. The local contribution of the full condition is therefore
\[
 \dim L_w-\dim H^0_w=2-1=1 .
\]

At every other prime of \(S\) the
invariant space is zero. A nonzero resulting class remains
nonzero on \(U_m\) by inflation--restriction, since the
coefficient module has no \(U_m\)-invariants. On inertia
away from \(\ell,w\) in \(\ker\eta\cap\ker P\) its cocycle
is zero, either by the unramified condition or by local
triviality and \(\kappa=1\) on that inertia. At \(w\) the
restriction to inertia is additive. Its \(\ell\)-primary
tame quotient is procyclic and the degree-\(\ell\) field
fixed by \(H_w\) is totally ramified there. Thus it vanishes
on the inertia in \(U_m\) at \(w\) as well. Its restriction
at \(B\) is a coboundary. These are exactly the local
conditions needed to descend the two classes to the frame
quotient.

\end{proof}

\input{fitted-arithmetic}

\paragraph{A characteristic polynomial on every test curve.}
To control points of the degree-zero algebra we first examine the characteristic
polynomial along a curve in characteristic \(\ell\). Let \(A\) be a complete DVR of
characteristic \(\ell\) with finite residue field and a local map from \(\mathcal O\) through
its residue field \(k\). Take a marked ordinary representation
over \(A\) reducing in the given basis to the diagonal model
\((P,1,P^{-1})\) on \(\ker\eta\). On \(G_{L^+}\) it factors
through the marked quotient. Indeed the killed inertia elements
away from \(\ell\) have \(\eta=P=1\), and the inertia condition
on the point kills these elements. Continuity kills the closed
normal subgroup they generate. Above \(\ell\) the ordered
inertial characters on the frame are \(\Psi^{w_i}\), now
viewed in \(A\), where \(w=(1,0,-1)\). They are one on the
kernel of \(P\) in the marked inertia.

\begin{proposition}[Characteristic polynomials on marked ordinary curves]\label{prop:marked-residual-polynomial}
Let \(A\) be a complete DVR of characteristic \(\ell\) with
finite residue field, with \(\mathcal O\to A\) factoring through
\(k\). Take a marked ordinary curve over \(A\) with residual diagonal
\((P,1,P^{-1})\) on \(\ker\eta\) and ordered ordinary characters
\((\Psi,1,\Psi^{-1})\) on the marked inertia. On the entire quadratic kernel the identity \eqref{eq:tag-h3} reads
\[
 \det(X-r(g))=(X-P(g))(X-1)(X-P(g)^{-1})\qquad(g\in\ker\eta).
\] In the
fixed marked basis the curve is upper or lower triangular,
and its diagonal on the positive component is prescribed
by the residual characters.
\end{proposition}

\begin{proof}
\proofstep{Applying the rank-three support result to the residual curve.}
On the rank-three frame let \(J\) be the image of the Hecke
ideal and \(I_E\) the ideal pulled back from the outer corner
on the exceptional locus. The intersection \(I=J\cap I_E\)
lies below every prime off the two triangular loci. This follows
from propagation on the type-I locus and the rank-two comparison
on the exceptional locus. Write \(\lambda\) for the ideal
generated by the maximal ideal of the coefficients and the
weight differences \(\psi_{vi}(s)-1\) on inertia with \(P(s)=1\).
Both \(J+\lambda\) and \(I_E+\lambda\) have radical the
degree-zero maximal ideal.

For a prime \(\mathfrak p\) the inclusion \(JI_E\subseteq I\) implies
\[
 I+\lambda\subseteq\mathfrak p\quad\Longrightarrow\quad J+\lambda\subseteq\mathfrak p\ \text{ or }\ I_E+\lambda\subseteq\mathfrak p .
\]
Both force the same degree-zero maximal ideal.
The curve kills \(\lambda\). If its kernel is off both triangular
loci, its degree-zero evaluations consequently equal their
residual constants. In the expansion of a characteristic
polynomial each cycle of off-diagonal entries has total weight
zero and vanishes at the residual point. Only products of the
prescribed diagonal entries survive.

\proofstep{The split characteristic polynomial on the fixed subgroup.}
On a triangular curve the diagonal entries are characters.
Divide each by the corresponding residual character. The
quotient takes values in \(1+\mathfrak m_A\). On marked inertia with \(P=1\), any character value \(u\) of the quotient satisfies its ordinary polynomial. Since the coefficient ring here is a domain,
\[
 (u-1)^3=0\quad\Longrightarrow\quad u=1 .
\]
This kills the local abelian generators and, by polarization, their conjugates. At each truncation it is a
character into a finite \(\ell\)-group, and the abelian
class-field condition of the frame gives a single nonzero
exponent \(h\) for all these quotients. Passing through
all truncations gives \(z^h=1\) in \(1+\mathfrak m_A\).

In an equal-characteristic DVR that group is torsion-free.
The prime-to-\(\ell\) part follows by factoring a difference
of powers and the \(\ell\)-part by Frobenius. Thus the three
diagonal characters are fixed in this case as well. If the
comparison was made with enlarged residue coefficients, extend
the test DVR by the corresponding finite residue field
extension, apply the result there and use injectivity on
coefficients. Pulling the polynomial back along the map from
\(G_L\) to the positive component of the frame gives
\[
       \det(X-r(g))=(X-P(g))(X-1)(X-P(g)^{-1})
                        \quad(g\in G_L).
       \tag{H3}\label{eq:tag-h3}
\]
Here and below \(r(g)\) denotes the matrix part on the positive
component. The subgroup \(G_L\) in \eqref{eq:tag-h3} is independent of \(A,r\).

\proofstep{Extend the polynomial identity to the original quadratic kernel.}
For \(g\in\ker\eta\), take \(d\ge1\) with \(g^d\in G_L\).
Let \(\lambda_i\) be the eigenvalues of \(T=r(g)\) in an
algebraic closure of the fraction field, counted with multiplicity.
Triangularizing \(T\) gives the factorization
\[
 \det(X^d-T^d)
   =\prod_i (X^d-\lambda_i^d)
   =\det(X-T)
       \prod_i\left(\sum_{a=0}^{d-1}X^{d-1-a}\lambda_i^a\right).
\]
The first polynomial belongs to \(k[X]\) by \eqref{eq:tag-h3} on \(g^d\),
so all \(\lambda_i\), and hence the coefficients of \(\det(X-T)\),
are algebraic over \(k\). This factorization needs no separability
of \(X^d-\lambda_i^d\). An algebraic element of the domain \(A\)
congruent to \(a\in k\) must equal \(a\), since the field it
generates over \(k\) embeds into the residue field under
reduction. This proves \eqref{eq:tag-h3} on all of \(\ker\eta\), even when
\(\ell\mid d\).

\proofstep{Triangularity in the marked basis.}
We need triangularity in the \emph{marked basis}. Put
\(g^*=cg^{-1}c\). The marked difference in the frame is the diagonal matrix
\[
 r(g_*)-r(g_*^*)=\operatorname{diag}(z,0,-z).
\]
The reduction of \(z\) is \(P(g_*)-P(g_*)^{-1}\ne0\).
Work over the fraction field. An invariant subspace is stable under the marked difference \(D\). The three eigenvalues are distinct, so their coordinate projectors are polynomials in \(D\):
\[
 e_\lambda=\prod_{\nu\ne\lambda}\frac{D-\nu I}{\lambda-\nu},\qquad \lambda,\nu\in\{z,0,-z\}.
\]

Let \(L_1,L_2,L_3\) be the marked coordinate lines, in that order.
A stable subspace is a sum of its images under the three projectors.
Polarization makes its orthogonal complement for the reversed pairing
stable as well, and the complementary planes of the coordinate lines are
\[
 L_1^\perp=L_1\oplus L_2,\qquad L_2^\perp=L_1\oplus L_3,\qquad
 L_3^\perp=L_2\oplus L_3 .
\]
Thus a stable outer line gives a full flag. A stable middle line or
outer plane gives the decomposition \(L_2\oplus(L_1\oplus L_3)\).
If the outer plane is reducible, its two coordinate projectors again
give a stable outer line and a flag. The other possibilities are an irreducible triple
and a block decomposition with irreducible outer pair.

In
the former case Brauer--Nesbitt applied to \eqref{eq:tag-h3} would identify
the irreducible triple with a sum of characters. In the
latter case the middle diagonal entry is an eigenvalue,
algebraic over \(k\) and congruent to one, so it is one.
Since the middle eigenvalue is one, divide the characteristic polynomial in \eqref{eq:tag-h3} by its common factor:
\[
 P_{\rm outer}(X)=\frac{P_{\rm triple}(X)}{X-1}.
\]
The resulting character polynomials identify the semisimplification of the outer pair by Brauer--Nesbitt. Thus every curve is upper or lower triangular.
The equalities of entries descend to \(A\) by injectivity.
Each diagonal entry is then an eigenvalue congruent to its
prescribed value, so equals it on \(\ker\eta\). Multiplication
by \(c\), whose matrix part is one, extends the diagonal
identity to the other coset.

\end{proof}

\paragraph{Finiteness and an integral point in the same orbit.}
The identities on curves can now be used on the original Galois quotient. For the
auxiliary field over \(\mathbb Q(i)\) put \(\Gamma=G_{\Q,S}\), the quotient
by the closed normal subgroup generated by inertia away
from the specified finite set. The surjection
\(G_\Q\to\Gamma\) identifies the image of \(G_K\) with
\(\ker\eta_\Gamma\). The preceding triangularity and diagonal
identities therefore hold on the host. In the criterion for a triangular point the diagonal identity is also
obtained by dividing diagonal characters on \(G_K\).
The inertia conditions make their quotients unramified at
each truncation, and \(\Z[i]\) has class number one.

\begin{construction}
\[
\begin{array}{ll}
\text{monomial on the curve}&\text{reason it vanishes}\\ \hline
\text{contains a weight-zero variable}&\text{the diagonal is the fixed residue diagonal}\\
\text{uses only nonzero weights}&\text{both signs occur to sum to zero}\\
 &\text{one strict sign vanishes by triangularity}
\end{array}
\]
\caption{The closed-fibre test for a nonconstant monomial of weight zero at the Gaussian host. Triangularity is available on each marked residual curve.}\label{fig:marked-host-zero-invariants}
\end{construction}

\begin{proposition}[The integral invariant algebra]\label{prop:marked-invariant-algebra}
Let \(R\) be the complete normalized graded ring of the
Gaussian host just described, with the triangularity and
diagonal identities on all its residual test curves.
The degree-zero subalgebra \(R_0\) is a finite \(\mathcal O\)-module.
For each characteristic-zero point extending the coefficient embedding,
let \(a_i\) be its values on homogeneous generators of the maximal ideal,
of degrees \(d_i\). There exist a finite extension \(O'\) of the coefficient DVR and a grading-torus scalar
\(t\in\operatorname{Frac}(O')^\times\) for which
\[
                  t^{d_i}a_i\in\mathfrak m_{O'}\qquad(\text{every }i).
\]
\end{proposition}

\begin{proof}
\proofstep{Finiteness by testing curves in the closed fibre.}
In an adapted power series
presentation over \(\mathcal O\) take variables \(z_i\) from
entries of \(r-r_0\) on the positive component, with their
entry weights \(w_i\), where \(r_0\) is the split point. They
belong to the maximal ideal. On every residual test curve
the zero-weight variables vanish, as do all variables of one
strict sign. Indeed, if a nonconstant monomial of weight zero has a factor of weight zero it vanishes on the curve. Otherwise its exponents \(a_i\) satisfy
\[
 \sum_i a_iw_i=0,\qquad \sum_i a_i>0 ,
\]
so at least one variable of each strict sign occurs. The variables of one of those signs all vanish on the curve.

Write \(\varpi\) for the coefficient uniformizer. Suppose that a prime ideal
\(\mathfrak p\) contains \(\varpi R\) but does not contain one of these
monomials \(b\). In the complete local domain \(R/\mathfrak p\) we have
\[
 \bar b=b+\mathfrak p\ne0 .
\]
We claim that a further prime
quotient of dimension one can be chosen in which the image of \(b\)
is still nonzero.

This follows by choosing parameters
with \(b\) as the last parameter in the domain, cutting by
the others and taking a minimal prime avoiding \(b\).
Normalization of the resulting complete one-dimensional
local domain gives a complete DVR with finite residue field.
This would be a test curve with \(b\ne0\), contrary to the vanishing on every test curve. Thus some power \(b^m\) belongs to \(\varpi R\) in the ambient graded ring. Extract its degree-zero component:
\[
 b^m\in\varpi R\quad\Longrightarrow\quad b^m\in(\varpi R)_0=\varpi R_0 .
\] The degree-zero monomials topologically generate \(R_0\), since their semigroup is finitely generated. The radical test on these generators thus gives
\[
 \sqrt{\varpi R_0}=\mathfrak m_{R_0}.
\] Noetherianity
gives a power containment. The special fiber has finite
length and the two adic topologies agree. Lifting a basis
of that fiber successively through powers of \(\varpi\)
shows that \(R_0\) is finite over \(\mathcal O\).

\proofstep{An interval of allowable rescalings.}
The characteristic-zero point is initially an evaluation of
the algebra spanned by homogeneous pieces. Write \(a_i\)
for the values of the \(z_i\), with valuation \(v\). Finiteness
of the local algebra \(R_0\) implies \(v>0\) for the values
of its maximal ideal, with \(v(0)=+\infty\). Thus \(v(a_i)>0\)
for weight zero. For \(w_i>0>w_j\), the monomial
\(z_i^{|w_j|}z_j^{w_i}\) has strictly positive valuation.
Consequently the finitely many strict constraints
\[
             -v(a_i)/w_i < b < v(a_j)/|w_j|
                \qquad(w_i>0>w_j,\ a_ia_j\ne0)
\]
have a common rational solution, also when only one sign
occurs by keeping the corresponding one-sided constraints.
Choose \(t\in\overline{\Q}_\ell^\times\) with \(v(t)=b\).
All the \(t^{w_i}a_i\) lie in the maximal ideal of one finite
coefficient DVR. Evaluation of power series converges there.

\proofstep{Evaluation of homogeneous series.}
To check agreement with the given point on a homogeneous
series, decompose it into finitely many monomials of its weight
times series of weight zero. Such generators exist by the
coordinatewise minimal exponent vectors and Dickson's lemma.
Evaluation is continuous on \(R_0\) by finiteness and the
adic power containment. The two evaluations on weight \(j\)
therefore differ exactly by \(t^j\). In matrix terms the
integral representation and the given point are conjugate
by \(\operatorname{diag}(t,1,t^{-1})\), with the semidirect
conjugation law on the negative coset. Both use the image of
the \emph{same} \(g_*\).
\end{proof}

\section{Descent from the marked field to the two ordinary levels}
\label{sec:host-descent}

The preceding appendix produced integral points on the larger marked field. Which
constituents can they supply at the smaller ordinary levels where we will evaluate
Frobenius traces? We determine the irreducible constituents
after restriction and keep track of the ordinary flags and
coefficient types in descent. The two possibilities for a
rank-three point give different packets in the length
calculation that follows. Lemma~\ref{lem:rank-three-host-restriction} controls which constituents remain irreducible on restriction. Proposition~\ref{prop:ordinary-host-attachment} returns the corresponding cuspidal systems to their original fields.

Retain the quadratic field \(K\), the marked element and the support conditions of Section~\ref{sec:marked-hosts}. Write \(H=G_K\) for the kernel of the quadratic character defining \(K/\mathbb Q\). At a type-I point the rank-three representation is irreducible on \(H\). At a type-II point for the general datum, let \(W\) be the outer plane and \(U\) the middle line. On \(H\) the preserved decomposition is
\[
 V=W\oplus U,\qquad \dim W=2,\quad\dim U=1,\quad W|_H\text{ irreducible}.
\] The length argument distinguishes the packets
coming from these two cases. We first determine the group on
which irreducibility actually holds.

\paragraph{Restriction of the rank-three points.}
Start with restriction, where an irreducible representation may acquire several
constituents. Use the integral representative of a type-I point and its
ordinary flag at \(B\). The residual trace on \(H\) is
\(P+1+P^{-1}\). Given an open subgroup \(V\) of \(H\), its normal core is the kernel of the finite permutation action:
\[
 U=\bigcap_{h\in H}hVh^{-1}=\ker\bigl(H\longrightarrow\operatorname{Perm}(H/V)\bigr).
\] In \(U\cap I_B\) choose \(u\)
with \(\psi(u)=1\) and \(\epsilon(u)\) of infinite order.
Indeed rational cyclotomic inertia surjects onto
\(\mathbb Z_\ell^\times\). Start with a lift of \(1+\ell\)
and take a power to enter the kernels of \(\eta,\psi\) and
the open core. In an ordinary basis the three eigenvalues
of \(u\) are \(\epsilon(u)^e,1,\epsilon(u)^{-e}\). A further
upper triangular change makes this matrix diagonal.

\begin{lemma}[Absolute irreducibility after restriction]
\label{lem:rank-three-host-restriction}
Let \(r\) be a type-I rank-three point with the integral
ordinary flag at \(B\) and residual trace \(P+1+P^{-1}\) on the
quadratic kernel \(H\) of Appendix~\ref{sec:linked-seeds}.
Assume the marked cyclotomic and local inertia data at that level.
The two restriction statements are
\[
 U\le H\text{ open}\ \Longrightarrow\ r|_U\text{ absolutely irreducible},
 \qquad r|_{\ker(\epsilon|_H)}\text{ absolutely irreducible}.
\] The associated
rank-two corner at the two-block points is treated separately
in the ordinary support comparison.
\end{lemma}

\begin{proof}
\proofstep{The isotypic summands after restriction.}
If restriction to \(U\) is reducible, it is semisimple because
the sum of the \(H\)-translates of a simple \(U\)-subspace
is the whole space, and a sum of simples is semisimple.
The isotypic summands form a single orbit. Since three
is prime, either \(U\) acts by scalars or there are three
distinct character lines permuted transitively. The first
case contradicts the spectrum of \(u\). In the second case
the lines are precisely the three eigenlines of \(u\).
Every matrix of \(H\) is monomial in this basis. The
separating element \(x\in D_B\cap H\) constructed above is
also upper triangular, so fixes each line. A transitive
subgroup of \(S_3\) contains a three-cycle.

For a lift \(g\)
of this cycle, \(gx^n\) has trace zero for every \(n\).
Reducing the trace in the original integral model and writing
\(a=P(g)\), \(b=P(x)\) gives
\[
         (ab^n)^2+ab^n+1=0\qquad(n=0,1,2).
\]
The three roots are distinct since \(b^2\ne1\), a
contradiction in every residue characteristic. Thus the
representation \(\rho\) is irreducible on every open
subgroup of \(H\), in particular on \(G_L\).

\proofstep{Restriction to the full cyclotomic kernel.}
It remains irreducible on \(H\cap\ker\epsilon\). Here is a
characteristic-zero argument for this further restriction.
Let \(G\) be the Zariski closure of \(\rho(H)\) over the
algebraically closed coefficient field. The preimages of
its finitely many components are closed and open, so its
identity component \(G^\circ\) acts irreducibly. Let \(U\) be the unipotent radical and \(V\) the representation space. Unipotence gives a nonzero fixed vector and normality gives stability, so irreducibility implies
\[
 0\ne V^U\subseteq V,\qquad G^\circ V^U=V^U\quad\Longrightarrow\quad V^U=V .
\] Hence the radical is trivial.
Schur's lemma makes the center scalar. As a connected
reductive group is the product of its connected center and
derived group, the derived group acts irreducibly. The
Zariski closure of \(\rho(H\cap\ker\epsilon)\) contains it since
commutators of the dense subgroup belong to that closure
and the commutator map is algebraic. This proves the claim.

\end{proof}

\begin{proposition}[Attachment on the original fields]
\label{prop:ordinary-host-attachment}
For the contributing type-I and outer-plane points, the
ordinary support comparison and cyclic descent yield the
cuspidal systems on the designated original fields, with the
specified shifts of weight and local level. In rank three the weight shift and multiplier are
\[
 r_3=\rho\otimes\epsilon^{-1},\qquad
 r_3(cgc)^{\mathsf T}Jr_3(g)=\epsilon(g)^{-2}J\quad(g\in G_K),
\]
where \(J\) reverses the marked basis. Each cyclic step uses irreducibility on the restricted group and the stipulated
local ordinary condition.
\end{proposition}

\begin{proof}
\proofstep{The ordinary cusps over \(L\).}
At an integral point, pass to a common finite coefficient
DVR with the seed before applying the support comparisons.
On \(G_{L^+}\) the point kills the required tame kernels
and has the prescribed flag characters on inertia, so it
defines a point of the framed ordinary quotient. In rank
three the restrictions of the positive-component matrices
agree with \(\rho|_{G_L}\). In rank two they agree with the
outer block at a type-II point. Choose a second omitted-prime pair whose omitted prime differs in each rank. Writing \(p_n,p_n'\) for the first and second choices in rank \(n\), the requirement is
\[
 p_2\ne p'_2,\qquad p_3\ne p'_3 .
\]
The comparison for the first pair gives the finite fiber
used in \eqref{eq:tag-h3}. Each run also puts the irreducible point of
the relevant rank on its own Hecke support.

A common
algebraic character twist makes its inertia weights a
classical specialization of the ordinary module. A
contributing ordinary eigenform then has attached
characteristic polynomials equal to those of the point
with that twist. A decomposition into proper cuspidal or Speh blocks would split the semisimple Galois representation at matching Frobenii into pieces of smaller rank. In particular its characteristic polynomial there would be the product
\[
 \det(1-X\rho^{\rm ss}(g))=\prod_i\det(1-X\rho_i(g)),\qquad \sum_i\dim\rho_i=\dim\rho,
\]
for the global block constituents \(\rho_i\). Irreducibility excludes this. Thus the
transfer is cuspidal. Equality at the common spherical
places and strong multiplicity one identify the two
transfers after adjusting the common twist. At either
omitted prime use the local comparison from the other
run. This gives the local support and size conditions of
the ordinary comparison wherever they are required.
Undoing an algebraic character twist translates the common
center of the size interval by its absolute-value exponent.

\proofstep{The weight shift of the rank-three system.}
In rank three write \(r_3=\rho\epsilon^{-1}\). On the negative
component the frame acts by inverse transpose relative to
the reversing matrix \(J\), and the matrix part at \(c\) is
one. On \(G_K\) this gives
\[
             r_3(cgc)^t J r_3(g)=\epsilon(g)^{-2}J .
\]
The twist thus gives the multiplier \(\epsilon^{1-3}\).
It preserves every comparison of ratios of graded
characters on a flag. The algebraic weights on inertia
become \((e-1,-1,-e-1)\).

\proofstep{The rank-two corner over \(\mathbb Q\).}
For general data the character on the middle line reduces to
one on \(H\). On inertia where \(P(h)\ne1\) it is a root of
the polynomial with roots \(\Psi(h),1,\Psi(h)^{-1}\). Only
the middle root reduces to one there. At a place above
\(\ell\), choose \(s\) in inertia with \(P(s)^2\ne1\).
When \(P(h)=1\), use multiplicativity on \(s\) and \(hs\). A representation matrix at \(h\) is recovered from those two by
\[
 B(h)=B(hs)B(s)^{-1} .
\]

At the other places
the inertia with \(P=\eta=1\) belongs to the prescribed tame
kernel. Thus the middle character is unramified everywhere.
It is one modulo the maximal ideal, and the class number
of the chosen \(K\) is prime to \(\ell\). The unramified
class field quotient forces the character to be one. The
determinant of the outer block has residue one, and is one
on every inertia group by the same three-root polynomial.
The same class group argument gives determinant one on \(H\).

\proofstep{Extension of the outer plane to the rational group.}
Write \(B(g)\) for this outer block and
\(D=\operatorname{diag}(-1,1)\). For a determinant-one
two by two matrix, inverse transpose conjugated by the
reversing matrix is conjugation by \(D\), as follows by
the two by two inverse formula. Thus polarization gives \(B(cgc)=D B(g)D^{-1}\).
With \(D^2=1\), the formula
\[
 W(hc^j)=B(h)D^j,\qquad
     h\in H,\ j\in\{0,1\},
\]
defines a continuous representation on \(H\rtimes\langle c\rangle\):
the product of the two matrices uses
\(D^j B(h') D^{-j}=B(c^jh'c^{-j})\).
Its determinant is one on \(H\) and \(-1\) at \(c\), hence is
\(\eta\).

It remains irreducible on \(G_L\). Otherwise the
normal restriction is semisimple and the marked element,
whose two residual eigenvalues differ, excludes scalar
restriction. There are two distinct lines, swapped by
some \(g\in H\). Then \(g\) and \(gg_*\) both have trace
zero. Reduction on the outer block forces
\[
 P(g)^2=P(gg_*)^2=-1,
\]
contradicting \eqref{eq:tag-h1}.
The rank-two run above consequently gives a regular
cuspidal system over \(L\) for this particular corner
after its algebraic twist.

\proofstep{Descent of the global system.}
Filter the abelian extension \(L/K\) by cyclic steps of
prime degree. Each intermediate field contains the
imaginary quadratic field and is abelian over \(\Q\),
hence is CM. At a step \(F'/F\) the given representation
exists on \(G_F\) and is irreducible even on \(G_{F'}\).
If \(\sigma\) generates the cyclic group, the good Frobenius polynomials of the attaching cusp \(\pi\) and of its conjugate agree. By strong multiplicity one this gives
\[
 {}^\sigma\pi\simeq\pi .
\] Cyclic
base change descends it to a cusp on \(F\). The infinity
parameter is regular algebraic because at a complex place the
local extension is trivial and the compatible parameters
agree. The Galois representation attached to this
descent agrees after restriction with the given one by
compatibility at split places, Chebotarev and
Brauer--Nesbitt.

Here degree-one primes over \(F\) are
dense for the required Frobenius comparison over \(F'\) since
the higher relative degrees have density zero.
Identify the two restrictions on \(G_{F'}\) using the
intertwiner. If the resulting representations are
\(r,\widetilde r\), then \(\widetilde r(g)^{-1}r(g)\)
commutes with their common action of \(G_{F'}\): both matrices
carry the \(h\)-action to the \(ghg^{-1}\)-action.
Schur's lemma gives a scalar \(\nu(g)\), and
\[
       r(g)=\nu(g)\widetilde r(g),\qquad
       \nu(gh)=\nu(g)\nu(h),\qquad \nu|_{G_{F'}}=1 .
\]
This character of the cyclic quotient gives the finite
Hecke twist that selects the specified descent.
The conjugate duality required of the descended cusp
follows by the same comparison at good places from the
given pairing. Iteration gives the rank-three polarized cusp
over \(K\) and the rank-two cusp for \(W|_{G_K}\). The
alternating rank-two pairing has multiplier \(\det W\), and
composing with \(W(c)^{-1}\) supplies the conjugate pairing.
\hfill\mbox{\cite{arthur-clozel,clozel-rajan,harris-lan-taylor-thorne}}

\proofstep{The additional rank-two descent.}
There is one further cyclic descent for the rank-two
system, from \(K\) to \(\Q\). The complex parameter has
distinct weights with integral difference \(2e>0\) and
finite determinant on the Galois side. Conjugation
invariance pairs the two characters of the complex Weil
group. With positive gap, the infinity parameter of the descent is the induction from one of the distinct conjugate Weil characters \(\xi,\xi^c\):
\[
 \operatorname{Ind}_{W_{\mathbb C}}^{W_{\mathbb R}}\xi,\qquad\xi\ne\xi^c .
\]

It therefore gives the discrete series up to twist. After the algebraic normalization
\(W(-e)\) the two cyclotomic exponents are \(0,-2e\),
the geometric convention for weight \(2e+1\). Write \(r_f\) for the classical newform representation before
choosing its descent twist. Polynomial comparison at degree-one primes
gives \(r_f|_{G_K}\simeq W(-e)|_{G_K}\). Applying the scalar
intertwiner calculation to this quadratic step gives
\[
 W(-e)\simeq r_f\otimes\nu,\qquad
 \nu\in\operatorname{Hom}(\operatorname{Gal}(K/\mathbb Q),
                      \{\pm1\})=\{1,\eta\}.
\] Keeping the particular
extension is needed for both off-diagonal Selmer summands.

\proofstep{The local condition through descent.}
Let \(\kappa_i\) be the graded characters of the flag and put
\(\delta_{ij}=\kappa_i/(\epsilon\kappa_j)\). In terms of their
images on the decomposition group, the two conditions read
\begin{equation}\label{eq:host-character-image-conditions}
 \begin{aligned}
 U_v         &\iff \delta_{ij}(G_{F_v})\ne\{1\} &&\text{for all }i,j,\\
 U_v^\infty &\iff |\delta_{ij}(G_{F_v})|=\infty &&\text{for all }i,j .
 \end{aligned}
\end{equation}
Jordan--Hölder makes the tests independent of the chosen flag. For a
finite local extension with subgroup \(H\subseteq G_{F_v}\),
\[
 [\delta_{ij}(G_{F_v}):\delta_{ij}(H)]\le [G_{F_v}:H]<\infty .
\]
Thus restriction preserves the second condition in both directions.
A nontrivial finite character can become trivial on \(H\), so the
first condition of \eqref{eq:host-character-image-conditions}
needs a separate argument on the chosen extension.

\proofstep{Ordinary flags under local restriction.}
First suppose \(U_v\) holds downstairs and the ordered
inertial characters are cyclotomic powers of strictly
decreasing exponents times finite characters of level
\(N\). The ordinary representation becomes crystalline
over the finite extension killing those inertia characters
of finite order. For the induction add one graded character
\(\kappa_j\) to a previously treated subrepresentation
\(V_0\). The representation
\(V=V_0\otimes\kappa_j^{-1}\) is de Rham downstairs by
finite descent and crystalline on that extension. Its
strictly positive cyclotomic exponents give
\(\operatorname{Fil}^0D_{\rm dR}(V)=0\).

Also \(H^0(V^*(1))=0\) by \(U_v\). Local duality and the Euler
characteristic formula give
\[
 \dim H^1(F_v,V)=[F_v:\Q_\ell]\dim V+\dim H^0(F_v,V).
\]
This equals \(\dim H_f^1(F_v,V)\) by the Bloch--Kato
exponential sequence. The extension class therefore
becomes crystalline where its ends are crystalline.
This proves the induction. The dimension calculation took
place downstairs, where the hypothesis on the characters
was imposed. \hfill\mbox{\cite{bloch-kato-finite}}

A forbidden equality of characters upstairs may be moved
to one of the frame places. The positive coset acts by
conjugation, the negative one by conjugation of the dual.
Conjugating transports an equality to the corresponding decomposition group. For duality the test is symmetric in the ordered pair of indices, since
\[
 \kappa_i^{-1}=\epsilon\,\kappa_j^{-1}
      \quad\Longleftrightarrow\quad
 \kappa_j=\epsilon\,\kappa_i .
\]
Both operations preserve potential crystallinity. The decomposition group of the resulting
frame place lies on the positive component by splitting in
\eqref{eq:tag-h2}. Its characters are precisely those of the restricted
ordinary flag, by comparison of characteristic polynomials
for all group elements. The even-level conclusion of
Section~\ref{sec:linked-components} at the potentially
crystalline point excludes the equality. This proves the
required persistence. In rank two the two exponents differ
by \(2e\ge2\), so the same assertion follows directly on
every open subgroup of inertia.

\proofstep{Nonresonance at the even level.}
There is a further use of the even-level result on the
original rank-three host. At a general-data place over \(\ell\), the local point factors through the quotient for the prescribed potentially crystalline type. If its ideal is \(I_{\rm pc}\) and \(x_v\) is the point, this says
\[
 I_{\rm pc}\subseteq\ker x_v,\qquad R_v\twoheadrightarrow R_v/I_{\rm pc}\xrightarrow{\ x_v\ }\overline{\mathbb Q}_\ell .
\] Its crystalline model on
the specified open subgroup is defined over a finite
coefficient extension by evaluating the local deformation
ring there. Conjugation and duality give the model at the
conjugate place.

If an equality occurred already over
\(K\), its restriction would occur over \(L\), and the
same frame test rules it out. A different auxiliary marked
element can be used for this test by diagonalizing the marked
difference with its distinct residues and using the paired
basis to retain the normalization at \(c\). The two
representations on \(G_K\) are isomorphic. On the
auxiliary field over \(\mathbb Q(i)\) one can test directly on the element \(x\)
at \(\ell\) of Section~\ref{sec:marked-hosts}. The three eigenvalues of the separator remain distinct in the residue field:
\[
 \{P(x),1,P(x)^{-1}\}\text{ has cardinality }3,\qquad \bar\epsilon(x)=1 .
\] An equality on two different lines would imply
\[
  \kappa_i=\epsilon\kappa_j
     \quad\Longrightarrow\quad
  \bar\kappa_i(x)=\bar\epsilon(x)\bar\kappa_j(x)=\bar\kappa_j(x)
       \qquad (i\ne j),
\]
contradicting distinctness.
For the same line the cyclotomic character is nontrivial
on local inertia. Conjugating the place preserves the test.
These arguments include \(e=1\).

For clarity, the local output glued from the two ordinary
runs gives \(U_{v'}\Rightarrow U_{v'}^\infty\) over \(L\).
For an attached cusp it gives a single real \(t\) such
that its local support at \(v'\nmid\ell\), and also at
\(v'\mid\ell\) satisfying \(U_{v'}^\infty\), consists of
characters \(\xi\) with
\[
           q_{v'}^{t-1/2}<|\xi(\varpi_{v'})|
                          <q_{v'}^{t+1/2}.
\]
The interval excludes ratios \(|\cdot|^{\pm1}\). Hence
the principal series is irreducible and its parameter
has zero monodromy. Persistence and the finite-index
criterion for infinite order descend the first assertion
about \(U^\infty\).

\proofstep{Descending the valuation interval.}
To descend the interval in a cyclic step, the upstairs
support characters on units must come from below by norm.
Away from \(\ell\) the downstairs inertia polynomial is
that of a list of global finite characters. Under the stated compatibility, equality of the inertia characteristic polynomials on the full upstairs inertia identifies the semisimplification. If \(\theta_i\) are the finite characters downstairs, at the places away from \(\ell\) this reads
\[
 r_{\rm up}|_{I_{v'}}^{\rm ss}\simeq\bigoplus_i\theta_i|_{I_{v'}} .
\] Above \(\ell\), when
the interval is needed, the representation is de Rham
and the upstairs Weil parameter is toral. Exactness of
potentially semistable comparison matches its characters
with the Weil characters of the graded lines. A de Rham
character under local reciprocity is locally algebraic.
Factor the locally algebraic character as \(\chi_{\rm alg}\chi_{\rm sm}\).

Under restriction from \(v\) to \(v'\), both factors are pulled back by the local norm \(N=N_{v'/v}\), so
\[
 (\chi_{\rm alg}\chi_{\rm sm})\circ N=(\chi_{\rm alg}\circ N)(\chi_{\rm sm}\circ N).
\]
Thus at either kind of place each upstairs character
extends by norm on units and extends on the full local
multiplicative group by choosing the value of a uniformizer. It can be
chosen by taking a root, accounting for its unit factor.
Equivalently the corresponding upstairs Weil character
extends to the downstairs Weil group. Local cyclic base
change restricts the Frobenius-semisimple parameter.
Monodromy downstairs is already zero by restriction.
Each isotypic summand upstairs is preserved downstairs,
since its character extends. Dividing by an extension,
one obtains a representation of the cyclic quotient,
which splits into characters in characteristic zero.
Thus the downstairs parameter is toral too.

For a place \(v'/v\) of residue degree \(f\), restriction of the local character is norm pullback by reciprocity:
\[
 \chi|_{G_{v'}}\ \longleftrightarrow\ \chi_v\circ N_{v'/v}.
\] For an inducing
character \(\alpha\), the norm of a uniformizer is a unit
times \(\varpi_v^f\), and the smooth character on units has
absolute value one. Consequently the upstairs bound reads
\[
 q_v^{f(t-1/2)}
   <|(\alpha\circ\mathrm N)(\varpi_{v'})|
      =|\alpha(\varpi_v)|^f
   <q_v^{f(t+1/2)} .
\]
Taking positive \(f\)-th roots gives exactly the original
bounds downstairs. Split places give the identity
comparison. These arguments iterate, and also transport
the rational rank-two comparison to an imaginary quadratic
field on which \(W\) remains irreducible, using the
newform's base change. \hfill\mbox{\cite{arthur-clozel,blggt-lg}}

These descents apply to the integral torus representatives
of Proposition~\ref{prop:marked-invariant-algebra}. The
following level calculation uses the original inertia
characters, so the resulting packet counts do not carry a
factor from the larger marked field.

\proofstep{Level on the original field.}
Away from \(\ell\) the type-I inertia is conjugate to
\(\psi,1,\psi^{-1}\) on \(H\). At \(\ell\) its flag has the same finite parts. Twisting by
\(\epsilon^{-1}\) therefore changes the ordered cyclotomic exponents by
\[
             (e,0,-e)\ \longmapsto\ (e-1,-1,-e-1).
\] For \(W\) at \(p\ne\ell\) the inertia
types on \(I_p\cap H\) are \(\psi^{\pm1}\). Each character
extends to \(I_p\), so the isotypic spaces are preserved.
If the subgroup has index two, the remaining generator
has scalar square on each such space, hence diagonalizes.
The extensions differ from the specified characters by
a power of \(\eta\). This gives the level
\(M_2=N|D_K|\).

At \(\ell\) a nonzero middle coordinate of the ordinary first line would intertwine the high inertia character \(\chi\) with the trivial one. If its value is \(a\), equivariance gives
\[
 (\chi(h)-1)a=0\quad(h\in I_\ell),\qquad \chi|_{I_\ell}\ne1\quad\Longrightarrow\quad a=0 .
\]
The first line therefore projects into the outer plane. They differ on an infinite-order
cyclotomic element after killing \(\psi\). In the plane
this inertia line is unique, since the second character
has exponent \(-e\), and it is stable under decomposition
because its character on inertia is the restriction of a
global character. For the plane, the two inertia exponents and the conductor condition on their finite factors \(\theta_i\) are
\[
 (e,-e),\qquad \operatorname{cond}(\theta_i)\mid M_2\quad(i=1,2).
\]

\proofstep{The fixed-vector bounds at the descended level.}
Here is the passage from these level data to fixed vectors.
Use \(M_3=N\) in rank three and \(M_2\) in rank two.
The ordinary de Rham jumps determine the archimedean
infinitesimal character with entries \((n-1)/2+m_i\)
in the algebraic Satake convention, where \(m_i\) are
the displayed exponents. At each finite place the local
parameter is toral with zero monodromy by the size and
nonresonance comparison. Each of its characters is trivial
on the inertia fixing \(\mu_{M_n}\). At \(\ell\) the smooth
parts of the de Rham characters give precisely this
condition.

At a place \(v\mid M_n\) over \(p\), the norm on local units satisfies
\[
 \mathrm N_{F_v/\mathbb Q_p}(1+M_n\mathcal O_v)
                    \subseteq 1+M_n\mathbb Z_p ,
\]
where \(F_v\) denotes the local field of the descended group.
Reciprocity and the inertia condition imply that each inducing
character is trivial on \(1+M_n\mathcal O_v\).
For the irreducible principal series, the induced-function
model has a vector supported on \(B K(M_n)_v\) taking
value one at the identity because the inducing character and
modulus are trivial on the intersection with the compact
subgroup. At the remaining places the unramified toral
parameter gives a spherical vector. Taking their tensor
product gives the required finite-adelic fixed vector.
For the rational packet the same argument applies after
irreducible quadratic base change wherever the packet
comparison requests it. Neither discriminants nor
principal levels of the enlarged field occur in \(M_n\).

\end{proof}

\paragraph{The rational adjoint Selmer groups.}
Having identified the possible constituents, we need to control their rational adjoint
deformations. For a regular algebraic polarized cusp over a CM field
\(F\), Thorne's adjoint vanishing theorem applies if the
attached representation is irreducible on
\(G_{F(\mu_{\ell^\infty})}\). It gives vanishing of
\(H_f^1(F^+,\operatorname{ad}r)\) for the extension of
the adjoint by the conjugate-dual polarization. In rank three
use \(F=K\) and the cusp attached above.

For the negative component at \(c\) the differential of inverse transpose contributes the minus sign. Modulo \(t^2\),
\[
 J^{-1}(1+tX)^{-t}J=1-tJ^{-1}X^tJ .
\] A
common character twist cancels in the adjoint. The
cyclotomic irreducibility proof at the start of this
section verifies the remaining hypothesis. The theorem thus
gives precisely the rational vanishing used for the type-I
host point. One can first pass to the enlarged marked frame, transport the ordinary and local genericity data
there, apply polarized adjoint vanishing, and return using restriction--corestriction in characteristic
zero. Direct application over \(K\) gives the same group.
\hfill\mbox{\cite{thorne-adjoint}}

\begin{lemma}[Adjoint Selmer vanishing for the descended systems]
\label{lem:descended-adjoint-vanishing}
With the applicable polarization and absolute irreducibility
on the cyclotomic kernel, let \(V\) be the rational adjoint of the descended system over the CM field \(F\),
extended to its maximal totally real subfield \(F^+\) by the polarization.
In the specified cases
\[
                            H_f^1(F^+,V)=0 .
\] The proof
handles the rank-two adjoint diagonal and off-diagonal
summands separately.
\end{lemma}

\begin{proof}
\proofstep{Strong irreducibility of the rational rank-two block.}
For the general-data corner it suffices to prove cyclotomic
irreducibility on \(K\). Inertia \(I\) at \(\ell\) lies
in \(H\) and has full cyclotomic image since \(K\) is
unramified at \(\ell\). Let \(H'=H\cap\ker\epsilon\). The preceding equality of cyclotomic images means that for each \(h\in H\) there is \(i\in I\) with the same image. Thus
\[
 i^{-1}h\in H',\qquad H=IH'.
\] If \(W|_{H'}\) were reducible,
the normal restriction argument would make it scalar or a
sum of two different characters. In the scalar case the
ordinary inertia line would be stable under \(I H'\).

In
the second case inertia permutes the two lines. Choose in
\(I\cap\ker\psi\) an element with infinite-order cyclotomic
value. An element that exchanged these two lines would have a matrix of the following shape in their basis:
\[
 \begin{pmatrix}0&a\\ b&0\end{pmatrix}^{\!2}=ab\,\mathrm{id}.
\]
Its eigenvalue ratio would then have square one, contrary to the infinite-order ratio of the element just chosen. Its two eigenlines are consequently the
two \(H'\)-lines. The ordinary inertia line must be one of
these, and again is stable under \(I H'\). Both alternatives
contradict irreducibility on \(H\). The rank-two cusp over
\(K\) therefore gives polarized adjoint vanishing there
by the same theorem.

\proofstep{The diagonal and the off-diagonal adjoint summands.}
Decompose the polarized endomorphisms relative to the
outer plane and the middle line. The outer endomorphisms
have this rank-two polarized action. Indeed the symmetric
pairing from the original frame on the outer plane is the
reversing matrix, equal up to scalar to the alternating
determinant pairing composed with \(W(c)^{-1}\).
On the one-dimensional middle block, transpose fixes every scalar. Thus the polarized action of \(c\) on its endomorphisms is the minus sign:
\[
 x\longmapsto-x^t=-x\qquad(x\in\operatorname{End}(\text{middle line})).
\]
This summand has character \(\eta\).

To see its finite Selmer group vanishes in characteristic zero,
observe that after restriction to \(K\) a finite class with trivial
coefficients is unramified at every finite prime (at
\(\ell\), a crystalline extension of the trivial character
by itself is unramified). The class group is finite,
and restriction is injective in characteristic zero.
The two off-diagonal spaces are interchanged by \(c\).
One restricts to \(W|_H\) and the other to \((W|_H)^\vee\).
On \(H\) the determinant pairing identifies these two modules.
Induction and the projection formula therefore give their sum as
\[
 \operatorname{Ind}_{H}^{G_\Q}(W|_H)
  \simeq W\otimes\operatorname{Ind}_{H}^{G_\Q}\mathbf 1
  \simeq W\oplus(W\otimes\eta).
\]
The second isomorphism uses that \(H\) has index two and the
coefficient field has characteristic zero.

For the newform \(f\) above put \(k=2e+1\). In Kato's
geometric Frobenius convention the representation
\(V_\lambda(f)\) is \(W(-e)\) after scalar extension.
Apply his Theorem~14.2(1) over the abelian field \(K\),
at the twist \(r=e\). Indeed \(1\le e\le k-1\) and
\(e\ne k/2\).

For every stable lattice the associated
torsion Selmer group is finite. A nonzero rational line of finite Selmer classes would give an infinite subgroup of the torsion Selmer group. Indeed for a line \(L\) over the coefficient field \(E\) and any full lattice \(\Lambda\subset L\),
\[
 L/\Lambda\simeq E/\mathcal O_E,\qquad |E/\mathcal O_E|=\infty .
\]
The torsion bound thus forces the rational finite Selmer group to vanish. Here the
integral global cohomology unramified off the finite set
of bad places and \(\ell\) is finitely generated, so its
image is a lattice. The local finite conditions are
compatible with the maps to torsion coefficients. Both
\(W\) and \(W\otimes\eta\) restrict to \(W|_H\). Finite
classes restrict to finite classes, and corestriction
composed with restriction is multiplication by two.
This proves vanishing on both off-diagonal spaces over
\(\Q\), including at \(e=1\). The direct-sum decomposition
now gives the entire polarized type-II vanishing.
\hfill\mbox{\cite{kato-modular,thorne-adjoint}}

\end{proof}

\paragraph{Fixed-inertia derivatives on the auxiliary field over \(\mathbb Q(i)\).}
The fixed-inertia condition on the Gaussian field requires a local derivative
calculation of its own. It will be used in the comparison with the split point.

\begin{lemma}[Fixed-inertia derivatives are finite]\label{lem:gaussian-fixed-inertia-finite}
For the marked Gaussian test and its prescribed ordinary
characters, let \(V\) denote the polarized adjoint and \(z\) the first-order
cocycle of a normalized dual-number deformation of the Gaussian test
representation. For every deformation satisfying the fixed-inertia equations,
\[
                       [z]\in H_f^1(\mathbb Q,V)
                              \ \subseteq\ H^1(\mathbb Q,V).
\]
\end{lemma}

\begin{proof}
A deformation to dual numbers with first derivative \(y\)
has matrix part \((1+\delta y(g))r(g)\), \(\delta^2=0\).
This defines a cocycle for the polarized adjoint. At an
away prime outside the bad set the deformation is
unramified. At a level prime outside the exceptional set the
local equations kill the derivative on inertia in
\(\ker\eta\cap\ker P\), a finite-index subgroup of inertia.

At an exceptional prime the dual-number deformation on
positive inertia is conjugate to the prescribed diagonal
action with its exact three characters. Differentiating
this equality of matrices expresses the cocycle there
as a coboundary, with a single conjugating derivative
for all inertia elements. For either finite-index inertia subgroup \(I'\subset I\), let
\(\xi\in H^1(I,V)\) be this local class over the coefficient field.
The preceding restriction calculation gives
\[
 [I:I']\,\xi
 =\operatorname{cor}_{I'}^I\operatorname{res}_{I'}^I(\xi)=0 .
\]
The positive integer \([I:I']\) is invertible over the characteristic-zero
coefficient field, so \(\xi=0\) on full inertia.

\proofstep{The finite local condition above the coefficient prime.}
At \(\ell\) restrict to the open normal subgroup of the
decomposition group killing \(\eta\) and the
\(\ell\)-local component of \(\psi\). On its inertia
the ordinary dual-number characters are exactly
\(\epsilon^{e},1,\epsilon^{-e}\). The separator
\(x^{\ell^{v_\ell(N)}}\) belongs to this subgroup since its
cyclotomic value is one modulo \(\ell^{v_\ell(N)+1}\)
by successive powering. In characteristic \(\ell\) the separator retains its distinct residual eigenvalues after powering, since
\[
 a^{\ell^k}-b^{\ell^k}=(a-b)^{\ell^k}\ne0\qquad(a\ne b).
\] Thus no two different
graded characters differ by \(\epsilon\) on this subgroup.
Viewed over the coefficient field, each dual-number line
is a self-extension of a crystalline character by an
unramified additive derivative, so is crystalline. To
add the next such two-dimensional block to the preceding
ones, the Hom space from the new block to the preceding
part has a filtration with characters of strictly
positive cyclotomic exponents.

Its dual twist has no
invariants by the separator. The local dimensions for this Hom space agree. Since the finite subspace embeds naturally in the full group,
\[
 H_f^1\subseteq H^1,\qquad \dim H_f^1=\dim H^1\quad\Longrightarrow\quad H_f^1=H^1 .
\]
The extension of crystalline blocks is therefore crystalline. This induction proves that the entire
self-extension is crystalline on the subgroup. Equivalently
the adjoint cocycle splits after tensoring with
\(B_{\rm cris}\) there. Averaging a splitting over the
finite quotient gives one over the original group. The
cocycle is therefore finite at \(\ell\) too. All matrices
and changes of basis may be taken over a finite
coefficient extension.

To see this for a continuous cocycle \(z:G\to V_{\overline{\Q}_\ell}\), fix a finite field of definition for the action. The sets \(H_L=z^{-1}(V_L)\), for the countably many finite extensions \(L\) of that field in \(\overline{\Q}_\ell\), are closed and cover \(G\). The cocycle identity makes each a subgroup. By Baire's theorem one \(H_L\) has interior, hence is open. Choose coset representatives \(g_1,\ldots,g_b\). For \(h\in H_L\),
\[
             z(g_i h)=z(g_i)+g_i z(h)\in z(g_i)+V_L .
\]
Adjoining the finitely many coordinates of the \(z(g_i)\) gives a common finite field of definition.
\hfill\mbox{\cite{bloch-kato-finite}}

\end{proof}

\begin{construction}
\[
\begin{array}{lll}
\text{branch}&\text{constituent descended}&\text{Gaussian application}\\ \hline
\text{type I}&\text{irreducible rank three}&\text{enters the trace product}\\
\text{type II}&\text{outer plane in rank two}&\text{excluded under }\mathrm{Xpp}
\end{array}
\]
\caption{The two possible branches at a two-sided point. The exclusion in the Gaussian application uses the residual vanishing assumption for the outer plane.}\label{fig:descent-trace-branches}
\end{construction}

The two branches in the table tell us which descended constituents to count. We can
now form the trace products on the original fields.

\begin{proposition}[Trace
products on the original packets]
\label{prop:original-packet-count}
For the indicated type-I and outer-plane points, the trace
packets and denominator allowances give the nonzero integer \(A_{J_0}(p)\) of \eqref{eq:local-odd-d}.
In the notation of that estimate its size satisfies
\[
 \log |A_{J_0}(p)|\ll_\psi
   \varphi([U:J_0])\bigl(\mathcal N(j+1)\bigr)^a\log p .
\] The Gaussian hypothesis excludes the separate
outer-plane case in that application. The power of
\(N(e+1)\) depends only on the original quadratic field and
fixed dimension.
\end{proposition}

\begin{proof}
\proofstep{The residual grading-torus tangent.}
The Gaussian tangent class is now zero by the type-I
vanishing. On either host a coboundary is an
infinitesimal conjugation by \(1+\delta X\). The
normalization of the marked difference forces \(X\)
diagonal because the off-diagonal entries of its commutator
with \(\operatorname{diag}(z,0,-z)\) must vanish, with
\(z\ne0\). Normalization at \(c\) imposes the adjoint condition on the surviving diagonal entries, leaving
\[
 X=-J^{-1}X^tJ,\qquad X=\operatorname{diag}(u,0,-u).
\] This is exactly the
infinitesimal grading action. At a point with nonzero
coordinates of both signs it accounts for the tangent
along its conjugacy orbit.

\proofstep{Which points enter the two products.}
A nonzero entry of each strict sign precludes both triangular
orientations. The marked coordinate test in
Section~\ref{sec:marked-hosts} then leaves type I and the
type-II outer plane. The rank-three trace estimate keeps both. For the case over the Gaussian field the
extra residual hypothesis used by the submission is vanishing
for the constrained weight \(-2\) classes, denoted
\(\mathrm{Xpp}\). In zero-based indices the coefficient on the entry \((2,0)\) is the one-dimensional module
\[
 k(\eta P^{-2}),\qquad g\cdot x=\eta(g)P(g)^{-2}x .
\] We explain exactly where this extra
arithmetic enters.

\proofstep{The residual class detected by an outer-block entry.}
In a type-II block the entry \((2,1)\) is identically zero.
If any \((2,0)\)-entry is nonzero, choose \(d=r(g_0)_{20}\)
of smallest valuation and put \(b(g)=r(g)_{20}/d\).
These quotients are integral. Multiplication on the positive
component gives, on reducing,
\[
  z(g)=\overline{b(g)}P(g)^{-1},\qquad
  z(gh)=z(g)+\eta(g)P(g)^{-2}z(h).
\]
For the extension of this equality to the negative component
use \(r(c)_{\rm mat}=1\). On the outer block the action of
\(c\) is inverse transpose with the reversing pairing, which
sends its lower left entry to its negative divided by its
determinant. That determinant is one modulo the maximal
ideal. This gives the displayed sign and \(z(c)=0\).
Continuity of the residual quotient follows by evaluating
the continuous entry modulo one power of the uniformizer
beyond \(v(d)\). Away from \(T_b\) the killed inertia
elements have zero entry.

At a place in \(T_b\), write
the first ordinary local line in the integral model with
generator \((1,h_1,h_2)^t\) congruent to the first coordinate
vector. For every positive local \(g\), let \(V(g)\) be
the eigenvalue on this vector minus \(r(g)_{22}\).
Reading the last row now gives the proportionality
\[
 d\,b(g)=h_2 V(g).
\] At the separating
element \(s\) of this place the residue of \(V(s)\) is
\(P(s)-P(s)^{-1}\ne0\). Hence \(h_2/d=b(s)/V(s)\)
is integral, and on reduction \(z(g)\) is a single
constant times \(1-P(g)^{-2}\) for all such \(g\).
This also applies at the exceptional places where the
local model is diagonal. Local coboundaries extend from
the positive kernel by averaging over its quotient of
order at most two. Thus \(z\) satisfies the indicated
residual Selmer conditions.

If the residue coefficients
were enlarged, apply each linear functional to the base
residue field before using the vanishing. Each scalar cocycle is a coboundary \(g\mapsto(\eta(g)P(g)^{-2}-1)a\). At conjugation this reads
\[
 0=z(c)=(-1-1)a=-2a .
\]
Since \(\ell\) is odd, \(a=0\) for every scalar functional, and therefore \(z=0\). This
contradicts \(b(g_0)=1\). The type-II block would therefore
itself be upper triangular. Under \(\mathrm{Xpp}\) only
type I contributes to the two-sided Gaussian points.

\proofstep{Exclusion in the Gaussian case.}
For the present auxiliary field over \(\mathbb Q(i)\) this residual vanishing follows
from the Bernoulli unit calculation in
Paragraph~\ref{par:dyadic-residual}. We give the conductor
check because the character here involves a square.
Use \(q_0\) for the bound on the conductor of the square
of the prime-to-\(\ell\) odd-conductor part of \(\psi\).
It is the fixed target conductor in the application over the Gaussian field.
Write \(\omega\) for the Teichmüller lift of the residual
cyclotomic character. Lift \(\eta P^2\) as the product of
the primitive characters inducing
\[
 \varphi=\omega^{2e}(\psi_{\rm odd}^{(\ell')})^2,
       \qquad \rho=\chi_4\psi_2^2 .
\]
Their conductors are coprime. The odd conductor \(f\) of
the first divides \(\ell q_0\). With \(d\) the order of
\(\ell\) modulo the least common multiple of four and
the order of \(\psi_{\rm odd}^{(\ell')}\), its values have
order dividing \(\ell^d-1\). Put \(r_0=v_2(\ell^d-1)\ge2\).

Since \(c=c_2(\psi)\ge4\), the value of \(\psi_2\) on
five has order \(2^{c-2}\). Squaring divides that order
by two, and \(\chi_4(5)=1\). Hence \(\rho\) has exact
conductor \(2^{c-1}\), including at \(c=4\).
Its product with \(\varphi\) is odd. The large-conductor
hypothesis gives
\[
        2^{c-1}\ge (\ell q_0)4^{r_0}\ge f4^{r_0}.
\]
The Bernoulli unit assertion applies to this product and to
its inverse. In the paragraph cited above the only
requirements for the Stickelberger argument are
unramified restriction on the kernel of the character
and the local coboundary condition at \(\ell\) to
ensure this unramifiedness there.

On
\(G_{\Q(\mu_{f2^{c-1}})}\) the inverse residual character
\(\eta P^{-2}\) is trivial. The cocycle restricts to an
unramified homomorphism because the intersection of inertia
with \(\ker\eta\cap\ker P\) has prime-to-\(\ell\) index
in inertia, and local coboundaries vanish on the character
kernel. Let \(z\) denote this restricted unramified homomorphism
and \(\Theta\) the smoothed Stickelberger element. Annihilation
and the character computation give
\[
           0=\Theta z=\chi(\Theta)z,\qquad \chi(\Theta)\in k^\times
\]
on the component being tested, with \(\chi\) its corresponding
residual character. Hence \(z=0\). The original cocycle factors
through the remaining abelian cyclotomic quotient. If \(g_0\)
has character value \(\chi_0(g_0)\ne1\) for its original
coefficient character \(\chi_0=\eta P^{-2}\), commutativity
gives
\[
   (\chi_0(g_0)-1)z(g)=(\chi_0(g)-1)z(g_0).
\]
This expresses the descended cocycle as a coboundary. This proves the required residual vanishing
on the actual host.
\hfill\mbox{\cite{washington-nonp,washington-cyclo}}

\proofstep{The trace bound for each admitted point.}
The quantitative contribution of an admitted point can
now be read at the original Frobenius. In the notation
of (D), the rank-three cusp has level \(N\) and exponents
\((e-1,-1,-e-1)\) at both complex embeddings of the fixed
\(K\). To bound its packet size polynomially in \(N(e+1)\),
use the finite cell complex at level three, the cover at level
\(3N\) and the Weyl dimension bound
in Section~\ref{sec:rank-four-deformation}. Clozel's
conjugation theorem preserves this packet and its algebraic
normalization. Let \(b\) be an algebraic Satake root at a good
place of norm \(q\). An invariant Galois lattice makes it
a unit away from the residue characteristic \(p\).
At \(p\), apply the spherical crystalline comparison and the
packet weak-admissibility estimate
(Sections~\ref{sec:local-global} and~\ref{sec:rank-four-deformation}):
\[
 v_p(q^{e-1}b^{-1})
     =v_p(q)\bigl(e-1-v_p(b)/v_p(q)\bigr)\ge0 .
\]
Indeed the maximal Hodge exponent is \(e-1\). This applies to
each coefficient embedding of each packet member.

One can also obtain the denominator bound on a
definite unitary group. After quadratic base change split
at the tested place, a character twist unramified there
makes the cusp conjugate self-dual. The Hecke operators
on the algebraic coefficient lattice bound the elementary
symmetric functions of the roots and their reciprocals.
For a monic polynomial \(X^n+c_1 X^{n-1}+\cdots+c_n\) over a nonarchimedean field, any root \(\alpha\) satisfies
\[
 |\alpha|\le\max_{1\le j\le n}|c_j|^{1/j}.
\]
Indeed a larger value would make the leading term strictly larger than every other term. Apply this bound both to the roots and to their reciprocals to obtain the slope inequalities. The shift in the weights cancels
the value of the twisting character and the split local
extension leaves \(q\) and the Satake multiset unchanged.
\hfill\mbox{\cite{clozel-motifs,borel-serre,blggt,blggt-lg}}

For the rank-two contribution one can also count the
classical newforms directly. On the rational inertia the
two characters away from \(\ell\) have conductor bounded
by \(M_2\) and the inertia image is finite. At \(\ell\)
the gap \(2e\) and the ordinary flag give potential
crystallinity by the flag argument, and the two smooth
Weil characters have the same conductor bound. A Tate twist is unramified for the Weil parameter used to compute the Artin conductor. In this normalization its exponent obeys
\[
 a(r\otimes |\cdot|^t)=a(r).
\] Thus by
local--global compatibility and the conductor formula for
newforms the conductor divides \(M_2^2\). The normalized
newforms of these conductors occur independently in
\(S_{2e+1}(\Gamma_1(M_2^2))\).

At modulus \(M_2^2\) the modular curve index is bounded by the square of that modulus:
\[
 [\mathrm{SL}_2(\mathbb Z):\Gamma_1(M_2^2)]\le (M_2^2)^2=M_2^4 .
\]
The valence dimension bound is linear in the weight times
this index, so the dimension is \(O((e+1)M_2^4)\). The
list is stable under conjugation. At a good \(p\) let
\(\alpha,\beta\) be the algebraic Hecke roots before the
twist. They are integral, have absolute value \(p^e\)
in every embedding and satisfy \(\alpha\beta=p^{2e}\xi\)
with \(\xi\) a root of unity.

If \(a=\alpha/p^e\) and
\(b=\beta/p^e\), then
\[
             p^e a=\alpha,\qquad p^e a^{-1}=\xi^{-1}\beta,
 \qquad      p^e b=\beta,\qquad p^e b^{-1}=\xi^{-1}\alpha .
\]
At residue degree \(h\), each of \(a^{\pm h},b^{\pm h}\)
is consequently integral after multiplication by \(p^{eh}\).
This factor also preserves the middle eigenvalue \(1\), and
clears the denominator in the trace.
The packet can also be counted from its irreducible restrictions to the quadratic field, the choices of roots at positive residue degree and the at most two descent characters for each fixed irreducible restriction.
The direct count gives a polynomial bound of the same
required form. \hfill\mbox{\cite{deligne-galois,carayol-local,deligne-weil}}

\proofstep{Counting isomorphism classes by torus orbits.}
For a fixed character and either type, equality of complete
systems gives equality of semisimple representations at
all good Frobenii and hence isomorphism over \(K\).
Assign a packet index to each selected conjugacy class. The complete system, rather than one trace, gives the implication
\[
 \text{same packet index}\quad\Longrightarrow\quad\text{same selected conjugacy class},
\]
proving injectivity for the count. Thus no
identification of such indices is needed. A common integer product over
both packets and the full character orbit is the integer
in (D). For the Gaussian case use just the first packet
and allow its two field embeddings, as in the case over the Gaussian field
product. In (D) the implied constant and exponent depend
only on the chosen quadratic field and the fixed dimension.
In particular neither the marked field, coefficient
ramification, omitted primes, nor the choice of ordinary
seed enters the power of \(N(e+1)\).

Section~\ref{sec:rank-three-length} supplies the integral
chart comparison for (L). The trace is a function of
weight zero. The length comparison uses one point
from each two-sided torus orbit, and the marked
bases make those representatives give distinct
restriction isomorphism classes in the original packets.
\end{proof}

\section{Integral length bounds from rank-three Frobenius traces}
\label{sec:rank-three-length}
\hypertarget{source-ordinary-rank-three-length}{}

The descent has identified the ordinary constituents in the two packets. We now return
to the original coefficient ring to measure the length they contribute.
The graded deformation directions give a bound on the
integral lengths in terms of valuations of Frobenius trace
differences in the two descended packets. The bounds enter the odd-prime cohomology window of
Appendix~\ref{sec:cohomology}. The localization kernel maps to the torsion derivations in \eqref{eq:local-length-1}. The final bound \eqref{eq:local-length-11} uses the valuation on the original coefficient field and one representative of each two-sided orbit.

Conjugation by \(D(t)=\operatorname{diag}(t,1,t^{-1})\) gives the grading directly on matrix units. With \((w_1,w_2,w_3)=(1,0,-1)\),
\[
 D(t)E_{ij}D(t)^{-1}=t^{w_i-w_j}E_{ij}.
\] An entry in row \(i\)
and column \(j\), numbered from zero, has weight \(j-i\).
The semidirect product describing the polarization acts on its
Lie algebra through \(X\mapsto-JX^{\mathsf T}J\) on the
quadratic coset, with \(J\) the reversal matrix. This
involution preserves the entry weights. Write \(r_0\) for the
split representation and \(g_*=g_1\) for the mark. The
matrix of \(r(c)\) is fixed. The normalization at the mark
diagonalizes
\[
              r(g_*)-r(cg_*^{-1}c^{-1})
              \quad\text{on the quadratic kernel.}
\]
The difference at \(r_0\) has
diagonal entries \(\Psi(g_*)-\Psi(g_*)^{-1},0,
\Psi(g_*)^{-1}-\Psi(g_*)\). These are distinct in the residue
field by the separator and oddness of \(\ell\).

\subsection{The integral cotangent comparison}

Which part of the integral cotangent module can record the desired length? The grading
isolates it. Let \(O\) be the coefficient DVR, with fraction field \(E\),
uniformizer \(\pi\) and \(v_\pi(\pi)=1\). If \(x:A\to O\) is
the origin of a graded ambient algebra, put \(J_x=\ker x\).
The lattice at weight \(-1\) is the image \(C_{-1}\) of
\(A_{-1}\) in \(J_x/\operatorname{sat}_O(J_x^2)\). The
equations \(F_i\) of weight \(-1\) give a submodule
\(\langle F_i\rangle\) of this lattice.

\begin{construction}
\[
\begin{array}{ccc}
\text{Smith summand}&\text{admissible image of its generator}&\text{length}\\ \hline
O/(\pi^d),\ d<s&\pi^{s-d}O/(\pi^s)&d\\
O/(\pi^d),\ d\ge s&O/(\pi^s)&s\\ \hline
&&\min(d,s)
\end{array}
\]
\caption{One Smith summand in the conormal bound, with \(s=v_\pi(b)\). A point derivation with values in \(O/(b)\) can see at most the smaller exponent.}\label{fig:cotangent-smith-factors}
\end{construction}

\begin{lemma}[Torsion derivations and the saturated conormal]
\label{lem:torsion-derivations-conormal}
Suppose an additive
injection from an \(O\)-module \(M\) sends a class \(m\) to
an \(O/(b)\)-valued point derivation, for \(0\ne b\in O\).
Suppose each derivation kills the \(F_i\) and is the
reduction of an \(O\)-valued point derivation of weight
\(-1\) on \(A\). If the quotient
\(C_{-1}/\langle F_i\rangle\) has finite length, then
\[
 \operatorname{length}_O M\le
 \operatorname{length}_O(C_{-1}/\langle F_i\rangle).
\]
\end{lemma}

\begin{proof}
Every integral point derivation at \(x\)
kills \(J_x^2\) by the product rule and its saturation
by torsion freeness of \(O\). It follows that the injection
factors through
\[
             \operatorname{Hom}_O(C_{-1}/\langle F_i\rangle,O/(b)).
             \tag{1}\label{eq:local-length-1}
\]
Indeed the scalars and every weight other than \(-1\)
already have zero derivative. If the domain lattice
quotient has finite length, the summand \(O/(\pi^d)\)
in its Smith decomposition gives a summand of length
\(\min(d,v_\pi(b))\) in \eqref{eq:local-length-1}. Counting both finite additive
groups in powers of the residue field cardinality gives
the desired length inequality even for an additive
injection. In fact if \((d_i)_i\) are the Smith exponents,
the target has length
\(\sum_i\min\{d_i,v_\pi(b)\}\le\sum_i d_i\).
\end{proof}

\begin{auxlemma}[Annihilation in large multidegree]\label{lem:rees-multidegree-annihilation}
Let \(R\) be a Noetherian \(\mathbb N^s\)-graded algebra generated over
\(R_{\boldsymbol0}\) by the pieces \(R_{\mathbf e_i}\) in the unit
directions. Put \(R_+=\prod_{i=1}^s (R_{\mathbf e_i})\).
Let \(H\subset R\) be a homogeneous ideal and \(J\subset R_{\boldsymbol0}\) an ideal.
If each prime containing \(H\) but not containing \(R_+\) contains \(J\),
there is an integer \(h\ge1\) such that
\[
 J^h(R/H)_{\mathbf d}=0\qquad(d_i\ge h\text{ for every }i).
\]
\end{auxlemma}

\begin{proof}
Every prime over \(H\) contains \(J R_+\), so \(J R_+\subseteq\sqrt H\).
By finite generation \((J R_+)^h\subseteq H\) for some \(h\).
A monomial of multidegree \(\mathbf d\) with every \(d_i\ge h\)
contains at least \(h\) factors from each unit direction. Consequently
\[
 R_{\mathbf d}\subseteq R_+^h,\qquad
 J^h R_{\mathbf d}\subseteq (J R_+)^h\subseteq H.
\]
\end{proof}

\begin{proposition}[Graded cotangent length bound]\label{prop:rank-three-conormal}
Let \(B=A_0\) be complete Noetherian local
and let \(A\) be generated over \(B\) by finitely many
elements of weights \(-2,-1,1,2\). Write \(I\) for the
ideal of the representation chart in \(A\). The ambient
point has a smooth formal chart at \(x\) over \(E\),
with the first row giving the tangent dimensions and
the second row the number of equations of \(I\).
The equations vanish at the section \(x\).
\[
\begin{array}{c|ccccc}
 j&-2&-1&0&1&2\\ \hline
 \text{ambient dimensions}&s&r&m&r'+1&s'+1\\
 \text{equations}&s&r+1&m&r'&s'
\end{array}
\tag{2}\label{eq:local-length-2}
\]
The equations in nonpositive weights span the rational
cotangent spaces of their weights.

Assume that \(B/(I\cap B)\)
is finite over \(O\). Over an algebraic closure of \(E\)
a point is two-sided if some negative and some positive
weight coordinate are nonzero. Over a fixed point other
than \(x\), at least one of the two strictly one-sided
sets with the same restriction to \(B\) must be empty.
At a two-sided point all tangent derivations of \(A/I\)
are generated by the torus tangent. A function \(a\in B\)
vanishes at \(x\) and is nonzero at every two-sided point.

For \(r\ge1\), these hypotheses give
\[
 \operatorname{length}_O(C_{-1}/\langle F_i\rangle)
       \le 2\sum_{\xi\in T}v_\pi\bigl(\xi(a)\bigr),
 \tag{3}\label{eq:local-length-3}
\]
where \(T\) contains exactly one point from each two-sided
torus orbit.
\end{proposition}

\begin{proof}
\proofstep{Slices transverse to the grading torus.}
A point of an orbit can be rescaled to the
slice \(f=1\) for a weight \(-2\) element \(f\) nonzero
on it. Use also a weight \(2\) element \(g\) nonzero on
it. Such even-weight elements are available by squaring.
Only finitely many \(f,g\) from generating sets are
needed. Inverting \(fg\) makes both \(f\) and \(g\) units.

On each slice all even-weight values follow from the
invariant values. Products of pairs of odd-weight values
follow in the same way. The only remaining ambiguity is the action of \(\mu_2\) on the slice. Thus for the tuple of these coordinates it is a common sign:
\[
 (y_i)_i\longmapsto(\varepsilon y_i)_i,\qquad\varepsilon\in\{1,-1\}.
\] The slices are finite over the invariant fibers
and reduced by the tangent hypothesis. This proves the
finiteness of \(T\). It also proves that each generic
invariant point on this locus contributes at most one
copy to any given weight in the unsquared equation quotient.

\proofstep{Squaring the equation list.}
For completeness, start the length calculation by replacing
the \(F_i\) with a Smith list. In the formal chart the linear terms of these equations are independent in their respective coordinates, and the extra relation starts in higher degree:
\[
 \operatorname{lin}(F_i)=\pi^{h_i}y_i\ (1\le i\le r),\qquad \operatorname{lin}(F_*)=0 .
\] The length on the left of
\eqref{eq:local-length-3} equals \(\sum_i h_i\). Over \(E\), eliminate the
weight-zero variables by the \(m\) independent equations
of weight zero. Denote the resulting formal chart map
by \(\phi\). If necessary, replace \(a\) by
\(a+F_i u\) for \(u\in A_1\) with nonzero image under
\(\phi\), so that \(\phi(a)\) is nonzero. This change
keeps every value in \eqref{eq:local-length-3}.

Use the three Rees directions over \(B\) given by
\[
        A_{-2},\qquad A_2,\qquad BF_*^2+A_{-2}^2A_2.
        \tag{4}\label{eq:local-length-4}
\]
In the power series chart write the variables as \(x_i\)
with weights \(w_i\). A finite set \(M_j\) of componentwise
minimal exponent vectors generates the solutions of
\(\sum_i w_i\alpha_i=j\) by addition of weight-zero exponents.
Indeed the minimal set is finite by Dickson's lemma, and any
nonminimal vector above \(\beta\in M_j\) has weight-zero difference
\(\alpha-\beta\). Assign each term of a series to one such \(\beta\).
This gives a finite sum \(\sum_{\beta\in M_j}x^\beta f_\beta\)
with each \(f_\beta\) of weight zero.

Similarly at weight \(2n\)
one can remove a weight-two submonomial for \(n\ge1\): use a
positive weight-two factor when one is present, or two factors
of weight one otherwise. Repeat on the remaining monomial.
Reverse all signs at weight \(-2n\). Grouping the terms shows
\[
                       A_{\pm2n}=A_{\pm2}^{\,n}.
\] Quotient the Rees ring by the
homogeneous ideal of coefficients killed by \(\phi\), and
write \(R\) for the quotient. Choose sections \(\mathbf t\)
consisting of \(F_i^2\) and the weight \(-2\) equations
in the first direction, the squares of the weight \(1\)
equations and the weight \(2\) equations in the second,
and \(F_*^2\) in the third. Regularity of \(a\) on this quotient means that multiplication by \(a\) gives the exact sequence
\[
 0\longrightarrow R\xrightarrow{\ \cdot a\ }R\longrightarrow R/aR\longrightarrow0 .
\]

\proofstep{Preserving the relevant leading coefficients.}
In large multidegree \(d\), the pieces of \(R/(\mathbf t)\)
are finite \(O\)-modules. To check the support, work away
from the ideal generated by products of sections of the
three directions. Choose a section outside the prime
in each direction. Any relation of \(I\cap B\) can
be written as a sum of homogeneous equation multiples.
For even-weight equations, multiply by enough of the
chosen sections to put each term in a Rees degree in
which the equation is a factor.

For an odd-weight equation, pass to its square in the even comparison and then use the implication valid in a prime ideal \(\mathfrak p\):
\[
 z^2\in\mathfrak p\quad\Longrightarrow\quad z\in\mathfrak p .
\]
A section with coefficient \(F_*^2\) in the first
direction can be exchanged with the section carrying
the same coefficient in the third direction after
multiplication by the chosen third section. Thus, with \(J=I\cap B\) and \(R_+\) the product of the
three ideals generated by unit-direction sections, we have
\[
 (\mathbf t)\subseteq\mathfrak p,\quad R_+\not\subseteq\mathfrak p
     \quad\Longrightarrow\quad J\subseteq\mathfrak p\cap B .
\]
This contraction holds on the integral Rees ring.

Suppose now that a characteristic-zero contracted prime
also kills \(A_{-2}A_2\). Nonvanishing of the first and second
sections implies that their coefficients survive outside
the nilradical in the corresponding fiber of \(A/I\).
Indeed, if a power of either coefficient belongs to the
fiber ideal, clearing denominators and making the same
section comparisons forces the power of the section
into the prime. This gives both a negative and a positive
point above the same invariant values. Both points are one-sided. Indeed, at a point over a field the vanishing of all products forces at least one side to vanish:
\[
 A_{-2}(z)A_2(z)=0\quad\Longrightarrow\quad A_{-2}(z)=0\ \text{or}\ A_2(z)=0 .
\] Scaling to the fixed point
keeps their invariant values.

For the fiber of \(x\),
the rational spanning in weight \(-1\) kills all weight
\(-1\) values on a negative one-sided point. In the
saturated square at weight \(-1\) every contribution
has a positive factor or a weight-zero factor vanishing
at \(x\). In weight \(-2\) every additional product has two factors of weight \(-1\). If \(z\) is the point under consideration, then for each such pair
\[
 (ab)(z)=a(z)b(z)=0\qquad(a,b\text{ of weight }-1).
\]
Spanning then also kills the weight \(-2\) values. Thus
the negative point cannot occur at \(x\). At every other
fixed point the one-sided hypothesis gives the exclusion.

\proofstep{Saturation in the three Rees directions.}
Apply Lemma~\ref{lem:rees-multidegree-annihilation} over the
coefficient DVR to \(H=(\mathbf t)\), using the integral prime
contraction just proved. For a fixed integer \(h\ge1\) this gives
\[
      J^h(R/(\mathbf t))_d=0\qquad(d_1,d_2,d_3\ge h).
\] Each nilpotent thickening of
the finite invariant quotient is also finite over \(O\),
by the filtration by powers of its defining ideal.
Adding the equation \(a=0\) leaves only special-fiber support in these
pieces by its assumed nonvanishing at the two-sided
points. These support tests apply to the homology of the
corresponding Koszul lists, which is killed by the list
of sections.

In large multidegrees write \(\chi\) for the alternating
sum of lengths of Koszul homology. The comparison is
\[
   \chi(\mathbf t;R/aR)_d
          \le \chi(a;(R/(\mathbf t))_d).
          \tag{5}\label{eq:local-length-5}
\]
Here \(a\) is regular on \(R\), so the left side equals
\(\chi(\mathbf t,a;R)_d\). Successively taking
kernel and cokernel of multiplication by each section
gives the inequality. The omitted terms have nonnegative
Euler characteristics in large degrees. This positivity,
also used for monotonicity under quotients, follows by
homogeneous prime filtration and induction on the list:
on a prime factor a nonvanishing first section is regular.

If it vanishes, the contribution is the difference of the
two shifted Euler characteristics for the tail. If the replacing section has multidegree \(e\) and is regular on the prime factor \(M\), the calculation uses the sequence
\[
 0\longrightarrow M(-e)\xrightarrow{\ \cdot s\ } M\longrightarrow M/sM\longrightarrow 0 .
\]
Thus only its degree enters the difference of shifted Euler characteristics. If the whole section degree vanishes on the
prime factor, the factor has zero pieces in large degrees
in that direction. For a degree-zero section one can
take the replacement to be \(1\).

\proofstep{The lower bound from the Smith exponents.}
On the left of \eqref{eq:local-length-5}, take leading coefficients in the chart \(\phi\) in the respective negative and positive weighted degrees
\[
 X_d=2d_0+4d_2,\qquad Y_d=2d_1+2d_2 .
\]
Each element in degree \(d\) has at least these totals:
\(F_*\) has weight \(-1\) and no linear term, so its
negative total is at least two. The leading coefficients
form a finite free \(O\)-module acted on by \(B\)
through \(x\). They span the full bidegree over \(E\).
To see the spanning, use representatives of the
individual chart axes in their weights. Group the
weight \(1\) axes in pairs and the weight \(2\) axes
singly on each side. Each third-direction factor takes
two negative blocks and one positive block. Within the
fixed bidegree the products of axis representatives
have independent leading terms by the filtration by
ordinary total degree. The resulting action factors through the vanishing of \(a\) at the point, namely
\[
 R\twoheadrightarrow R/aR\qquad(a(x)=0),
\]

On this quotient \(F_i^2\) acts by \(\pi^{2h_i} y_i^2\).
Splitting off its scalar by Koszul additivity contributes
\(2h_i\) times the Euler characteristic modulo \(\pi\)
of the list with this section omitted. The degree
pieces are torsion-free, so the successive
\(\pi\)-quotients count the same monomials. Positivity
also applies to the list with all scalars removed.
The Hilbert series for the two leading degrees is
\[
 \frac{1}{(1-X)^r(1-X^2)^s(1-Y)^{r'+1}(1-Y^2)^{s'+1}}.
\]

The remaining sections multiply it by
\((1-X^2)^{r-1+s}(1-Y^2)^{r'+s'}(1-X^4Y^2)\).
After cancellation the first numerator factors are
\((1+X)^{r-1}(1+Y)^{r'}\). Separate the two degrees in the numerator by adding and subtracting \(Y^2\):
\[
 1-X^4Y^2=(1-Y^2)+Y^2(1-X^4).
\]
The second summand has bounded \(X\)-degree after
cancellation. The first summand has eventual coefficient
\(2^{r+r'-1}\). Thus the left side of \eqref{eq:local-length-5} is at least
\(2^{r+r'}\sum_i h_i\).

\proofstep{The upper bound from two-sided orbits.}
On the right of \eqref{eq:local-length-5} the kernel and cokernel lengths on
the \(O\)-torsion cancel. On the free part the Euler
characteristic is the valuation of the determinant of
multiplication by \(a\). It is the sum over the geometric
invariant points with the multiplicities in the generic
degree piece. At these points multiplication by a negative and by a positive even-weight unit permits shifts between the corresponding weight pieces: if \(u\) is such a unit of weight \(w\),
\[
 M_j\xrightarrow[\sim]{\ \cdot u\ } M_{j+w},\qquad m'\longmapsto u^{-1}m'\ \text{for the inverse}.
\] Each module in \eqref{eq:local-length-4}
then becomes its full weight piece.

The generic Rees
piece is a quotient of the corresponding weight piece
with ordinary equations \(F_i^2,F_*^2,E_i^2\), the
weight \(2,-2\) equations, and the weight-zero
equations, where \(E_i\) are the weight \(1\) equations.
Replacing an equation \(q\) by \(q^2\) at most doubles
each multiplicity bound, by the two pieces of the
filtration by \(q\), with one grade shifted. The
unsquared orbit bound above gives multiplicity at
most \(2^{r+1+r'}\) after the replacements. The factor
of two in \eqref{eq:local-length-3} follows on comparing with the lower bound.
\label{cor:ordinary-rank-three-length-1}

\end{proof}

\subsection{The split point and the localization kernel}

To apply the cotangent comparison we must identify the part of the localization kernel
it detects. Begin at the split point with rational coefficients, then pass to finite
torsion depth.

\begin{lemma}[Rational vanishing at the normalized split point]
\label{lem:rank-three-split-vanishing}
For the split representation at the Gaussian fixed-inertia chart of
Proposition~\ref{prop:gaussian-rank-three-chart} or the general
crystalline chart of Proposition~\ref{prop:rank-three-crystalline-local-base}, the normalized rational point derivations satisfy 
\[
 \operatorname{Der}_{x,j}^{\rm norm}=0\qquad(j=-2,-1,0),
\]
 where the subscript denotes the weight at the split point \(x\). In these
weights the equations therefore span the rational conormal
spaces in the hypotheses of Proposition~\ref{prop:rank-three-conormal}.
\end{lemma}

\begin{proof}
Write \(\Gamma\) for the host Galois group, and
\(\Psi=\epsilon_\ell^{\,e}\psi\). On the quadratic kernel
the weight \(j\) adjoint space of \(r_0\) has scalar action
\(\Psi^j\). On \(\Gamma\) its eigencharacters are
\(\Psi^j\) and \(\eta\Psi^j\), on the \(+1\) and \(-1\)
eigenspaces, respectively, of \(-JX^{\mathsf T}J\).
The involution sends \(E_{ij}\) to \(-E_{2-j,2-i}\) for \(0\le i,j\le2\). A two-element orbit contributes one eigenvector of each sign. A fixed matrix position contributes one negative eigenvector. Counting the orbits in each weight gives
\[
 \begin{array}{c|rrrrr}
 j&-2&-1&0&1&2\\ \hline
 \dim (\mathrm{Ad}\,r_0)_j&1&2&3&2&1\\
 \dim (\mathrm{Ad}\,r_0)_j^{+}&0&1&1&1&0\\
 \dim (\mathrm{Ad}\,r_0)_j^{-}&1&1&2&1&1
 \end{array}
\]

The rational nonpositive-weight groups used at the split
point vanish by the following rank-one argument. At
weight zero restrict to the quadratic kernel, so that
both characters are trivial. The local condition gives
an unramified additive character everywhere. For the
Gaussian chart this follows from vanishing on the
residual-kernel inertia and finite index there, and
from the unramified weight-zero condition at \(\ell\).
For the crystalline chart, inertia at \(\ell\) lies in
the quadratic kernel and the finite local conditions on
the trivial character are unramified. Class field theory
makes the resulting homomorphism zero, since the class
group is finite. Restriction to an open subgroup is
injective on rational cohomology by corestriction.

\proofstep{The negative-weight strict local condition.}
For negative \(j\), write \(j=-n\) with \(n>0\). The corresponding characters have the form
\[
 \epsilon_\ell^{ej}\theta=\epsilon_\ell^{-en}\theta,\qquad\theta\text{ finite}.
\]
At \(\ell\) the Gaussian local condition kills the class;
a crystalline tangent at the split point has finite local
condition, which is zero for a negative Tate twist of a
finite character. Indeed its de Rham quotient by
\(\mathrm{Fil}^0\) and its invariant space are zero.
Restriction to a finite field killing the finite local
type gives the same finite condition by averaging over
the finite quotient. Away from \(\ell\) an unramified
class with these coefficients vanishes because Frobenius
minus one on its inertia invariants is invertible. Even the full local rational cohomology vanishes here.

For this one-dimensional character space \(V\) away from \(\ell\), the local Euler characteristic formula and duality give
\[
 \dim H^1_v(V)=\dim H^0_v(V)+\dim H^0_v(V^\vee(1))=0 .
\] Thus the
global class becomes everywhere locally split on a finite
abelian extension \(L\) killing \(\theta\) and containing
\(\mu_\ell\). The field is totally complex since
\(\ell\) is odd. Take \(S\) over \(L\) to contain the
places above \(\ell\) and the ramification of the cocycle
and coefficients. The classes under consideration are
unramified outside such a finite \(S\).

\proofstep{Rational vanishing from the height and duality inputs.}
Use the finiteness part of the height estimate of
Section~\ref{sec:cohom-height} at the fixed field \(L\)
and exponent \(en+1\). For a \(G_{L,S}\)-module \(M\) write \(\operatorname{Sha}^i(M)\)
for the kernel of localization in degree \(i\), with all
places of \(S\) included. Poitou--Tate gives
\[
 \operatorname{Hom}\!\left(\operatorname{Sha}^1(\Q_\ell/\Z_\ell(-en)),
                                     \mathbb Q/\mathbb Z\right)
       \simeq \operatorname{Sha}^2(\Z_\ell(en+1)).
\]
The height input makes the right side finite, so the strict
torsion kernel on the left has finite order.
Extra places of \(S\) preserve the required finiteness bound
on this kernel: in the degree-two localization sequence
the new terms are local terms, where the class in the
kernel has zero image.

Equivalently an everywhere locally
split cocycle in degree one factors through the group
unramified outside the places above \(\ell\).
Finite coefficient extension preserves this finiteness.
A nonzero rational class would give torsion classes of
unbounded order: multiply a rational class by
arbitrary inverse powers of the uniformizer and project
to the divisible coefficient module. Let \(V\) be the rational coefficient representation and \(T\) its lattice. The sequence of coefficients identifies the kernel of projection on cohomology:
\[
 \ker\bigl(H^1(V)\to H^1(V/T)\bigr)=\operatorname{im}\bigl(H^1(T)\to H^1(V)\bigr).
\]
The image is a lattice in its rational span by finiteness of the residual global cohomology. Local splitting persists
under projection. This contradicts finiteness of the
strict torsion kernel.

For a cocycle with values in an
algebraic coefficient closure, a finite extension contains
its values. To see this, use the countable family of finite
extensions containing the field of definition of the action.
Their preimages under the cocycle are closed subgroups,
and Baire's theorem gives one with open preimage. For a cocycle \(z\) let \(H\) be the open preimage obtained in this way. If \(g_1,\ldots,g_r\) represent \(G/H\), adjoining the finitely many coordinate values of \(z(g_i)\) suffices by
\[
 z(g_i h)=z(g_i)+g_i z(h)\qquad(h\in H).
\] This proves the nonpositive-weight vanishing
over both fields of coefficients.

An infinitesimal normalized point at \(r_0\) now has
zero derivative in each of these weights. Its matrix
derivative times \(r_0(g)^{-1}\) is a cocycle, by the
product law with square-zero coefficients. A coboundary
is infinitesimal conjugation by \(1+\varepsilon X\).
Fixing \(r(c)\) and the normalization at \(g_*\) imposes the simultaneous conditions
\[
 X=-JX^{\mathsf T}J,\qquad X\text{ diagonal}.
\] Such an
\(X\) is a multiple of \(\operatorname{diag}(1,0,-1)\),
which centralizes the split polarized representation.
Since matrix entries separate infinitesimal deformations,
the chart derivative is zero. Dualizing, the equations
span the rational cotangent spaces in \eqref{eq:local-length-2}. In the
Gaussian power series presentation this can also be
tested over \(O\): a rational vector in the annihilator
would have an integral nonzero multiple annihilating
the equations and would give a normalized point over
\(O[\varepsilon]/(\varepsilon^2)\).
\hfill\mbox{\cite{bloch-kato-finite,milne-duality}}\label{cor:ordinary-rank-three-length-2}

\end{proof}

\paragraph{Transfer at every torsion depth.}
Rational vanishing deals with the generic fibre. A length comparison requires a map
that detects classes even at finite depth.

\begin{proposition}[Integral transfer at all coefficient depths]
\label{prop:rank-three-integral-transfer}
The weight-minus-one localization kernel with coefficients
\((E/O)(\Psi^{-1})\), on either of the two designated hosts,
injects at each torsion depth into the derivations in condition
\textup{\eqref{eq:local-length-1}}. Lifting uses the Gaussian chart or the crystalline
local base, respectively. Let \(G\) be this kernel over the coefficient ring \(O\) with
uniformizer \(\pi\) and absolute ramification index \(e_O\). A uniform
length bound from \eqref{eq:local-length-3} passes to the full kernel:
\[
 (\forall h\ge1,\ \operatorname{length}_O G[\pi^h]\le L)
       \quad\Longrightarrow\quad \operatorname{length}_O G\le L .
\]
The normalized bound is divided by \(e_O\) to use \(v_\ell\).
\end{proposition}

\begin{proof}
\proofstep{The weight-minus-one cocycle.}
Let \(W=(E/O)(\Psi^{-1})\) and \(b\in O\setminus\{0\}\).
For brevity put \(W[b]=(O/(b))(\Psi^{-1})\), using the
isomorphism \(u\mapsto u/b\bmod O\). A class in the
rank-one localization kernel killed by \(b\) has a
representative valued in \(W[b]\). In fact if \(bz\)
is the coboundary of \(u\in W\), choose a vector \(u'\)
with \(bu'=u\) and subtract its coboundary. The class
in finite-coefficient global cohomology is unique, since
\(\Psi(g_*)^{-1}-1\) is a unit and thus \(W^\Gamma=0\).

On an inertia subgroup on which \(\Psi=1\), the action on the coefficient is trivial. Thus a coboundary on that subgroup satisfies
\[
 (da)(i)=i\cdot a-a=0,
\]
and the finite representative vanishes there. At a decomposition group where some
\(\Psi(g)^2-1\) is a unit, a trivializing vector in \(W\)
for this representative is killed by \(b\). Indeed
its product with \(\Psi(g)^{-1}-1\) is killed by \(b\).
Consequently the local class already vanishes with
coefficients in \(W[b]\).

The weight \(-1\) adjoint embedding of \(W[b]\) uses
\(T=E_{10}-E_{21}\), where \(E_{ij}\) is the elementary
matrix in row \(i\) and column \(j\). Write
\(g_*'=cg_*^{-1}c^{-1}\). Normalize the cocycle \(z\)
by making
\[
    L(z)=\Psi(g_*)z(g_*)-\Psi(g_*')z(g_*')
\]
zero. Subtract the coboundary of
\(L(z)/(\Psi(g_*)^{-1}-\Psi(g_*))\), which is defined
since the denominator is a unit. Also \(z(c)=0\)
because \(c^2=1\), \(\Psi(c)=1\) and \(2\) is a unit.
The representation
\[
             (1+\varepsilon z(g)T)r_0(g),
       \qquad \varepsilon^2=0,\quad b\varepsilon=0,
             \tag{6}\label{eq:local-length-6}
\]
therefore satisfies the two normalization conditions.
Write \(D_{10}(g)\) for the coefficient of \(\varepsilon\) in the \((1,0)\) entry of \eqref{eq:local-length-6}. Since \(T_{10}=1\), the entry calculation gives
\[
    D_{10}(g)=\Psi(g)z(g),\qquad
    z(g)=\Psi(g)^{-1}D_{10}(g)\quad(g\in\Gamma).
\]
Thus the resulting derivatives determine the normalized cocycle, and addition of cocycles adds the derivatives entrywise.

\proofstep{Integral lifting on the Gaussian chart.}
For the Gaussian chart quotient by the torsion-depth
local conditions at \(\ell\) and at the exceptional
tame places where \(\Psi(g)^2-1\) is a unit on some
element of the decomposition group. These are precisely
the places where weight \(-1\) must be a local coboundary
in \eqref{eq:local-length-6}. At the other tame places the representation
kills the residual-kernel inertia, on which
\(\Psi=1\). The normalized representation therefore
defines a point of its universal algebra. In the ambient power series ring centered at the split point, let \(X_i\) be its coordinates.

For each coordinate of weight \(-1\) choose \(d_i\in O\) lifting the given derivative in \(O/(b)\); set \(d_i=0\) in every other weight. The point derivation
\[
  D(f)=\sum_i d_i\,\frac{\partial f}{\partial X_i}(0)
\]
is integral. Modulo \((b)\) it gives the original infinitesimal point, so \(D(F)\equiv0\pmod b\) for every equation \(F\) of the chart. The original
localization group on the absolute Galois group
descends to \(G_{\Q,S}\), with \(S\) the primes dividing
\(\ell N\), because its cocycles vanish on all inertia
outside \(S\). Thus \eqref{eq:local-length-1} bounds every \(b\)-torsion
subgroup by the same cotangent length.

\proofstep{Integral lifting on the crystalline chart.}
For the general-data chart the global condition kills
the residual-kernel inertia at each place away from
\(\ell\). The datum makes \(\Psi\) trivial on those
subgroups. For a class \(z\) in this localization kernel the condition at each prime dividing \(\ell N\) is on the full group:
\[
 \operatorname{res}_{D_p}[z]=0\qquad(p\mid\ell N).
\] Otherwise it is zero already on
inertia. In either case the finite representative vanishes
on the subgroup to be killed, including above primes
ramified only in the quadratic field.

At \(\ell\) the mark's unit separator gives a trivializing vector in \(O/(b)\). Lift a conjugating matrix for this vector to a weight \(-1\) matrix \(Y\) over \(O\), choosing its sign to recover the given coboundary. The integral local lift is
\[
 (1+\varepsilon Y)r_0(g)(1+\varepsilon Y)^{-1}
    =r_0(g)+\varepsilon\bigl(Yr_0(g)-r_0(g)Y\bigr),
       \qquad \varepsilon^2=0
\]
on the positive component, with the polarized action extending it. Being conjugate to the split point, it lies on the crystalline local quotient. The ambient global chart is a relative
power series algebra over that quotient, as described
below. Sending its remaining coordinates to arbitrary
integral lifts of their derivatives in weight \(-1\)
gives an integral point derivation of the ambient algebra.
Its reduction comes from \eqref{eq:local-length-6} and vanishes on every
global equation. This is the lift required by \eqref{eq:local-length-1}.
Lifting on the crystalline local quotient first is
necessary for this use of saturation.

\proofstep{The union over all coefficient depths.}
The divisible module \(W\) is the union of its
finite torsion submodules, and a continuous cocycle
into it has finite image by compactness. Thus the
localization group is a union of its \(\pi^h\)-torsion
subgroups. For the localization group \(G\) the torsion filtration and its common length bound \(L\) satisfy
\[
 G_h=G[\pi^h]\subseteq G_{h+1},\qquad \operatorname{length}(G_h)\le L\quad(h\ge1).
\] Additivity gives
\[
 \operatorname{length}_O(G_{h+1}/G_h)
 =\operatorname{length}_O G_{h+1}-\operatorname{length}_O G_h .
\]
Every strict inclusion contributes at least one; the sum of these increments is at most \(L\). The chain stabilizes and its union \(G\) has length at most \(L\). Dividing by the
ramification index of \(O/\Z_\ell\) takes the bound
in \eqref{eq:local-length-3} from \(v_\pi\) to \(v_\ell\). The procedure
keeps the packet bound independent of torsion depth
and of the residue-field and coefficient degrees.
\label{cor:ordinary-rank-three-length-3}

\end{proof}

\subsection{The Gaussian inertia chart}

The integral comparison is ready to use once the local presentation is known. We
construct it first for the Gaussian point, where the deformation fixes inertia.

\begin{proposition}[The Gaussian fixed-inertia presentation]
\label{prop:gaussian-rank-three-chart}
Under the Gaussian mark and its local conditions, the graded
presentation has the dimensions and equations of \textup{\eqref{eq:local-length-2}},
with a finite degree-zero ring as required for
Proposition~\ref{prop:rank-three-conormal}. Let \(n_{-1}\) be the number of weight-minus-one variables and \(G\) the
localization kernel of Proposition~\ref{prop:rank-three-integral-transfer}.
The coordinate map in \eqref{eq:local-length-6} gives
\[
                         n_{-1}=0\quad\Longrightarrow\quad G=0 .
\]
\end{proposition}

\begin{proof}
\proofstep{The fixed-inertia local problem.}
Use \(S=\{v:v\mid\ell N\}\), the Auxiliary field over \(\mathbb Q(i)\) quotient
\(\Gamma=G_{\Q,S}\), and the one mark from
Section~\ref{sec:marked-hosts}. At \(\ell\) the local
representation preserves a flag with the three ordered
inertia characters \(\Psi,1,\Psi^{-1}\). At an exceptional
tame place it is diagonal up to a change of basis
congruent to the identity, with the prescribed inertia
characters. On the other tame decomposition groups the inertia subgroup killed by the representation is prescribed by
\[
 N_v=\ker(\eta|_{I_v})\cap\ker(P|_{I_v}),\qquad P=\bar\Psi .
\] This is the subgroup on which the
coefficient character is trivial there. In both tame
cases \(I_v/N_v\) has order prime to \(\ell\).
The local determinant or polarization condition is the
one inherited from the semidirect product.

\proofstep{Lifting the local conditions across an extension.}
These local conditions lift across a graded Artinian
small extension. Here the kernel is killed by the
maximal ideal and is supported in one weight. At
\(\ell\), use the eigenspaces of the separator to
choose the flag basis, with the lines labelled by their
reductions. Let \(T\) be the lifted separator and \(\alpha_i\) its labelled eigenvalues. The projector onto its \(i\)-th eigenspace is integral, because
\[
 e_i=\prod_{j\ne i}\frac{T-\alpha_j}{\alpha_i-\alpha_j},\qquad \alpha_i-\alpha_j\in O^\times\quad(i\ne j).
\]
Thus the necessary changes of flag basis with the prescribed entry weights lift. In the flag basis lift the torus
action by taking the given inertia characters and
lifting a Frobenius in the semidirect torus. The
procyclic unramified quotient permits this choice,
also when \(\eta\) acts on the torus. Lift the
upper triangular entries in order of root height.
At each step let \(u(g,h)\) be the multiplication error in
the small-extension kernel of that root weight \(j>0\).

Associativity of the lifted products gives the two-cocycle identity for their additive defect:
\[
 g\,u(h,k)-u(gh,k)+u(g,hk)-u(g,h)=0 .
\] The separator \(x\) has
\(\bar\epsilon_\ell(x)=\eta(x)=1\) and \(P(x)^2\ne1\).
On \(\operatorname{Hom}(V_j,k(1))\) it acts by \(P(x)^{-j}\),
which is nontrivial for both possible positive weights \(1,2\).
Therefore
\[
 H^2(\Q_\ell,V_j)\simeq H^0(\Q_\ell,V_j^\vee(1))^\vee=0.
\]
Write \(u=da\) for a continuous cochain \(a\). Subtracting
\(a\) from the selected root lifts removes the multiplication
error while keeping the other weights fixed. At each exceptional tame place, the separator kills invariants and dual invariants in nonzero weights. Local duality and the Euler characteristic in residue characteristic different from \(\ell\) then give 
\[
 h^2=0,\qquad h^1=h^0+h^2=0,
\]
 where \(h^i\) denotes the local dimension in the selected nonzero weight. A canonical
eigenbasis there diagonalizes the inertial action,
and the decomposition group preserves these lines
with the action prescribed by \(\eta\).

Lifting its
Frobenius in the torus and changing back to the original
basis gives the desired lift. At the remaining tame
places lift the finite prime-to-\(\ell\) action of
\(I_v/N_v\) by averaging. Its lifts are conjugate by
a matrix in the kernel. Let \(r_I\) be the lifted inertia action and \(F\) a first Frobenius lift. Reduction identifies the two inertia maps, since the residual representation already respects conjugation:
\[
 \overline{F r_I(i) F^{-1}}=\overline{r_I(\mathrm{fr}\,i\,\mathrm{fr}^{-1})}.
\] The averaging argument conjugates
one to the other by a kernel matrix \(C\). Replacing \(F\)
by \(CF\) gives
\[
         (CF)r_I(i)(CF)^{-1}
                    =r_I(\mathrm{fr}\,i\,\mathrm{fr}^{-1}).
\]
This is the required Frobenius-inertia relation. The
extension of the finite inertia group by the unramified
procyclic quotient is generated by inertia and one
Frobenius with those relations. All the linear
corrections in the kernel can be taken in the chosen
weight, because the residual adjoint action preserves
the grading.

\proofstep{The tangent and obstruction dimensions.}
We give the dimension computation controlling the global
relations. Fix a weight \(j\), let \(d=\dim V_j\),
let \(d_\infty=\dim V_j^{c=1}\), and write \(g_i,h_{v,i}\)
for the dimensions of the global and local cohomology
groups. For the space of permitted local cocycle classes the inclusion is
\[
 L_v\subseteq H^1(\mathbb Q_v,V_j).
\] At a tame place
\(\dim L_v=h_{v,0}\). Indeed restriction to inertia
of a cocycle vanishing on \(N_v\) is a coboundary by
averaging. Let \(T\) be Frobenius on \(V_j^{I_v}\). Rank-nullity identifies the dimension of its coinvariants with that of its invariants:
\[
 \dim\operatorname{coker}(T-1)=\dim\ker(T-1).
\]

At exceptional places the same
formula holds by the specified inertia and the
vanishing of \(H^1\) in nonzero weights.
At \(\ell\), \(L_\ell=0\) for \(j<0\), is unramified
for \(j=0\), and is all of \(H^1\) for \(j>0\).
The tangent assertion for the flag follows directly
from the ordered characters. A negative-weight
entry comes only from changing the flag basis, a
weight-zero entry can vary its unramified character,
and an upper-weight cocycle gives a flag-preserving
deformation with fixed diagonal.

Let \(t_j\) be the dimension of the Selmer tangent
space and \(o_j\) the dimension of its obstruction
space. If \(\lambda\) is the rank of the localization
map to \(\bigoplus_v H^1(\Q_v,V_j)/L_v\), then
\[
 \begin{split}
 t_j&=g_1-\lambda,\\
 o_j&=\sum_v (h_{v,1}-\dim L_v)-\lambda
        +g_2-\sum_v h_{v,2}.
 \end{split}
 \tag{7}\label{eq:local-length-7}
\]
The first part of \(o_j\) is the degree-one localization
cokernel. The rest is the kernel of localization in
degree two in the cochain complex recording a global
multiplication error and its local trivializations.
The degree-two localization map is surjective by the
Poitou--Tate end sequence.

In fact
\(H^0(\Q,V_j^\vee(1))=0\) for each \(j\): for
nonzero \(j\) use the preceding separator, and for
\(j=0\) use an element of inertia at \(\ell\) with
nontrivial residual cyclotomic value. The global
Euler characteristic is
\(g_1+ d_\infty=g_0+g_2+d\).
The local Euler characteristic is
\(h_{v,1}=h_{v,0}+h_{v,2}+d\,1_{v=\ell}\).
Subtracting the two lines of \eqref{eq:local-length-7} gives
\[
           t_j-o_j=g_0-d_\infty+\dim L_\ell-h_{\ell,0}.
           \tag{8}\label{eq:local-length-8}
\]
For \(j\ne0\) the global and \(\ell\)-local invariants
are zero by the separator. In positive weights
\(h_{\ell,1}=d\) because \(h_{\ell,2}=0\).
At weight zero \(g_0=d_\infty=1\).
Consequently the five defects, in increasing weight,
are \(0,-1,0,1,1\).

For clarity, these are bounds for a presentation of
the normalized deformation problem. The residual
representation is fixed, the flags in the local
conditions have the marked residual labels, and the
inertia conditions are invariant under a change of
frame. At the global origin any infinitesimal
cocycle can have its value at \(c\) and the
off-diagonal derivative of the marked difference removed
by a coboundary. If \(z\) is the adjoint cocycle and
\(dY(g)=g\cdot Y-Y\), then
\[
 z(c)+c\cdot z(c)=0,\qquad d\bigl(-z(c)/2\bigr)(c)=z(c).
\]
Subtract this coboundary first. The remaining normalization
uses a \(c\)-invariant matrix \(Y\).

In the eigenbasis of the
marked operator with eigenvalues \(\lambda_i\), let \(D=\operatorname{diag}(\lambda_i)\). To realize the required off-diagonal change \(Z\) in the derivative solve
\[
 [D,Y]_{ij}=(\lambda_i-\lambda_j)Y_{ij}=Z_{ij},\qquad
 Y_{ij}=\frac{Z_{ij}}{\lambda_i-\lambda_j}\quad(i\ne j).
\]
The denominators are units at the mark. The prescribed change belongs to the invariant off-diagonal space, on which this commutator is invertible. The remaining conjugators centralize the split
point. Thus normalized tangents contain
one representative of every Selmer class.

The cohomology of every open subgroup of \(G_{\Q,S}\)
with finite \(\ell\)-torsion coefficients is finite in
degree one. This is the finite-dimensional tangent
condition for pro-representability of the framed
deformation problem. The frame supplies a complete
Noetherian universal ring. Its matrix entries,
after subtracting their split values, give
homogeneous variables. Choose homogeneous entries whose classes form a basis of the relative cotangent space modulo the maximal ideal. With \(r\) such entries \(x_i\), the resulting presentation over the coefficient ring \(O\) has the map
\[
 O[[X_1,\ldots,X_r]]\twoheadrightarrow R,\qquad X_i\longmapsto x_i .
\] That choice generates by complete
Nakayama. The variables have weights in
\([-2,2]\). The ideal of relations is homogeneous,
as follows by applying the torus character with an
indeterminate unit and comparing weights in each
finite quotient.

\proofstep{The relation count for the completed global presentation.}
Here is the relation bound underlying this presentation.
For a graded small extension, choose a set-theoretic
lift of the matrices and local homomorphism lifts.
Let \(I\) denote the square-zero extension kernel, \(c\) the global multiplication error and \(b_v\) the local correction at \(v\). In the mapping-cone convention of \eqref{eq:local-length-7}, associativity and the local lifting equations give
\[
  dc=0,\qquad \operatorname{res}_v(c)-db_v=0 .
\]
If the global lift changes by \(\alpha\), let \(\lambda_v\) be the difference between the two permitted local lifts, read as a local tangent cocycle. Its class lies in the prescribed local subspace \(L_v\). The change of the error pair is
\[
   \bigl(d\alpha,(\operatorname{res}_v\alpha-\lambda_v)_v\bigr).
\]
A local change of basis gives the special case \(\lambda_v=d\beta_v\).

In the relative complex with the specified local tangent conditions, all these changes give the same obstruction class, as reflected in the quotient by \(L_v\) in \eqref{eq:local-length-7}. If a functional \(u\) on the extension kernel kills the obstruction class \(o\), its pushout has obstruction
\[
 o_{\rm pushout}=u_*(o)=0 .
\]
Hence there is a global lift with the required local conditions. The normalization at \(c\)
and the mark can be restored there by a kernel
conjugation as above. Apply this to the dual of the
space of minimal relations of one weight in the
power series presentation. The calculation can be
made in a finite quotient inducing the same space
of minimal relations, by Artin--Rees. A functional
giving a lift would let the universal ring map back
to the corresponding nonzero extension of its
Artinian quotient.

If the two maps out of the power series ring differ by \(\delta_i\in I\) on the variables, then for each minimal relation \(F\) their difference is
\[
 F(\mathbf X+\boldsymbol\delta)-F(\mathbf X)
     =\sum_i (\partial_i F)(\mathbf X)\,\delta_i=0.
\]
The Taylor expansion ends after its linear part because \(I^2=0\). For the minimal relations each linear coefficient lies in the maximal ideal, so 
\[
 (\partial_i F)\delta_i\in\mathfrak m I=0\qquad\text{for every }i.
\]
 The section would thus kill the selected relation even in the extension where its class is nonzero, a contradiction. Hence the dual
of the minimal relation space injects into
the obstruction space. This proves the bound of
\(o_j\) relations in weight \(j\).

If \(n_j\) is the number of chosen variables, then
\(n_j\ge t_j\). By \eqref{eq:local-length-8} the relation bounds can be
padded by zero equations to
\[
 (n_{-2},\ n_{-1}+1,\ n_0,\ n_1-1,\ n_2-1).
\]
There is at least one variable in each positive
weight by \eqref{eq:local-length-8}. In case \(n_{-1}=0\) every
localized class has zero image in the coordinate
derivations in \eqref{eq:local-length-6}, so the desired length is zero.
Otherwise these variables and equations have
exactly \eqref{eq:local-length-2}. Take \(A\) to consist of finite
sums of homogeneous series in the ambient power
series ring. Its split formal chart has the
specified weight dimensions. Homogeneous projection
shows the ideal of the quotient is generated in
\(A\) by the homogeneous equations. The invariant
ring in the quotient is finite by the Gaussian
fixed-mark argument, and it also makes the
rescaling to integral points possible. This gives
the chart to which \eqref{eq:local-length-3} applies.
\hfill\mbox{\cite{milne-duality}}\label{cor:ordinary-rank-three-length-4}

\end{proof}

\subsection{The crystalline chart for general data}

For the general point the local calculation at the coefficient prime uses a
crystalline chart. In the rank-three trace estimate the order of \(\psi\) is prime to \(\ell\),
so its values can be taken over an unramified
coefficient DVR. Normalized length is invariant
under further coefficient extension. Work over
the general datum with \(\ell\nmid D_K h_K\),
where \(h_K\) is the class number. Away from \(\ell\)
kill \(I_v\cap\ker\eta\cap\ker P\); the split lift
kills it because the finite inertia character has
order prime to \(\ell\). We may work with the
finite set of places above \(\ell N|D_K|\), using
the normal closure of the killed subgroups in the
global deformation problem.

\paragraph{The base at \(\ell\).}
To keep the codimension count integral we must control saturation at the coefficient
prime. The following elementary calculation is the local ingredient.

\begin{auxlemma}[A coefficient test for saturation]\label{lem:crystalline-coefficient-test}
Let \((B,\mathfrak n)\) be a local ring with \(\mathfrak n^h=0\).
Suppose \(P,y\in B[t]/(t^n)\) satisfy \(Py=0\) and the image of \(P\)
in \((B/\mathfrak n)[t]/(t^n)\) is \(t^b\) times a unit, where
\(0\le b<n\). Evaluation at the origin then satisfies
\[
                       n>hb\quad\Longrightarrow\quad y(0)=0\quad\text{in }B .
\]
\end{auxlemma}
\begin{proof}
Write \(y=\sum_{i<n}y_i t^i\). We prove
\[
  i<n-vb\quad\Longrightarrow\quad y_i\in\mathfrak n^v
  \qquad(0\le v\le h).
\]
For \(v=0\) the assertion holds in \(B\). Suppose it holds at \(v<h\)
and truncate modulo \(t^{n-vb}\). The truncated coefficients of \(y\)
lie in \(\mathfrak n^v\), so \(P\) acts on their classes in
\(\mathfrak n^v/\mathfrak n^{v+1}\) by its residue \(t^b u\), with \(u\)
invertible. Cancelling \(u\) in the truncated polynomial module gives
\[
 t^b\sum_{i<n-vb}[y_i]t^i=0\pmod {t^{n-vb}}.
\]
Consequently \([y_i]=0\) for \(i<n-(v+1)b\), as required.
At \(v=h\), \(0<n-hb\) puts \(y_0\) in \(\mathfrak n^h=0\).
\end{proof}

\begin{construction}
\[
\begin{array}{ccc}
 B &\longleftarrow& B[t]/(t^n)\\[-2pt]
 \downarrow\ {\scriptstyle\bmod\,\mathfrak n}&&
    \downarrow\ {\scriptstyle\bmod\,\mathfrak n}\\
 E&\longleftarrow&E[t]/(t^n)\\
 \multicolumn{3}{c}{\text{horizontal arrows evaluate at }t=0}\\[2pt]
 \multicolumn{3}{c}{\ker\bigl(B[t]/(t^n)\longrightarrow
       B\times_E E[t]/(t^n)\bigr)=t\mathfrak n[t]/(t^n)}
\end{array}
\]
\caption{The test algebra joining the given point to the unramified curve. Evaluation and reduction have common target in the residue field. Formal smoothness lifts the compatible pair across the resulting nilpotent kernel.}\label{fig:crystalline-thickening-test}
\end{construction}

\begin{proposition}[The saturated crystalline local chart]
\label{prop:rank-three-crystalline-local-base}
For the general-data host with \(\ell\nmid D_K h_K\) and
\(\psi\) of order prime to \(\ell\), kill
\(I_v\cap\ker\eta\cap\ker P\) on inertia away from \(\ell\).
Use the prime through
the split point in the radical of the ordered-inertia equations
saturated by \(\ell D_0D_1\). Its formal characteristic-zero chart is crystalline, including at
the split point. The dimensions \(d_j\) of the homogeneous coordinate
spaces in \eqref{eq:rank-three-crystalline-chart} are
\[
                (d_{-2},d_{-1},d_0,d_1,d_2)=(1,2,3,4,2).
\] On the selected global chart the test-curve
comparison gives a finite degree-zero algebra, and the
two-sided characteristic-zero orbits have integral representatives.
\end{proposition}

\begin{proof}
\proofstep{The ordered-inertia equations and their saturation.}
Let \(U_\ell\) be the unrestricted framed
local ring and \(H=\ker\eta\cap\ker\tau\) in
the local group, where \(\tau\) is the finite
character in \(\Psi|_{I_\ell}=\epsilon_\ell^e\tau\).
For \(s\in H\cap I_\ell\) put \(q_i(s)=\Psi(s)^{1-i}\),
\(0\le i\le2\), and \(S_i(s)=r(s)-q_i(s)1\).
Impose the characteristic polynomial with roots
\(q_0(s),q_1(s),q_2(s)\), and impose
\[
          S_0(s_0)S_1(s_1)S_2(s_2)=0
          \qquad(s_i\in H\cap I_\ell).
\]

Choose \(s\in H\cap I_\ell\) of infinite
cyclotomic image. Its three \(q_i(s)\) are
distinct. With \(c_i=\prod_{j\ne i}(q_i(s)-q_j(s))\),
the spectral projector onto the \(i\)-th
eigenline is \(c_i^{-1}\prod_{j\ne i}S_j(s)\)
over the fraction field. Take \(\mathrm{fr}\in H\)
generating the unramified quotient and let
\(T_i\) be the trace of
\(\prod_{j\ne i}S_j(s)r(\mathrm{fr})\).
The two adjacent differences are
\[
       D_i=c_{i+1}T_i-
            \epsilon_\ell(\mathrm{fr})c_iT_{i+1}
       \qquad(i=0,1).
\]
The local ideal is the prime through the split
point of the radical of the ordinary equations
saturated by \(\Delta=\ell D_0D_1\).
We explain the use of saturation in both directions.

\proofstep{The ordered crystalline equations after inverting the adjacent factors.}
In characteristic zero the ordinary equations give
an ordered flag. For the first line and last
quotient, use a variable inertia element in
the first and last factors of the triple product
respectively and the chosen \(s\) in the other two.
The projectors isolate the first line and the last quotient of the ordered flag. On those one-dimensional graded pieces the operator thus has the two scalar restrictions
\[
 q_0\,\mathrm{id}\quad\text{on the first line},\qquad q_2\,\mathrm{id}\quad\text{on the last quotient}.
\] The middle value follows from the trace.
These flags are unique with the ordered characters
because the eigenvalues at \(s\) are distinct.
The decomposition group preserves the flag with
the conjugate-dual action on the quadratic coset.
On \(H\), write \(\nu_1,\nu_2,\nu_3\) for the unramified twists. The ordered graded characters restrict to
\[
 \nu_1\epsilon_\ell^e,\qquad\nu_2,\qquad\nu_3\epsilon_\ell^{-e}.
\]
In the ordered basis \(T_i/c_i\) is the \(i\)-th
Frobenius value.

Fix a finite local characteristic-zero coefficient
algebra \(B\) on the saturated quotient. One can
factor its point through a finite local algebra
\(B'\) of that quotient such that any element of
\(B'\) killed by \(\Delta\) maps to zero in \(B\).
Indeed in the generic local ring of the saturated
quotient, \(\Delta\) is a nonzerodivisor. Apply
Artin--Rees to its principal submodule and the
powers of the maximal ideal \(\mathfrak m\).
Writing \(R\) for this local ring, choose \(c\) from Artin--Rees,
so \(\mathfrak m^h\cap\Delta R\subseteq\mathfrak m^{h-c}\Delta R\)
for \(h\ge c\). Since \(\Delta\) is a nonzerodivisor,
\[
 \Delta b\in\mathfrak m^h\quad\Longrightarrow\quad
 b\in\mathfrak m^{h-c}.
\]
Choose \(h\) large enough that this last ideal maps to zero in \(B\). Taking
\(B'\) to be this \(\mathfrak m^h\)-quotient
gives the assertion. Its residue field is
finite over the coefficient fraction field,
since the point has values in a finite
algebra.

\proofstep{The vanishing required by the saturation.}
On \(H\), a ratio of adjacent characters is an
unramified twist of \(\epsilon_\ell^e\).
For a positive Tate exponent \(a\), the
quotient of \(H^1\) by its finite subspace
is dual to the finite \(H^1\) of the Tate
dual, by local duality. For \(a>1\) the Tate dual has exponent \(1-a<0\). Its finite local group vanishes, so local duality gives
\[
 H^1/H_f^1=0\qquad(a>1).
\] For \(a=1\), write the coefficient as
\(\epsilon_\ell\otimes\theta\) with \(\theta\) unramified.
The Tate dual is \(\theta^\vee\). Its finite
\(H^1\) is the Frobenius coinvariant quotient.
Dualizing that quotient makes the possible nonfinite
part the kernel of \(\theta(\mathrm{fr})-1\).
In terms of the adjacent characters this scalar is
\[
 \theta(\mathrm{fr})-1
  = \frac{(\kappa_i/\kappa_j)(\mathrm{fr})
                          -\epsilon_\ell(\mathrm{fr})}
          {\epsilon_\ell(\mathrm{fr})}.
\]
The denominator is a unit.

Over \(B'\) a class in this
kernel is killed by the indicated factor of \(\Delta\),
hence is killed by \(\Delta\), so its image in \(B\)
is zero by the choice of \(B'\). This gives finite
extensions on the two adjacent subquotients.
For the outer character ratio the positive exponent is at least two, so the preceding local dimension calculation gives
\[
 2e\ge2\quad\Longrightarrow\quad H^1/H_f^1=0
\]
for that ratio.
There is no additional obstruction for the
full upper triangular representation. Once
the upper two-dimensional block is crystalline,
tensor it by the inverse of the last character.
Its two successive characters have positive
exponents. The finite \(H^1\) maps onto the
finite \(H^1\) of its last quotient by exactness
and the dimension formula for the de Rham
quotient. The remaining outer-line class is
finite as just proved. Thus the full point
over \(B\) is crystalline on \(H\).

\proofstep{Recovering the crystalline quotient, including the split point.}
Conversely take a crystalline deformation on \(H\)
over a local Artin coefficient algebra reducing
to the split lift. The Frobenius after a full
residue-field period has three distinct slopes
at the split lift and hence three separate
eigenlines over the Artin algebra. Order the three Frobenius eigenlines \(D^{(1)},D^{(2)},D^{(3)}\) by increasing slope. The initial subspaces for which crystallinity must be checked are
\[
 D_1=D^{(1)},\qquad D_2=D^{(1)}\oplus D^{(2)} .
\] To verify this with
nilpotent coefficients, regard the filtered
Frobenius module as a vector space over the
coefficient field. Its multiset of Hodge
weights is the split multiset repeated
according to the length of the Artin
algebra. The Newton number of each initial
eigenline sum is the sum of the corresponding
lowest Hodge weights. The induced filtration on a subspace of this dimension has
at least the sum of the smallest Hodge weights, counted with
those multiplicities.

If \(D_i\) is either initial eigenline
sum, and \(h_{\min}(D_i)\) is that minimum, then
\[
 h_{\min}(D_i)\le t_H(D_i)\le t_N(D_i)=h_{\min}(D_i).
\]
The middle inequality is weak admissibility of the full
object on the Frobenius-stable \(D_i\). Equality, together
with the inherited inequalities on all its subobjects, makes
\(D_i\) weakly admissible. The crystalline equivalence
gives the desired flag on the representation.
The graded terms are rank one over the
Artin algebra, with fixed Tate exponents
and unramified factors on \(H\). Freeness
can be checked by reduction and the
underlying vector space dimension.

On the overlap of two base changes the two pulled-back ordered flags have the same defining property. Uniqueness therefore gives the descent compatibility, with \(p_1,p_2\) the two projections:
\[
 p_1^*\mathcal F=p_2^*\mathcal F .
\] Its finite inertial
characters lift those at the split point,
because the order of a finite inertia
quotient is invertible over a characteristic-zero
coefficient field. This proves the ordinary
equations on every such deformation.

For a single crystalline point as above,
\(\Delta\) need not be a unit. To see that
it kills the saturated ideal, let \(\Delta^k a\)
lie in the ordinary ideal and use the following
deformation of the test algebra. Take the local algebra \(B\) with residue coefficient field \(E\) and maximal ideal \(\mathfrak n\) of nilpotence index at most \(h\):
\[
 B/\mathfrak n=E,\qquad \mathfrak n^h=0 .
\] Choose \(n>2kh\). On
\(C=E[t]/t^n\), vary the split lift by
an unramified torus character with value
\(\operatorname{diag}(1+t,1,(1+t)^{-1})\)
on a generator of the unramified quotient.
On \(\mathrm{fr}\) the power of \(1+t\)
has nonzero exponent. Each adjacent difference either has nonzero constant term or differs from a nonzero multiple of \(t\) by higher terms.

For the latter case in characteristic zero use
\[
 (1+t)^a-1\equiv at\pmod{t^2},\qquad a\ne0 .
\]
Its order in \(t\) is therefore one. The image of \(\Delta\)
is \(t^j\) times a unit with \(j\le2\).
The local crystalline functor at the split
lift is formally smooth. In the filtered
Frobenius description one lifts the invertible
Frobenius and the filtration by direct summands
across a nilpotent extension. Across the strict successive layers the Hodge and Newton numbers are additive. Writing \(D_i\) for the successive quotients, the equality starts with
\[
 t_H(D)=\sum_i t_H(D_i)=\sum_i t_N(D_i)=t_N(D).
\]
Strictness and weak admissibility of the layers give the subobject inequalities as well. Equivalence with
crystalline representations recovers the action
and one lifts the frame on its free coefficient
module. The fixed finite type on descent
is unchanged across the nilpotent extension.
This is the local smoothness of the crystalline
quotient theorem.

On the unramified
quadratic kernel the conjugate-dual descent
datum is retained by order-two averaging
on the characteristic-zero lifting and
change-of-basis spaces. The given \(B\)-point and the unramified twist over
\(C\) reduce to the same split lift. They therefore give
a point over \(B\times_E C\). The coefficient map
\[
       B[t]/t^n\longrightarrow B\times_E C,\qquad
             f(t)\longmapsto (f(0),\bar f(t))
\]
is surjective with kernel \(t\mathfrak n[t]/(t^n)\).
The kernel is nilpotent. Formal smoothness with the
polarized descent retained gives a lift across this map.

\proofstep{Keeping the integral local equation after saturation.}
Let \(P\) be the image of \(\Delta^k\)
and \(y\) the image of \(a\) in this
thickening. The ordinary equations give
\(Py=0\) and \(P\bmod\mathfrak n\) is
\(t^{jk}\) times a unit. At the two possible positive weights the nilpotence bound stays strictly below the truncation order:
\[
 j\le2\quad\Longrightarrow\quad hjk\le2kh<n .
\] Lemma~\ref{lem:crystalline-coefficient-test} with \(b=jk\) gives \(y(0)=0\). Evaluation at \(t=0\) recovers the given \(B\)-point, so that point kills \(a\), as required for every element of the saturated ideal.
The saturation consequently gives the
crystalline formal germ at the split lift,
including tangent points over dual numbers.
The completed germ is smooth and therefore
has a single component through that point.
Its closure defines the integral local base.

Conjugation congruent to the identity preserves
the split component. Indeed the action by the
universal congruence matrix on a domain
component lands in a formal power series domain
over that component. The preimage of its zero
prime under the conjugation map is contained in
the original component prime by specialization
to the identity. If \(P\) denotes the original component prime and \(Q\) the contracted zero prime, this gives a chain of primes
\[
 Q\subseteq P,\qquad P\text{ minimal}\quad\Longrightarrow\quad Q=P .
\] In particular the conjugate of
the split representation over the integral
dual numbers factors through the chosen local
quotient, as used in the transfer.

\proofstep{The local weight directions on the selected component.}
With characteristic-zero adjoint coefficients, the dimension at the split point is
\[
 \dim H^1_f(V_j)=h^0_j+\dim\bigl(D_{\rm dR}(V_j)/\mathrm{Fil}^0\bigr).
\]
Changing the actual frame contributes
\(\dim V_j-h^0_j\). The de Rham quotient has
dimension \(\dim V_j\) in positive weights and
zero in other weights. This follows from the
Tate twists \(\epsilon_\ell^{ej}\) on the
weight spaces, with only finite character
factors. The local framed smooth chart therefore
has dimensions
\[
                    (1,2,3,4,2)
                    \tag{9}\label{eq:rank-three-crystalline-chart}
\]
in increasing weight. Homogeneous matrix
entries give compatible axes. In the formal
completion over the fraction field, their
differentials span because a tangent annihilating
all entries gives the split deformation.

Choose
a homogeneous basis among them and change formal
coordinates by the inverse function theorem.
This gives the graded local chart, with the
same weights on the ambient integral
presentation. The integral component injects into
the smooth completed chart, since its defining
prime is the contraction of the prime of that
chart. Matrix entries map to homogeneous functions.
Thus for a relation with homogeneous components \(F_j\), projection to any fixed weight in the completed chart yields
\[
 \sum_j F_j=0\quad\Longrightarrow\quad F_i=0\quad\text{in weight }i .
\] By injectivity on the
integral component, the component ideal is closed
under these weight projections. The integral
relative presentation inherits this completed
grading.
\hfill\mbox{\cite{bloch-kato-finite,colmez-fontaine,kisin-crystalline}}\label{cor:ordinary-rank-three-length-5}

\proofstep{Inertia on test curves.}
The ordinary flag at the generic point of
the same component has the full prescribed
inertia characters. On \(H\cap I_\ell\)
these are given by the ordinary equations.
For a character ratio \(\delta\) trivial on \(H\cap I_\ell\), the remaining inertia datum factors through the finite quotient
\[
 I_\ell\twoheadrightarrow I_\ell/(H\cap I_\ell)\xrightarrow{\ \delta\ }\operatorname{im}(\delta).
\]
Specialization at the split lift fixes this finite factor.
More explicitly the eigenspaces of \(r(s)\)
at the separating element \(s\) label the
flag across any specialization in which
the \(q_i(s)\) stay distinct, including
the split lift over the fraction field.
The ratios at other inertia elements
have finite order and hence are constant
on the component on this labelled
chart once their value at the split
point is fixed. The ordered flag with
its prescribed action is a closed
condition on the projective flag
scheme over the integral local base.

Let \(Z\) denote that closed flag locus and \(S\) its integral base. Properness makes the image closed, so
\[
 \eta_S\in\operatorname{im}(Z\to S)\quad\Longrightarrow\quad\operatorname{im}(Z\to S)=S .
\]
Each field specialization thus has a flag after extending its residue field.
This also follows by extending the
generic flag over valuation rings.
For a map to a complete local DVR, write \(A\) for the
matrix of the separator. Its distinct reduced eigenvalues
lift by Hensel's lemma to roots \(a_i\) with unit differences.
The projectors
\[
                   e_i=\prod_{j\ne i}
                        \frac{A-a_j}{a_i-a_j}
\]
give the integral eigenlines. The ordered flag after field
extension is the sum of these lines in the prescribed residual
order, so it descends to the original valuation ring. This supplies the marked
ordinary condition even when the DVR
kills the coefficient uniformizer.

\proofstep{Applying the residual polynomial comparison to this chart.}
Now apply the finite-end statement of
Section~\ref{sec:marked-hosts} to
each such characteristic-\(\ell\) DVR
point of the crystalline global ring.
Away from \(\ell\), the inertia elements to be killed form the intersection
\[
 \ker(\eta|_{I_v})\cap\ker(\Psi|_{I_v}) .
\]
For the general datum this agrees with the prescribed residual kernel. The statement
makes the characteristic polynomial
split on a fixed subgroup of finite
index. The finite-index argument in
that section makes it split on the
whole quadratic kernel, and the mark
then gives triangularity with the
split diagonal. Thus every diagonal
coordinate minus its split value, and
every balanced monomial with a
positive and a negative entry, vanishes
on every test curve.

As on the auxiliary field over \(\mathbb Q(i)\), the curve test puts each such generator \(z\) in the radical of the uniformizer ideal. Equivalently, in the indicated local ring \(R\),
\[
 z\in\sqrt{\varpi R}\quad\Longleftrightarrow\quad z^m\in\varpi R\ \text{for some }m\ge1 .
\] Complete Nakayama gives
finiteness of the degree-zero algebra
over \(O\) on the selected chart.
The same valuation interval as in
the Gaussian case gives an integral
representative of every two-sided
characteristic-zero orbit. Here the
fact that this point is on the
crystalline quotient supplies the
finite condition at \(\ell\) for
its infinitesimal deformations.
\label{cor:ordinary-rank-three-length-6}

\end{proof}

\paragraph{The relative global presentation.}
With the integral base in hand we can count the remaining global variables and
relations relative to it.

\begin{lemma}[The normalized relative presentation]
\label{lem:relative-rank-three-presentation}
Over the integral local base of Proposition~\ref{prop:rank-three-crystalline-local-base},
the completed relative presentation after normalization at
\(c,g_*\) has the counts of \eqref{eq:local-length-10}. In increasing order of
weights \(-2,-1,0,1,2\), its numbers of coordinates and equations are
\[
 \begin{aligned}
 (n_j)_j&=(a_{-2}+1,a_{-1}+2,a_0+3,a_1+4,a_2+2),\\
 (r_j)_j&=(a_{-2}+1,a_{-1}+3,a_0+3,a_1+3,a_2+1).
 \end{aligned}
\]
These counts, its rational chart and the integral derivative lift
apply to the same quotient.
\end{lemma}

\begin{proof}
\proofstep{Separating the local and relative global presentations.}
Before normalizing at \(c,g_*\), the relative
deformation problem fixes a framed representation
at \(\ell\) on the chosen local quotient. In the
global presentation it records the local
conjugacy comparing the two frames. For a
residual weight space \(V_j\), write
\(g_i,h_{v,i}\) as in \eqref{eq:local-length-7}, and let \(\lambda\)
be the rank of localization into the full
\(H^1\) at \(\ell\) and the quotient by the
unramified conditions at the other places.
Those away conditions have dimension \(h_{v,0}\)
by averaging over the residual inertia quotient.
At \(\ell\) a global relative tangent has
zero cohomology class. Its local trivializing
vector can still vary by local invariants.

After taking the global changes of frame
into account, the relative tangent dimension is
\[
                a_j=g_1-\lambda+h_{\ell,0}-g_0 .
\]
Here the cancellation of the common change of
frame can also be seen by setting the local
conjugacy equal to the identity at \(\ell\)
and allowing the global frame to vary. For
the obstruction space the relative complex gives
\[
 h_{\ell,1}+\sum_{v\ne\ell}(h_{v,1}-h_{v,0})
                  -\lambda+g_2-\sum_v h_{v,2}.
\]
To justify the last sum, the general datum
gives an element with \(\eta=-1\),
\(P=1\) and \(\bar\epsilon_\ell=1\), from
ramification of \(K\) away from \(\ell N\).
Multiplying it by \(c\) gives an element
acting trivially on all \(V_j\) but with
cyclotomic value \(-1\). On the Tate-twisted dual the element just constructed acts as minus the identity. Hence in characteristic zero
\[
 v\in H^0(\mathbb Q,V_j^\vee(1))\quad\Longrightarrow\quad v=-v\quad\Longrightarrow\quad v=0 .
\]
Poitou--Tate then gives the claimed surjectivity in degree two.
The local and global Euler formulas now give
obstruction dimension \(a_j+\dim V_j^{c=1}\).

\proofstep{The relative equations and their weights.}
The relative ring is a quotient of a power
series algebra over the local crystalline
quotient with \(a_j\) additional variables
of weight \(j\). Indeed the relative tangent
coordinates lift a basis modulo the
local base and its maximal ideal. The
small-extension computation for the Gaussian
chart applies with the lift at \(\ell\)
prescribed by this base. One records its
change of frame in degree one of the
relative complex. The tame lift still
uses the finite inertia quotient.
For each weight \(j\), write \(\mathrm{Rel}_j\) for the space of relative minimal relations and \(\mathrm{Obs}_j\) for the corresponding relative obstruction space. The small-extension test gives
\[
 \mathrm{Rel}_j^\vee\hookrightarrow\mathrm{Obs}_j,\qquad\dim\mathrm{Rel}_j\le\dim\mathrm{Obs}_j .
\] This proves the
claimed relation bound over the local base.
Represent the additional coordinates by
matrix entries, so their weights stay in
\([-2,2]\). Zero equations pad each weight
to the bound.

\proofstep{The eight normalization equations.}
There are eight normalization equations.
Applied at \(c\), the relation \(c^2=1\) imposes the following symmetry on the matrix \(Y\) of the quadratic coset:
\[
 Y=JY^{\mathsf T}J .
\] The entries in positions
\((0,0),(1,1),(0,1),(1,0),(0,2),(2,0)\)
therefore suffice to set \(Y=1\).
Once this is done the marked difference is
\(Z-JZ^{\mathsf T}J\) with \(Z=r(g_*)\).
The off-diagonal equations reduce to the
\((0,1)\) and \((1,0)\) entries, one of
each weight \(1,-1\). Thus the normalization
adds \((1,2,2,2,1)\) equations in
increasing weight. The combined chart has
\[
\begin{array}{c|ccccc}
j&-2&-1&0&1&2\\ \hline
\text{ambient dimensions}&
a_{-2}+1&a_{-1}+2&a_0+3&a_1+4&a_2+2\\
\text{equations}&
a_{-2}+1&a_{-1}+3&a_0+3&a_1+3&a_2+1 .
\end{array}
\tag{10}\label{eq:local-length-10}
\]

We retain the entire integral local base
when forming the relative power series
ring. Its formal chart over the fraction
field is obtained by adjoining the same
relative variables to \eqref{eq:rank-three-crystalline-chart}. Passing to the
finite sums of homogeneous components
commutes with the quotient by the
homogeneous equations by projection on
each weight. Thus the smoothness,
equation counts and integral derivative
lifting required for \eqref{eq:local-length-3} and \eqref{eq:local-length-1} refer
to the same chart.
\hfill\mbox{\cite{milne-duality}}\label{cor:ordinary-rank-three-length-7}

\end{proof}

We now have the local and global models used in
Proposition~\ref{prop:rank-three-conormal}. The Gaussian
model and the crystalline model are used in separate
applications. In each one the bound is read on the original
coefficient ring.

\subsection{Frobenius trace bounds for the original Galois representations}

Return to the representations before passage to the marked field. The comparison on
either chart now expresses the length bound in terms of traces in the corresponding
packet.

\begin{proposition}[A length bound in terms of Frobenius trace differences]
\label{prop:rank-three-original-trace-bound}
Use the packet over the original field and the normalized Gaussian
or crystalline chart of Propositions~\ref{prop:gaussian-rank-three-chart}
and~\ref{prop:rank-three-crystalline-local-base}, respectively, with the
relative presentation of Lemma~\ref{lem:relative-rank-three-presentation}.
The normalized
length \(s_\ell(\psi,e)\) satisfies \textup{\eqref{eq:local-length-11}}: its bound is
twice the sum of valuations of the selected Frobenius trace
differences over representatives of the two-sided orbits.
For one character kernel \(J_0\), its rank-three contribution \(S(J_0)\) at the selected test prime satisfies 
\[
 S(J_0)\le 2v_\ell(A_{J_0}(p)),\qquad p\ne\ell.
\]

The Gaussian large-conductor application uses its separate
inert-prime packet. The packet exponents are those of the
original hosts.
\end{proposition}

\begin{proof}
We check the remaining point conditions of \eqref{eq:local-length-3}.
On a fixed point, the restriction to the
quadratic kernel is diagonal. Each of its
three characters differs from the corresponding
split character by a character unramified at
every finite place of \(K\) and congruent
to one. Here are the inertia checks for
that statement. At a tame place of residual
kernel type the diagonal character factors
on inertia through a group of order prime
to \(\ell\). Let \(m\) be the order of the finite inertia quotient. Reduction distinguishes the character values because \((m,\ell)=1\), so
\[
 \zeta,\zeta'\in\mu_m,\quad\bar\zeta=\bar\zeta'\quad\Longrightarrow\quad\zeta=\zeta'.
\] At a tame place of prescribed
inertia type the local eigenbasis has its
residual labels. The same argument using
the local eigenbasis at \(\ell\) gives the
ordered labels there.

In the general-data chart, on the inertia subgroup where the ordered polynomials are imposed, let \(\chi_i\) denote the three diagonal characters and \(\chi_i^0\) their prescribed values. Take a separating element \(h\) and label the three distinct roots \(\lambda_i=\chi_i(h)=\chi_i^0(h)\) in the fixed frame. Work over the characteristic-zero field of this fixed point. For any \(g\) in that subgroup, equality of the traces prescribed at \(h^a g\) gives
\[
 \sum_{i=1}^3 \lambda_i^a\bigl(\chi_i(g)-\chi_i^0(g)\bigr)=0
              \qquad (a=0,1,2).
\]
The determinant of these three equations is \(\prod_{i<j}(\lambda_j-\lambda_i)\ne0\). Thus each labeled value agrees, even when \(g\) itself has coinciding roots.

By unramified class field theory the quotient character \(\xi\)
factors through the class group of \(K\). Its image is finite and
congruent to one, and hence has \(\ell\)-power order. Thus
\[
 |\operatorname{im}\xi|\mid h_K,\qquad
 |\operatorname{im}\xi|=\ell^u,\qquad \ell\nmid h_K
 \quad\Longrightarrow\quad \xi=1.
\] For the Gaussian
case \(h_K=1\). The fixed value at complex
conjugation determines the representation
on the other coset. Hence the only
fixed point is \(x\). This suffices for
the one-sided condition away from \(x\).

\proofstep{The tangent criterion for the two-sided points.}
At a two-sided point of the Gaussian chart,
the block exclusion in Section~\ref{sec:host-descent}
gives an irreducible restriction to \(K\).
For general data, the two possibilities for the restriction are
\[
 V\text{ irreducible of dimension }3,\qquad\text{or}\qquad V=W_{\rm outer}\oplus U_{\rm middle}
\]
with an irreducible outer plane and a middle line. In both charts
restrict the marked integral representative
at \(\ell\) to obtain the finite adjoint
local condition. For the Gaussian chart
this was the fixed-inertia derivative
calculation of Section~\ref{sec:host-descent}.

For the general datum it follows on the
crystalline local quotient by the
crystalline extension criterion, after
the same finite base change that
trivializes the local type. Away from
\(\ell\), the inertia condition gives
the finite local condition by restriction
and averaging. Thus the adjoint class
of any tangent to a two-sided point
vanishes by the rational Selmer results
of that section. The constraints from complex conjugation and the marked difference leave only the following line of coboundary matrices:
\[
 X=u\operatorname{diag}(1,0,-1).
\]
It produces exactly the torus tangent.

\proofstep{Continuity and the completed integral chart.}
These point and derivative arguments apply
on the graded algebra. A degree-zero
point is integral on the finite invariant
quotient. Nonzero products of opposite
weights give the strict interval of
allowable torus rescalings, so the full
matrix coordinates can be taken in the
maximal ideal of a finite coefficient
DVR away from their split values. A first
order lift of this point has the same
extension to the complete chart. For
each generator of nonzero weight its value
is the rescaled integral value plus a
nilpotent term.

In a dual-number ring, a product of \(n\) rescaled
values \(a_i+\delta b_i\) has the expansion
\[
 \prod_{i=1}^n(a_i+\delta b_i)
   =\prod_{i=1}^n a_i+
       \delta\sum_{i=1}^n b_i\!\prod_{j\ne i}a_j .
\]
The \(a_i\) are in the maximal ideal, and the finitely
many \(b_i\) have one fixed lower valuation bound. Both
coefficients therefore tend to zero uniformly in such products
as \(n\) increases. At the
finite invariant quotient, a point
derivation on its local coefficient
algebra is continuous by this same
product rule on powers of the maximal
ideal. Hence completing the chart
gives a normalized dual-number
deformation. Its matrix entries still
separate the derivations. This is the
chart version of the torus tangent
condition in \eqref{eq:local-length-3}.

\proofstep{Applying the conormal bound to the selected trace.}
Take \(a\) to be the trace at the
selected Frobenius minus the split
trace. It has weight zero and
vanishes at \(x\). The packet bounds
of Section~\ref{sec:host-descent}
show it is nonzero at each of
the two-sided points. Apply \eqref{eq:local-length-3}
and the integral transfer \eqref{eq:local-length-1} at
every torsion depth. With \(T\)
one representative of each
two-sided orbit, the normalized
rank-one length satisfies
\[
 s_\ell(\psi,e)
 \ \le\ 2\sum_{\xi\in T}
   v_\ell\!\left(\operatorname{tr}r_\xi(\mathrm{Frob}^{-1})
                  -\operatorname{tr}r_0(\mathrm{Frob}^{-1})\right).
 \tag{11}\label{eq:local-length-11}
\]
An intertwiner on the quadratic
kernel is diagonal in the marked
bases. For an irreducible triple,
polarization then makes it a scalar
times a grading-torus element.
For the plane plus line, the
middle diagonal entry can be
changed to the square root of the
product of the two outer entries.
This preserves the intertwiner.
So in either case the injection
from orbits to the packet indices
in \eqref{eq:local-odd-d} applies to precisely the
representatives in \eqref{eq:local-length-11}.

\proofstep{Normalized valuations on the original packets.}
For the rank-three trace estimate sum \eqref{eq:local-length-11} over the
characters with one kernel. At a test prime \(p\ne\ell\) the common congruence integer bounds the sum over characters with this kernel by
\[
 2v_\ell(A_{J_0}(p)).
\]
Two fixed distinct test primes
give the bound at all admissible
\(\ell\), as in the arithmetic sum
following \eqref{eq:local-odd-l}. The cardinality
and valuation factors are at the
original levels \(N\) and \(N|D_K|\).
Once \(K\) and the pair of test
primes are fixed, the exponent
depends on neither \(N,e\) nor the
residual prime or coefficient
ring.

For the case over the Gaussian field, the fixed
Gaussian field and the inert
prime \(p\le 12\ell N\) give
\(s_\ell(\psi,1)\log\ell\ll
(\ell N)^C\) at the large
dyadic conductor, using only the
rank-three packet. The smaller
conductor goes through the
height estimate as in the
odd-prime count. These deductions
use the height and marked
support inputs of the previous
sections. The enlargement of the
marked field enters neither
packet exponent.
\label{cor:ordinary-rank-three-length-8}
\end{proof}


\section{Integral symplectic bases and transport of ordinary flags in rank four}
\label{sec:additional-rank-four}
In the even-order argument we must connect a mixed Hilbert seed to a rational target
in rank four. A comparison of their residual representations alone would lose the
positions of their integral ordinary flags. We choose the frame using a marked element
with distinct residual eigenvalues. The cup calculation and
the transport of flags give the three-link path from the
rational target to the mixed Hilbert seed used in
Paragraph~\ref{par:count-two-propagation}. The target normalization
and cup lift also enter the four-link construction of
Appendix~\ref{sec:count-three-chain}.

For the mixed Hilbert seed of Appendix~\ref{par:count-ii-tower},
write \(r_2\) for its rank-two representation and \(k_0\)
for the real quadratic inducing field. Lemma~\ref{lem:rank-four-seed-diagonalization} fixes a symplectic eigenbasis on the induced seed. Proposition~\ref{prop:rational-rank-four-torus} treats the target lattice separately before their two sets of local flags are compared.

\begin{lemma}[Integral diagonalization of the mixed seed]
\label{lem:rank-four-seed-diagonalization}
In odd residual characteristic, with the marked element and
determinant identity of the mixed seed, one can choose an
integral eigenbasis at \(\tau^2\) in which the residual
rank-two representation is a nonsplit upper triangular
extension. In an integral eigenbasis of the induced lattice at \(\tau\), its symplectic form has matrix 
\[
 J_{\rm ind}=uJ_{\rm std},\qquad u\in O^\times,
\]
 over the coefficient DVR \(O\).
\end{lemma}

\begin{proof}
\proofstep{Use the two rank-two eigenprojectors to choose the extension orientation.}
At the normalizing element
\(\tau^2\) the two residual eigenvalues are distinct.
The eigenprojectors therefore preserve an integral lattice.
The residual trace is the sum of the two characters, so
\(\operatorname{tr}(E_{00}r_2(g)E_{11}r_2(h))\) reduces to zero
for all \(g,h\) on the inducing subgroup. Thus the minima of the
valuations of the two off-diagonal entries sum to a positive
integer. Irreducibility makes both minima finite. Scale the
two basis vectors relatively so the upper-right minimum is
zero. After the relative scaling, the sum of the two minimum valuations is the same positive integer as before. If \(m_{12},m_{21}\) denote the scaled minima,
\[
 m_{12}=0,\qquad m_{12}+m_{21}>0\quad\Longrightarrow\quad m_{21}>0,
\]
so the lower-left entry vanishes residually. The
upper-right entry survives on some element, and diagonalization
at \(\tau^2\) gives a nonsplit residual extension.

\proofstep{Induce the basis and compute the alternating form.}
Denote the
eigenvalues there by \(a^2,b^2\), choosing signs to match the
residual labels and \(ab=\mu(\tau)\) for the similitude
character \(\mu\). This is possible at the odd residual prime
by the determinant identity. On induction across the two cosets
the matrix of \(\tau\) is
\(\left(\begin{smallmatrix}0&\operatorname{diag}(a^2,b^2)\\
I_2&0\end{smallmatrix}\right)\). The columns of
\[
 P=\begin{pmatrix}0&a&0&a\\ b&0&-b&0\\0&-1&0&1\\1&0&1&0\end{pmatrix}
\]
are integral eigenvectors of eigenvalues \(b,-a,-b,a\).
Its determinant is a unit. For the alternating block form
\(\operatorname{diag}(J_2,abJ_2)\), with
\(J_2=\left(\begin{smallmatrix}0&1\\-1&0\end{smallmatrix}\right)\),
the new frame is symplectic up to the unit \(-2ab\).
\end{proof}

We now use this frame to compute the entries and flags of the
seed. We then normalize the rational target and join the two
sets of residual data.

\paragraph{The integral flags of the mixed seed.}
\label{par:count-two-seed-flags}
Keep the eigenbasis just constructed. To compare the seed with the target we first
need the positions of the seed's ordinary flags in these coordinates.

\begin{lemma}[Residual entries and ordinary flags of the mixed seed]
\label{lem:rank-four-mixed-flags}
Use the eigenbasis of Lemma~\ref{lem:rank-four-seed-diagonalization}.
Let \(r\) be the induced representation in that basis. Write \(\operatorname{supp}_{\rm off}(\bar r)\) for the pairs \((i,j)\), \(i\ne j\), for which \(\bar r_{ij}(g)\ne0\) for some \(g\). Then 
\[
 \operatorname{supp}_{\rm off}(\bar r)\subseteq\{(1,0),(3,2),(3,0),(1,2)\}.
\]
 The \((1,0)\) entry is nonzero
somewhere or both \((3,0)\) and \((1,2)\) are nonzero somewhere.
At every prescribed \(\ell\)-adic place, an integral lower
unipotent symplectic frame for the ordinary flag has \((2,0)\)
entry equal to \(\pm1\) and \((2,1)\) entry in the maximal ideal.
The inertia deletion conditions of the mixed seed hold in this
frame. All entries and comparisons may be read over one finite
coefficient extension.
\end{lemma}

\begin{proof}
\proofstep{Compute the residual entries before and after the other coset.}
For \(g\in G_{Fk_0}\), write \(X(g),Y(g)\) for the
residual upper-right entries of the two rank-two
constituents on the original two cosets. The two
diagonal characters of both constituents are
\(1,\bar\eta\), since \(\eta\) comes from \(\Q\).
Put \(\gamma=b/a\) in the basis change \(P\) above.
Multiplying by \(P^{-1}\) on the left and \(P\) on the
right gives
\[
\begin{array}{c|rrrr}
\text{entry} &(1,0)&(3,2)&(3,0)&(1,2)\\ \hline
\text{value}&
(\bar\gamma X-Y)/2 &-(\bar\gamma X-Y)/2&
(\bar\gamma X+Y)/2 &-(\bar\gamma X+Y)/2 .
\end{array}
\]
All other off-diagonal entries are zero. The diagonal
entries have the marked residual order. On the other
coset in \(G_F\), multiplication by the diagonal matrix
of \(\tau\) preserves the same support and gives the
remaining marked characters. Since \(X\) is nonzero
somewhere, either the \((1,0)\) entry is nonzero at some
element, or both \((3,0)\) and \((1,2)\) are nonzero.
On reducing, the scalar twist therefore has no effect on the tested entries: for its character \(\chi\),
\[
 \overline{\chi(g)r(g)}=\bar\chi(g)\bar r(g)=\bar r(g).
\]

\proofstep{Normalize the eigenlines in each ordering of the local constituents.}
At a place above \(\ell\) the decomposition group is
inside \(G_{Fk_0}\). Its two constituents evaluate
\(r_2\) at the distinct places of \(k_0\) above \(\ell\)
because \(\tau\) exchanges those places. The residual
ordinary lines have opposite characters. Choose an element of the local decomposition group at which
\(\bar\eta\ne1\), and write its matrix as
\(\left(\begin{smallmatrix}a_0&b_0\\c_0&d_0\end{smallmatrix}\right)\).
The lattice of the triangular residual extension has
\(a_0\equiv1\), \(d_0\equiv\eta\), and \(c_0\equiv0\).

For an eigenvector \((p,q)^{\mathsf t}\) of eigenvalue \(\lambda\),
\[
       (\lambda-a_0)p=b_0 q,\qquad
       (\lambda-d_0)q=c_0 p.
\]
If \(\bar\lambda=\bar\eta\), the first multiplier is a unit,
so \(|p|\le|q|\) and the line has a generator \((u,1)\).
If \(\bar\lambda=1\), the second multiplier is a unit and
\(|c_0|<1\), so the generator has the form \((1,w)\) with
\(|w|<1\).

In the first ordering, these lines are \((u,1)\) in
the first constituent and \((1,w)\) in the second.
The largest exponent is \(k+t\) on the first line.
Applying \(P^{-1}\) and normalizing gives first flag
column \((1,\gamma u,-1,\gamma u)^{\mathsf t}\).
The second line is proportional to
\((w,-1,w,1)^{\mathsf t}\).
Write \(v_1=(1,\gamma u,-1,\gamma u)^{\mathsf t}\) for the normalized first column. Subtract its first-coordinate multiple from the second generator and normalize the second coordinate:
\[
 v_2=\frac{(w,-1,w,1)^{\mathsf t}-wv_1}{-1-\gamma uw}.
\]

In the other ordering, the largest exponent is on
the second constituent. The corresponding first flag
column is \((1,-u,1,u)^{\mathsf t}\), and the other
line is proportional to
\((w,\gamma,-w,\gamma)^{\mathsf t}\).
Subtract \(w\) times the first column and divide by
\(\gamma+uw\) for the second flag column.
The ultrametric inequality computes both divisors in the two flag normalizations. Since \(|\gamma|=1\) and \(|uw|<1\),
\[
 |-1-\gamma uw|=1,\qquad |\gamma+uw|=1 .
\]

\proofstep{Complete the pair of columns to a symplectic flag.}
In either case the two columns span a Lagrangian
plane, since the original lines lie in separate
symplectic summands. Write the columns as
\((1,p,q,r)^{\mathsf t},(0,1,s,t)^{\mathsf t}\).
The Lagrangian equation is \(t+ps=q\). Complete them
with \((0,0,1,-p)^{\mathsf t}\) and
\((0,0,0,1)^{\mathsf t}\); the resulting lower
unipotent matrix is symplectic by the same equation.
The first three columns span the orthogonal of the
first line. Thus the frame is adapted to the full
stable flag, with characters of exponents
\(k+t,k,0,-t\) before the scalar twist.

In the two orderings the required entries \(q\) and \(s\) of the normalized matrix satisfy
\[
 q\in\{1,-1\},\qquad |s|<1
\]
in positions \((2,0)\) and \((2,1)\), respectively.
The two constituents' ordinary characters reduce,
in this order, to \(\bar\eta,1,\bar\eta,1\).
They coincide with the marked characters because
\(\overline\Psi\) is trivial on the decomposition
group and the diagonal order was fixed at \(\tau\).

\proofstep{Check the scalar twist on the deleted inertia.}
With \(\vartheta=\epsilon_\ell^{e+1}\omega^{-1}\eta\psi_0\),
the square of the scalar twist is
\[
 \delta^2=\epsilon_\ell^{e+1-k}\omega^{-1}\psi_0\phi_0^{-1}
 \quad\text{on }G_F.
\]
On the deep tower \(\psi_0\) has prime-to-\(\ell\) order.
Above \(\ell\), \(\phi_0\) is one on the decomposition
group and \(\omega\) is the Teichm\"uller part of
\(\epsilon_\ell\). On inertia in the cyclotomic kernel,
\(\psi_0\) reduces to one by the prescribed residual
inertia. Its prime-to-\(\ell\) order makes it one.

The two possible square roots are separated by reduction. Thus
\[
 (\delta-1)(\delta+1)=0,\qquad\delta\equiv1\quad\Longrightarrow\quad \delta=1
\]
at these places. On open inertia where
\(\psi_0=\omega=1\), the exponent of \(\delta^2\) is
\(e+1-k\), which gives the doubled weights used
above. At a prescribed inertia element away from
\(\ell\) with trivial residual labels, the tame
tower condition kills \(\psi_0\), and the prime-to-
\(\ell\) character \(\eta\) is one by reduction. The
quadratic character \(\phi_0\) is unramified on the
tower and \(\epsilon_\ell,\omega\) are unramified away
from \(\ell\). The two constituents are trivial on these inertia elements because the type is \(1\oplus\eta\) and \(\eta=1\) on them. Thus \(\delta^2=1\) there and the same reduction test gives
\[
 \delta^2=1,\quad\bar\delta=1\quad\Longrightarrow\quad\delta=1 .
\]
These are the required inertia conditions for the
integral seed and its reduction.

\proofstep{Choose one finite coefficient extension.}
All these matrices descend to the integers of one
finite extension of \(\Q_\ell\). For the representation,
write \(\overline{\Q}_\ell\) as a countable union of
finite fields over \(\Q_\ell\). Their intersections
with the matrix entries cut out closed subgroups
covering the compact Galois group. By the Baire category theorem at least one of these closed subgroups has nonempty interior. Calling it \(H\), compactness then gives
\[
 H\text{ open in }G,\qquad [G:H]<\infty .
\] Adjoining the entries on finitely
many coset representatives, and the finitely many
flag entries, then gives the desired field.
The displayed unit and strict norm bounds persist
under this coefficient choice. Thus the reduced seed
has the required support and nonzero marked entries
over a common residual coefficient extension.
\label{cor:cohomology-count-two-seed-1}
\end{proof}

\paragraph{The initial symplectic labels and flags.}
\label{par:count-two-labels}
We now change endpoints. At the rational target we again start in the marked
eigenbasis, so that later changes of lattice can be read on the same root spaces.

\begin{lemma}[The rational flag in the marked eigenbasis]
\label{lem:rational-rank-four-marking}
For the rational datum, there is an integral symplectic
eigenbasis diagonalizing the normalizing element \(\tau\).
In this basis the ordinary flag admits a lower unipotent
symplectic frame \(N\). Denote its four independent lower entries by
\(N_{10},N_{20},N_{30},N_{21}\). Let \(S_K\) consist of the degrees
of nonzero off-diagonal entries of the global matrices \(r(g)\),
\(g\in G\), over the fraction field \(K\). Write \(S^{(0)}\)
for the degrees of the global and independent frame entries
that are units on this initial lattice. Then
\[
 -w_\Psi\in S_K,\qquad
 S_K\cap\{w_\beta,-w_\alpha,-w_\vartheta\}\ne\varnothing,
 \qquad S^{(0)}\subseteq\{-w_\Psi\}.
\]
The initial residual global matrices are diagonal. Only a frame
entry can contribute to \(S^{(0)}\).
\end{lemma}

\begin{proof}
\proofstep{Place the ordinary flag in a chart with unit leading minors.}
For the rational point use the residual order
\((\bar\vartheta,\bar\Psi,\bar\alpha,1)\).
The designated Frobenius at \(\ell\) separates the
characters of the two spaces on the indices
\(\{0,2\}\) and \(\{1,3\}\). The residue of the
ordinary line has its character from the first
space and lies in that space. The two-dimensional
step of the Lagrangian flag splits across the two
spaces, by the Frobenius projectors, with a line in
each. The two spaces are paired by the alternating
form. If the first line has nonzero coordinate zero,
the Lagrangian plane projects isomorphically to the
first two coordinates. Indeed if its other line had
coordinate one zero it would be spanned by the
vector of index three, which pairs nontrivially
with the first line.

If instead coordinate zero vanishes, choose the signed symplectic permutation inducing the exchanges
\[
 (0\ 2)(1\ 3).
\]
This is allowed because the original first line
then has character both \(\bar\vartheta\) and
\(\bar\alpha\) on the marked decomposition group.
Their equality says \(\bar\Psi=1\) on the whole group,
and the other two characters coincide there as well.
The permutation acts on the global basis, with the
order of the given local inertia characters retained.
In either order the flag has unit first two leading
minors. The third is a unit by the relation between
the inverse matrix and its symplectic adjoint.

\proofstep{Diagonalize the marked global matrix without changing the residual flag.}
Let \(\lambda_i\) be the four Hensel lifts of the distinct
eigenvalues of the matrix \(R(\tau)\) in the complete
coefficient ring. The projectors
\[
 e_i=\prod_{j\ne i}\frac{R(\tau)-\lambda_j I}
                              {\lambda_i-\lambda_j}
\]
are integral because every denominator is a unit.
Applying \(e_i\) to the corresponding standard basis vector
gives an eigenvector reducing to that vector. For the
alternating form \(\langle\, , \,\rangle\) and the similitude
\(s(\tau)\), the eigenvector identity is
\[
  (\lambda_i\lambda_j-s(\tau))\langle v_i,v_j\rangle=0 .
\]
Only the paired indices \(0,3\) and \(1,2\) can have a nonzero
pairing, by their distinct residual eigenvalues. Each of
those pairings is a unit, congruent to the old pairing.
Rescale one vector of each pair to recover the symplectic
form with a change of basis congruent to the identity.

In this basis the unit leading minors give Gaussian
factorization of the flag frame as \(N U\), with \(N\) lower
unipotent and \(U\) upper triangular. Write \(J\) for the antidiagonal matrix of the form and \(\mu\) for the similitude of the frame \(A=NU\). The group automorphism \(\iota(X)=J^{-1}X^{-t}J\) preserves the lower unipotent and the upper triangular subgroups.

Since \(A^t J A=\mu J\), it gives a second factorization:
\[
 NU=A=\mu\iota(A)=\iota(N)\bigl(\mu\iota(U)\bigr).
\]
Consequently \(N^{-1}\iota(N)=U(\mu\iota(U))^{-1}\) is both lower unipotent and upper triangular. It is the identity, so \(\iota(N)=N\), equivalently \(N^tJN=J\). The upper factor preserves the ordinary flag
and its characters. Modulo the uniformizer the first two
columns follow their separate Frobenius spaces, and the
perpendicular conditions determine the two later steps.
Thus the only possible residual entries of \(N-I\) have
degree \(-w_\Psi\).

\proofstep{Locate the needed roots over the fraction field.}
If the rational representation is irreducible, the
algebra spanned by its matrices has all matrix-unit
positions in this eigenbasis. Let \(P_i\) be the diagonal projectors in this eigenbasis.
An arrow \(j\to i\) means that some action matrix \(r(g)\)
has \(r(g)_{ij}\ne0\). Then
\[
 P_i r(g)P_j=r(g)_{ij}E_{ij}.
\]
A proper terminal component of the directed graph would span
a proper stable coordinate subspace, so irreducibility gives
a path between each ordered pair of vertices. Repeated multiplication along a path removes each intermediate index:
\[
 E_{i_k i_{k-1}}\cdots E_{i_1 i_0}=E_{i_k i_0}.
\]
The projectors are polynomials in \(r(\tau)\). Hence every
product just obtained is in the \(K\)-linear span of the \(r(g)\):
products of group matrices are again group matrices. A nonzero
\((i,j)\)-coordinate in this span forces \(r(g)_{ij}\ne0\)
for some \(g\), with no assertion about its valuation.

In the block-extension case, the original matrices
preserve the space on indices \(1,3\), and the
diagonalizing eigenvectors in that block lie in the
same space. The change of basis is block triangular
and this block still spans the full matrix algebra.
On this block the matrix-unit calculation gives
\[
    (1,3),(3,1)\text{ in the matrix support over }K
       \quad\Longrightarrow\quad \{w_\Psi,-w_\Psi\}\subset S_K .
\]
 The nonsplit extension has a
nonzero cross-block entry. Otherwise the other
coordinate block gives a stable complement and
conjugating it back gives a stable complement
which is the graph of a map to the first block.
Such a graph is exactly the split extension
excluded in the rational datum. With rows \(1,3\) and columns \(0,2\), the four degrees are
\[
  \deg(1,2)=w_\beta,\quad \deg(3,0)=-w_\vartheta,\quad
  \deg(1,0)=\deg(3,2)=-w_\alpha .
\]

These entries give the two assertions about \(S_K\).
The initial global reduction is the given diagonal model and the
change by Hensel projectors is the identity on its reduction,
up to the preceding simultaneous permutation. The flag calculation
gives the asserted inclusion for \(S^{(0)}\).
\label{cor:cohomology-count-two-normalization-1}
\end{proof}

\paragraph{The torus normalization of the target.}
\label{par:count-two-normalization}
The rational eigenbasis fixes the labels but leaves a diagonal freedom. We use it to
reduce the target lattice to the shapes needed for interpolation.

\begin{proposition}[The two rational lattice shapes]
\label{prop:rational-rank-four-torus}
Assume the odd-root localization vanishing for the marked rational
datum and the basis of Lemma~\ref{lem:rational-rank-four-marking}.
There is a diagonal symplectic similitude, after adjoining
a square root of the uniformizer, giving an integral model
whose reduced global matrices \(r\) and reduced lower frame \(N\)
satisfy the following bounds. Here \(S\) counts the degrees of
nonzero reduced global entries off the diagonal and of the nonzero entries among the four
independent frame coordinates after the change of lattice:

\[
 S\subseteq\{w_\beta,-w_\Psi,-w_\alpha,-w_\vartheta\},
       \qquad r_{12}\not\equiv0,\qquad N_{21}=0.
\] Either degree \(-w_\Psi\) occurs in the residual
global matrices or it occurs only through \(N_{20}\ne0\),
in which case \(r_{12}\) cannot have the form
\(c(d_2-d_1)\) on the marked decomposition group.
\end{proposition}

\begin{proof}
\proofstep{Maximize the two torus coordinates subject to integrality.}
Use the normalized rational datum above in its small
coefficient model over a complete
discrete valuation ring with uniformizer \(\varpi\).
Use additive valuation normalized by
\(v(\varpi)=1\). Retain the four independent lower frame entries
\((1,0),(2,0),(3,0),(2,1)\); the symplectic equations
express the other lower entries in these. For each degree \(w=(w_1,w_2)\) occurring in a nonzero global or retained frame entry, define the minimum
\[
 k_w=\min\{v(a):a\ne0\text{ a global or retained frame entry of degree }w\}.
\] The given integral model makes
this a minimum of nonnegative integers. All global
off-diagonal entries, and the frame entries other than the
possible degree \(-w_\Psi\), have positive valuation.

Conjugation by the proposed diagonal with weights
\((a+b,b,a,0)\) changes a degree-\(w\) coefficient of
valuation \(k\) into one of valuation \(k+(w_1a+w_2b)/2\).
Thus integrality of every entry to be checked is equivalent
to requiring
\[
             w_1a+w_2b+2k_w\ge0
                    \qquad\text{for each degree }w\text{ with }k_w<\infty .
\]
The initial lattice makes \((0,0)\) feasible. Put \(k_w=+\infty\)
if there is no entry of degree \(w\). The two assertions about
\(S_K\) in Lemma~\ref{lem:rational-rank-four-marking} give
\[
 b\le 2k_{-w_\Psi},\qquad
 a\le e b+2k_{(-1,e)}\quad\text{for some }e\in\{-1,0,1\}.
\]
Each right side for its indicated degree is finite. Thus the feasible
integer pairs attain a maximum in \(b\), and on that slice in \(a\).

Take \(p^2=\varpi\) and use
the diagonal symplectic similitude with entries
\((p^{a+b},p^b,p^a,1)\). Subscripts of the
four ordered positions correspond to
\(((1,1),(0,1),(1,0),(0,0))\).

For the diagonal matrix \(D\) just defined, let \(w\) be the difference
of the two position degrees at an entry \(x_{ij}\). With \(v_\varpi(\varpi)=1\),
\[
 (DxD^{-1})_{ij}=p^{w_1a+w_2b}x_{ij},\qquad
 v_\varpi((DxD^{-1})_{ij})=v_\varpi(x_{ij})+(w_1a+w_2b)/2.
\]
The defining inequalities are exactly nonnegativity of these valuations
for the matrix and frame coordinates. For an original unit \(u\) and
a nonzero coordinate \(x=u\varpi^h\), this gives the integral element
\[
           p^{w_1a+w_2b}x=u p^{2h+w_1a+w_2b}
                     \in\mathcal O_{K(p)}.
\] Conjugation
preserves the diagonal, the ordinary flag equations
and the normalizing matrix at \(\tau\). This also
preserves the small-model condition and continuity.
Multiplication by finitely many fixed powers of \(p\)
only shifts the adic indices by fixed amounts.

\proofstep{Find the roots that prevent an increase in the maximizer.}
Reduction after conjugation gives the exact identity
\[
 S=\{w:k_w<\infty,\ w_1a+w_2b+2k_w=0\}.
\]
Indeed the minimum is attained in the discrete valuation ring,
and precisely entries of new valuation zero survive. Call these
degrees tight. In all subsequent root and flag equations, \(r\)
and \(N\) denote the reductions on this new lattice. For any tight root other
than \(-w_\Psi\), the value of \(w_1a+w_2b\)
is strictly negative. For \(-w_\Psi\) tight it gives
\(b\ge0\). The failed moves \((1,0)\) and \((0,1)\)
at the maximum require a tight root in each of
\[
 \{-w_\alpha,w_\beta,-w_\vartheta\},\qquad
 \{-w_\Psi,-w_\beta,-w_\vartheta\}.
\]
For either move the slack \(w_1a+w_2b+2k_w\) decreases
by at most one, so a failed move must meet a tight constraint.

At a tight root the integer \(w_1a+w_2b\) is even. The two
diagonal weights have linear forms \(a-b\) and \(a+b\)
up to sign. Tightness of either makes \(a+b\) even.
If neither blocking root is diagonal, the first gives
\(a\equiv0\pmod2\) and the second gives \(b\equiv0\pmod2\).
In either case the slack of a diagonal constraint is even. The failed moves \((-1,1)\) and \((1,1)\)
therefore require a tight root in each of
\[
 \{-w_\Psi,w_\alpha,-w_\beta\},\qquad
 \{-w_\Psi,-w_\alpha,-w_\vartheta\}.
\]
The only inequalities that drop by two in these moves
are diagonal roots, for which the preceding slack is
even. Hence slack one cannot block one of these moves.

For brevity put \(A=w_\alpha\), \(B=w_\beta\),
\(P=w_\Psi\), \(T=w_\vartheta\) in the following
table. The tight set is contained in one of the
listed sets and contains the two marked roots in its
row.
\[
\begin{array}{c|l}
\text{two tight roots}&\text{containing set}\\ \hline
B,-P & \{B,-P,-A,-T\}\\
-A,-P & \{-A,-P,-T\}\\
-T,-P & \{-T,-P\}\\
-A,-B & \{-A,-B,-P,-T\}\\
-T,-B & \{-T,-B,-P\}\\
-T,A & \{-T,A,-B,-P\}
\end{array}
\]
To check the table, suppose first that \(B\) is
tight, so \(b<a\) and \(-B\) is excluded.
The third failed move forces \(-P\) or \(A\).
If \(-P\) is absent, \(A\) gives \(b<a<0\), excluding
every root for the second move. Thus \(-P\) is
tight and \(a>b\ge0\), giving the first row.
If \(B\) is absent and \(-B\) is present, the first
move forces \(-A\) or \(-T\).

With \(-A\) tight
one has \(b>a>0\), giving row four. Otherwise the inequalities supplied by \(-T\) and \(-B\) give
\[
 0<a+b<2b\quad\Longrightarrow\quad b>0,
\]
where the strict upper bound uses \(a<b\).
According as \(A\) is absent or present one obtains
row five or six. If neither \(B\) nor \(-B\) is tight,
the second and third moves require either \(-P\) or
both \(-T,A\). These last two exclude \(-A,T,P\)
by their signs, giving row six. With \(-P\) tight,
\(b\ge0\) and the first move requires \(-A\) or \(-T\).
If \(-A\) is present, the signs exclude \(A,T,P\),
giving row two. In the case where \(-P\) is tight and \(-A\) is absent, \(-T\) must be present. The remaining test reads
\[
 A\text{ present}\ \Longrightarrow\ \text{row six},\qquad A\text{ absent}\ \Longrightarrow\ \text{row three}.
\]

\proofstep{Use localization to exclude unsupported root cocycles.}
Use the odd-root localization vanishing for this marked
rational datum from the rank-four cohomology calculation.
If \(r_{ij}\) has degree in \(\{\pm A,\pm B,\pm T\}\) and
only the diagonal terms in its multiplication law, then
\(c_{ij}=r_{ij}/d_j\) satisfies
\[
    c_{ij}(gh)=c_{ij}(g)+
                  (d_i/d_j)(g)c_{ij}(h).
\]
If its local restriction is a coboundary, the localization
statement makes it a global coboundary. In the marked
eigenbasis, \(c_{ij}(\tau)=0\), while \(d_i(\tau)/d_j(\tau)\ne1\).
This forces the global coboundary constant to be zero.
The same reasoning on the designated local Frobenius
forces a local constant to vanish whenever its root
character is nontrivial there.

If \(-A\) is absent, \(r_{10}=r_{32}=0\). Hence the two
intermediate terms of the \((3,0)\) entry vanish:
\[
 r_{31}(g)r_{10}(h)+r_{32}(g)r_{20}(h)=0.
\]
In the first-column flag equation the intermediate term involving \(N_{10}\) vanishes. The remaining identity is
\[
 r_{30}=(d_0-d_3)N_{30}\qquad(N_{10}=0).
\]
Thus the odd-root argument kills \(r_{30}\) and \(N_{30}\)
and excludes \(-T\) from the tight set.

If both \(B,P\)
are absent, \(r_{12}=r_{13}=0\); similarly the \((1,0)\)
entry has no intermediate terms, and its local equation
is \(r_{10}=(d_0-d_1)N_{10}\). Hence \(r_{10}=N_{10}=0\).
In this two-term situation \((r^{-1})_{10}=0\) as well,
since \(r(g^{-1})=r(g)^{-1}\). The symplectic adjoint formula
then gives \(r_{32}=0\). This excludes \(-A\) when \(B,P\)
are absent. The two exclusions leave only the first row
of the tight-root table.

\proofstep{Distinguish a global entry from an ordinary frame entry.}
The remaining off-diagonal degrees of the reduced matrices form a subset of
\[
 \{B,-P,-A,-T\}.
\] The frame has \(N_{21}=0\), since
its degree \(-B\) is absent. Tightness in degree \(B\)
comes from the global \((1,2)\) entry, since the frame
is lower triangular. The first lattice shape occurs
when degree \(-P\) appears in the global matrices.

In the second shape degree \(-P\) is tight only
through \(N_{20}\ne0\). Row two of the first-column
equation forces \(d_2=d_0\) on the marked decomposition
group. The global \((1,2)\) entry has the two-term
multiplication law. If it satisfied
\[
                      r_{12}=c(d_2-d_1)
\]
on this group for a scalar \(c\), normalized odd-root
vanishing would make it zero everywhere. Its nonzero
entry gives the required second lattice shape. \label{cor:cohomology-count-two-normalization-2}
\end{proof}

\paragraph{A locally split cup lift at auxiliary primes.}
\label{par:count-two-cup}
To join the two endpoints we next construct the required extension classes. The
residual path begins over the tame real field \(F_0\).
Write \(\Gamma=G_{F_0}\), let \(k\) be the common residual
coefficient field, and put
\[
 A=d_1/d_0,\qquad B=d_2/d_0 ,
\]
where \(d_i\) are the four ordered residual characters. Their
symplectic products satisfy \(d_0d_3=d_1d_2\). The normalizing
element \(\tau\) has distinct diagonal entries.

In particular
\(A(\tau)\ne1\) and \((AB)(\tau)\ne1\). The normalized target
has support in the diagonal and the entries
\[
 (1,0),\ (2,0),\ (3,0),\ (3,1),\ (3,2),\ (1,2).
\]
Its local lower unipotent symplectic frame \(N_v\) has
\((N_v)_{21}=0\). Put \(y(g)=r(g)_{20}/d_0(g)\).
Multiplication in that support gives
\[
 y(gh)=y(g)+B(g)y(h),\qquad y(\tau)=0 .
\]
For \(s\) in the marked decomposition group \(G_v\), the first
column of \(r(s)N_v=N_v(N_v^{-1}r(s)N_v)\) gives
\[
                    y(s)=(1-B(s))n_v,\qquad n_v=(N_v)_{20}.
\]
This uses the prescribed first diagonal character in the local
upper triangular matrix. It applies even when \(y\) is zero and
the local entry \(n_v\) is nonzero.

\begin{lemma}[The two rational shapes after restriction]\label{lem:count-two-target-restriction}
In the first lattice shape of Proposition~\ref{prop:rational-rank-four-torus},
the two entry functions on the restricted group satisfy
\[
 r_{20}|_{G_{F_0}}\not\equiv0,\qquad r_{12}|_{G_{F_0}}\not\equiv0 .
\]
In the second shape the global \((2,0)\) entry is zero, the normalized
marked frame has nonzero \((2,0)\) coordinate, and a \((1,0)\) entry is
nonzero on \(G_{F_0}\).
\end{lemma}

\begin{proof}
The fitted field \(F_0\) is abelian over \(\Q\). For an entry
whose normalized coordinate is a cocycle \(f\) of
character \(\theta\), suppose \(f\) vanishes on \(G_{F_0}\).
For any \(g\in G_\Q\), normality puts
\(g\tau g^{-1}\) in \(G_{F_0}\), and the cocycle
identity gives
\[
              \begin{aligned}
 0=f(g\tau g^{-1}) &= f(g)+\theta(g)f(\tau)
                       +\theta(g)\theta(\tau)f(g^{-1})\\
                   &=(1-\theta(\tau))f(g),
 \end{aligned}
\]
Since \(1-\theta(\tau)\ne0\), the restriction of a
nonzero normalized cocycle to \(G_{F_0}\) is nonzero.
The \((2,0)\) coordinate has character \(d_2/d_0\). For the
\((1,2)\) coordinate, no intermediate product contributes to
the \((1,2)\) entry, so dividing by \(d_2\) gives a cocycle
of character \(d_1/d_2\). The two coordinates therefore survive
in the first shape. The other allowed pattern is distinguished by the global representation entry and the local frame entry
\[
 r(g)_{20}=0\ (\text{all }g),\qquad n=N_{20}\ne0 .
\]

In this support the normalized \((1,0)\) entry is a
cocycle, since its only additional product would use
the zero \((2,0)\) entry. If it vanishes on the tower
it vanishes over \(\Q\) by the displayed argument.
On the rational marked decomposition group the
first-column equations in rows two and one then
give \(d_2=d_0\) and
\[
             r(s)_{12}=\frac{N_{10}}{n}(d_2(s)-d_1(s)).
\]
This contradicts the non-coboundary in the second
lattice shape. Thus the restricted point has a
nonzero \(-w_\alpha\) entry. The frame transport calculation in
Paragraph~\ref{par:count-two-restriction} gives
\(n\mapsto (d_2(g)/d_0(g))n\) when the global \((2,0)\) entry
vanishes. Thus the other coordinate required in the second
shape survives as well.
\end{proof}

\begin{lemma}[A prescribed tame ramification direction]\label{lem:cup-ramification-direction}
Let \(\ell\) be an odd prime, \(F_0\) a number field, and \(M\) a
finite continuous \(\mathbb F_\ell\)-representation of \(G_{F_0}\).
Fix a nonzero \(e_0\in M\), a finite set \(S\) containing the places
above \(\ell\), the infinite places and the ramification of \(M\),
and a finite Galois extension of \(\Q\) containing the normal closure
of \(F_0\), a splitting field of \(M\), and \(\mu_\ell\). There is a
continuous cocycle \(c\colon G_{F_0}\to M\) unramified outside \(S\)
and the places above finitely many rational primes splitting completely
in that extension. Its local class is zero at every \(v\in S\).
On the new inertia groups its image satisfies 
\[
 c(I_v)\subseteq\mathbb F_\ell e_0\quad(v\text{ new}),\qquad c(I_v)\ne0\quad\text{for at least one new }v.
\]

\end{lemma}

\begin{proof}
\proofstep{Specify the new local ramification.}
Take a set \(P\) of places over sufficiently many distinct
new rational primes, including all the places over each of those
primes. Chebotarev permits arbitrarily large \(|P|\).
Let \(M'=\operatorname{Hom}(M,\mu_\ell)\) with its natural Galois
action. At \(v\in P\) both modules are trivial and the residue
cardinality is one modulo \(\ell\). The quotient of
\(H^1(F_{0,v},M)\) by its unramified subspace identifies with the
homomorphisms from the cyclic exponent-\(\ell\) quotient of tame
inertia to \(M\). This follows from the tame presentation
\(\operatorname{Frob}\,t\,\operatorname{Frob}^{-1}=t^{q_v}\).
Requiring the image to lie in \(\mathbb F_\ell e_0\) gives a
one-dimensional subspace of this quotient.
\proofstep{Bound the dual obstruction.}
In the sum of local groups over \(S\cup P\) take the subgroup whose component at each place is
\[
 U_v=\begin{cases}H^1_{\rm unr}(G_v,M)&v\in P,\\0&v\in S.\end{cases}
\]
By local Tate duality its annihilator imposes unramifiedness at
each place of \(P\). By Poitou--Tate exactness the annihilator of
the image of \(H^1(G_{F_0,S\cup P},M)\) is the image of
\(H^1(G_{F_0,S\cup P},M')\). Intersect this latter image with
\(U^\perp\). Any representative global class is unramified at
\(P\), since the local inertia actions are trivial there and
coboundaries vanish on those inertia groups.

It inflates from
\(H^1(G_{F_0,S},M')\), by factoring the cocycle through the quotient
that kills the normal closure of the new inertia groups.
This fixed cohomology group is finite. Let \(L\) be the
localization image modulo \(U\), in the quotient of the sum
\(V_{\rm loc}\) of local groups over \(S\cup P\). The annihilator
description gives
\[
 \operatorname{codim}_{V_{\rm loc}/U}L
 =\dim L^\perp\le \dim_{\mathbb F_\ell}H^1(G_{F_0,S},M').
\]
\proofstep{Obtain a class in the prescribed directions.}
Let \(D_P\subseteq V_{\rm loc}/U\) be the direct sum of the
prescribed ramification lines and zero entries at \(S\).
Then \(\dim D_P=|P|\). Choose
\(|P|>\dim_{\mathbb F_\ell}H^1(G_{F_0,S},M')\).
The dimension formula gives
\[
 \dim(D_P\cap L)\ge |P|-\operatorname{codim}_{V_{\rm loc}/U}L>0.
\]
Lift a nonzero element of this intersection to a global class.
At the new places the inertia restriction is independent of its
representative and is nonzero on some inertia element, which acts
trivially on \(M\). At an old place, a representative is a
coboundary on the full decomposition group, so is zero on any
subgroup on which \(M\) acts trivially.
\hfill\mbox{\cite{milne-duality}}\label{cor:cohomology-count-two-chain-1}
\end{proof}

\begin{proposition}[The marked auxiliary symplectic datum]\label{prop:count-two-auxiliary-datum}
Take the residual characters and the marked \((2,0)\) target
column obtained by restriction in
Lemma~\ref{lem:count-two-target-restriction}. For these data,
one can choose an auxiliary representation \(r'\) with the same diagonal.
Writing \(r,N_v\) for the target matrices and frames and \(N'_v\) for the auxiliary frames, one has
\[
 r'_{20}=r_{20},\qquad (N'_v)_{20}=(N_v)_{20}\quad(v\text{ marked}),
                         \qquad r'_{10}\not\equiv0 .
\] It is diagonal at
\(\tau\), kills the required old inertia and is locally conjugate
to the marked diagonal at the prescribed local places. The new
ramification is supported above finitely many completely split
rational primes as in Lemma~\ref{lem:cup-ramification-direction}.
\end{proposition}

\begin{proof}
\proofstep{Lift the new root together with its cup correction.}
Apply Lemma~\ref{lem:cup-ramification-direction} with \(M\)
carrying coordinates \((x,z)\)
and action
\[
 \rho(g)(u,v)=\bigl(A(g)u,\ (AB)(g)v+y(g)A(g)u\bigr).
\]
The second component of \(\rho(g)\rho(h)(u,v)\) is
\[
 (AB)(gh)v+
 A(gh)\bigl(B(g)y(h)+y(g)\bigr)u .
\]
This is the second component of \(\rho(gh)(u,v)\) by the cocycle
identity for \(y\). Let \(k_{\rm d}\) be the finite subfield generated
by the diagonal characters' images. The span \(V_y\) over
\(k_{\rm d}\) of the finite image of \(y\) is finite, since \(y\) is
continuous on a compact group with discrete target. The chosen additive module is finite even when \(k\) is infinite, since it is the sum of the two finite pieces in the action:
\[
 M=k_{\rm d}\oplus V_y,\qquad |M|=|k_{\rm d}|\,|V_y| .
\]

Take \(e_0=(1,0)\).
All originally prescribed ramification is retained by taking
\(S\) large enough to contain it and using the same construction
with the additional condition that the rational primes split in
the field cut out by the residual diagonal, the quadratic CM
character and the mod-\(\ell\) cyclotomic character. On any old
inertia on which the former representation is required to be
the identity both the diagonal and \(y\) are trivial. The local
coboundary constructed there is zero on that subgroup.

The cocycle coordinates obey
\[
\begin{aligned}
 x(gh)&=x(g)+A(g)x(h),\
 z(gh)&=z(g)+(AB)(g)z(h)+y(g)A(g)x(h).
\end{aligned}
\]
On \(G_v\) write the local coboundary as
\[
 x(s)=(1-A(s))m_v,\qquad
 z(s)=(1-(AB)(s))o_v-y(s)A(s)m_v .
\]
Subtract a global coboundary to make \(x(\tau)=z(\tau)=0\).
At \(\tau\) the action is diagonal because \(y(\tau)=0\), and
neither of its two eigenvalues is one. Thus the subtraction is
possible and just changes the local constants. At the inertia element where \(x\ne0\), the action is trivial, so the subtraction leaves \(x\) unchanged. The same argument works with finitely
many fixed \(y_j\), by replacing \(z\) by a tuple with the indicated
action in each coordinate.

\proofstep{Produce the symplectic matrices and local flags.}
Set
\[
\begin{aligned}
 U(x,y,t)&=
 \begin{pmatrix}
 1&0&0&0\\ x&1&0&0\\ y&0&1&0\\ t&y&-x&1
 \end{pmatrix},\\
 U(x,y,t)U(x',y',t')&=U(x+x',y+y',t+t'+yx'-xy') .
\end{aligned}
\]
For the alternating form pairing positions \(0,3\) and \(1,2\)
with the same sign, \(U\) is symplectic and conjugation by
\(\mathcal D(g)=\operatorname{diag}(d_i(g))\) scales its three
coordinates by \(A(g),B(g),(AB)(g)\).
Let
\[
 R_1(g)=U(x(g),y(g),2z(g)-x(g)y(g))\mathcal D(g).
\]
Writing \(c=2z-xy\), the two cocycle equations give
\[
 c(gh)-c(g)-(AB)(g)c(h)
   =y(g)A(g)x(h)-x(g)B(g)y(h).
\]
Together with the identities for \(x,y\) this is exactly the
multiplication law for the unipotent factor of
\(R_1(g)R_1(h)\). Thus \(R_1\) is a representation with
the prescribed similitude and diagonal.

It is diagonal at \(\tau\), kills the required inertia, and has a
nonzero \((1,0)\) entry. At \(v\) take
\[
                         N_{1,v}=U(m_v,n_v,2o_v-m_vn_v).
\]
On \(G_v\), conjugating \(\mathcal D(s)\) by this matrix gives
\(R_1(s)\). In the third unipotent coordinate both sides give
\((1-(AB)(s))(2o_v-m_vn_v)+(B(s)-A(s))m_vn_v\).
This proves the flag conditions with both the full local
diagonal and the smaller inertia quotient. It also retains
the given \((2,0)\) matrix and frame entries.
\label{cor:cohomology-count-two-chain-2}
\end{proof}

\paragraph{Removing the auxiliary ramification and transporting the flags.}
\label{par:count-two-restriction}
The cup lift introduced tame ramification at the new primes. Before comparing the flags we restrict to a compositum of real cyclotomic subfields on which the new inertia cocycles vanish. The root classes needed for the path must survive this same restriction.

\begin{auxlemma}[Restriction of a normalized root cocycle]
\label{lem:root-cocycle-ell-power}
Let \(k\) be a field of characteristic \(\ell\), \(G\) a group,
\(\theta:G\to k^\times\) a character and \(u:G\to k\) a nonzero
cocycle for the action by \(\theta\). Suppose \(u(\tau)=0\) and
\(\theta(\tau)\ne1\) for some \(\tau\in G\). For any normal subgroup \(H\) the nonvanishing test is 
\[
 \exp(G/H)\mid\ell^a\quad\Longrightarrow\quad u|_H\ne0.
\]

\end{auxlemma}

\begin{proof}
Choose \(g\) with \(u(g)\ne0\). When \(\theta(g)=1\), replace \(g\)
by \(\tau g\). Its cocycle value is \(\theta(\tau)u(g)\ne0\) and
its character value is \(\theta(\tau)\ne1\).
For the resulting \(g\) the cocycle identity and the Frobenius map give
\[
 u(g^{\ell^a})=u(g)\sum_{j=0}^{\ell^a-1}\theta(g)^j
 =u(g)\frac{\theta(g)^{\ell^a}-1}{\theta(g)-1}
 =u(g)(\theta(g)-1)^{\ell^a-1}\ne0 .
\]
The element \(g^{\ell^a}\) belongs to \(H\).
\end{proof}

\begin{proposition}[Restriction and transport of the marked flags]\label{prop:count-two-tame-return}
For each new rational prime \(w\) of the auxiliary datum of
Proposition~\ref{prop:count-two-auxiliary-datum}, adjoin the real
cyclic degree-\(\ell\) subfield of \(\mathbb Q(\mu_w)\) to \(F_0\).
Let \(F_1\) be the resulting compositum. For each of these new
places \(v\) of \(F_0\), with tame inertia cocycle \(z_v\), restriction gives
\[
                          z_v|_{I_v\cap G_{F_1}}=0 .
\]
Passing to a further finite real cyclotomic layer preserves the
required nonzero root entries, the diagonal normalization and the
common \((2,0)\) frame coordinates of the target and auxiliary datum.
\end{proposition}

\begin{proof}
\proofstep{Kill the added tame inertia.}
For a new rational prime \(w\), every completion \(F_{0,v}\)
above \(w\) is \(\Q_w\). At its inertia, the old \(y\) and the
diagonal characters are trivial, and the new cocycle is a
homomorphism to an additive group killed by \(\ell\).
Write \(I_v\) for inertia and \(z\) for this restricted
additive cocycle. Since \(w\ne\ell\), wild inertia and the
prime-to-\(\ell\) factors of tame inertia have zero image. The
remaining factor is procyclic, so \(z\) factors as
\begin{equation}\label{eq:count-two-tame-quotient}
 I_v\longrightarrow I_v^{\mathrm{tame},(\ell)}\simeq\mathbb Z_\ell
  \longrightarrow\mathbb Z_\ell/\ell\mathbb Z_\ell
  \xrightarrow{\ \bar z\ }(k,+).
\end{equation} The cyclic
degree-\(\ell\) subfield of \(\Q(\mu_w)\) realizes that
quotient, since \(w\equiv1\pmod\ell\) and \(\Q(\mu_w)\) is
totally ramified at \(w\). Adjoin these subfields for all
new \(w\) to obtain \(F_1\). This extension is real and abelian
over \(\Q\); it is unramified over \(F_0\) at \(\ell\) and the
quotient of absolute groups over \(F_0\) has exponent dividing
\(\ell\).

By \eqref{eq:count-two-tame-quotient} every new inertia
restriction vanishes on \(G_{F_1}\). For \(F'=F_1\) take \(a=1\);
after adjoining a finite layer of the real cyclotomic
\(\mathbb Z_\ell\)-extension choose \(a\) so that the exponent
of \(\operatorname{Gal}(F'/F_0)\) divides \(\ell^a\). Then
\[
 \tau^{\ell^a}\in G_{F'},\qquad
 d_i(\tau^{\ell^a})-d_j(\tau^{\ell^a})
       =(d_i(\tau)-d_j(\tau))^{\ell^a}\ne0\quad(i\ne j).
\]
The restricted diagonal characters therefore still separate the labels.

\proofstep{Preserve the nonzero root entries.}
For each normalized cocycle \(u\) the normalizing element satisfies
\(u(\tau)=0\) and \(\theta(\tau)\ne1\). The quotient by the new absolute
group has exponent dividing a power of \(\ell\), so
Lemma~\ref{lem:root-cocycle-ell-power} preserves a nonzero value on
that group.
The same two-by-two multiplication applies to the retained
entries whenever the indicated block support is used.
Thus the restricted cocycle can be nonzero even if the original inertia element lies outside the subgroup.

\proofstep{Transport the local flags.}
At a place of a deeper layer, take a conjugator \(g\) over
the earlier field carrying its marked predecessor to the
place in question. Transport the flag by \(r(g)N_v\) and
take its lower unipotent factor. For \(A=r(g)N_v\) the successive determinants in the flag calculation are
\[
 \Delta_0(A)=1,\qquad \Delta_m(A)=\det(A_{ij})_{0\le i,j<m},
\]
with indices starting at zero throughout.

In the support
of the first line, the first row has no off-diagonal terms and
\(A_{21}=0\). Gaussian elimination therefore gives
\[
 \Delta_{i+1}/\Delta_i=d_i(g)\quad(0\le i\le2),\qquad
 \Delta_4/\Delta_3=\det(A)/(d_0(g)d_1(g)d_2(g)).
\]
These four pivots are units. Hence the decomposition into a lower
unipotent and an upper triangular factor is unique.
The symplectic adjoint reverses products and preserves
upper triangular and lower triangular types. Applying it
to the factorization shows by uniqueness that the lower
unipotent factor is symplectic. A different conjugator
differs on the right by an element of the original
decomposition group and hence gives the same flag.
In the transported frame the \((2,0)\) entry is
\[
             \frac{r(g)_{20}+d_2(g)(N_v)_{20}}{d_0(g)} .
\]
It is identical for the target and the auxiliary datum.
The \((2,1)\) entries remain zero. The local upper
triangular matrices have the original characters pulled
back, since conjugation inside the earlier Galois group
does not change those characters.

The same Gaussian construction carries the integral
target to each marked place. The leading principal
minors are units since their reductions are the
pivots just computed. Its flag characters are the
original rational characters restricted to inertia.
The rational point kills the inertia on the residual
kernel outside its marked set. At an original exception allowed by the target, the inertia and tower conditions give
\[
 \bar\Psi=1,\qquad\kappa=1\quad\Longrightarrow\quad\psi_0=1
\]
on the relevant inertia. Hence the rational point's other inertia condition applies there as well.
On inertia above \(\ell\) in the cyclotomic kernel,
the trivial residual labels and the prime-to-\(\ell\)
order of \(\psi_0\) on the deep tower give all four
ordered inertia characters equal to one.
\label{cor:cohomology-count-two-chain-3}
\end{proof}

\paragraph{The two Levi blocks and the three lines.}
\label{par:count-two-lines}
We can now read the path from its block decompositions. At each step the endpoints agree on the diagonal blocks. The middle step changes which
pair is held fixed.
\begin{construction}
\begin{tikzpicture}[>=Stealth,node distance=7mm,every node/.style={inner sep=2pt}]
\node (target) {Restricted rational target};
\node[below=of target] (aux) {Auxiliary datum};
\node[below=of aux] (proj) {One-root projection};
\node[below=of proj] (seed) {Mixed Hilbert seed};
\draw[->] (target) -- node[right=3mm] {\(\{0,2\},\{1,3\}\)} (aux);
\draw[->] (aux) -- node[right=3mm] {\(\{0,1\},\{2,3\}\)} (proj);
\draw[->] (proj) -- node[right=3mm] {\(\{0,2\},\{1,3\}\)} (seed);
\end{tikzpicture}
\caption{Block orders for the three affine lines of Proposition~\ref{prop:count-two-three-lines}. In the indicated order the endpoint differences belong to the square-zero off-block ideal.}\label{fig:analytic-three-lines}
\end{construction}

\begin{proposition}[Three lines joining the rational target to the seed]\label{prop:count-two-three-lines}
In the even-order case, after the restriction of
Proposition~\ref{prop:count-two-tame-return} and at a deep layer
where the mixed seed is available, the rational target and that
seed are joined by affine interpolation from target to auxiliary datum,
then from the auxiliary datum to its one-root projection, and finally
from that projection to the seed. Write \(D_0,D_1,D_2,D_3\) for the four marked data in that order.
There are families \(\mathcal D_i\) over \(k[t]\), affine in the matrix
and frame coordinates, with
\[
           \mathcal D_i(0)=D_i,\qquad \mathcal D_i(1)=D_{i+1}
                         \qquad(0\le i<3).
\]
They preserve the marked flags, similitude,
diagonal normalization and the prescribed inertia characters.
\end{proposition}

\begin{proof}
\proofstep{Project to the one-root datum.}
In the even-order case the character \(B\) is trivial on the
decomposition groups above \(\ell\). On the deeper tower,
project the auxiliary global matrices to their diagonal
blocks in the order \(\{0,1\},\{2,3\}\). Do the same
to the local frames. These projections respect products
on the block lower triangular algebra; the symplectic
form pairs the two blocks, so the projected matrices and
frames still satisfy the same similitude and flag
identities. The projected residual representation has
only the off-diagonal entries of degree \(-w_\alpha\),
including the nonzero entry in position \((1,0)\).
In its projected frame the local triangular matrix is
diagonal, because the representation and frame are
lower triangular as well.

At each marked place, fix the frame's \((2,0)\) entry using the seed value \(e_v\):
\[
 (N_v)_{20}=e_v .
\]
Right-multiply the projected frame by \(U(0,e_v,0)\).
On the marked decomposition group the first and third diagonal characters agree. Cancelling this common character in the paired-product identity gives
\[
 d_0=d_2,\qquad d_0d_3=d_1d_2\quad\Longrightarrow\quad d_1=d_3 .
\] Thus the frame multiplier commutes with the
diagonal local matrix. It preserves the \((1,0)\) entry
and inserts \(e_v\) in position \((2,0)\).

\proofstep{Construct the three lines.}
The first line joins the restricted target to the
auxiliary datum. The second joins the auxiliary datum
to this one-root projection with the modified frame.
The third joins the projection to the reduced seed.
For the first and third lines, order the blocks as
\(\{0,2\},\{1,3\}\). Both endpoints have the same
diagonal blocks, by the common \((2,0)\) entry, the
paired \((3,1)\) entry determined by the alternating
form, and the fixed diagonal characters. In the third line the two endpoints have identical readings in the global and the frame coordinates at position \((2,0)\):
\[
 (\text{global entry},\text{marked frame entry})=(0,e_v).
\] The second line uses
\(\{0,1\},\{2,3\}\), with the same diagonal blocks by
the common \((1,0)\) entry and its paired \((3,2)\)
entry.

In each order, all differences are in the
square-zero ideal of strictly off-block matrices.
The adjoint preserves this ideal, since the form
pairs the two ordered blocks. To verify the interpolation, write any interpolated matrix
as \(M_0+t\Delta M\), with \(\Delta M\) in the square-zero ideal.
For two such matrices multiplication gives
\[
 (M_0+t\Delta M)(N_0+t\Delta N)
 =M_0N_0+t(M_0\Delta N+\Delta M N_0)
 =(1-t)M_0N_0+tM_1N_1 .
\]
The adjoint is linear and preserves the ideal, so the same
calculation applies to both the product identities and the
similitude identity. Applying the equality to the individual basis columns \(e_i\) reads the marked flag along each line:
\[
 R(s)(N_v e_i)=N_vD_v(s)e_i .
\] The
interpolants are defined over \(k[t]\) and preserve
the normalization and all marked characters.

The original target vertex uses the two cases of the
lattice and restriction calculation above. The auxiliary
vertex retains the nonzero \(-w_\Psi\) if present,
and otherwise has the same nonzero marked frame entry
as the target and the new \(-w_\alpha\) entry. The
third vertex uses the nonzero \(-w_\alpha\) entry and
the seed's nonzero frame entry. For the seed, let \(S_{\ne0}\) be the set of weights of nonzero entries in its prescribed support. The two allowed alternatives are
\[
 -w_\alpha\in S_{\ne0}\qquad\text{or}\qquad \{w_\beta,-w_\vartheta\}\subseteq S_{\ne0}.
\]
In the second alternative the \(w_\alpha\)-character must also be nontrivial on the marked local group.

These choices put the auxiliary fields before the
bounds on the deeper cyclotomic tower. The finite
set of ramified auxiliary primes and \(F_1\) are fixed
using the target over \(F_0\). The cocycle bound \(C\)
is taken at this tame base \(F_1\). Choose the deeper layer only after obtaining \(C\) over \(F_1\). If \(d_{\rm layer}\) denotes its degree, impose both the depth condition and
\[
 d_{\rm layer}\ge30(C+1)^2
\]
before constructing the mixed seed on that layer. The cup lift, the ramification removal and the
transport of the target are independent of this later
seed. Its frame entry is used only in the last two lines.
\label{cor:cohomology-count-two-chain-4}
\end{proof}

\endgroup
